% draft snapshot: fixed photon mass; sections as of the snapshot time
% preamble.tex -- shared macros for the paper (input by main.tex before \begin{document})
\documentclass[11pt,a4paper]{amsart}
\usepackage[utf8]{inputenc}
\usepackage[T1]{fontenc}
\usepackage{amsmath,amssymb,amsthm,mathrsfs}
\usepackage[margin=1in]{geometry}
\usepackage{enumitem}
\usepackage{xcolor}   % needed for the colour expressions in hyperref's options below
\usepackage[colorlinks=true,linkcolor=blue!50!black,citecolor=blue!50!black,urlcolor=blue!50!black]{hyperref}
\usepackage{cleveref}
\crefname{hypothesis}{Hypothesis}{Hypotheses}
\Crefname{hypothesis}{Hypothesis}{Hypotheses}
\crefname{convention}{Convention}{Conventions}
\Crefname{convention}{Convention}{Conventions}

% ---------- theorem environments (numbered within sections) ----------
\theoremstyle{plain}
\newtheorem{theorem}{Theorem}[section]
\newtheorem{proposition}[theorem]{Proposition}
\newtheorem{lemma}[theorem]{Lemma}
\newtheorem{corollary}[theorem]{Corollary}
\theoremstyle{definition}
\newtheorem{definition}[theorem]{Definition}
\newtheorem{hypothesis}[theorem]{Hypothesis}
\newtheorem{convention}[theorem]{Convention}
\theoremstyle{remark}
\newtheorem{remark}[theorem]{Remark}
\newtheorem{example}[theorem]{Example}

% ---------- sets, fields, operators ----------
\newcommand{\R}{\ensuremath{\mathbb R}}
\newcommand{\C}{\ensuremath{\mathbb C}}
\newcommand{\Z}{\ensuremath{\mathbb Z}}
\newcommand{\N}{\ensuremath{\mathbb N}}
\newcommand{\E}{\ensuremath{\mathbb E}}
\newcommand{\cS}{\ensuremath{\mathcal S}}
\newcommand{\cD}{\ensuremath{\mathcal D}}
\newcommand{\cK}{\ensuremath{\mathcal K}}
\newcommand{\cI}{\ensuremath{\mathcal I}}
\newcommand{\cH}{\ensuremath{\mathcal H}}
\newcommand{\cV}{\ensuremath{\mathcal V}}
\newcommand{\cF}{\ensuremath{\mathcal F}}
\newcommand{\cT}{\ensuremath{\mathcal T}}
\newcommand{\cZ}{\ensuremath{\mathcal Z}}
\newcommand{\cR}{\ensuremath{\mathcal R}}
\newcommand{\cM}{\ensuremath{\mathcal M}}
\newcommand{\cB}{\ensuremath{\mathcal B}}
\newcommand{\cA}{\ensuremath{\mathcal A}}
\newcommand{\cW}{\ensuremath{\mathcal W}}
\newcommand{\cG}{\ensuremath{\mathcal G}}
\newcommand{\cQ}{\ensuremath{\mathcal Q}}
\newcommand{\cC}{\ensuremath{\mathcal C}}
\newcommand{\cL}{\ensuremath{\mathcal L}}
\newcommand{\sK}{\ensuremath{\mathsf K}}
\newcommand{\sS}{\ensuremath{\mathsf S}}
\newcommand{\sW}{\ensuremath{\mathsf W}}
\newcommand{\amp}{\ensuremath{\mathsf a}}
\newcommand{\Wc}{\ensuremath{\mathsf W}}
\DeclareMathOperator{\supp}{supp}
\DeclareMathOperator{\dist}{dist}
\DeclareMathOperator{\Var}{Var}
\DeclareMathOperator{\Cov}{Cov}
\DeclareMathOperator{\vertex}{vert}
\DeclareMathOperator{\alive}{alive}
\DeclareMathOperator{\Tr}{Tr}
\DeclareMathOperator{\diag}{diag}
\DeclareMathOperator{\curl}{curl}
\renewcommand{\Re}{\operatorname{Re}}
\renewcommand{\Im}{\operatorname{Im}}
\newcommand{\eps}{\ensuremath{\varepsilon}}
\newcommand{\Lam}{\ensuremath{\Lambda}}
\newcommand{\ND}{\ensuremath{\mathrm{ND}}}
\newcommand{\Oh}{\ensuremath{O_h}}
\newcommand{\Oth}{O(3)}
\newcommand{\Euc}{\ensuremath{\R^3\rtimes O(3)}}
\newcommand{\Cub}{\ensuremath{\R^3\rtimes O_h}}
\newcommand{\dstar}{\ensuremath{d^{*}}}
\newcommand{\ie}{i.e.\ }
\newcommand{\eg}{e.g.\ }
\newcommand{\lmu}{\ensuremath{\mu_0}}            % the fixed photon mass
\newcommand{\mus}{\ensuremath{\mu_*}}            % the reference photon mass of the frozen family
\newcommand{\dm}{\ensuremath{\delta m^2}}        % the mass counterterm

% ---------- lattices, cubes, states ----------
\newcommand{\Lk}{\ensuremath{\Lambda_k}}                 % \eps_k Z^3
\newcommand{\CJ}{\ensuremath{C_J}}                       % the cube [-MJ,MJ)^3
\newcommand{\CK}{\ensuremath{C^{K}_{J}}}                 % King's RG cube (as a state convention)
\newcommand{\XJ}{\ensuremath{X_J}}                       % C_J \cap \Lambda_k
\newcommand{\TJ}{\ensuremath{T_J}}                       % the three top face layers
\newcommand{\NDK}{\ensuremath{\ND^{K}}}                  % King's convention (scalar killed on T_J)
\newcommand{\NDa}{\ensuremath{\ND^{(a)}}}                % comparison convention (scalar alive on X_J)
\newcommand{\avK}[1]{\ensuremath{\langle #1\rangle^{K}_{k,J}}}
\newcommand{\ava}[1]{\ensuremath{\langle #1\rangle^{(a)}_{k,J}}}
\newcommand{\av}[1]{\ensuremath{\langle #1\rangle}}
\newcommand{\ZK}{\ensuremath{Z^{K}}}                     % King's printed state functional
\newcommand{\ZKinf}{\ensuremath{Z^{K}_{\infty}}}
\newcommand{\ZKJ}{\ensuremath{Z^{K}_{J}}}
\newcommand{\Ginf}{\ensuremath{G_\infty}}                % the tube object (Leg W)
\newcommand{\ZLIM}[1]{\ensuremath{Z^{#1}_{\infty}}}      % LIM state of a continuum subsequence
\newcommand{\Dtube}{\ensuremath{\cD_{\sigma_+,H}}}       % OS0's Wick tube
\newcommand{\Wick}[1]{\ensuremath{{:}\phi^2{:}(#1)}}
\newcommand{\Fcurv}[1]{F(#1)}
\newcommand{\norm}[1]{\ensuremath{\lVert #1\rVert}}
\newcommand{\abs}[1]{\ensuremath{\lvert #1\rvert}}
\newcommand{\wnorm}[1]{\ensuremath{\lvert #1\rvert_{w}}}
\newcommand{\fournorm}[1]{\ensuremath{\lvert #1\rvert_{4}}}
\newcommand{\ip}[2]{\ensuremath{\langle #1,#2\rangle}}

% ---------- named hypotheses (always typeset through these macros) ----------
\newcommand{\hyp}[1]{\ensuremath{(\mathrm{#1})}}
\newcommand{\EC}{\ensuremath{(\mathrm{E}_{C})}}
\newcommand{\ECK}{\ensuremath{(\mathrm{E}_{C})^{K}}}
\newcommand{\ECa}{\ensuremath{(\mathrm{E}_{C})^{(a)}}}
\newcommand{\Wone}{\ensuremath{(\mathrm{W1})}}
\newcommand{\Wtwo}{\ensuremath{(\mathrm{W2})}}
\newcommand{\Tall}{\ensuremath{(\mathrm{T})}}
\newcommand{\Tone}{\ensuremath{(\mathrm{T1})}}
\newcommand{\Ttwo}{\ensuremath{(\mathrm{T2})}}
\newcommand{\Tthree}{\ensuremath{(\mathrm{T3})}}
\newcommand{\BALm}{\ensuremath{(\mathrm{BAL}\text{-}m)}}
\newcommand{\RC}{\ensuremath{(\mathrm{R}_{C})}}
\newcommand{\RW}{\ensuremath{(\mathrm{R}_{W})}}
\newcommand{\RCfr}{\ensuremath{(\mathrm{R}_{C})^{\mathrm{fr}}}}
\newcommand{\RWfr}{\ensuremath{(\mathrm{R}_{W})^{\mathrm{fr},(a)}}}
\newcommand{\TWK}{\ensuremath{(\mathrm{TW}^{K})}}        % item 51: the Wick tube on King's cube
\newcommand{\FAK}{\ensuremath{(\mathrm{FA}^{K})}}        % item 53: Form (A) on King's cube (optional extension)
\newcommand{\Emu}{\ensuremath{(\mathrm{E}_{\mu})}}
\newcommand{\EmuK}{\ensuremath{(\mathrm{E}_{\mu})^{K}}}
\newcommand{\EmuI}{\ensuremath{(\mathrm{E}_{\mu})^{\cI}}}
\newcommand{\EmuW}{\ensuremath{(\mathrm{E}_{\mu})^{W}}}
\newcommand{\NEGz}{\ensuremath{(\mathrm{NEG}_{0})}}
\newcommand{\NEGs}[1]{\ensuremath{(\mathrm{NEG}_{0})[#1]}}
\newcommand{\NEGlim}{\ensuremath{(\mathrm{NEG}_{\lim})}}
\newcommand{\CTmu}{\ensuremath{(\mathrm{CT}_{\mu})}}
\newcommand{\CTmufour}{\ensuremath{(\mathrm{CT}_{\mu})^{(4)}}}
\newcommand{\hypAlph}{\ensuremath{(\alpha)}}   % named hypothesis (alpha); not \Alph, which is LaTeX's counter formatter
% \Alph is LaTeX's counter formatter (\Alph{enumi}); the sections also write \Alph for the hypothesis (alpha).
% Dispatch on the next token: a brace means the counter formatter, anything else the hypothesis.
\makeatletter
\makeatother
\newcommand{\CONT}{\ensuremath{(\mathrm{CONT}_{\mu,J})}}
\newcommand{\Kfr}{\ensuremath{(\mathrm{K}^{\mathrm{fr}}_{\cI})}}
\newcommand{\Hrelunif}{\ensuremath{(\mathrm{H0rel})^{\mathrm{unif}}}}
\newcommand{\Kposunif}{\ensuremath{(\mathrm{K}\text{-}\mathrm{pos})^{\mathrm{unif}}_{N}}}
\newcommand{\KfrWunif}{\ensuremath{(\mathrm{K}^{\mathrm{fr}}_{W})^{\mathrm{unif}}}}
\newcommand{\TR}{\ensuremath{(\mathrm{TR})}}
\newcommand{\TRW}{\ensuremath{(\mathrm{TR})^{W}}}
\newcommand{\Bcell}{\ensuremath{(\mathrm{B}_{\mathrm{cell}})}}
\newcommand{\SCH}{\ensuremath{(\mathrm{SCH})}}
\newcommand{\UNIQ}{\ensuremath{(\mathrm{UNIQ})}}
\newcommand{\PBIR}{\ensuremath{(\mathrm{PB}_{\mathrm{IR}})}}
\newcommand{\PBmus}{\ensuremath{(\mathrm{PB}_{\mus})}}
\newcommand{\NTtau}{\ensuremath{(\mathrm{NT}_{\tau})}}
\newcommand{\hypR}{\ensuremath{(\mathrm{R})}}
\newcommand{\hypRK}{\ensuremath{(\mathrm{R})^{K}}}
\newcommand{\CFK}{\ensuremath{(\mathrm{C}\text{-}\mathrm{FK})}}
\newcommand{\Bminus}[1]{\ensuremath{(\mathrm{B}^{-})[#1]}}
\newcommand{\Itwo}{\ensuremath{(\mathrm{I2})}}
\newcommand{\Ithree}{\ensuremath{(\mathrm{I3})}}
\newcommand{\Ifour}{\ensuremath{(\mathrm{I4})}}
\newcommand{\UVS}{\ensuremath{(\mathrm{UVS}^{+}_{\infty})}}

% ---------- Osterwalder--Schrader axiom labels ----------
\newcommand{\OSz}{\ensuremath{\mathrm{OS0}}}
\newcommand{\OSo}{\ensuremath{\mathrm{OS1}}}
\newcommand{\OSt}{\ensuremath{\mathrm{OS2}}}
\newcommand{\OSth}{\ensuremath{\mathrm{OS3}}}
\newcommand{\OSf}{\ensuremath{\mathrm{OS4}}}
\newcommand{\Ezz}{\ensuremath{\mathrm{E0}''}}

% ---------- small conveniences ----------
\newcommand{\defeq}{:=}
\newcommand{\inv}{\ensuremath{^{-1}}}
\newcommand{\half}{\ensuremath{\tfrac12}}
\newcommand{\lmt}[1]{\ensuremath{\lim_{#1}}}
\allowdisplaybreaks
\providecommand{\cP}{\mathcal{P}}

% ---------- visible placeholder for values or lists still changing (draft only) ----------
\newcommand{\moving}[1]{\textcolor{orange!80!black}{[#1]}}

% preamble_rot.tex -- additional macros for the rotation part (sections r01..r07).
% To be \input by main.tex right after preamble.tex. No macro here clashes with preamble.tex
% (checked against every \newcommand there and against the macro names used in sec/s01..s23).
% Conventions: the shell mass is written \varkappa (never \kappa, which is the tube constant of s10--s11);
% the source box of the heavy shell is \Hbox (never H, the tube half-width of s11); the unnormalised torus
% functional is \Zun (never \cZ, the characteristic functional of s09); the pair space of real sources is \cE.

% ---------- geometry of the rotation part ----------
\newcommand{\thz}{\ensuremath{\theta_0}}                       % arctan(3/4)
\newcommand{\Rgen}{\ensuremath{\mathcal R_0}}                  % the generating set {R_{mu nu}(+-theta_0)}
\newcommand{\Gz}{\ensuremath{G_0}}                             % the group generated by \Rgen (dense in SO(3))
\renewcommand{\DJ}{\ensuremath{D_J}}                         % King's torus R^3/(50MJ Z)^3 (renew: T1 fontenc defines \DJ as the letter Dj, unused here)
\newcommand{\GamJ}{\ensuremath{\Gamma_J}}                      % its fundamental cell [-25MJ,25MJ)^3
\newcommand{\Qb}[1]{\ensuremath{Q_{#1}}}                       % the cube [-b,b]^3 (as Q_R in s02)
\newcommand{\Nr}[2]{\ensuremath{N_{#1}(#2)}}                   % the closed l^infty r-neighbourhood
\newcommand{\Bbar}[1]{\ensuremath{\bar B_{#1}}}                % closed Euclidean ball about 0

% ---------- walled tori ----------
\newcommand{\fc}{\ensuremath{\mathfrak c}}                     % a walled configuration (J',w) or (J',psi,n)
\newcommand{\fs}{\ensuremath{\mathfrak s}}                     % a shell configuration (\fc,b',L)
\newcommand{\Aw}[1]{\ensuremath{\cA(#1)}}                      % admissible walled configurations for the set #1
\newcommand{\Wv}{\ensuremath{\cW}}                             % walled value \Wv^a_\fc(f,g)
\newcommand{\Wenv}{\ensuremath{\cW^{\mathrm{env}}}}            % the lower envelope (curvature sector)
\newcommand{\Wenva}{\ensuremath{\cW^{a,\mathrm{env}}}}         % the lower envelope in the scheme a
\newcommand{\Tw}[2]{\ensuremath{T^{#1}_{#2}}}                  % walled ratio T^a_{\fc,\Lambda}(f,g)
\newcommand{\Zun}{\ensuremath{\mathsf Z}}                      % unnormalised torus functional \Zun_{\Lambda,T}(s)
\newcommand{\cE}{\ensuremath{\mathcal E}}                      % \cE_\rho = pairs of real sources supported in \Bbar\rho
\newcommand{\env}{\mathrm{env}}

% ---------- the heavy shell ----------
\newcommand{\Hbox}{\ensuremath{\mathsf H}}                     % the source box (a closed cube)
\newcommand{\Hboxp}{\ensuremath{\mathsf H'}}                   % the source box of the exterior-shell configuration
\newcommand{\Psik}{\ensuremath{\Psi_k}}                        % two-level expectation at spacing eps_k
\newcommand{\Psio}{\ensuremath{\Psi^0}}                        % its continuum limit
\newcommand{\tail}{\operatorname{tail}}
\newcommand{\Dchi}{\ensuremath{D_\chi}}                        % the dead set {chi>0}
\newcommand{\tauGB}{\ensuremath{\tau}}                         % the photon-switch constant tau(\fs)
\newcommand{\srw}{\ensuremath{\mathsf a}}                      % sup_z |z|_infty G_{Z^3}(0,z)
\newcommand{\vol}{\operatorname{vol}}
\newcommand{\Icap}{\ensuremath{I^{\mathrm{cap}}}}              % the capped stopping-scale integral of (HS)
\newcommand{\cJ}{\ensuremath{\mathcal J}}                      % the ratio \cJ_\varkappa(z) in the proof of the shell theorem
\newcommand{\Mdisc}{\ensuremath{M_\circ}}                      % the disc constant (not the block integer M)
\newcommand{\sfa}{\ensuremath{\mathfrak a}}                    % \sfa(eta_0), \sfb(eta_0): the two-sided bound's errors
\newcommand{\sfb}{\ensuremath{\mathfrak b}}
\newcommand{\Xdiff}{\ensuremath{\varXi}}                       % the entire difference \Xdiff = Psi_K - Psi_infty (\Xi itself is LaTeX's letter)

% ---------- named hypotheses of the rotation part (always typeset through these macros) ----------
\newcommand{\Kz}{\ensuremath{(\mathrm{K}_0)}}                  % the torus comparison, curvature sector
\newcommand{\Ka}{\ensuremath{(\mathrm{K})_a}}                  % the torus comparison, all sources, scheme a
\newcommand{\TIz}{\ensuremath{(\mathrm{TIGHT}_0)}}             % the closing hypothesis, curvature sector
\newcommand{\TIa}{\ensuremath{(\mathrm{TIGHT})_a}}             % the closing hypothesis, scheme a
\newcommand{\TItwo}{\ensuremath{(\mathrm{TIGHT}_2)}}           % the two-point form
\newcommand{\REz}{\ensuremath{(\mathrm{RES}_0)}}               % heavy-shell decoupling at one configuration, curvature
\newcommand{\REa}{\ensuremath{(\mathrm{RES})_a}}               % the same, scheme a
\newcommand{\REcont}{\ensuremath{(\mathrm{RES}^0)}}            % its continuum form
\newcommand{\SHz}{\ensuremath{(\mathrm{SH}_0)}}                % shell-kill decoupling, curvature
\newcommand{\SHa}{\ensuremath{(\mathrm{SH})_a}}                % the same, scheme a
\newcommand{\HS}{\ensuremath{(\mathrm{HS})}}                   % the heavy-shell estimates (two-sided stopping-scale bound + kappa-Cauchy property)
\newcommand{\DSh}{\ensuremath{(\mathrm{DS})}}                  % the dead-shell chain for the exterior shell
\newcommand{\hypF}{\ensuremath{(\mathrm F)}}                   % one-sided bound of King's type at one spacing
\newcommand{\hypFz}{\ensuremath{(\mathrm F_0)}}                % its curvature restriction
\newcommand{\hypV}{\ensuremath{(\mathrm V)}}                   % vacuum constants cancel
\newcommand{\hypCc}{\ensuremath{(\mathrm C)}}                  % source-linear constants cancel
\newcommand{\Wadm}[1]{\ensuremath{(\mathrm W)_{#1}}}           % admissibility of the scheme #1 = (W1)^(a)+(W2)^(a)
\newcommand{\Wsh}{\ensuremath{(\mathrm W^\sharp)}}             % the general-source scheme hypothesis
\newcommand{\Einf}{\ensuremath{(\mathrm E_\infty)}}            % existence of the dead-shell torus value (not consumed)


\title[OS axioms for King's cube state at fixed photon mass]{Osterwalder--Schrader axioms for King's cube state of the three-dimensional abelian Higgs model at fixed photon mass}
\author{}
\date{}

\begin{document}

\begin{abstract}
King constructed the continuum, infinite-volume limit of the three-dimensional $U(1)$ Higgs model along Neumann--Dirichlet cubes at block factor $5$ \cite{K1,K2}. For his printed state at fixed photon mass we prove the Osterwalder--Schrader axioms: regularity $\Ezz$ with absolute constants, invariance under the full Euclidean group, reflection positivity about every plane, symmetry, and the measure with its ergodic decomposition, on whose components the cluster axiom holds in all its forms. The existence of King's cube limits, the existence clause of his Theorem~2.1, is proved from his renormalisation-group analysis in a ratio form in which the vacuum energies cancel; rotation invariance is proved by comparing King's cubes with periodic tori carrying a smooth mass wall on two lattices whose orientations differ by $\arctan(3/4)$. With real scalar sources the same holds jointly for the curvature and the renormalised Wick square on a tube of small Wick amplitudes, without a Euclidean clause. The curvature results hold on King's whole coupling range; the Wick results hold under a one-sided sign clause of the Wick leg, which we prove on explicit coupling regions. The limit of vanishing photon mass is left to later work.
\end{abstract}

\maketitle

% s01 Introduction and main results
\section{Introduction and main results}

\subsection{King's construction}

King \cite{K1,K2} constructed the three-dimensional $U(1)$ Higgs model, scalar electrodynamics with a
massive photon, as a continuum and infinite-volume limit of lattice models. The lattice theory lives on
$\Lk=\eps_k\Z^3$ with spacings $\eps_k=5^{-k}$; its action has a photon $A$ in Feynman form with mass
$\lmu>0$, an $N$-component scalar $\phi$ with quartic coupling $\lambda>0$ and bare mass $m>0$, minimally
coupled with charge $e>0$, and one mass counterterm $\dm_{\eps_k}$. The counterterm is King's free-boundary
one: it is computed with infinite-lattice propagators, so it is one number per spacing, the same at every
site of every box and every torus. The ultraviolet problem is settled in \cite{K1} with Ba{\l}aban's
renormalisation group \cite{Ba1,Ba2,Ba3,Ba4} at block factor $5$: on a finite volume the effective
actions converge as $k\to\infty$. The infinite-volume problem is treated in \cite{K2} with the correlation
inequalities of Brydges, Fr\"ohlich and Seiler \cite{BFS1}. King's finite-volume states are
Neumann--Dirichlet states on the cubes $C_J=[-MJ,MJ)^3$, $M\ge1$ his block integer: the photon gradient
terms are kept only inside the cube (Neumann), and the scalar field vanishes on a boundary layer
(Dirichlet). The limit is iterated, first $k\to\infty$ at fixed $J$, then $J\to\infty$; the second limit
exists because the cube states are monotone in $J$ on a cone of observables, by \cite{BFS1}.

This paper proves the Osterwalder--Schrader axioms for King's printed state at fixed photon mass, with
full Euclidean invariance, and the corresponding statements for the joint field of the curvature and the
renormalised Wick square on a tube of small Wick amplitudes. It is self-contained in the following sense.
The existence of King's cube limits, the existence clause of his Theorem~2.1 \cite{K2}, is proved in the
existence part of the paper (\cref{sec:k01}, and \moving{existence section with scalar sources} for the
version with scalar sources): King's route, the comparison of two spacings at a common block scale, is run
in a ratio form in which the vacuum energies cancel, so that the convergence of the vacuum energies, which
his printed argument uses and which is not available on the cube, is not needed. The passage from the cube
limits to the axioms in \cite{K2} is asserted at three places, and at each we proceed differently. First,
King's Neumann--Dirichlet cube form is not reflection symmetric about any plane, and its natural
symmetrisation is not reflection positive (an exact counterexample at $e=0$, \Cref{rem:s07:R1}), so his
reflection positivity goes through a modified box state; that this box state is reflection positive is his
Proposition~3.3, which he derives from \cite[Cor.~5.4]{BFS1}, whose hypotheses do not cover the finite-box
Gaussian of the modified action. Proposition~3.3 is true, and it is proved here (\Cref{lem:s07:RPbox})
rather than derived from \cite[Cor.~5.4]{BFS1}. Second, the inequality by which he passes from the box
states back to the cube limit holds in the reverse direction for the class of observables his argument
needs: weakly, at every coupling, for the box state with the site-independent counterterm, and strictly at
$e=0$ (\Cref{lem:s07:R}); it is replaced by a two-sided sandwich between cube states with a margin of order
$\lmu^{-1}$ (\Cref{lem:s07:L,lem:s07:U}). Third, the continuum subsequences through which he reaches the
limit theory are built here (\Cref{rem:s12:K9}). With scalar sources, finally, his Theorem~2.1 needs a
vacuum counterterm that his Wick ordering does not contain; we work with a scheme $a$ and prove the
existence of the cube limits with sources in that scheme.

The authors' earlier paper \cite{SymPhase} constructs torus limits of the same model at block factor $2$
and proves the Osterwalder--Schrader axioms with cubic symmetry for them; its abstract measure theory (the
reconstruction of the measure, its ergodic decomposition and the weak clustering of the components) is used
here as it stands, cited by label. Rotation invariance is proved in the rotation part of this paper, whose
main theorem precedes the main theorem of the paper: King's cubes are compared with periodic tori carrying
a smooth mass wall, the walled tori on two lattices whose orientations differ by $\arctan(3/4)$ are
compared, and King's state is identified as the lower envelope of the walled tori through a heavy-shell
decoupling estimate on King's block-spin flow with the shell mass inside the matter covariance.

\subsection{The two cube conventions}

King's cube state $\NDK_{k,J}$ (\Cref{def:s02:Kcube}) has its photon variables on $X_J=C_J\cap\Lk$ and its
scalar alive on $\{x\in X_J:\ x_\nu<MJ-\eps_k\ \text{for all }\nu\}$, that is, killed on the three top
face layers $T_J$ of the cube. The comparison cube state $\NDa_{k,J}$ (\Cref{def:s02:acube}) has the
scalar alive on all of $X_J$. The Wick-tube lemmas and the lattice comparisons of the rotation part are
written for the comparison convention, while King's printed state is on his own. The two are interlaced:
for observables of the cone class, with supports at $\ell^\infty$-distance $>1$ from the complement of the
cube,
\[
\avK{\Theta}\ \le\ \ava{\Theta}\ \le\ \langle\Theta\rangle^{K}_{k,J+1},
\]
and reversely for the opposite class (\Cref{lem:s04:G}). Hence both families have the same $J$-limit at
fixed $k$; the comparison family's fixed-cube limits inferior and superior are squeezed to King's printed
state $\ZKinf$ as $J\to\infty$, and every continuum-subsequence limit of the comparison family equals
$\ZKinf$ (\Cref{cor:s04:G2}). What does not transfer is the existence of the fixed-$J$ limit $k\to\infty$
itself: on one convention it does not imply it on the other. This is why the existence theorems are proved
on King's convention, where the printed state lives and where King's block-spin flow runs; the Wick tube,
proved on the comparison convention, is carried to King's convention by a transfer theorem
(\Cref{thm:s13:TWK}).

\subsection{The main theorem}

Sources are real $2$-forms $f_j\in C^\infty_c(\R^3;\Lambda^2)$ with coefficients $z_j$ and real scalar test
functions $h_j\in C^\infty_c(\R^3)$ with coefficients $w_j$; $F(f)$ is the curvature paired with $f$,
$W^a(h)$ the renormalised Wick square in the scheme $a$, and $g(z)=\sum_j(\Re z_j)f_j$,
$k(\zeta)=\sum_j(\Re w_j)h_j=k_+-k_-$, with the norms $\norm{\cdot}_2$, $\fournorm{\cdot}$,
$\wnorm{\cdot}$ and the tube $\Dtube$ of \Cref{def:s02:sources,def:s02:tube}. The Schwinger functions
$S_n$, $S_{n,p}$ are the Taylor coefficients of the limit functionals at $0$. A scheme $a$ is admissible
if $a=z$, the reference scheme, or $z_k-a_k$ converges (\Cref{def:s03:adm}). The following summary is
restated with every clause in full and proved as \Cref{thm:s13:main}.

\begin{theorem}[Main Theorem, summary]\label{thm:s01:main}
Let $e>0$, $\lambda,m,\lmu>0$ and $M\ge1$ be arbitrary.

\medskip\noindent\textbf{(A) The curvature field.} King's printed state
$\ZKinf(f,0)=\lim_{J}\lim_{k}\avK{e^{F(f)}}$ exists for every real $2$-form $f$, extends to an entire
function of the coefficients of every finite family, and satisfies:
\begin{enumerate}[label=\textup{(\roman*)},leftmargin=*]
\item \textup{(\OSz)} $\abs{\ZKinf(\sum_jz_jf_j,0)}\le e^{\norm{g(z)}_2^2}$ for all $z\in\C^N$, with no
dependence on $\lmu$, $e$, $\lambda$, $m$ or $M$, and
$\abs{\ZKinf(\sum_jz_jf_j,0)}\le\exp\big(\tfrac12\lmu^{-2}\norm{\sum_j(\Re z_j)\dstar f_j}_2^2\big)$;
for real test functions $\abs{S_n(f_1,\dots,f_n)}\le(\sqrt2e)^n(n!)^{1/2}\prod_j\norm{f_j}_2
\le(e\pi/2)^n(n!)^{1/2}\prod_j\fournorm{f_j}$; the $S_n$ are tempered distributions and satisfy the
regularity axiom \Ezz\ with $s=4$, the constant $\sigma_n=(e\pi)^n(n!)^{1/2}$ for complex
$f_j\in\cS$, and every $\beta>\tfrac12$.
\item \textup{(\OSo)} Invariance under the full Euclidean group $\Euc$: under all translations of
$\R^3$, under the cubic group $\Oh$ about every point, and under all rotations, $F$ transforming as a
$2$-form; the $S_n$ are Euclidean covariant.
\item \textup{(\OSt)} Reflection positivity about every plane $\{x\cdot n=c\}$, $n$ a unit vector,
$c\in\R$, in matrix form for exponentials with arbitrary complex curvature coefficients and in polynomial
form.
\item \textup{(\OSth)} $S_n$ is real and symmetric, and $S_n=0$ for odd $n$.
\item \textup{(The measure)} For every even $r\ge14$ there is exactly one Borel probability measure
$\mu_F$ on $\cH^{-r}(\R^3;\Lambda^2)$ with moments $S_n$ and characteristic functional
$\ZKinf(\mathrm i\,\cdot,0)$ on $\cS$; its image on $H^{-r}_{\rm loc}$ is determined by the characteristic
functional and by the moments. The laws of the localised lattice fields $\chi F_{\Lk}$ under $\NDK_{k,J}$
converge weakly on $H^{-r}_{\rm per}(Q')$, $Q'$ a box containing $\supp\chi$, along the full sequence
$k\to\infty$ at every fixed $J\ge J_0(\chi)$, and then as $J\to\infty$, to the law of $\chi F$ under
$\mu_F$.
\item \textup{(\OSf\ on the ergodic components)} $\mu_F$ is $\Euc$-invariant and $\mu_F=\int\mu_\xi\,d\nu(\xi)$
with translation-ergodic components; $\nu$-almost every component is $\Euc$-invariant, reflection positive
about every plane, clusters in the mean and under diagonal separation, has a unique joint vacuum in every
direction, satisfies \OSf\ in all its forms for every choice of the time direction, and has real symmetric
moments with a factorial growth bound with component constants; the decomposition is $\Oh$-invariant, each
component being a mixture of $k\in\{1,2,3,6\}$ ergodic measures. The vanishing of the odd moments is not
claimed for a component.
\end{enumerate}

\medskip\noindent\textbf{(B) The joint field on the Wick tube.} Let the couplings satisfy the sign clause
\Tthree\ of \Cref{hyp:s03:T3}, $\sigma_1(M_1)>0$ for some $M_1\in(0,m)$; in particular, let them lie in
one of the explicit regions of \Cref{prop:s14:T3}. Let $a$ be an admissible scheme. Then for every
family of sources supported in $Q_R=[-R,R]^3$ the joint functional of $(F,W^a)$ has a limit $\Ginf$ on
the tube $\Dtube$, with
\[
\sigma_+=s_1,\qquad H=\tfrac{s_1}{16},\qquad \kappa=\tfrac{C_U}{s_1},
\]
and $K_-$ determined by the cell constant $B_{\rm cell}$ (\Cref{eq:s11:constants}), where
$s_1\in(0,\sigma_1(M_1))$ is the amplitude of the extensive bound and $C_U$ its constant
(\Cref{eq:s10:I2,prop:s10:U}); the constants depend on $(e,\lambda,m,\lmu,M,a)$ and on the auxiliary choices $M_1$,
$\chi_0$, $s_1$, and on nothing else, and they are not uniform in $\lmu$. The joint Schwinger functions
$S_{n,p}$ satisfy \textup{(i)--(vi)} in the following form:
\begin{enumerate}[label=\textup{(\roman*)},leftmargin=*]
\item $\abs{\Ginf(\zeta)}\le\exp\big(\norm{g}_2^2+\kappa\norm{k_+}_{L^1}+K_-\wnorm{k_-}\big)$ on
$\Dtube$, and for real test functions
$\abs{S_{n,p}}\le e^{K'\sigma_0}(\sqrt2e)^n(n!)^{1/2}(e/\sigma_0)^p\,p!\prod_i\norm{f_i}_2\prod_j\wnorm{h_j}$
with $\sigma_0=H/4$ and $K'=\kappa\pi^2+K_-$ (a factor $2^{n+p}$ for complex test functions); \Ezz\ holds
with $s=4$ and every $\beta>1$.
\item Translations and $\Oh$ about every point hold as identities on the tube, $W^a$ transforming as a
scalar; there is no Euclidean clause.
\item Reflection positivity about every coordinate plane for families whose Wick amplitudes satisfy
$\amp(\zeta_i)<\sigma_W=\sigma_+/8$, in particular for all purely imaginary Wick coefficients, and the
polynomial form.
\item $S_{n,p}=0$ for odd $n$.
\item[\textup{(v), (vi)}] As in \textup{(A)} for the measure on $\cH^{-r}(\R^3;\Lambda^2\oplus\R)$,
$r\ge14$ even, with the scheme-shifted lattice field in the convergence clause, and with the weak form of
\OSf\ (ergodicity under $\R^3$ and clustering in coordinate directions) on $\nu$-almost every component;
there is no Euclidean clause.
\end{enumerate}
The same statements hold for the tube object $G^K_\infty$ of King's printed state on his own convention
(\Cref{thm:s13:TWK,prop:s13:WK}).
\end{theorem}

The $\lmu$-free bound in (A)(i) is absolute: it depends on none of the parameters of the model. The
$\lmu^{-2}$ bound is the finite-volume Gaussian bound carried to the limit. Part (B) is the headline form
of the Wick leg; the extension of (B) to arbitrary Wick amplitudes, under the hypothesis \Wtwo\ at every
amplitude, with the only Euclidean clause of the Wick leg, is stated separately as \Cref{thm:s23:FormA}
among the open problems.

\subsection{The massless limit}

The photon mass $\lmu$ is fixed throughout: King's construction, and every statement above, is at
$\lmu>0$, and the constants of the Wick leg are not uniform in $\lmu$. The limit of vanishing photon mass,
along a family whose counterterm and scheme are frozen at a reference mass, is not treated in this paper;
it is listed among the open problems as later work.

\subsection{The coupling range}

The curvature leg holds on all of King's range: $e>0$ and $\lambda,m,\lmu>0$ arbitrary. Every smallness
condition in \cite{K1,K2,Ba1,Ba2,Ba3,Ba4} is a cap on the running couplings at a stopping scale that is
chosen after the point $(e,\lambda,m,\lmu)$ and the datum, so the quantifier over the couplings is the one
over the sources, and no sentence ``$e,\lambda$ sufficiently small'' occurs
(\Cref{lem:s14:CW,rem:s14:quantifier}). No relation $\lmu\le m$ is used or needed. The only coupling
condition in the paper is the one-sided sign clause \Tthree\ of the Wick leg, $\sigma_1(M_1)>0$ for some
$M_1\in(0,m)$; it fails exactly when $m$ minimises King's free-boundary bare coefficient $\alpha_\infty$ on
$(0,m]$, and to one-loop order it holds if and only if
\[
e^2<e_c^2(\lambda,m,\lmu)=\max\Big\{\frac{8\pi(m+\lmu)^2}{m+2\lmu}\Big(1+\frac{\lambda(N+2)}{2\pi m}\Big),\
4\pi(m+\lmu)\Big(1+\frac{\lambda(N+2)}{\pi m}\Big)\Big\}
\]
(\Cref{eq:s14:ec2}), so that a half-line of couplings fails it at one loop. Beyond one loop, \Tthree\ is
proved on explicit coupling regions, at $\lmu=m$ on a region containing $0\le\lambda,e^2\le0.1m$ and at
every $\lmu\in(0,m]$ on explicit boxes (\Cref{prop:s14:T3}); the Wick leg is stated on those regions, and
more generally wherever \Tthree\ holds. The mass clause in Ba{\l}aban's upper bound, a cap on the
dimensionless scalar mass at the stopping scales where his Proposition~3.1 is used, is part of that
proposition's statement and is one more one-sided cap on the stopping scale, not a condition on the
couplings (\Cref{rem:s14:BALm}).

\subsection{What is proved where}

The inputs of the main theorem are collected in \Cref{sec:s03}: the existence theorems for King's cube
limits without and with scalar sources (\Cref{hyp:s03:EC,hyp:s03:W1}), proved in the existence part; the
rotation theorem (\Cref{hyp:s03:RC}), proved in the rotation part; Ba{\l}aban's torus stability bounds and
the $c$-number lemma in the form the Wick leg uses (\Cref{hyp:s03:T1,hyp:s03:T2}); and the coupling
condition \Tthree. Everything else is proved in the sections between. Open are only the optional
extensions: the Wick sector at arbitrary amplitudes (\Wtwo, Form~(A), \Cref{thm:s23:FormA}), the relation
of King's state to the torus states of \cite{SymPhase} and its uniqueness and ergodicity, the limit of
vanishing photon mass, and non-triviality; they are collected in the final section.

\subsection{Structure of the paper}

The model, King's cube, the comparison cube, the sources and the standing lemmas (the correlation
inequalities of \cite{BFS1}, with the two inequalities that are only sketched there proved in full, Gaussian
domination, cutting and the elementary facts) come first, then the inputs with their provenance. The
interlacing of the two cube conventions (\Cref{lem:s04:G,cor:s04:G2}) and its three geometry margins
follow, and then the existence part: King's existence theorem on the cube (\Cref{sec:k01}) and its version
with scalar sources. The curvature sector is then built up in the order of the axioms: regularity with the
absolute constants (\Cref{thm:s05:C,thm:s05:CK}), translations, the cubic group and the rotation clause
(\Cref{thm:s06:S1K,thm:s06:O3}), reflection positivity through the modified box and its passage to the
limit with \OSth\ and the polynomial form (\Cref{lem:s07:RPbox,thm:s08:OS2C,prop:s08:OS3,prop:s08:poly}),
and the measure with its ergodic decomposition (\Cref{thm:s09:S2K}). The Wick tube is treated in two
sections, the torus inputs and the extensive bound (\Cref{prop:s10:RP,prop:s10:U}) and the uniform tube
bound with its limits (\Cref{thm:s11:W,lem:s11:TWtube}); continuum subsequences and the states they
define, with no existence input, form a section of their own (\Cref{thm:s12:LIM}). The rotation part
follows: the geometry of the two lattices and the walled tori, the admissible schemes and the full object,
the torus comparison and its reduction to one-sided bounds at one spacing, the closing estimate and the
invariance theorems, and the heavy shell in two sections; a short section collects the corrections to
King's papers that it uses. The main theorem is assembled in \Cref{thm:s13:main}, and the coupling range
is established in \Cref{lem:s14:CW,prop:s14:T3}. The paper ends with the open problems.

% s02 The model, King's cube and the standing lemmas
\section{The model, King's cube and the standing lemmas}

\subsection{Couplings, lattices and the action}

Throughout, $e>0$ and $\lambda,m,\lmu>0$ are the charge, the quartic coupling, the bare scalar mass and the photon mass; no further condition on the couplings is imposed, since every smallness condition of \cite{K1,K2,Ba1,Ba2,Ba3,Ba4} is a cap on the running couplings at a stopping scale chosen after the point $(e,\lambda,m,\lmu)$ and after the datum (the section on the coupling range). The block factor is $5$: the spacings are $\eps_k\defeq5^{-k}$, the lattices $\Lk\defeq\eps_k\Z^3$ with the standard frame $(e_\alpha)$, and $M\in\N$ is King's large-block integer, fixed once and for all. King's cubes are $\CJ\defeq[-MJ,MJ)^3$, $J\in\N$, and $\XJ\defeq\CJ\cap\Lk$ (the spacing understood from the context). A \emph{King torus} is $\R^3/(\ell\Z)^3$ with $\ell\in M\N$ even; $\Lk$ descends to it, with $\ell\,5^k$ sites per side.

The scalar field is $\phi(x)=r_xe^{\mathrm i\vartheta_x}\in\C$ and the photon field a real $A_\alpha(x)$ for every site $x$ and direction $\alpha$; the bond $\langle x,x+\eps e_\alpha\rangle$ carries $A_\alpha(x)$ and the reversed bond $-A_\alpha(x)$. King's action \cite[(2.5)]{K2} at spacing $\eps$ is
\begin{equation*}
\begin{aligned}
S&=\half\sum_{\alpha,\beta}\sum_x\eps\,\big(A_\alpha(x+\eps e_\beta)-A_\alpha(x)\big)^2+\frac{\lmu^2}2\sum_x\eps^3|A(x)|^2
+\half\sum_\beta\sum_x\eps\,\big|e^{\mathrm ie\eps A_\beta(x)}\phi(x+\eps e_\beta)-\phi(x)\big|^2\\
&\quad+\sum_x\eps^3\Big[\half\big(m^2+\dm_\eps\big)r_x^2+\lambda r_x^4\Big],
\end{aligned}
\end{equation*}
the photon in Feynman form with mass $\lmu$, and the hopping term of a bond $\langle x,y\rangle$, $y=x+\eps e_\beta$, equal to $\frac\eps2(r_x^2+r_y^2)-\eps r_xr_y\cos(\vartheta_x-e\eps A_\beta(x)-\vartheta_y)$: a \emph{diagonal part} and a \emph{cosine part}. The counterterm $\dm_\eps$ is King's free-boundary counterterm \cite[p.~330]{K2}, fixed by his two-point graphs of order at most four in $(e,\lambda)$ at zero momentum with the propagators of the infinite lattice $\eps\Z^3$: it is \emph{one number per spacing}, the same at every site of every box and torus. The Wick constant is the free one, $\sW_\eps\defeq(-\Delta_{\eps\Z^3}+m^2)^{-1}(0,0)$. This \emph{free convention} is used on cubes and tori alike unless said otherwise; King's own torus convention appears only inside the torus inputs of the Wick leg.

\subsection{The states}

\begin{definition}[Neumann--Dirichlet boxes, killed sites, the comparison cube]\label{def:s02:acube}
For a finite set $X\subset\Lk$ the \emph{Neumann--Dirichlet state} $\langle\cdot\rangle^{\ND}_{k,X}$ is the normalised measure $Z^{-1}e^{-S}\prod_{x\in X}d^2\phi(x)\prod_{x\in X,\alpha}dA_\alpha(x)$ with $\phi=0$ at every site outside $X$ (Dirichlet condition: a bond leaving $X$ contributes only its diagonal part $\frac\eps2r_x^2$ at its end $x\in X$), photon gradient terms only for pairs $x,x+\eps e_\beta$ both in $X$, and the photon mass term on $X$ (Neumann condition). For $U\subset X$ the state \emph{with the matter killed on $U$} replaces the single-site measure $d^2\phi(x)$ times its weight by the point mass at $\phi(x)=0$ for $x\in U$; every hopping term with an endpoint in $U$ then reduces to its diagonal part at the other endpoint, and the photon variables are unchanged. The \emph{comparison cube} is $\NDa_{k,J}\defeq\langle\cdot\rangle^{\ND}_{k,\XJ}$, written $\ava{\cdot}$: the scalar is alive on all of $\XJ$, with Dirichlet planes $x_\mu=-MJ-\eps_k$ and $x_\mu=MJ$.
\end{definition}

\begin{definition}[King's RG cube]\label{def:s02:Kcube}
Let $\alive^K_J\defeq\{x\in\XJ:\ x_\mu<MJ-\eps_k\ \forall\mu\}$ and $\TJ\defeq\XJ\setminus\alive^K_J$, the three top face layers. \emph{King's RG cube} $\NDK_{k,J}$, written $\avK{\cdot}$, is the Neumann--Dirichlet state on $\XJ$ with the matter killed on $\TJ$: the photon variables $A_\alpha(x)$, $x\in\XJ$, their gradient terms for pairs in $\XJ$ and their mass term on $\XJ$ are those of the comparison cube; the scalar is alive on $\alive^K_J$ and $\phi=0$ on every other site of $\Lk$, in particular on $\TJ$. Its Dirichlet planes are $x_\mu=-MJ-\eps_k$ and $x_\mu=MJ-\eps_k$, at separation $2MJ$; the Neumann image planes of the photon, $x_\mu=-MJ-\eps_k/2$ and $x_\mu=MJ-\eps_k/2$, are the same in both conventions.
\end{definition}

The two conventions differ exactly by the matter on $\TJ$; the photon Gaussian is the same. King's RG cube is the box \cite[(3.3)--(3.4)]{K2} with the Dirichlet planes of \cite[(3.9)]{K2} and the photon reflection planes of \cite[(3.7)]{K2}, the geometry of his renormalisation group; his Theorem 3.2 and his modified box \cite[(3.38)]{K2} are written in the comparison convention. The two families are interlaced in Section~4; the existence theorems of this paper are posed on $\NDK$, while the rotation part and the cell lemmas of the Wick leg are written for the comparison family. The \emph{periodic state} $\langle\cdot\rangle_{k,T}$ on a King torus $T$ has the action summed over the finite graph $\Lk/(\ell\Z)^3$, in the free convention.

\paragraph{BFS measures.} A \emph{BFS measure} on a finite set $V$ of sites is a probability measure
\[
Z^{-1}\exp\Big(\sum_{\{x,y\}\in B}c_{xy}\,r_xr_y\cos\big(\vartheta_x-\vartheta_y-\ell_{xy}(A)\big)\Big)\prod_{x\in V}d\rho_x(r_x)\,d\vartheta_x\;d\gamma(A),
\]
with $B$ a set of pairs of sites, $c_{xy}\ge0$, each $\rho_x$ a finite measure on $[0,\infty)$ with $\int e^{\alpha r^2}d\rho_x<\infty$ for all $\alpha\ge0$ (point masses at $0$, that is killed sites, included), $\gamma$ a non-degenerate centred Gaussian measure on $\R^m$ (the photon variables), and linear functionals $\ell_{xy}$. Every state above is a BFS measure: $c_{xy}=\eps$ for the bonds present, the single-site measures contain $e^{-\eps^3[\frac12(m^2+\dm_\eps)r^2+\lambda r^4]}$ and the diagonal parts $\frac\eps2r_x^2$ of all bonds at $x$ (those leading out of the box or into a killed site included); the photon form is positive definite because $\lmu>0$, and the exponential moments are finite because $\lambda>0$. So is every measure obtained from such a state by multiplying some $c_{xy}$ or some photon gradient terms by $t\in[0,1]$, adding $s|\ell(A)|^2$ ($s\ge0$) to the action, multiplying single-site measures by $e^{\eps^3q(x)r^2}$ with real $q$, or killing sites. A BFS measure is \emph{of King type} if $c_{xy}\in[0,\eps]$ and each $\rho_x$ is a point mass at $0$ or has the density $r\,e^{-\eps^3[\frac12(m^2+\dm_\eps)r^2+\lambda r^4]-\eps^3w(x)r^2-\frac\eps2n_xr^2}$ with $w(x)\ge0$, $n_x\in\{0,\dots,6\}$; every state above, and every measure obtained from one by the operations just listed with $q\le0$, is of King type. Every state above is invariant under the \emph{charge conjugation} $(A,\vartheta)\mapsto(-A,-\vartheta)$, which changes the sign of every field strength and fixes $r$.

\subsection{Sources, schemes and the objects}

\begin{definition}[sources, schemes, classes, transport]\label{def:s02:sources}
A \emph{curvature source} is a real $2$-form $f\in C^\infty_c(\R^3;\Lambda^2)$, a \emph{scalar source} a real $h\in C^\infty_c(\R^3)$. For the plaquette $p=(x;\alpha<\beta)$ with centre $c(p)$ let $F(p)\defeq\eps^{-1}(A_\alpha(x)+A_\beta(x+\eps e_\alpha)-A_\alpha(x+\eps e_\beta)-A_\beta(x))$ and
\[
F_{\Lk}(f)\defeq\sum_p\eps^3f_{c(p)}(e_\alpha,e_\beta)F(p),\qquad
\sS_{\Lk}(h)\defeq\sum_x\eps^3h(x),\qquad
\Wick{h}\defeq\sum_x\eps^3h(x)\big(|\phi(x)|^2-2\sW_\eps\big),
\]
the sums over the plaquettes (corners in the box), respectively the sites, of the region. A \emph{scheme} is a real sequence $a=(a_k)$, $a_k$ depending only on $\eps_k$ and the couplings; $W^a(h)\defeq\sum_x\eps^3h(x)(|\phi(x)|^2-2\sW_\eps-a_k)$. For a finite family $\zeta=(z_j,f_j;w_j,h_j)_j$ with complex coefficients put $\Theta(\zeta)\defeq\exp(\sum_jz_jF(f_j)+\sum_jw_j\Wick{h_j})$ and, for a Neumann--Dirichlet state on a box $X$ (killed sites allowed),
\[
Z^a_{k,X}(\zeta)\defeq\Big\langle\exp\Big(\sum_jz_jF(f_j)+\sum_jw_jW^a(h_j)\Big)\Big\rangle^{\ND}_{k,X}
=e^{-\sum_jw_ja_k\sS_{\Lk}(h_j)}\big\langle\Theta(\zeta)\big\rangle^{\ND}_{k,X};
\]
$Z^a_{k,J}\defeq Z^a_{k,\XJ}$ for the comparison cube, $Z^{a,K}_{k,J}$ for King's RG cube, and $N^a_{k,X}(0,q)\defeq e^{-(a_k+2\sW_\eps)\sS(q)}\langle e^{\sum_x\eps^3q(x)r_x^2}\rangle^{\ND}_{k,X}$ for one real scalar source. At $h=0$ the multiplier is $1$ and $a$ is dropped: $Z_{k,J}(z)\defeq\ava{e^{\sum_jz_jF(f_j)}}$, $Z^K_{k,J}(z)\defeq\avK{e^{\sum_jz_jF(f_j)}}$ for $z\in\C^n$. The \emph{cone} is $\cK\defeq\{(f,g):f\text{ a real $2$-form},\ g\ge0\}$ with the observable $\Theta(\mathrm if,g)$ (\emph{class $\cI$}); \emph{class $\cD$} consists of the observables $\Theta(f,-g)$, $(f,g)\in\cK$. Transport: $(O\cdot h)(x)\defeq h(O^{-1}x)$, $(O\cdot f)_x(u,v)\defeq f_{O^{-1}x}(O^{-1}u,O^{-1}v)$ for $O\in\Oth$; $\tau_bh\defeq h(\cdot-b)$, $(\tau_bf)_x\defeq f_{x-b}$. For $\nu\in\{1,2,3\}$, $c\in\R$ let $\rho=\rho_{\nu,c}$, $x\mapsto x-2(x_\nu-c)e_\nu$; $\theta h\defeq h\circ\rho$, $\theta f\defeq\rho^*f$, $(\rho^*f)_x(u,v)\defeq f_{\rho x}(\rho u,\rho v)$ (the components $f_{\mu\nu}$, $\mu\ne\nu$, change sign); $\theta\zeta^*\defeq(\bar z_j,\rho^*f_j;\bar w_j,h_j\circ\rho)_j$, and $\zeta\sqcup\zeta'$ is the union of two families. The \emph{amplitude} is $\amp(\zeta)\defeq\sum_j|\Re w_j|\,\norm{h_j}_\infty$. The \emph{reference scheme} is $z=a^\sharp$, $a^\sharp_k\defeq e^2X_{\eps_k}$, $X_\eps$ the coefficient of $e^2$ in the perturbative expansion of $\langle|\phi(0)|^2\rangle$ on $\eps\Z^3$ at $\lambda=0$ with the order-$e^2$ part of $\dm_\eps$ inserted.
\end{definition}

By charge conjugation $\langle\Theta(\mathrm if,g)\rangle=\langle\cos F(f)\,e^{\Wick{g}}\rangle$ is real for every state above, and $Z^a_{k,X}(-z,w)=Z^a_{k,X}(z,w)$.

\begin{remark}[schemes]\label{rem:s02:scheme}
The Wick leg is stated for the reference scheme $z$, or for an admissible scheme $a$ (\cref{def:s03:adm}: $z_k-a_k$ converges, hence is bounded). If the fixed-cube limits exist in $(0,\infty)$ for both $a$ and $z$ at a real scalar source, the scheme is admissible: for $0\le h\in C^\infty_c$ with $\int h>0$ supported well inside some $\CJ$, $Z^z_{k,J}(0,h)/Z^a_{k,J}(0,h)=e^{-(z_k-a_k)\sS_{\Lk}(h)}$ converges to a number in $(0,\infty)$, and $\sS_{\Lk}(h)\to\int h>0$. What the Wick leg uses is only $\sup_k|z_k-a_k|<\infty$.
\end{remark}

\begin{definition}[norms and the tube]\label{def:s02:tube}
For a family with sources $f_1,\dots,f_n$, $h_1,\dots,h_p$ supported in $Q_R\defeq[-R,R]^3$ and $\zeta=(z,w)\in\C^{n+p}$ put $g(\zeta)\defeq\sum_i(\Re z_i)f_i$ and $k(\zeta)\defeq\sum_j(\Re w_j)h_j=k_+-k_-$ (pointwise positive and negative parts). Lattice norms: $\norm{g}_{2,k}^2\defeq\sum_p\eps_k^3g_{c(p)}(e_\alpha,e_\beta)^2$ over the plaquettes of $\Lk$, $\norm{q}_{1,k}\defeq\sum_{x\in\Lk}\eps_k^3|q(x)|$. Continuum norms: $\norm{g}_2^2\defeq\sum_{\alpha<\beta}\norm{g_{\alpha\beta}}_{L^2}^2$; $\wnorm{q}\defeq\sup_x(1+|x|^2)^2|q(x)|$; $\fournorm{q}\defeq\sup_{x,|\gamma|\le4}(1+|x|^2)^2|D^\gamma q(x)|$ and $\fournorm{g}\defeq(\sum_{\alpha<\beta}\fournorm{g_{\alpha\beta}}^2)^{1/2}$, as in \cite{SymPhase}. For $\sigma,H>0$ the \emph{tube} is
\[
\cD_{\sigma,H}\defeq\big\{\zeta\in\C^{n+p}:\ \norm{k_+(\zeta)}_\infty<\sigma/2,\ \ \norm{k_-(\zeta)}_\infty<H/2\big\}.
\]
Since $s\mapsto\norm{(\sum_js_jh_j)_\pm}_\infty=\sup_x(\pm\sum_js_jh_j(x))_+$ is a supremum of convex functions of $s=\Re w$, hence convex and continuous, $\cD_{\sigma,H}$ is the preimage under $\zeta\mapsto\Re w$ of the intersection of two open convex sublevel sets: open, convex, hence connected; it contains $0$ and every $\zeta$ with purely imaginary $w$; $\norm{k(\zeta)}_\infty<\min(\sigma,H)/2$ implies $\zeta\in\cD_{\sigma,H}$; its real points form an open subset of $\R^{n+p}$. The tube of the Wick leg is $\Dtube$, with $\sigma_+$ and $H$ fixed there.
\end{definition}

\paragraph{The objects.} For real $2$-forms $f_1,\dots,f_n$ let $J_0(f)$ be the least $J$ such that every support lies at $\ell^\infty$-distance $>1$ from $\R^3\setminus\CJ$. King's \emph{cube limits} and \emph{printed state} are
\[
\ZKJ(z)\defeq\lim_{k\to\infty}Z^K_{k,J}(z)\quad(J\ge J_0(f)),\qquad
\ZKinf\Big(\sum_jz_jf_j,0\Big)\defeq\lim_{J\to\infty}\ZKJ(z),
\]
whenever the limits exist: their existence at real $z$ is \cref{hyp:s03:EC}, and their existence at every $z\in\C^n$, entire in $z$, is derived from it in the curvature-sector section. The object is scheme-free. For general sources in a scheme $a$ the \emph{tube object} is $G_J(\zeta)\defeq\lim_kZ^a_{k,J}(\zeta)$, $\Ginf(\zeta)\defeq\lim_JG_J(\zeta)$, $\zeta\in\Dtube$, whenever these limits exist (the Wick-leg sections); where the full object $Z^a_\infty\defeq\lim_J\lim_kZ^a_{k,J}$ of the optional extension (\cref{thm:s23:FormA}) exists it agrees with $\Ginf$ on the tube, being the same iterated limit.

\subsection{Elementary facts}

\begin{lemma}[elementary facts]\label{lem:s02:E}
\emph{(E1)} For $\eps\le1$ and any translate $\Gamma$ of $\eps\Z^3$: $\sum_{x\in\Gamma}\eps^3(1+|x|^2)^{-4}\le C_0\defeq(5/2)^4\pi^2/8<49$ and $\sum_{x\in\Gamma}\eps^3(1+|x|^2)^{-2}\le C_0'\defeq(5/2)^2\pi^2<62$.
\emph{(E2)} $\norm{q}_{L^1}\le\pi^2\wnorm{q}$, $\norm{q}_{L^2}\le\frac{\pi}{2\sqrt2}\fournorm{q}$, $\norm{q}_\infty\le\wnorm{q}\le\fournorm{q}$, $\norm{g}_2\le\frac{\pi}{2\sqrt2}\fournorm{g}$; $\wnorm{\cdot}$ is a norm, so $\wnorm{k(\zeta)}\le\sum_j|\Re w_j|\wnorm{h_j}$ and $\wnorm{k_\pm}\le\wnorm{k}$.
\emph{(E3)} Sampling is evaluation at points: it commutes with pointwise positive and negative parts, and a sampled sup is at most the continuum sup, so $\zeta\in\cD_{\sigma,H}$ for the continuum sources implies the same bounds for the sampled ones. For continuous compactly supported $g,q$: $\norm{g}_{2,k}\to\norm{g}_2$, $\norm{q}_{1,k}\to\norm{q}_{L^1}$, and for every $k$: $\norm{g}_{2,k}^2\le C_0\fournorm{g}^2$, $\norm{q}_{1,k}\le C_0'\wnorm{q}$.
\emph{(E4)} $n^{n/2}\le e^{n/2}(n!)^{1/2}$.
\end{lemma}

\begin{proof}
(E1) For $y$ in the cube of side $\eps$ centred at $x\in\Gamma$, $|x-y|^2\le\frac34$ and $|y|^2\le2|x|^2+2|x-y|^2\le2|x|^2+\frac32$, so $1+|y|^2\le\frac52(1+|x|^2)$ and $(1+|x|^2)^{-r}\le(5/2)^r(1+|y|^2)^{-r}$. Integrate over the cube and sum over the disjoint cubes: $\sum_{x\in\Gamma}\eps^3(1+|x|^2)^{-r}\le(5/2)^r\int_{\R^3}(1+|y|^2)^{-r}dy$, with $\int(1+|y|^2)^{-4}=4\pi\cdot\frac\pi{32}=\frac{\pi^2}8$ and $\int(1+|y|^2)^{-2}=4\pi\cdot\frac\pi4=\pi^2$. (E2) $\norm{q}_{L^1}\le\wnorm{q}\int(1+|y|^2)^{-2}$, $\norm{q}_{L^2}^2\le\fournorm{q}^2\int(1+|y|^2)^{-4}$; the rest is immediate. (E3) Riemann sums for the convergences; the plaquette centres of one orientation form a translate of $\eps_k\Z^3$, so (E1) bounds $\sum_{p:(\alpha,\beta)}\eps_k^3g_{c(p)}(e_\alpha,e_\beta)^2$ by $C_0\fournorm{g_{\alpha\beta}}^2$, and the three orientations sum to $C_0\fournorm{g}^2$; the last bound is (E1) with $r=2$. (E4) is $n^n\le e^nn!$.
\end{proof}

\paragraph{Lattice-to-continuum.} On a box $X$ let $\ip AB_k\defeq\sum_{x\in X}\sum_\alpha\eps^3A_\alpha(x)B_\alpha(x)$, and for a real $2$-form $g$ whose non-vanishing plaquettes have their corners in $X$ define $\dstar_kg$ by $F_{\Lk}(g)=\ip{\dstar_kg}{A}_k$. Each component $(\dstar_kg)_\gamma$ is a backward difference quotient of the sampled components $g_{\gamma\beta}$, $\beta\ne\gamma$; for smooth compactly supported $g$ it converges uniformly to $(\dstar g)_\gamma$, up to the orientation sign convention of the components, at the rate $O(\eps_k)$ of Taylor's theorem. Hence
\begin{equation}\label{eq:s02:riemann}
\norm{\dstar_kg}_{2,k}\to\norm{\dstar g}_{L^2},\qquad \sup_k\norm{\dstar_kg}_{2,k}<\infty,\qquad \norm{g}_{2,k}\to\norm{g}_2\qquad(k\to\infty).
\end{equation}

\subsection{Correlation inequalities}

For a BFS measure, a real linear functional $F$ of $A$ and $g\ge0$ on $V$ put $\Theta_{F,g}\defeq e^{\mathrm iF}e^{\sum_xg_xr_x^2}$, $\Theta^-_{F,g}\defeq e^{F}e^{-\sum_xg_xr_x^2}$, and $\langle X;Y\rangle\defeq\langle XY\rangle-\langle X\rangle\langle Y\rangle$. Since $F(f)$ is a real linear functional of $A$ and $e^{\Wick{g}}$ a positive constant times $e^{\sum_x\eps^3g(x)r_x^2}$, class $\cI$ consists of multiples of observables $\Theta_{F,g}$ and class $\cD$ of multiples of observables $\Theta^-_{F,g}$.

\begin{lemma}[correlation inequalities]\label{lem:s02:BFS}
Let $\langle\cdot\rangle$ be a BFS measure, $F$ a real linear functional of $A$, $g\ge0$, $\Theta=\Theta_{F,g}$, $\Theta^-=\Theta^-_{F,g}$. Then $\langle\Theta\rangle$ is real, and
\begin{enumerate}[label=\textup{(\alph*)},nosep]
\item $\langle\Theta;|\ell(A)|^2\rangle\le0$ for every complex linear functional $\ell$ of $A$; hence adding $s|\ell(A)|^2$ to the action, $s\ge0$, does not decrease $\langle\Theta\rangle$, and deleting a photon gradient term does not increase it;
\item $\langle\Theta;r_xr_y\cos(\vartheta_x-\vartheta_y-\ell_{xy}(A))\rangle\ge0$; hence decreasing $c_{xy}$ does not increase $\langle\Theta\rangle$;
\item $\langle\Theta;r_x^2\rangle\ge0$; hence multiplying $\rho_x$ by $e^{-sr^2}$, $s\ge0$, does not increase $\langle\Theta\rangle$, and killing a set $U$ of sites does not increase it, provided $g=0$ on $U$ or the measure is of King type;
\item for $\Theta^-$ the three inequalities and the three monotonicity statements hold with the opposite signs;
\item if the pairs in $B$ and the Gaussian $\gamma$ do not couple the fields at the sites of $X\subset V$ (and the photon variables attached to them) to the others, the measure is a product, and for an observable depending only on the fields in $X$ the expectation equals the expectation in the $X$-factor.
\end{enumerate}
\end{lemma}

\begin{proof}
(a)--(c) are \cite[Thms.~6.1, 6.2(1),(2)]{BFS1}, whose hypotheses are those of a BFS measure: an arbitrary finite lattice, $c_{xy}\ge0$, finite single-site measures on $[0,\infty)$ with $\int e^{\alpha r^2}d\rho_x<\infty$ for all $\alpha\ge0$ (a point mass at $0$ is one), an arbitrary centred Gaussian with covariance $C\ge0$, and $g\ge0$ (so $g=0$ is included). Theorem 6.1 states $\langle P(r)\cos(n\cdot\vartheta+f\cdot A);Q(r)\cos(m\cdot\vartheta+h\cdot A)\rangle\ge0$ for polynomials $P,Q$ in the $r_x$ with positive coefficients; (a), (b) follow as in the printed proof of 6.2(1),(2): $e^{\sum g_xr_x^2}$ is a limit of such polynomials, $e^{\mathrm iF}$ may be replaced by $\cos F$ because the expectation and both observables are invariant under $(A,\vartheta)\mapsto(-A,-\vartheta)$ (which also gives the realness), and $|\ell(A)|^2=\ell_1(A)^2+\ell_2(A)^2$ with $\ell_r(A)^2=\lim_{t\downarrow0}t^{-2}[1-\cos(t\ell_r(A))]$; (c) is Theorem 6.1 with $Q=r_x^2$ and $m=0$, $h=0$. Each monotonicity statement follows by differentiating in the parameter, every interpolating measure being a BFS measure. Killing: if $g=0$ on $U$, set $c_{xy}=0$ on every bond at $U$ (not increasing, by (b)); the measure is then a product (e) of the matter at $U$, which carries no observable and cancels, and of the state with the matter killed on $U$ (the diagonal parts at the other ends of the bonds into $U$ are already in their single-site measures). For King type in general, multiply $\rho_x$, $x\in U$, by $e^{-s\eps^3r^2}$ (not increasing) and let $s\to\infty$ by Lemma~\ref{lem:s02:Kprime}.

(d) is stated in \cite[Thm.~6.2(3),(4)]{BFS1} for the photon and the hopping inequality with a sketched proof, and its $r^2$-inequality is \cite[(3.28)]{K2}, stated there without proof; we prove all three. Let $\mu$ be the unnormalised measure and $\omega'=(r',\vartheta',A')$ an independent copy of $\omega=(r,\vartheta,A)$; then $(\int d\mu)^2\langle\Theta^-;Y\rangle=\iint\Theta^-(\omega)(Y(\omega)-Y(\omega'))\,d\mu(\omega)d\mu(\omega')$. Write $r=u-v$, $r'=u+v$ (so $u_x\ge|v_x|$), $\vartheta=\vartheta_++\vartheta_-$, $\vartheta'=\vartheta_+-\vartheta_-$, $A=a_++a_-$, $A'=a_+-a_-$. Then $a_\pm=(A\pm A')/2$ are independent with the common centred Gaussian law of covariance $C/2$, and the radial law $\prod_xd\rho_x(u_x-v_x)d\rho_x(u_x+v_x)$ on $\{u\ge|v|\}$ is invariant under $v_x\mapsto-v_x$ at each site. With $\chi_{xy}\defeq\vartheta_{+,x}-\vartheta_{+,y}-\ell_{xy}(a_+)$ and $\psi_{xy}\defeq\vartheta_{-,x}-\vartheta_{-,y}-\ell_{xy}(a_-)$ the unprimed phase is $\chi+\psi$, the primed one $\chi-\psi$, and
\begin{align*}
r_xr_y\cos(\chi+\psi)+r'_xr'_y\cos(\chi-\psi)&=2(u_xu_y+v_xv_y)\cos\chi\cos\psi+2(u_xv_y+v_xu_y)\sin\chi\sin\psi,\\
r_xr_y\cos(\chi+\psi)-r'_xr'_y\cos(\chi-\psi)&=-2(u_xv_y+v_xu_y)\cos\chi\cos\psi-2(u_xu_y+v_xv_y)\sin\chi\sin\psi,
\end{align*}
$r_x^2-r_x'^2=-4u_xv_x$, $\ell(A)^2-\ell(A')^2=4\ell(a_+)\ell(a_-)$ for real $\ell$, and $\Theta^-(\omega)=e^{F(a_+)}e^{F(a_-)}e^{-\sum_xg_x(u_x^2+v_x^2)}e^{2\sum_xg_xu_xv_x}$. Expand $\exp(\sum_bc_b[\text{first line}])$ and $e^{2\sum g_xu_xv_x}$ in power series (absolutely convergent at fixed spacing). Every term of the integrand is a sign $\sigma$ times a monomial in $(u,v)$ with non-negative coefficient, times the even radial factor $e^{-\sum g_x(u_x^2+v_x^2)}$, times $V(\vartheta_+,a_+)V(\vartheta_-,a_-)$ with one real function $V$ (a product of $e^{F(\cdot)}$, of factors $\cos\chi_b$ or $\sin\chi_b$, and of $\ell(\cdot)$ for the photon slot). The photon integral of $V(\vartheta_+,a_+)V(\vartheta_-,a_-)$ factorises because $a_+,a_-$ are independent and identically distributed. The angular integral factorises by the following argument: per site $d\vartheta\,d\vartheta'=2\,d\vartheta_+d\vartheta_-$, and $[0,2\pi)^2$ in the $(\vartheta,\vartheta')$-variables is a fundamental domain of the lattice $L$ generated by $(\pi,\pi)$ and $(\pi,-\pi)$ in the $(\vartheta_+,\vartheta_-)$-plane, which contains $(2\pi\Z)^2$ with index $2$; $V(\vartheta_+,\cdot)V(\vartheta_-,\cdot)$ is $L$-invariant, because the simultaneous shift of $\vartheta_{+,x}$ and $\vartheta_{-,x}$ by $\pi$ at one site changes every $\chi_b$ and $\psi_b$ with $b\ni x$ by $\pi$ and leaves $\cos\chi_b\cos\psi_b$ and $\sin\chi_b\sin\psi_b$ unchanged; hence $\int_{[0,2\pi)^2}VV\,d\vartheta\,d\vartheta'=2\int_{L\text{-domain}}VV\,d\vartheta_+d\vartheta_-=\int_{[0,2\pi)^2}VV\,d\vartheta_+d\vartheta_-=(\int_0^{2\pi}V)^2$. The radial integral of a monomial vanishes if its degree in some $v_x$ is odd, by the $v\mapsto-v$ symmetry, and is non-negative otherwise, since $u_x\ge|v_x|$. So every term has the sign $\sigma$: $\sigma=-1$ for $Y=r_x^2$ (from $-4u_xv_x$) and for the hopping term (second line), $\sigma=+1$ for $Y=\ell(A)^2$ (for complex $\ell$, real and imaginary parts). Hence $\langle\Theta^-;r_x^2\rangle\le0$, $\langle\Theta^-;r_xr_y\cos(\cdot)\rangle\le0$, $\langle\Theta^-;|\ell(A)|^2\rangle\ge0$. Monotonicity and killing follow as in (a)--(c), with (d)'s hopping inequality in place of (b) and Lemma~\ref{lem:s02:Kprime}, which covers $\Theta^-$, for King type. (e) is clear.
\end{proof}

\begin{lemma}[current representation and Gaussian domination]\label{lem:s02:GD}
Let $\langle\cdot\rangle$ be a BFS measure whose Gaussian has covariance $C$. Integrating out the matter at fixed $A$, the matter integral is $\Psi(A)=\sum_qw(q)\cos\ip{\kappa(q)}A$, $q$ running over the integer-valued divergence-free currents on $B$, $w(q)\ge0$, $\sum_qw(q)<\infty$, $\ip{\kappa(q)}A\defeq\sum_{xy}q_{xy}\ell_{xy}(A)$. Let $\E^\sharp$ be the probability measure on currents with weights proportional to $w(q)e^{-\ip{\kappa(q)}{C\kappa(q)}/2}$. Then for $v\in\R^m$, $s\in\R$,
\[
\langle e^{s\ip vA}\rangle=e^{s^2\ip v{Cv}/2}\,\E^\sharp\cos\big(s\ip{Cv}{\kappa}\big)\le e^{s^2\ip v{Cv}/2}.
\]
\end{lemma}

\begin{proof}
Each bond factor expands as $e^{c_{xy}r_xr_y\cos t}=\sum_{q\in\Z}I_q(c_{xy}r_xr_y)e^{\mathrm iqt}$ with $I_q\ge0$, $I_{-q}=I_q$, $\sum_qI_q(c)=e^c$; after the radial integrals the multiple series converges absolutely, dominated by $\prod_xe^{\frac12\sum_yc_{xy}r_x^2}$, integrable against the single-site measures. Integrating the angles forces $\sum_yq_{xy}=0$ at every site (at a killed site $r_x=0$ and $I_q(0)=\delta_{q0}$, so the currents vanish on its bonds); pairing $q$ with $-q$ gives $\Psi$ with $w(q)\ge0$. For a centred Gaussian $\gamma$ with covariance $C$ and $a\in\C^m$, $\int e^{\ip aA}d\gamma=e^{\ip a{Ca}/2}$; with $a=sv\pm\mathrm i\kappa$ and the average of the two signs, $\int e^{s\ip vA}\cos\ip\kappa A\,d\gamma=e^{s^2\ip v{Cv}/2-\ip\kappa{C\kappa}/2}\cos(s\ip v{C\kappa})$. Sum over $q$ with the weights $w(q)$, divide by the same sum at $s=0$ (positive), and use $|\cos|\le1$.
\end{proof}

The lemma applies to every state of this paper at every $k$, $J$, $\lmu>0$, with or without a real scalar source absorbed in the single-site measures (the photon covariance is unchanged).

\begin{lemma}[killing]\label{lem:s02:Kprime}
Let $\langle\cdot\rangle$ be a BFS measure of King type, $U$ a set of alive sites, $\chi:U\to(0,\infty)$, and $\langle\cdot\rangle_{U,s}$ the measure with $\rho_x$ multiplied by $e^{-s\eps^3\chi(x)r_x^2}$ for $x\in U$. Let $\Theta$ be continuous with $|\Theta(A,\phi)|\le e^{\ell(A)}e^{C\sum_xr_x^2}$ for a real linear functional $\ell$ and $C\ge0$. Then $\langle\Theta\rangle_{U,s}\to\langle\Theta\rangle_U$ as $s\to\infty$, $\langle\cdot\rangle_U$ the measure with $\rho_x$ replaced by the point mass at $0$ for $x\in U$ (the state with the matter killed on $U$).
\end{lemma}

\begin{proof}
The unnormalised density is $e^{-Q(A)/2}dA\prod_{x\text{ alive}}w_x(r_x)r_xdr_xd\vartheta_x\prod_{b}e^{c_br_xr_y\cos(\cdot)}\prod_{x\in U}e^{-s\eps^3\chi(x)r_x^2}$, $Q$ positive definite, killed sites carrying no variables. Substitute $r_x=(s\eps^3\chi(x))^{-1/2}u_x$ for $x\in U$; the Jacobians cancel between numerator and denominator. \emph{Domination for $s\ge s_0$:} at $x\in U$, $w_x(r_x)\le e^{\eps^3|\dm_\eps|r_x^2/2}$ (the other exponents in $w_x$ are $\le0$), $|\Theta|$ contributes $e^{Cr_x^2}$, and each of the at most six bond weights at $x$ is at most $e^{\frac\eps2r_x^2+\frac\eps2r_y^2}$ ($0\le c_b\le\eps$); in the variable $u_x$ these factors together are $\le e^{u_x^2/2}$ once $s\ge s_0(\eps,m,\dm_\eps,C,\min\chi)$, so the $U$-factor is at most $e^{-u_x^2/2}u_x$, integrable on $[0,\infty)\times[0,2\pi)$; at an alive $y\notin U$ the factor is at most $w_y(r_y)e^{(C+3\eps)r_y^2}r_y$, integrable since $\lambda>0$; the photon factor $e^{\ell(A)}e^{-Q(A)/2}$ is integrable since $Q$ is positive definite. \emph{Pointwise limit:} at fixed $(u,\vartheta,A,r_{V\setminus U})$, $r_x\to0$ for $x\in U$, so $\Theta\to\Theta|_{\phi_U=0}$ by continuity, every bond weight with an endpoint in $U$ and every $w_x(r_x)$, $x\in U$, tends to $1$, and the $U$-factor tends to $e^{-u_x^2}u_x\,du_xd\vartheta_x$, whose integral cancels between numerator and denominator. By dominated convergence the ratio tends to $\int\Theta|_{\phi_U=0}d\mu_U/\int d\mu_U$, where $\mu_U$ has the cosines at $U$ removed, the weights $w_y$ (which already contain all diagonals at $y$) and the photon unchanged, and the point mass at $U$: that is $\langle\cdot\rangle_U$, with a positive denominator.
\end{proof}

\begin{lemma}[cutting]\label{lem:s02:cut}
Let $\langle\cdot\rangle_Y$ be a Neumann--Dirichlet box state on $\Lk$, possibly with killed sites, or a periodic state on a King torus, and let $X$ be a set of sites of $Y$ such that $X$ and its nearest neighbours lie in one fundamental cell (a clause in torus language: for a box state it has no content beyond $X\subset Y$). Let $\langle\cdot\rangle^{\ND}_X$ be the Neumann--Dirichlet state on $X$ whose single-site measures (killed sites and the diagonal parts of the bonds leaving $X$ included) are those of $\langle\cdot\rangle_Y$. Let $(f,g)\in\cK$ be such that every plaquette and every site at which $f$, respectively $g$, does not vanish lies in $X$. Then, at every lattice spacing,
\[
\langle\Theta(\mathrm if,g)\rangle^{\ND}_X\le\langle\Theta(\mathrm if,g)\rangle_Y,\qquad
\langle\Theta(f,-g)\rangle^{\ND}_X\ge\langle\Theta(f,-g)\rangle_Y,
\]
and the same for $Z^a$. The first inequality uses Lemma~\ref{lem:s02:BFS}(a),(b),(e) only; the second uses (d),(e). If no site of $X$ is killed, $\langle\cdot\rangle^{\ND}_X$ is the state of Definition~\ref{def:s02:acube}.
\end{lemma}

\begin{proof}
Delete from $Y$ every photon gradient term coupling a site of $X$ to a site outside $X$, and the cosine parts of all hopping terms between $X$ and its complement, keeping their diagonal parts $\frac\eps2(r_x^2+r_y^2)$. By Lemma~\ref{lem:s02:BFS}(a),(b) this does not increase the first expectation, by (d) it does not decrease the second, and every intermediate measure is a BFS measure. The result is a product (Lemma~\ref{lem:s02:BFS}(e)) whose $X$-factor is $\langle\cdot\rangle^{\ND}_X$: the diagonal parts of the cut bonds are exactly the Dirichlet terms and the remaining photon terms the Neumann ones of $X$; the observable depends on the $X$-variables only. The free Wick constant, $\dm_\eps$ and the scheme multiplier are the same on both sides.
\end{proof}

With $Y=X_{J+1}$, $X=\XJ$ this is King's monotonicity in the volume \cite[Thm.~3.2, (3.31)--(3.33)]{K2} between comparison cubes, in the form his proof proves: with one Wick constant for all $J$, the free subtraction King himself switches to at the start of that proof \cite[p.~333]{K2}. With the Dirichlet Wick constant of \cite[(2.13)]{K2}, which increases with $J$, the $c$-number factor moves against the monotone factor and the monotonicity does not follow; that form is not used. The box version with killed sites and the monotonicity of King's RG cubes are in Section~4.

\subsection{Algebra of the scheme, covariance, reflection, continuity}

\begin{lemma}[algebra of the scheme]\label{lem:s02:alg}
Let $a$ be a scheme and $Z^a$ the functional of a state above on a box at spacing $\eps_k$.
\begin{enumerate}[label=\textup{(\alph*)},nosep]
\item For real coefficients $w$ the multiplier $e^{-\sum_jw_ja_k\sS(h_j)}$ is a positive constant depending only on $k$, $w$ and $\sS_{\Lk}(h_j)$. Hence the inequalities of Lemmas~\ref{lem:s02:BFS} and~\ref{lem:s02:cut} between two states on the same lattice with the same sources hold for $Z^a$ in the same direction.
\item[(d)] $|Z^a(\zeta)|\le Z^a(\Re\zeta)$, and $\ln Z^a$ is convex in the real coefficients.
\item[(e)] $Z^a(f,h)\,Z^a(-f,-h)\ge1$ for real $(f,h)$.
\item[(f)] For real $f$ and real $k_1,k_2\ge0$: $Z^a(f,k_1-k_2)\le Z^a(2f,-2k_2)^{1/2}\,Z^a(0,2k_1)^{1/2}$.
\end{enumerate}
\end{lemma}

\begin{proof}
(a) is immediate. (d) $|\Theta(\zeta)|=\Theta(\Re\zeta)$ and the multiplier has modulus $e^{-\sum_j(\Re w_j)a_k\sS(h_j)}$; convexity is H\"older's inequality. (e) $\langle e^{X}\rangle\langle e^{-X}\rangle\ge1$ by Cauchy--Schwarz, and the multipliers cancel. (f) Cauchy--Schwarz gives $\langle e^{F+\Wick{k_1}-\Wick{k_2}}\rangle\le\langle e^{2F-2\Wick{k_2}}\rangle^{1/2}\langle e^{2\Wick{k_1}}\rangle^{1/2}$, and $e^{-a_k\sS(k_1-k_2)}=(e^{2a_k\sS(k_2)})^{1/2}(e^{-2a_k\sS(k_1)})^{1/2}$.
\end{proof}

\begin{lemma}[passive covariance]\label{lem:s02:pass}
Let $O\in\Oth$, $\Lambda$ a lattice of spacing $\eps$ with frame $(f_\alpha)$, and $X$ a box or a torus for $\Lambda$. Equip $O\Lambda$ with the frame $(Of_\alpha)$ and $OX$ with the transported boundary conditions. Then $\langle\Theta(z\,O\cdot f,w\,O\cdot h)\rangle_{O\Lambda,OX}=\langle\Theta(zf,wh)\rangle_{\Lambda,X}$ and $\sS_{O\Lambda}(O\cdot h)=\sS_\Lambda(h)$, likewise for families and for $Z^a$. On a torus the periodic model does not change if the frame $(f_\alpha)$ is replaced by $(\pm f_{\pi(\alpha)})$, $\pi$ a permutation.
\end{lemma}

\begin{proof}
The map $\phi'(Ox)=\phi(x)$, $A'_\alpha(Ox)=A_\alpha(x)$, components taken in the frames $(f_\alpha)$ and $(Of_\alpha)$, is a bijection of the configuration spaces preserving Lebesgue measure; the action is written in the frame, and $\dm_\eps$, $\sW_\eps$, $a_k$ depend on $\eps$ only; plaquettes, centres and field strengths are transported, and $(O\cdot f)_{Oc}(Of_\alpha,Of_\beta)=f_c(f_\alpha,f_\beta)$. A permutation of the frame relabels the components; reversing $f_\alpha$ replaces $A_\alpha(x)$ by $-A_\alpha(x-\eps f_\alpha)$, a bijection of the torus variables preserving the action, which changes the sign of both $F(p)$ and $f_{c(p)}(f_\alpha,f_\beta)$ and preserves the set of plaquette centres.
\end{proof}

\begin{lemma}[the reflection sandwich at fixed spacing, comparison convention]\label{lem:s02:refl}
Let $\sigma_3=\diag(1,1,-1)$ and let $(f,g)\in\cK$ be supported at $\ell^\infty$-distance $>1$ from $\R^3\setminus C_{J-1}$. Then for every $k\ge0$
\[
Z^a_{k,J-1}(\mathrm if,g)\le Z^a_{k,J}(\mathrm i\,\sigma_3\!\cdot\!f,\sigma_3\!\cdot\!g)\le Z^a_{k,J+1}(\mathrm if,g),
\]
and the reversed inequalities hold for $Z^a(f,-g)$.
\end{lemma}

\begin{proof}
$\sigma_3$ preserves $\Lk$ and $\sS_{\Lk}$ and maps $\XJ$ onto $X_+\defeq(\CJ+\eps_ke_3)\cap\Lk$. By Lemma~\ref{lem:s02:pass}, $Z^a_{k,J}(\mathrm i\sigma_3\cdot f,\sigma_3\cdot g)$ is the functional with sources $(f,g)$ of the Neumann--Dirichlet state on $X_+$ in the frame $(e_1,e_2,-e_3)$. Written in the frame $(e_\alpha)$, with $A'_3(x)=-A_3(x-\eps e_3)$, this is the BFS measure $R$ in which $\phi,A_1,A_2$ live on $X_+$ and $A_3$ on $X_+-\eps e_3=\XJ$, with Dirichlet conditions for $\phi$ on $X_+$, photon gradient and mass terms within these sets, and the hopping term in direction $3$ the standard one of the bond $\langle y-\eps e_3,y\rangle$; the plaquette terms $(x;\alpha,3)$ pick up two sign changes and keep their centres, so the observable is $\Theta(\mathrm if,g)$ in the standard frame. $R$ is obtained from the Neumann--Dirichlet state on $X_{J+1}$ by deleting photon gradient terms and hopping cosines and factorising, since $X_+\cup\XJ\subset X_{J+1}$, and the state on $X_{J-1}$ is obtained from $R$ in the same way, since $X_{J-1}\subset\XJ\cap X_+$; the diagonal parts of the cut bonds are the Dirichlet terms of the smaller state, and the sources stay inside $X_{J-1}$. Lemma~\ref{lem:s02:BFS}(a),(b),(e), respectively (d),(e), give the directions, and Lemma~\ref{lem:s02:alg}(a) carries them to $Z^a$.
\end{proof}

\begin{lemma}[continuity of convex functions]\label{lem:s02:baire}
Let $U$ be a non-empty open convex subset of a Fr\'echet space and $\Phi:U\to\R$ a convex Borel function. Then $\Phi$ is continuous on $U$.
\end{lemma}

\begin{proof}
$U$ is a Baire space and $U=\bigcup_{j\in\N}\{\Phi\le j\}$, so some $S\defeq\{\Phi\le j\}$ is non-meagre in $U$; being Borel, $S$ has the Baire property \cite{Kec95}: $O\setminus\mathcal M\subset S$ for a non-empty open $O\subset U$ and a meagre $\mathcal M$. Pick $x_0\in O$ and a balanced open neighbourhood $W$ of $0$ with $x_0+W+W\subset O$. For $w\in W$ and $y\defeq x_0+w$, $y+W\subset O$, and the point reflection $\iota(x)\defeq2y-x$ maps $y+W$ onto itself and preserves meagre sets; so $S\cap(y+W)$ and $\iota(S\cap(y+W))$ are comeagre in $y+W$ and intersect: there is $u\in S$ with $2y-u\in S$, and $\Phi(y)\le\frac12(\Phi(u)+\Phi(2y-u))\le j$. Hence $\Phi\le j$ on $x_0+W$. A convex function on an open convex set bounded above near one point is bounded above near every point $y\in U$ (choose $z\in U$ with $y$ strictly between $x_0$ and $z$, and take the convex combinations of $x_0+W$ with $z$), and a convex function locally bounded above is continuous \cite{RS1}.
\end{proof}

% s03 The inputs
\section{The inputs}\label{sec:s03}

This section states, in the form in which the later sections consume them, the results on which the theorems of this paper rest, and says where each is proved. Three are theorems of this paper proved in later parts: the existence of King's cube limits, without and with scalar sources (\cref{hyp:s03:EC,hyp:s03:W1}; the existence part, \cref{sec:k01} and \moving{existence section with scalar sources}), and the rotation theorem (\cref{hyp:s03:RC}; the rotation part, \moving{rotation part}). Two are Ba{\l}aban's torus estimates in the form used by the Wick leg (\cref{hyp:s03:T1,hyp:s03:T2}). One is a condition on the couplings, the sign clause \Tthree\ (\cref{hyp:s03:T3}), proved on explicit coupling regions in \cref{prop:s14:T3}; it is the only condition on the couplings anywhere in the paper. The hypotheses of the optional extension of the Wick leg to arbitrary amplitudes (\Wtwo, \RW, \FAK) are stated at the end of the section; no clause of the main theorem uses them. From the literature, the paper uses the correlation inequalities of \cite{BFS1}, the chessboard estimate of \cite{FILS78}, the reconstruction theorems of \cite{OS73,OS75}, Ba{\l}aban's torus analysis \cite{Ba1,Ba2,Ba3,Ba4} and King's papers \cite{K1,K2}, each at the place where it is applied; the general lemmas on the measure, its ergodic decomposition and weak clustering are taken from the authors' earlier paper on the torus limits \cite{SymPhase} (\cref{rem:s03:companion-inputs}).

\subsection{Objects and conventions}

The two cube conventions of Definitions~\ref{def:s02:Kcube} and~\ref{def:s02:acube} are King's convention $\NDK$ (the scalar field killed on the three top face layers $\TJ$ of the cube $\CJ$, the convention of the state King prints) and the comparison convention $\NDa$ (the scalar alive on all of $\XJ$), in which the tube results (Theorem~\ref{thm:s11:W}, Lemma~\ref{lem:s11:TWtube}) and the lattice comparisons of the rotation part are written. The two are interlaced by Lemma~\ref{lem:s04:G}. Which convention a statement is posed on is part of it:
\begin{itemize}[leftmargin=2em]
\item The existence theorems \cref{hyp:s03:EC,hyp:s03:W1} are posed on King's convention $\NDK$, on which King's renormalisation-group flow runs (\cref{rem:k01:cube}). By Corollary~\ref{cor:s04:G2} the existence of the fixed-cube limits on one convention does not imply it on the other; the comparison family's fixed-cube limits are reached through that corollary, and the statements proved for the comparison family under the assumption that they exist (\cref{def:s03:conv}) serve as templates for King's convention.
\item The rotation theorem \cref{hyp:s03:RC} is proved for the comparison family $\NDa$ and reaches King's printed state $\ZKinf$ through Corollary~\ref{cor:s04:G2}.
\item \Cref{hyp:s03:T1,hyp:s03:T2} and \Tthree\ are statements about King's tori (Definition~\ref{def:s10:torus}) and about $c$-numbers.
\end{itemize}
Throughout, $Z^{a}_{k,J}(z,w)$ denotes the generating functional of Definition~\ref{def:s02:sources} with curvature coefficients $z$ and Wick coefficients $w$ in the scheme $a$, on the convention in question (written $Z^{a,K}_{k,J}$ when the convention $\NDK$ is to be stressed); $z=a^\sharp$ is the reference scheme of Remark~\ref{rem:s02:scheme}; $\avK{\,\cdot\,}$ is King's cube expectation at spacing $\eps_k$.

\begin{definition}[convergence of the comparison family at fixed cubes]\label{def:s03:conv}
The comparison family is \emph{convergent at fixed cubes} if for every finite family $f_1,\dots,f_n$ of real $2$-forms in $C^\infty_c(\R^3;\Lambda^2)$, every $J\ge J_0(f)$ and every $x\in\R^n$ the limit $\lim_k\ava{e^{\sum_jx_jF(f_j)}}$ exists and lies in $(0,\infty)$.
\end{definition}
This is the clause of \cref{hyp:s03:EC} read for the family $\NDa$. It is not a result of this paper and is not needed for the main theorem: the statements below that assume it are proved for the comparison family as templates, and their forms for King's printed state are proved under \cref{hyp:s03:EC}; the continuum-subsequence section (\cref{thm:s12:LIM}) removes the dependence on the fixed-cube limits altogether.

\subsection{Existence of King's cube limits}

\begin{theorem}[King's cube limits, curvature]\label{hyp:s03:EC}
For every finite family $f_1,\dots,f_n$ of real $2$-forms in $C^\infty_c(\R^3;\Lambda^2)$, every $J\ge J_0(f)$, where $J_0(f)$ is the least $J$ such that every $\supp f_j$ lies at $\ell^\infty$-distance greater than $1$ from $\R^3\setminus\CJ$, and every $z\in\R^n$, the limit
\[
\lim_{k\to\infty}\avK{e^{\sum_j z_j\Fcurv{f_j}}}
\]
exists and lies in $(0,\infty)$.
\end{theorem}
Proved in \cref{sec:k01} (\cref{thm:k01:EC}), from King's renormalisation-group analysis of the cube.

\begin{remark}[on \cref{hyp:s03:EC}]\label{rem:s03:EC}
(i) \emph{Existence clause only.} \Cref{hyp:s03:EC} is used only through the existence of the fixed-cube limit at real coefficients: no constants clause, no uniformity in $J$ and no locally uniform convergence in $z$ enter from it. Every bound used on the limit is proved here, with absolute constants, in Theorem~\ref{thm:s05:C}; the monotonicity in the cube, the $J$-limit, the entire extension and the translations are proved in Lemma~\ref{lem:s04:G}, Theorem~\ref{thm:s05:C}, Lemma~\ref{lem:s05:TWC} and Theorem~\ref{thm:s06:O3}.
(ii) \emph{Origin.} This is the existence clause of King's Theorem~2.1 \cite{K2} at zero scalar source, for the state he prints. The proof in \cref{sec:k01} follows King's route, the comparison of two spacings at a common block scale, in a ratio form in which the vacuum energies cancel; what it uses of Ba{\l}aban's block-factor-$5$ analysis \cite{Ba1,Ba2,Ba3,Ba4} is stated there. The mass clause of \cite[Prop.~3.1]{Ba2}, upstream of it, is part of that proposition's statement and is cited under it (Remark~\ref{rem:s14:BALm}); it is not a condition of this paper.
(iii) \emph{Consumers.} The whole curvature sector (Theorems~\ref{thm:s05:C}, \ref{thm:s06:O3}, \ref{thm:s08:OS2C}, \ref{thm:s09:S2K}), the identification of the LIM states with King's state (Corollary~\ref{cor:s12:K}) and the main theorem (Theorem~\ref{thm:s13:main}).
\end{remark}

\begin{remark}[the cube]\label{rem:s03:cube-identification}
The cube on which the existence theorem is proved is King's RG cube $\CK$ of \cref{def:s02:Kcube}, with King's killing of the scalar on the top layers $\TJ$; it is the only cube on which King's block-spin flow exists (\cref{rem:k01:cube}). Every placement of a comparison-cube result on King's cube goes through Lemma~\ref{lem:s04:G} and its corollaries.
\end{remark}

\begin{theorem}[King's cube limits with scalar sources]\label{hyp:s03:W1}
Let $z=a^\sharp$ be the reference scheme and let $s_W\defeq s_1$ be the amplitude of Proposition~\ref{prop:s10:U}. For every finite family of real $2$-forms $f_j\in C^\infty_c(\R^3;\Lambda^2)$ and real functions $h_i\in C^\infty_c(\R^3)$, every $J\ge J_0(f,h)$ and every real coefficient vector $(z,w)$ whose positive scalar part satisfies $\sup\big(\sum_iw_ih_i\big)_+\le s_W$ (the negative scalar part is unrestricted), the limit $\lim_{k\to\infty}Z^{z}_{k,J}(z,w)$ exists and lies in $(0,\infty)$.
\end{theorem}
Proved in \moving{existence section with scalar sources}, by the route of \cref{sec:k01} with the scheme $c$-number carried along.

\begin{definition}[admissible schemes]\label{def:s03:adm}
A scheme $a$ (Remark~\ref{rem:s02:scheme}) is \emph{admissible} if $a=z$ or $\lim_k(z_k-a_k)$ exists. For an admissible scheme \cref{hyp:s03:W1} holds for $Z^a_{k,J}$ in place of $Z^z_{k,J}$ (Lemma~\ref{lem:s02:alg}(a)), and $\sup_k|z_k-a_k|<\infty$.
\end{definition}

\begin{remark}[on \cref{hyp:s03:W1}]\label{rem:s03:W1}
(i) At zero Wick coefficient \cref{hyp:s03:W1} is \cref{hyp:s03:EC}.
(ii) \emph{Origin.} \Cref{hyp:s03:W1} is \emph{not} King's Theorem~2.1 with scalar sources: as printed, in King's Wick ordering, that statement fails, and the state with scalar sources exists only after a vacuum counterterm, which is what the scheme supplies (Remark~\ref{rem:s02:scheme}). The threshold on the positive scalar part is the one at which the positive-source bound of the Wick leg is available (Proposition~\ref{prop:s11:W1prime}); what happens above it is the subject of the optional extension (\cref{hyp:s03:W2}).
(iii) One instance of \cref{hyp:s03:W1} is named separately because the statements on continuum subsequences need it alone: \Bcell\ (Hypothesis~\ref{hyp:s12:Bcell}; one cube, one negative Wick source).
(iv) \emph{Consumers.} The Wick tube (Proposition~\ref{prop:s10:U}, Theorem~\ref{thm:s11:W}), its transfer to King's convention (Theorem~\ref{thm:s13:TWK}), the Wick sector of Theorem~\ref{thm:s12:LIM}, and Form~(B) of the main theorem (Theorem~\ref{thm:s13:main}, Proposition~\ref{prop:s13:WK}).
\end{remark}

\subsection{The torus inputs of the Wick leg}

\begin{proposition}[torus stability at two scalar masses]\label{hyp:s03:T1}
For the scalar masses $M=m$ and $M=M_1$ (the auxiliary mass fixed in Proposition~\ref{prop:s10:U}) there are constants $E_\pm(M)<\infty$, a tail $K_1$ and, for each $k\ge K_1$, a $J_*(k)$ such that the source-free partition functions of King's tori $T$ (Definition~\ref{def:s10:torus}) at spacing $\eps_k$ and at scalar mass $M$ satisfy, in King's torus convention,
\[
e^{-E_-(M)\abs{T}}\le Z^{(M)}_{k,T}\le e^{E_+(M)\abs{T}}\qquad(k\ge K_1,\ J\ge J_*(k)).
\]
\end{proposition}
These are Ba{\l}aban's lower and upper stability bounds, \cite[(1.14)]{Ba1} and \cite[(1.3)]{Ba2}, at a fixed scalar mass, read at the two masses $m$ and $M_1$ and in King's torus normalisation of the counterterm; the form used here is established in \moving{existence part: Ba{\l}aban's torus bounds at block factor $5$}. Proposition~3.1 of \cite{Ba2}, behind the upper bound, is stated under a cap on the dimensionless scalar mass at the scales at which it is applied; it is cited under that cap, one more one-sided cap on the stopping scale and no condition on the bare mass (Remark~\ref{rem:s14:BALm}). Only a tail in $k$ and, for each $k$, large $J$ are used (Proposition~\ref{prop:s10:U}, Theorem~\ref{thm:s11:W}).

\begin{lemma}[the $c$-number input]\label{hyp:s03:T2}
Let $\alpha^\infty_k(M)=M^2+\dm_k(M)$ be King's free-boundary bare coefficient at scalar mass $M$ and spacing $\eps_k$, in the reference scheme $z$. For the auxiliary mass $M_1$ of Proposition~\ref{prop:s10:U}, the quantities
\[
\sigma_1(M_1)\defeq\liminf_{k\to\infty}\tfrac12\bigl(\alpha^\infty_k(m)-\alpha^\infty_k(M_1)\bigr)
\qquad\text{and}\qquad \beta_1(M_1)
\]
are finite, where $\beta_1(M_1)$ is the second constant entering Proposition~\ref{prop:s10:U}.
\end{lemma}
Proved in \moving{existence part: the $c$-number lemma}.

\begin{remark}[on \cref{hyp:s03:T2}]\label{rem:s03:T2}
(i) This is a statement about lattice Feynman integrals: the difference of two mass counterterms at the same photon mass. The limit $\alpha_\infty(M)=\lim_k\alpha^\infty_k(M)$ and its description as a function of the couplings are discussed in Lemma~\ref{lem:s14:CW}.
(ii) \emph{Consumers.} Proposition~\ref{prop:s10:U} (the constants $\sigma_+=s_1$, $H$, $\kappa$ of the tube) and the coupling range, Lemma~\ref{lem:s14:CW}.
\end{remark}

\begin{hypothesis}[\Tthree: the threshold]\label{hyp:s03:T3}
$\sigma_1(M_1)>0$ for some $M_1\in(0,m)$.
\end{hypothesis}

\begin{remark}[on \Tthree]\label{rem:s03:T3}
(i) \Tthree\ is the only condition on the couplings anywhere in this paper, and it is one-sided: it fails exactly on the set $E^<_W$ of couplings $(e,\lambda,m,\lmu)$ at which $m$ minimises $M\mapsto\alpha_\infty(M)$ on $(0,m]$ (Lemma~\ref{lem:s14:CW}). At one loop it holds if and only if
\[
e^2<e_c^2(\lambda,m,\lmu)=\max\Bigl\{\frac{8\pi(m+\lmu)^2}{m+2\lmu}\Bigl(1+\frac{\lambda(N+2)}{2\pi m}\Bigr),\ 4\pi(m+\lmu)\Bigl(1+\frac{\lambda(N+2)}{\pi m}\Bigr)\Bigr\},
\]
$N$ the number of real components of the scalar field (equation~\eqref{eq:s14:ec2}); so at one loop a half-line of values of $e^2$ fails. Beyond one loop, \Tthree\ is \emph{proved} on explicit coupling regions (Proposition~\ref{prop:s14:T3}): at $\lmu=m$ on a region containing $0\le\lambda,e^2\le0.1m$, and at every $\lmu\in(0,m]$ on explicit boxes; the main theorem is stated on the region where it holds. The two-sided variant and the set $E_W\subseteq E^<_W$ are described in Lemma~\ref{lem:s14:CW}.
(ii) \emph{Consumers.} Proposition~\ref{prop:s10:U} and Lemma~\ref{lem:s14:CW}.
\end{remark}

\subsection{The rotation theorem}

For a family $f$ of real $2$-forms and $J\ge J_0(f)$ let $L^f_J$ be the set of locally uniform limits on $\C^N$ of subsequences of $(a^{(a)}_{k,J})_k$, $a^{(a)}_{k,J}(z)=\ava{e^{\sum_jz_jF(f_j)}}$ (the comparison family; by Lemma~\ref{lem:s05:C2} and Montel's theorem the set is non-empty).

\begin{theorem}[rotation invariance of the comparison family's limit states]\label{hyp:s03:RC}
For every finite family $f$ of real $2$-forms in $C^\infty_c(\R^3;\Lambda^2)$ and every choice $\zeta_J\in L^f_J$, $J\ge J_0(f)$, the limit $Z^{\mathrm{LIM}}_\infty(\sum_jz_jf_j,0)\defeq\lim_{J\to\infty}\zeta_J(z)$ exists locally uniformly on $\C^N$, is entire and does not depend on the choices. The functional $Z^{\mathrm{LIM}}_\infty(\cdot,0)$ is invariant under the Euclidean group $\Euc$: $Z^{\mathrm{LIM}}_\infty(\gamma\cdot f,0)=Z^{\mathrm{LIM}}_\infty(f,0)$ for every Euclidean motion $\gamma$ and every complex coefficient vector, $F$ transforming as a $2$-form.
\end{theorem}
Proved in the rotation part \moving{rotation part}: King's cubes are compared with periodic tori carrying a smooth mass wall, the walled tori on two lattices whose orientations differ by $\arctan(3/4)$ are compared, and King's state is identified as the lower envelope of the walled tori.

\begin{remark}[on \cref{hyp:s03:RC}]\label{rem:s03:RC}
(i) The theorem is posed for the comparison family and for the $J$-limit of its subsequential fixed-cube limits, not for the fixed-cube limits themselves; it does not give \cref{hyp:s03:EC}. Under \cref{hyp:s03:EC}, Corollary~\ref{cor:s04:G2} identifies $Z^{\mathrm{LIM}}_\infty(\cdot,0)$ with King's printed state $\ZKinf(\cdot,0)$, which is thereby $\Euc$-invariant (Theorem~\ref{thm:s06:O3}).
(ii) \emph{Consumers.} The $O(3)$ clause of the curvature sector (Theorem~\ref{thm:s06:O3}), the Euclidean clause of the measure's ergodic components (Theorem~\ref{thm:s09:S2K}(c)), and the $O(3)$ form of the LIM theorem (Theorem~\ref{thm:s12:LIM}).
\end{remark}

\begin{remark}[the Wick tube on King's convention]\label{rem:s03:TWK}
Theorem~\ref{thm:s11:W} is proved for the comparison convention $\NDa$. Its transfer to King's convention --- the $(k,J)$-uniform tube bound~\eqref{eq:s11:W} for $Z^{a,K}_{k,J}$ and the fixed-cube limits on the tube --- is Theorem~\ref{thm:s13:TWK}, with the kill step of Remark~\ref{rem:s04:F}, the cell statement of Lemma~\ref{lem:s04:E} and the direction check for the negative leg as its convention-sensitive steps; Form~(B) of the main theorem then holds on King's convention (Proposition~\ref{prop:s13:WK}).
\end{remark}

\subsection{Hypotheses of the optional extension}

The following statements are hypotheses of the optional extension of the Wick leg to arbitrary amplitudes (\cref{thm:s23:FormA}) only; no clause of \cref{thm:s13:main} uses them.

\begin{hypothesis}[\Wtwo: the Wick sector at arbitrary amplitudes]\label{hyp:s03:W2}
Let $a$ be an admissible scheme. (a) For every finite family of real $2$-forms and real functions, every $J\ge J_0(f,h)$ and \emph{every} real coefficient vector $(z,w)$, the limit $\lim_kZ^a_{k,J}(z,w)$ exists and lies in $(0,\infty)$. (b) For every $q\in C^\infty_c(\R^3)$ with $q\ge0$,
\[
\sup_{J\ge J_q}\ \lim_{k\to\infty}Z^{a}_{k,J}(0,q)<\infty ,
\]
where $J_q$ is the least $J$ with $\supp q$ at $\ell^\infty$-distance greater than $1$ from $\R^3\setminus\CJ$.
\end{hypothesis}

\begin{remark}[on \Wtwo]\label{rem:s03:W2}
(i) Clause (a) is \cref{hyp:s03:W1} with the threshold on the positive scalar part removed. On the cone of amplitudes $\sup q\le s_1<\sigma_1(M_1)$, clause (b) is proved in this paper in $L^1$ form (Corollary~\ref{cor:s11:W2prime}, on the comparison convention; on $\NDK$ in $\limsup$ form by Remark~\ref{rem:s04:F}). Above the threshold both clauses are open; the exact obstruction is the free-energy bound $\Bminus{u}$ at amplitude $u$ above the threshold (Hypothesis~\ref{hyp:s23:Bminus}): the $\eps$-uniform stability of King's torus model with its bare coefficient shifted by $-2u$.
(ii) \emph{Consumers.} None in the headline. The headline Wick result is Form~(B) on the tube $\Dtube$, where only the proved cone instance enters, through Theorem~\ref{thm:s11:W}; \Wtwo\ is a hypothesis of Form~(A) (Theorem~\ref{thm:s23:FormA}) and of \RW.
\end{remark}

\begin{hypothesis}[\RW: rotation invariance of the Wick sector at arbitrary amplitudes]\label{hyp:s03:RW}
For an admissible scheme $a$ satisfying \Wtwo, the full object $Z^a_\infty=\lim_J\lim_kZ^a_{k,J}$ of the comparison family (whose existence and entireness follow from \Wtwo\ as in Lemma~\ref{lem:s08:W1}) is invariant under every Euclidean motion for all complex coefficients, $F$ transforming as a $2$-form and $W^a$ as a scalar.
\end{hypothesis}

\begin{remark}[on \RW]\label{rem:s03:RW}
\RW\ enters only in the Euclidean clause of Form~(A) (Theorem~\ref{thm:s23:FormA}). The headline Wick result, Form~(B) on the tube, has no Euclidean clause and does not assume it. The rotation part proves the curvature sector only.
\end{remark}

\begin{hypothesis}[\FAK: Form~(A) on King's convention]\label{hyp:s03:FAK}
On King's convention $\NDK$, under \Wtwo\ read for $\NDK$: the limit functional $Z^{a,K}_\infty$ is entire in all coefficients; the lattice symmetries $\Oh$ about every point hold for the joint field at all amplitudes; reflection positivity holds at arbitrary complex coefficients; the measure exists off the tube; and the Euclidean clause of the Wick sector holds for the King-family object.
\end{hypothesis}

\begin{remark}[on \FAK]\label{rem:s03:FAK}
\FAK\ is the King-convention counterpart of Form~(A) as the comparison-convention results give it. It is consumed only in the optional extension (Theorem~\ref{thm:s23:FormA}) and is listed in the open problems; the headline makes no use of it.
\end{remark}

\subsection{Hypotheses stated in later sections}

The statements on continuum subsequences (Theorem~\ref{thm:s12:LIM}) are proved under three hypotheses stated there, \Bcell\ (Hypothesis~\ref{hyp:s12:Bcell}; the cell bound, an instance of \cref{hyp:s03:W1}), \SCH\ (Hypothesis~\ref{hyp:s12:SCH}; $a=z$ or $\sup_k\abs{z_k-a_k}<\infty$) and \UNIQ\ (Hypothesis~\ref{hyp:s12:UNIQ}; any two continuum subsequences give the same LIM state). They are hypotheses of those statements only: all three are proved for the objects of the main theorem (Proposition~\ref{prop:s12:ID}), and none is a hypothesis of Theorem~\ref{thm:s13:main}. The optional extension (Theorem~\ref{thm:s23:FormA}) uses $\Bminus{u}$ (Hypothesis~\ref{hyp:s23:Bminus}), the obstruction behind \Wtwo\ above the threshold. The limit of vanishing photon mass is not treated in this paper.

\subsection{Results of the earlier paper used as inputs}

\begin{remark}[results of the torus paper used as proved]\label{rem:s03:companion-inputs}
The following results of the authors' earlier paper on the torus limits \cite{SymPhase} are used as proved and cited by their labels there; none is re-proved here, and each is used only under the hypotheses it carries there, which are matched to King's state at the places where they are applied (Theorem~\ref{thm:s09:S2K}, Proposition~\ref{prop:s14:T3}): its abstract chain from a limit functional with the properties of Theorem~\ref{thm:s05:C} and the reflection positivity of Theorem~\ref{thm:s08:OS2C} to the measure, its ergodic decomposition and weak clustering --- thm:D:decomp, lem:D:realize, lem:D:hermite, prop:discharge, thm:D:mean, thm:D:diag, prop:D:recon, prop:D:fixed(iii), cor:D:schwinger(iv), add:Ohdecomp, add:OS4, p:cor:rotcor, prop:C-tight, prop:C-RP(a), lem:C-FP and the lemmas of its Appendix~G1 (the two-loop majorants behind Proposition~\ref{prop:s14:T3}); Theorem~main(ii); the lemmas on the flat state and the current representation, lem:flat, lem:C-flatlimit, lem:C-current, lem:C-Fsector.
Of \cite{BFS1}, Theorems~6.1 and~6.2(1),(2) are cited as printed; items (3),(4) of its Theorem~6.2, a sketch in print, are a lemma of this paper (Lemma~\ref{lem:s02:BFS}). The chessboard estimate is Theorem~4.1/4.2 of \cite{FILS78}, applied after its hypotheses are checked (Proposition~\ref{prop:s10:RP}). Nothing of \cite{K1,K2} is used as a result: King's existence theorem is proved in \cref{sec:k01} and \moving{existence section with scalar sources}, and his reflection-positivity argument for the modified box and his passage to subsequences are replaced by proofs in this paper.
\end{remark}

% s04 The two cube conventions: interlacing and margins
\section{The two cube conventions: interlacing and margins}

King's printed state lives on King's convention $\NDK_{k,J}$ (\cref{def:s02:Kcube}); the rotation part and the Wick-tube estimates of this paper use the comparison convention $\NDa_{k,J}$ (\cref{def:s02:acube}). We show that the two families are interlaced in $J$ (\cref{lem:s04:G}), draw the consequences for the limits (\cref{cor:s04:G1,cor:s04:G2}), and show that the comparisons of a cube with a reflected cube, a larger box and a tiling by cells carry over to King's convention with the printed margins (\cref{lem:s04:C,lem:s04:D,lem:s04:E}).

\subsection{The two conventions and the support condition}

Throughout, $\eps=\eps_k=5^{-k}$, $\XJ=\CJ\cap\Lk=\{-MJ,\dots,MJ-\eps\}^3$, $\TJ=\{x\in\XJ:\ x_\mu=MJ-\eps\text{ for some }\mu\}$ and $\alive^K_J=\XJ\setminus\TJ=\{-MJ,\dots,MJ-2\eps\}^3$. Both states carry the same photon variables $A_\alpha(x)$, $x\in\XJ$, with the Neumann form of $\XJ$, so their photon Gaussians are identical and every plaquette observable is the same function of the shared variables. They differ exactly by the matter on $\TJ$: alive in the comparison convention (scalar Dirichlet planes $x_\mu=-MJ-\eps$, $x_\mu=MJ$), killed in King's (planes $x_\mu=-MJ-\eps$, $x_\mu=MJ-\eps$, $2MJ$ apart, a multiple of every block scale; cf.\ \cref{rem:s03:cube-identification}). A bond with one end alive and the other killed or outside the site set contributes its diagonal part $\frac\eps2r_x^2$ at the alive end, so an alive site carries six diagonals in either convention. In the proof of his Theorem~3.2 and in his box form (3.38), King \cite{K2} works in the comparison convention: every site of the region is alive there.

For $(f,g)$ with $f$ a real $2$-form and $g\ge0$ a scalar, both in $C^\infty_c$, we write as in \cref{def:s02:sources} $\Theta(\mathrm if,g)=e^{\mathrm iF(f)+W(g)}$ (class $\cI$), $\Theta(f,-g)=e^{F(f)-W(g)}$ (class $\cD$), $W(g)=\sum_x\eps^3g(x)(|\phi(x)|^2-2\sW_\eps)$ (free Wick constant), and $Z^a_{k,J}(\mathrm if,g)=e^{-a_k\sS(g)}\av{\Theta(\mathrm if,g)}$, $Z^a_{k,J}(f,-g)=e^{a_k\sS(g)}\av{\Theta(f,-g)}$, $\sS(g)=\sum_x\eps^3g(x)$; a superscript $K$ ($Z^{a,K}_{k,J}$) marks King's convention. All these expectations are real. The sources lie inside $\XJ$ if
\begin{equation}\label{eq:s04:S}\tag{S}
\begin{gathered}
\text{every plaquette $p=(x;\alpha,\beta)$ of $\Lk$ with $f_{c(p)}(e_\alpha,e_\beta)\ne0$ has its four corners in $\XJ$, and}\\
\text{every site $x\in\Lk$ with $g(x)\ne0$ lies in $\alive^K_J=\XJ\setminus\TJ$.}
\end{gathered}
\end{equation}
If $\supp f\cup\supp g\subset\{y:\dist_\infty(y,\R^3\setminus\CJ)>1\}$, as every hypothesis of this paper demands of its sources, then \eqref{eq:s04:S} holds: a site with $g(x)\ne0$ has $-MJ+1<x_\mu<MJ-1\le MJ-\eps$, so $x\in\alive^K_J$; a plaquette with $f_{c(p)}\ne0$ has its centre at $\ell^\infty$-distance $>1$ from $\R^3\setminus\CJ$ and its corners within $\eps/2\le\frac12$ of it, hence in $\XJ$. Also \eqref{eq:s04:S} for $X_{J_0}$ implies \eqref{eq:s04:S} for every $J\ge J_0$.

\subsection{Cut-outs and the interlacing lemma}

\begin{lemma}[cut-outs from a box with killed sites]\label{lem:s04:G3}
Let $\av{\cdot}_Y$ be the state of \cref{def:s02:acube} on a finite set $V\subset\Lk$ of sites with the matter killed on $T\subset V$, and $X\subset V$ any subset. Let $(f,g)$, $g\ge0$, satisfy \textup{(S$_X$)}: every plaquette with $f\ne0$ has its four corners in $X$ and every site with $g\ne0$ lies in $X$ ($g$ on $X\cap T$ is immaterial). Then deleting from $\av{\cdot}_Y$ the photon gradient terms coupling $X$ to $V\setminus X$ and the cosine parts of the hopping terms between $X$ and $V\setminus X$ yields a product whose $X$-factor is $\av{\cdot}^{(a),X\cap T}_X$, the state of \cref{def:s02:acube} on $X$ with the matter killed on $X\cap T$; consequently
\[
\av{\Theta(\mathrm if,g)}^{(a),X\cap T}_X\ \le\ \av{\Theta(\mathrm if,g)}_Y,\qquad
\av{\Theta(f,-g)}^{(a),X\cap T}_X\ \ge\ \av{\Theta(f,-g)}_Y,
\]
and the same for $Z^a$. In particular, \textup{(a)} if $X\cap T=\emptyset$ the factor is $\av{\cdot}^{(a)}_X$; \textup{(b)} if $X=\prod_\mu\{a_\mu,\dots,b_\mu\}$ is a box and $X\cap T=\{x\in X:\ x_\mu=b_\mu\text{ for some }\mu\}$, the factor is the state of King's convention on $X$.
\end{lemma}

\begin{proof}
Each deletion is in the cone direction: deleting a photon gradient term does not increase $\av{\Theta_{F,g}}$ and does not decrease $\av{\Theta^-_{F,g}}$ (\cref{lem:s02:BFS}(a),(d)), and setting a hopping coefficient $c_{xy}$ to $0$ likewise (\cref{lem:s02:BFS}(b),(d)); every intermediate measure is a BFS measure. After all deletions nothing couples the $X$-variables to the rest, so the measure is a product (\cref{lem:s02:BFS}(e)). In the $X$-factor the photon is the Neumann photon of $X$; a site of $X\cap T$ carries its point mass; an alive $x\in X\setminus T$ carries the single-site weight times $\frac\eps2r_x^2$ for each of its six bonds, since in $\av{\cdot}_Y$ every bond at an alive $x$ contributes $\frac\eps2r_x^2$ at $x$ whether its other end is alive in $V$ (the diagonal part of $\frac\eps2|e^{\mathrm ie\eps A}\phi_y-\phi_x|^2$), killed, or outside $V$ ($\phi=0$ there), and the cut keeps all diagonals; the cosines are those of the bonds within $X$ between alive sites. That is the state of \cref{def:s02:acube} on $X$ killed on $X\cap T$. By (S$_X$) the observable is a function of the $X$-factor's variables (at $x\in X\cap T$ the factor $e^{\eps^3g(x)r_x^2}$ is $1$ in both states), and the $c$-numbers $e^{\mp2\sW_\eps\sS(g)}$, $e^{\mp a_k\sS(g)}$ are the same positive constants on both sides (\cref{lem:s02:alg}(a); the counterterm is site-independent). Hence $\av{\Theta}^{(a),X\cap T}_X$ is the expectation of $\Theta$ in the fully deleted measure, which is $\le\av{\Theta}_Y$ on $\cI$ and $\ge$ on $\cD$.
\end{proof}

The clause of \cref{lem:s02:cut} that $X$ and its neighbours lie in one fundamental cell is a torus clause; for a box it has no content, since a bond leaving $V$ already carries its Dirichlet diagonal in $\av{\cdot}_Y$. Every sandwich $C_{J-1}\subset\cdot\subset C_{J+1}$ of this paper, read in King's convention, is an instance of \cref{lem:s04:G3}.

\begin{lemma}[interlacing of the two cube families]\label{lem:s04:G}
Fix $k\ge0$, $J\ge1$ and $(f,g)$, $g\ge0$, satisfying \eqref{eq:s04:S} for $\XJ$. Then
\[
\avK{\Theta(\mathrm if,g)}\ \le\ \ava{\Theta(\mathrm if,g)}\ \le\ \langle\Theta(\mathrm if,g)\rangle^{K}_{k,J+1},\qquad
\avK{\Theta(f,-g)}\ \ge\ \ava{\Theta(f,-g)}\ \ge\ \langle\Theta(f,-g)\rangle^{K}_{k,J+1},
\]
and the same for the scheme functionals $Z^a_{k,J}(\mathrm if,g)$, $Z^a_{k,J}(f,-g)$ of the three states.
\end{lemma}

\begin{proof}
\emph{Step 0 (one observable, one multiplier).} By \eqref{eq:s04:S} the plaquettes with non-zero coefficient have all corners in $\XJ$, whose photon variables are present in all three states, so $F(f)$ is one real linear functional $F$ of the shared variables; the sites with $g\ne0$ lie in $\alive^K_J$, alive in all three states; $\sW_\eps$, $\dm_\eps$, $\sS(g)$ are the same for the three regions. Hence, with $g'_x:=\eps^3g(x)\ge0$, $\Theta(\mathrm if,g)=e^{-2\sW_\eps\sS(g)}\Theta_{F,g'}$ and $\Theta(f,-g)=e^{2\sW_\eps\sS(g)}\Theta^-_{F,g'}$, where $\Theta_{F,g'}=e^{\mathrm iF}e^{\sum_xg'_xr_x^2}$, $\Theta^-_{F,g'}=e^{F}e^{-\sum_xg'_xr_x^2}$ are the observables of \cref{lem:s02:BFS}, and the scheme multiplier $e^{\mp a_k\sS(g)}$ is one positive constant for all three states (\cref{lem:s02:alg}(a)). It suffices to treat $\av{\Theta_{F,g'}}$, $\av{\Theta^-_{F,g'}}$.

\emph{Step 1 (kill step on $\cI$).} For $s\ge0$ let $\av{\cdot}_s$ be $\ava{\cdot}$ with the single-site measure at each $x\in\TJ$ multiplied by $e^{-s\eps^3r_x^2}$, a BFS measure, and $\varphi(s):=\av{\Theta_{F,g'}}_s=N(s)/D(s)$, $N(s)=\int\Theta\,e^{-s\eps^3\sum_{\TJ}r_x^2}d\mu_0$, $D(s)=\int e^{-s\eps^3\sum_{\TJ}r_x^2}d\mu_0$, $\mu_0$ the unnormalised measure at $s=0$. The $s$-derivative of either integrand is bounded in modulus by $\eps^3\sum_{\TJ}r_x^2$ times the integrand at $s=0$, integrable against the quartic weight $e^{-\lambda\eps^3r^4}$ and the Gaussian; so differentiation under the integral sign is justified, $D(s)>0$, and
\[
\varphi'(s)=-\eps^3\sum_{x\in\TJ}\big(\av{\Theta\,r_x^2}_s-\av{\Theta}_s\av{r_x^2}_s\big)=-\eps^3\sum_{x\in\TJ}\av{\Theta;r_x^2}_s\ \le\ 0
\]
by \cref{lem:s02:BFS}(c). So $\varphi(s)\le\varphi(0)=\ava{\Theta}$. Since $\Theta_{F,g'}$ is continuous with $|\Theta_{F,g'}|=e^{\sum_xg'_xr_x^2}$ and $\ava{\cdot}$ is of King type, \cref{lem:s02:Kprime} gives $\varphi(s)\to\av{\Theta}_{\TJ}$, the state with the matter killed on $\TJ$, which is $\avK{\cdot}$. Hence $\avK{\Theta}\le\ava{\Theta}$.

\emph{Step 2 (cut step).} Apply \cref{lem:s04:G3} with $V:=X_{J+1}$, $T:=T_{J+1}$ (so $\av{\cdot}_Y=\langle\cdot\rangle^K_{k,J+1}$), $X:=\XJ$: $\XJ\subset X_{J+1}$, and $\XJ\cap T_{J+1}=\emptyset$ since every site of $T_{J+1}$ has a coordinate $M(J+1)-\eps\ge MJ>MJ-\eps$. (At $k=0$, $M=1$ the outward neighbours of $\XJ$ lie on $T_{J+1}$ and are killed in $\av{\cdot}_Y$; harmless, since the bond from an alive $x\in\XJ$ to such a neighbour already carries $\frac\eps2r_x^2$ at $x$, the Dirichlet diagonal of the comparison cube.) By \cref{lem:s04:G3}(a) the $\XJ$-factor is $\ava{\cdot}$, and \eqref{eq:s04:S} is (S$_{\XJ}$): $\ava{\Theta(\mathrm if,g)}\le\langle\Theta(\mathrm if,g)\rangle^K_{k,J+1}$, $\ava{\Theta(f,-g)}\ge\langle\Theta(f,-g)\rangle^K_{k,J+1}$.

\emph{Step 3 (kill step on $\cD$).} With $\av{\cdot}_s$ as in Step~1 and $\varphi^-(s):=\av{\Theta^-_{F,g'}}_s$ (differentiation justified as before, now with $|\Theta^-|\le e^{F(A)}$ integrable against the Gaussian), $(\varphi^-)'(s)=-\eps^3\sum_{x\in\TJ}\av{\Theta^-;r_x^2}_s\ge0$ by \cref{lem:s02:BFS}(d), so $\varphi^-(s)\ge\varphi^-(0)=\ava{\Theta^-}$; as $|\Theta^-_{F,g'}|\le e^{F(A)}$ with $F$ real linear, \cref{lem:s02:Kprime} (with $\ell=F$, $C=0$) gives $\varphi^-(s)\to\avK{\Theta^-}$. Hence $\avK{\Theta(f,-g)}\ge\ava{\Theta(f,-g)}$.

\emph{Step 4 ($Z^a$).} The multiplier of Step~0 is one positive constant.
\end{proof}

Nothing in the proof uses a reflection symmetry of either cube (neither has one), the counterterm beyond its site-independence, or the renormalisation group. Since the observable vanishes at the killed sites in every kill step of this section, each such step also follows with no limit: set $c_{xy}=0$ on every bond at the set $U$ to be killed, one bond at a time (\cref{lem:s02:BFS}(b),(d)); the measure is then a product (\cref{lem:s02:BFS}(e)) whose $U$-factor cancels, and the other factor is the state with $U$ killed.

\subsection{Monotonicity of the K family and the limits}

\begin{corollary}[the K family is monotone in $J$]\label{cor:s04:G1}
Under \eqref{eq:s04:S} for $\XJ$: $\avK{\Theta(\mathrm if,g)}\le\langle\Theta(\mathrm if,g)\rangle^K_{k,J+1}$ and $\avK{\Theta(f,-g)}\ge\langle\Theta(f,-g)\rangle^K_{k,J+1}$, and the same for $Z^{a,K}$; hence for all $J\ge J_0$ whenever the sources satisfy \eqref{eq:s04:S} for $X_{J_0}$.
\end{corollary}

\begin{proof}
Chain the two inequalities of \cref{lem:s04:G}. (For the comparison family this is \cref{lem:s04:G3}(a) with $V=X_{J+1}$, $X=\XJ$, \ie the first inequality of \cref{lem:s02:cut}.)
\end{proof}

The proof of Theorem~3.2 of \cite{K2} first replaces the $J$-dependent Dirichlet Wick subtraction (2.13) by the free, $J$-independent one and then proves the monotonicity in $J$ in the comparison convention (\cref{lem:s04:G3}(a)); by \cref{cor:s04:G1} that statement holds between cubes of King's convention too, with the kill step (Step~1) not printed there. The form with the Dirichlet constant follows in neither convention and is not used here; at $g=0$ no Wick constant enters.

\begin{corollary}[the limits]\label{cor:s04:G2}
Let $f=(f_1,\dots,f_N)$ be real $2$-forms in $C^\infty_c$ with $\supp f_j\subset\{\dist_\infty(\cdot,\R^3\setminus C_{J_0})>1\}$; for $J\ge J_0$, $z\in\C^N$ put $a^K_{k,J}(z):=\avK{e^{\sum_jz_jF(f_j)}}$, $a^{(a)}_{k,J}(z):=\ava{e^{\sum_jz_jF(f_j)}}$, $g(z):=\sum_j(\Re z_j)f_j$. At the cone points $z=\mathrm it$, $t\in\R^N$, both are real with modulus $\le1$, and \eqref{eq:s04:S} holds for $\XJ$.
\begin{enumerate}[label=\textup{(\roman*)},leftmargin=*]
\item \textup{(One $J$-limit at fixed $k$; no hypothesis.)} For every $k$ and $t$, $J\mapsto a^K_{k,J}(\mathrm it)$ and $J\mapsto a^{(a)}_{k,J}(\mathrm it)$ are nondecreasing on $J\ge J_0$, bounded by $1$ and interlaced; hence $\lim_Ja^K_{k,J}(\mathrm it)=\lim_Ja^{(a)}_{k,J}(\mathrm it)\in[-1,1]$. The same holds for $\av{\Theta(\mathrm if,g)}$, $g\ge0$, with a common limit in $(-\infty,\infty]$.
\item \textup{(King's cube limits.)} By \cref{hyp:s03:EC}, posed on King's convention, for every such family, every $J\ge J_0$ and every $t\in\R^N$, $\lim_ka^K_{k,J}(t)$ exists in $(0,\infty)$. Then: ($\alpha$) $a^K_{k,J}\to\ZKJ$ locally uniformly on $\C^N$, $\ZKJ$ is entire and $|\ZKJ(z)|\le e^{\norm{g(z)}_2^2}$; ($\beta$) $\ZKJ\to\ZKinf$ locally uniformly on $\C^N$, $\ZKinf$ is entire and $|\ZKinf(z)|\le e^{\norm{g(z)}_2^2}$; ($\gamma$) for all $J\ge J_0$ and $t\in\R^N$,
\[
\ZKJ(\mathrm it)\ \le\ \liminf_ka^{(a)}_{k,J}(\mathrm it)\ \le\ \limsup_ka^{(a)}_{k,J}(\mathrm it)\ \le\ Z^K_{J+1}(\mathrm it),
\]
hence $\lim_J\liminf_ka^{(a)}_{k,J}(\mathrm it)=\lim_J\limsup_ka^{(a)}_{k,J}(\mathrm it)=\ZKinf(\mathrm it)$.
\item \textup{(The subsequential limits of the comparison family.)} Let $\cS=(k_m)$ be an increasing sequence of spacings along which $a^{(a)}_{k_m,J}\to Z^{\cS}_J$ locally uniformly on $\C^N$ for every $J\ge J_0$ and every such family (a continuum subsequence of the comparison family, \cref{def:s12:contsub}; they exist with no hypothesis, \cref{lem:s12:LIM0}). Then $Z^{\cS}_J\to\ZKinf$ locally uniformly on $\C^N$ as $J\to\infty$: every subsequential fixed-cube limit of the comparison family, followed by $J\to\infty$, is King's printed state. In particular, if the fixed-cube limits of the comparison family exist along the full sequence, $\lim_J\lim_ka^{(a)}_{k,J}=\ZKinf$.
\end{enumerate}
\end{corollary}

\begin{proof}
(i) Monotonicity is \cref{cor:s04:G1} for the K family and \cref{lem:s04:G3}(a) for the comparison family; \cref{lem:s04:G} interlaces them; two nondecreasing sequences with $a_J\le b_J\le a_{J+1}$ have the same limit, finite or $+\infty$.

(ii)($\alpha$) This is the first step of the proof of \cref{thm:s05:C}(b) run on the K family. (1) $a^K_{k,J}$ is entire (Gaussian decay in $A$, $e^{-\lambda\eps^3r^4}$ decay in $r$ at the alive sites). (2) The local bound of \cref{lem:s05:C2} holds for the K family with the same constant: $\Re\sum_jz_jF(f_j)=F(g(z))$, so $|a^K_{k,J}(z)|\le\avK{e^{F(g(z))}}$; \cref{lem:s02:GD}(a) with $s=1$ and $v:=d^*_kg(z)$ (defined by $F_{\Lk}(g(z))=\ip{v}{A}_k$) applies to any BFS measure whose Gaussian has covariance $C$ and gives $\av{e^{\ip{v}{A}}}\le e^{\ip{v}{Cv}/2}$; the K cube is a BFS measure whose Gaussian is the Neumann photon Gaussian of $\XJ$, the same as the comparison cube's, so its covariance is the $C$ of \cref{lem:s05:C1}, a statement about this $C$ and the plaquettes with corners in $\XJ$: $\ip{v}{Cv}_k\le2\norm{g(z)}_{2,k}^2$. Hence $|a^K_{k,J}(z)|\le e^{\norm{g(z)}_{2,k}^2}\le e^{3C_0\fournorm{g(z)}^2}$ uniformly in $k\ge0$, $J\ge J_0$, by (E3) of \cref{lem:s02:E}: the plaquette centres of each of the three orientations form a translate of $\eps_k\Z^3$, each contributing at most $C_0\fournorm{g}^2$ to $\norm{g}_{2,k}^2$ by (E1), and the three orientations give $3C_0$. Nothing in this chain sees the matter beyond the state being a BFS measure. (3) \Cref{hyp:s03:EC} gives convergence at every real point. (4) $\R^N$ is a set of uniqueness for entire functions on $\C^N$; by Montel's theorem every subsequence has a locally uniformly convergent sub-subsequence, all limits are entire and agree on $\R^N$, hence everywhere \cite{Hor90}. So $a^K_{k,J}\to\ZKJ$ locally uniformly, $\ZKJ$ is entire, and the bound passes to the limit, $|\ZKJ(z)|\le\lim_ke^{\norm{g(z)}_{2,k}^2}=e^{\norm{g(z)}_2^2}$ (\cref{lem:s02:E}), a $J$-free bound.

($\beta$) By \cref{cor:s04:G1} at every $k$ and ($\alpha$), $J\mapsto\ZKJ(\mathrm it)$ is nondecreasing and $\le1$, so $\ZKinf(\mathrm it):=\lim_J\ZKJ(\mathrm it)$ exists; $\{\ZKJ\}_{J\ge J_0}$ is locally uniformly bounded by the $J$-free bound of ($\alpha$); $(\mathrm i\R)^N$ is a maximal totally real subspace, hence a set of uniqueness for entire functions; Montel's theorem and the identity theorem give $\ZKJ\to\ZKinf$ locally uniformly, $\ZKinf$ entire, $|\ZKinf(z)|\le e^{\norm{g(z)}_2^2}$. (This is the proof of \cref{thm:s05:C}(b) with \cref{lem:s02:cut} replaced by \cref{cor:s04:G1}.)

($\gamma$) \cref{lem:s04:G} at every $k$, with $k\to\infty$ at the two ends (limits by ($\alpha$)) and $\liminf$, $\limsup$ in the middle, gives the sandwich; ($\beta$) gives the limits.

(iii) \cref{lem:s04:G} along $\cS$, with the full-sequence limits of ($\alpha$) at the ends, gives $\ZKJ(\mathrm it)\le Z^{\cS}_J(\mathrm it)\le Z^K_{J+1}(\mathrm it)$, so $\lim_JZ^{\cS}_J(\mathrm it)=\ZKinf(\mathrm it)$ by ($\beta$). $\{Z^{\cS}_J\}_J$ is locally bounded by the $J$-free bound $e^{\norm{g(z)}_2^2}$ (\cref{lem:s05:C2} passed to the limit), so by Montel's theorem and the uniqueness set $(\mathrm i\R)^N$ it converges locally uniformly to an entire function (\cref{lem:s12:LIM1}), which agrees with $\ZKinf$ on $(\mathrm i\R)^N$, hence on $\C^N$. If the fixed-cube limits exist along the full sequence, that sequence is a continuum subsequence.
\end{proof}

Statements (ii)--(iii) also hold with the two conventions exchanged, for $J\ge J_0+1$, when the comparison family is convergent at fixed cubes (\cref{def:s03:conv}; read \cref{lem:s04:G} as $a^{(a)}_{k,J-1}\le a^K_{k,J}\le a^{(a)}_{k,J}$); the exchanged (iii) needs \cref{lem:s12:LIM0,lem:s12:LIM1} for the K family, whose proofs use only the local bound and the photon-only bound $\av{F(u-u')^2}\le2\norm{u-u'}_{2,k}^2$, both transferred by the covariance identity in ($\alpha$).

\emph{What does not follow.} At fixed $J$, \cref{lem:s04:G} yields exactly ($\gamma$) and nothing on whether $\lim_ka^{(a)}_{k,J}(\mathrm it)$ exists: the bounds $\ZKJ(\mathrm it)$, $Z^K_{J+1}(\mathrm it)$ differ in general (the cubes differ by $M$, so the gap does not shrink with $k$), and nothing constrains the oscillation of $a^{(a)}_{k,J}(\mathrm it)$ in $k$ between them. So the existence of the fixed-cube limits on one convention does not imply it on the other by these inputs, and likewise with scalar sources: the existence theorem must be proved on the family on which the renormalisation-group construction is carried out; the other is a comparison family with the same infinite-volume state by (ii)--(iii). No counterexample is claimed.

\subsection{Three margins in King's convention}

Let $\sigma_3=\diag(1,1,-1)$, $X:=\XJ$, $X_+:=\sigma_3X=(\CJ+\eps e_3)\cap\Lk$ (third coordinates $\{-MJ+\eps,\dots,MJ\}$), $\sigma_3\TJ=\{x\in X_+:\ x_1=MJ-\eps\text{ or }x_2=MJ-\eps\text{ or }x_3=-MJ+\eps\}$ and $A_\sigma:=\sigma_3(\alive^K_J)=\{x:\ x_1,x_2\in\{-MJ,\dots,MJ-2\eps\},\ x_3\in\{-MJ+2\eps,\dots,MJ\}\}$.

\begin{lemma}[the reflection sandwich]\label{lem:s04:C}
Let $(f,g)\in\cK$ be supported at $\ell^\infty$-distance $>1$ from $\R^3\setminus C_{J_0}$.
\textup{(a)} If $M\ge2\eps_k$ (\ie $k\ge1$, or $M\ge2$), then for every $J\ge J_0+1$
\[
Z^{a,K}_{k,J-1}(\mathrm if,g)\ \le\ Z^{a,K}_{k,J}(\mathrm i\,\sigma_3\!\cdot\!f,\sigma_3\!\cdot\!g)\ \le\ Z^{a,K}_{k,J+1}(\mathrm if,g).
\]
\textup{(b)} For every $k\ge0$, $M\ge1$ and $J\ge J_0+2$ the same holds with $J\mp2$ in place of $J\mp1$.
\textup{(c)} For $k=0$, $M=1$ the one-step sandwich \textup{(a)} is not obtainable by deletions and kills: $A_\sigma$ contains alive sites ($x_3=MJ$) killed in King's cube of index $J+1$, and $\alive^K_{J-1}$ contains sites ($x_3=-MJ+\eps$) killed in the reflected state. Whether the inequality itself holds there is left open.
\end{lemma}

\begin{proof}
\emph{Step 1 (transport).} By \cref{lem:s06:PK}, $Z^{a,K}_{k,J}(\mathrm i\sigma_3\!\cdot\!f,\sigma_3\!\cdot\!g)=e^{-a_k\sS(g)}\av{\Theta(\mathrm if,g)}_R$, with $R$ the state on $X_+$ for the frame $(e_1,e_2,-e_3)$ killed on $\sigma_3\TJ$ ($\sS(\sigma_3\cdot g)=\sS(g)$). In the frame $(e_\alpha)$, as in \cref{lem:s02:refl}, $R$ is the BFS measure with $A_1,A_2$ on $X_+$ and $A_3$ on $X_+-\eps e_3=X$ (gradient pairs within these sets, masses on them), matter alive exactly on $A_\sigma$ with point masses on $\sigma_3\TJ$ and $\phi=0$ outside $X_+$, all six diagonals at each alive site, and the standard cosine on every bond with both ends in $A_\sigma$ (the flipped bond with $A'_3(x)=-A_3(x-\eps e_3)$ gives the standard cosine of $\langle x-\eps e_3,x\rangle$); $F(f)$, $W(g)$ are the standard ones (the plaquettes $(x;\alpha,3)$ pick up two sign changes and keep their centres). $R$ is of King type. The sources satisfy \eqref{eq:s04:S} for $X_{J-1}$ (resp.\ $X_{J-2}$), so by the inclusions below Step~0 of \cref{lem:s04:G} applies to all states compared.

\emph{Step 2 (upper end).} If $A_\sigma\subset\alive^K_{J'}$ and $X_+\cup X\subset X_{J'}$: from $\langle\cdot\rangle^K_{k,J'}$ kill every site of $\alive^K_{J'}\setminus A_\sigma$ (\cref{lem:s02:BFS}(c), \cref{lem:s02:Kprime}, which allows the already killed $T_{J'}$), leaving exactly the cosines of bonds within $A_\sigma$, \ie $R$'s; then delete every photon gradient term except the pairs within $X_+$ for $A_1,A_2$ and within $X$ for $A_3$ (\cref{lem:s02:BFS}(a); these are gradient pairs of $X_{J'}$). The result is a product (\cref{lem:s02:BFS}(e)) whose factor on $R$'s variables is $R$ (photon, alive set with all six diagonals by the bookkeeping of \cref{lem:s04:G3}, cosines); the other factor cancels. Hence $\av{\Theta(\mathrm if,g)}_R\le\langle\Theta(\mathrm if,g)\rangle^K_{k,J'}$. The inclusions hold iff $MJ\le MJ'-2\eps$ (direction $3$; the rest is trivial): for $J'=J+1$ iff $M\ge2\eps$, for $J'=J+2$ always.

\emph{Step 3 (lower end).} If $\alive^K_{J''}\subset A_\sigma$ and $X_{J''}\subset X\cap X_+$: from $R$ kill every site of $A_\sigma\setminus\alive^K_{J''}$ (\cref{lem:s02:BFS}(c), \cref{lem:s02:Kprime}) and delete every photon gradient term coupling an $X_{J''}$-variable to another variable of $R$ (\cref{lem:s02:BFS}(a)). The $X_{J''}$-factor has the Neumann photon of $X_{J''}$ (gradient pairs within $X_+\cap X_{J''}=X\cap X_{J''}=X_{J''}$) and the matter alive on $\alive^K_{J''}$ with all diagonals, killed on $T_{J''}$: it is $\langle\cdot\rangle^K_{k,J''}$; the rest cancels. Hence $\langle\Theta(\mathrm if,g)\rangle^K_{k,J''}\le\av{\Theta(\mathrm if,g)}_R$. The inclusions hold iff $-MJ''\ge-MJ+2\eps$: for $J''=J-1$ iff $M\ge2\eps$, for $J''=J-2$ always.

\emph{Step 4.} Steps 2--3 with $J\pm1$ give (a), with $J\pm2$ give (b); the multipliers agree (\cref{lem:s02:alg}(a)). For (c): at $k=0$, $M=1$ the top of $A_\sigma$ in direction $3$ is $J$ while that of $\alive^K_{J+1}$ is $J-1$, and the bottom of $\alive^K_{J-1}$ is $-J+1=-MJ+\eps$, the killed layer of $R$; deletions and kills never revive a killed site.
\end{proof}

In the curvature sector ($g=0$) the limit sentences of \cref{thm:s06:S1K} run on (a) for $k\ge1$ or on (b) for all $k$, with $k\to\infty$ by \cref{cor:s04:G2}(ii) and $J\to\infty$ by \cref{cor:s04:G1}; with scalar sources they need the K-family tube limits, supplied by \cref{thm:s13:TWK} (\cref{rem:s04:wick}).

\begin{lemma}[the upper box needs no change]\label{lem:s04:D}
Let $\Omega\subset B$ be lattice boxes at spacing $\eps$ such that $B$ contains every site within $\ell^\infty$-distance $N$ lattice steps of $\Omega$, $N\ge1$ (the hypothesis of \cref{lem:s07:U}, whose mass condition forces $N\ge1$), and let $\av{\cdot}^K_B$ be the state of King's convention on $B$ (killed on its top face layers $T_B$). Then \textup{(a)} $\Omega\cap T_B=\emptyset$; \textup{(b)} the second part of \cref{lem:s07:U} holds with $\av{\cdot}^K_B$ in place of the comparison state on $B$, with the same proof, and its first part is a photon statement, unchanged; \textup{(c)} consequently the sandwich in the proof of \cref{thm:s08:OS2C} holds for King's cubes with the printed cubes and margins \eqref{eq:s07:margin}: for every $k$ and every $\zeta$ admitted there, $Z^{a,K}_{k,J}(\zeta)\le\Phi_k(\zeta)\le Z^{a,K}_{k,J+2n}(\zeta)$, $\Phi_k$ the functional of the box state of \cref{def:s07:box} on $\Omega_{J,k}$; the lower end is \cref{lem:s07:L} with Step~1 of \cref{lem:s04:G} prepended. (By \cref{cor:s04:G1} the cube of index $J+2n+1$ would also do; it is not needed.)
\end{lemma}

\begin{proof}
(a) For $x\in\Omega$ and each $\mu$, $x+N\eps e_\mu\in B$, so $x_\mu<\max_Bx_\mu$ and $x\notin T_B$. (b) The proof of \cref{lem:s07:U} starts from the state on $B$ and removes all hopping cosines touching $B\setminus\Omega$; the matter on $\Omega$ is then the Dirichlet matter of $\Omega$ (diagonals of the bonds leaving $\Omega$ kept, each such bond ending in $B$ as $N\ge1$), coupled to $A$ only through the bonds within $\Omega$, and the sites of $B\setminus\Omega$ decouple. For $\av{\cdot}^K_B$ the cosines touching $T_B$ are already absent; remove the remaining cosines touching $B\setminus\Omega$ (\cref{lem:s02:BFS}(b) on $\cI$, (d) on $\cD$). By (a) every site of $\Omega$ is alive, each bond at $x\in\Omega$ carries $\frac\eps2r_x^2$ at $x$ (its other end alive, killed or outside), and the cosines within $\Omega$ are present; so the matter on $\Omega$ is the same Dirichlet matter, and the sites of $B\setminus\Omega$, alive or killed, decouple and cancel. The rest of the proof concerns only the photon law of $B$, the same in both conventions; the first part uses only the photon forms and a cutoff supported in $B$. (c) Lower end: \cref{lem:s07:L} gives $\langle Y\rangle^{(a)}_{C_J}\le\langle Y\rangle_{\Omega_{J,k}}$ on $\cI$ (its proof lands on the Dirichlet matter of $C_J$, the comparison convention), and Step~1 of \cref{lem:s04:G} gives $\langle Y\rangle^K_{C_J}\le\langle Y\rangle^{(a)}_{C_J}$, \eqref{eq:s04:S} holding since the sources of the enlarged family and their reflections are supported at $\ell^\infty$-distance $>1$ from $\R^3\setminus C_{J_0}$, $C_{J_0}\subset C_J$. Upper end: (b) with $B=X_{J+2n}$, whose hypothesis is the second margin of \eqref{eq:s07:margin}, a statement about photon sites. The multipliers agree (\cref{lem:s02:alg}(a)).
\end{proof}

For the cell bound we use the cells of \cref{def:s11:cells}: $\ell_c=2ML_c\ge10$; the cells $c=\prod_\mu[w_\mu,w_\mu+\ell_c)$ of the tiling $\cT_s$, $s\in\{0,1\}^3$, $w_\mu\in W_\mu(s_\mu)$, $W_\mu(0)=\ell_c\Z$, $W_\mu(1)=ML_c+\ell_c\Z$ (so $w_\mu,w_\mu+\ell_c,MJ\in\Z$); the good sets $G_s=\{x:\dist(x_\mu,W_\mu(s_\mu))>2\ \forall\mu\}$ and cutoffs $\chi_s$ with $\supp\chi_s\subset G_s$; $X_c:=c\cap\XJ$; $J_1(R)=\min\{J:MJ\ge R+\ell_c\}$ for sources supported in $Q_R=[-R,R]^3$; $B_{\rm cell}$, defined on the comparison cell $C_{L_c}$. $N^K_{k,C_J}$ is the normalised scalar-source functional of $\avK{\cdot}$ and $N^{(a),T}_{k,X}$ that of the comparison state on $X$ killed on $T\subset X$.

\begin{lemma}[the cell lemmas hold as printed]\label{lem:s04:E}
\textup{(a)} For a tiling $\cT_s$ and any lattice function $q\ge0$ on $\XJ$: $N^K_{k,C_J}(0,-q)\le\prod_iN^{(a),X_i\cap\TJ}_{k,X_i}(0,-q|_{X_i})$, over the site sets $X_i=c\cap\XJ$ of the cells meeting $\CJ$.
\textup{(b)} For a full cell $c\subset\CJ$: $X_c\cap\TJ\ne\emptyset$ iff $w_\mu+\ell_c=MJ$ for some $\mu$.
\textup{(c)} Let $q\ge0$ be any lattice function supported in $Q_R$, $J\ge J_1(R)$, $s\in\{0,1\}^3$ and $c\in\cT_s$ with $(\chi_sq)|_{X_c}\not\equiv0$. Then $-MJ<w_\mu$ and $w_\mu+\ell_c<MJ-2$ for every $\mu$: $c$ is a full cell inside $\CJ$ with $X_c\cap\TJ=\emptyset$, and $(\chi_sq)|_{X_c}$ is supported at $\ell^\infty$-distance $>2$ from $\R^3\setminus c$.
\textup{(d)} Hence \cref{lem:s11:W4,lem:s11:W5,lem:s11:W67} are unchanged and \cref{prop:s11:W8} holds for $\avK{\cdot}$ as printed, under its hypotheses and with its constant $K_-$: $\ln N^K_{k,C_J}(0,-q)\le K_-\wnorm{q}$ for every $k$, every $J\ge J_1(R)$ and every lattice $0\le q\le H$ supported in $Q_R$.
\textup{(e)} (Direction.) For a cell $c$ and $\emptyset\ne T\subset X_c$: $N^{(a),T}_{k,c}(0,-\tilde q)\ge N^{(a)}_{k,c}(0,-\tilde q)$ for $\tilde q\ge0$ vanishing on $T$. So $B_{\rm cell}$, a bound on a comparison cell, does not bound a source-carrying cell of King's type, and \textup{(c)} is load-bearing.
\end{lemma}

\begin{proof}
(a) The proof of \cref{lem:s11:W3} deletes the cross-cell photon gradient terms and hopping cosines, each in the $\cD$ direction (\cref{lem:s02:BFS}(d)); by \cref{lem:s02:BFS}(e) and the factor identification of \cref{lem:s04:G3} with $V=\XJ$, $T=\TJ$, $X=X_i$, the product has the factors $N^{(a),X_i\cap\TJ}_{k,X_i}$; the observable $e^{-\sum\eps^3qr^2}$ and the multiplier $e^{a_k\sS(q)}=\prod_ie^{a_k\sS(q|_{X_i})}$ factorise, and at the sites of $\TJ$ the observable factor is $1$ in every state. (b) The sites of a full cell have $x_\mu\in\{w_\mu,\dots,w_\mu+\ell_c-\eps\}$, and $x\in\TJ$ iff $w_\mu+\ell_c-\eps=MJ-\eps$ for some $\mu$. (c) Let $x\in X_c$ with $(\chi_sq)(x)\ne0$. Then $x\in G_s$; since $w_\mu,w_\mu+\ell_c\in W_\mu(s_\mu)$ and $w_\mu\le x_\mu<w_\mu+\ell_c$, this gives $w_\mu+2<x_\mu<w_\mu+\ell_c-2$; and $q(x)\ne0$ gives $|x_\mu|\le R$. Hence $w_\mu+\ell_c<R+\ell_c-2\le MJ-2$ (as $MJ\ge R+\ell_c$) and $w_\mu>-R-\ell_c+2\ge-MJ+2$. So $c\subset(-MJ,MJ)^3$, $w_\mu+\ell_c\ne MJ$, hence $X_c\cap\TJ=\emptyset$ by (b), and $\dist_\infty(x,\R^3\setminus c)=\min_\mu\min(x_\mu-w_\mu,w_\mu+\ell_c-x_\mu)>2$. (d) In the chain of \cref{prop:s11:W8} the Hölder step of \cref{lem:s11:W67} is an inequality within one state, convention-free; then (a) with the tiling $\cT_s$ and $q_s:=8\chi_sq$. A factor with $q_s|_{X_i}\equiv0$ equals $1$ whatever is killed on $X_i$ (this covers the clipped cells at the faces of $\CJ$ and any cell flush with them); a factor with $q_s|_{X_i}\not\equiv0$ is, by (c), the comparison-cell functional $N^{(a)}_{k,c}(0,-q_s|_c)$ of a full cell $c\subset\CJ$ with $q_s|_c$ supported at distance $>2$ from $\R^3\setminus c$ and $0\le q_s\le8H=H'$: exactly the hypotheses of \cref{lem:s11:W5} (hence of \cref{lem:s11:W4}: translation to $C_{L_c}$, $B_{\rm cell}$, \cref{prop:s11:W1prime} on the comparison box), as printed; \cref{lem:s11:W67} sums the cells. (A cell flush with a face, $w_\mu+\ell_c=MJ$, can meet $\supp q$ but never $\supp(\chi_sq)$; the chain feeds the cutoff sources, not $q$, to (a).) (e) Step~3 of \cref{lem:s04:G} applied to the comparison cell with $U=T$: $s\mapsto\av{\Theta^-}_s$ is nondecreasing (\cref{lem:s02:BFS}(d)) and tends to the killed state (\cref{lem:s02:Kprime}).
\end{proof}

\begin{remark}[one-line kill steps]\label{rem:s04:F}
\Cref{prop:s11:W2} (cube $\le$ torus on $\cI$), \cref{prop:s11:W1prime} (box $\le$ cube $\le$ torus), the lower end of \cref{lem:s07:L}, and the comparison of the cube with the box state of \cref{def:s07:box} on the cone (\cref{rem:s07:K8}) are each proved for the comparison cube; for King's cube prepend Step~1 of \cref{lem:s04:G} ($\avK{\cdot}\le\ava{\cdot}$ on $\cI$), and for the $\cD$ versions (cube $\ge$ torus for negative sources) Step~3. The same applies to a box of King's type (a killed face layer) in \cref{prop:s11:W1prime}. Nothing beyond \cref{lem:s04:G} is needed.
\end{remark}

\begin{remark}[the Wick sector]\label{rem:s04:wick}
With scalar sources $g\ge0$, \cref{cor:s04:G2}(i) holds as stated. The analogues of \cref{cor:s04:G2}(ii)--(iii) on the Wick tube need the fixed-cube limits and the $(k,J)$-uniform bound \eqref{eq:s11:W} for both families. For King's family the bound is the chain of \cref{thm:s11:W}(i) run on $\NDK$, whose convention-sensitive steps are the cube-versus-torus comparison (\cref{cor:s11:W2prime}, via \cref{rem:s04:F}) and the cell bound (\cref{lem:s04:E}); the assembly of that chain and the existence of the K-family tube limits are carried out in \cref{thm:s13:TWK}, after which the tube statements hold on King's convention (\cref{rem:s11:Kconvention}, \cref{prop:s13:WK}).
\end{remark}


\part*{The existence part}
% k01 King's existence theorem on the cube
% local macros for the named statements of this section and the next (to be moved to preamble.tex)
\providecommand{\moving}[1]{\textcolor{orange!80!black}{[#1]}}
\providecommand{\Hrel}{\ensuremath{(\mathrm{H0rel})_{C}}}
\providecommand{\Kpos}{\ensuremath{(\mathrm{K}\text{-}\mathrm{pos})_{N}}}
\providecommand{\BALint}{\ensuremath{(\mathrm{BAL}_{5})[\mathrm{int}]}}
\providecommand{\RENorb}{\ensuremath{(\mathrm{REN}_{\mathrm{orb}})}}
\providecommand{\CUTM}{\ensuremath{(\mathrm{CUT}_{M})}}
\providecommand{\Kh}{\ensuremath{(\mathrm{K}_{h})}}
\providecommand{\Uplus}{\ensuremath{(\mathrm{U}^{+})}}
\providecommand{\Bh}{\ensuremath{(\mathrm{B}^{h})}}
\providecommand{\Rzero}{\ensuremath{(\mathrm{R}_{0})}}
\providecommand{\Wsh}[1]{\ensuremath{(\mathrm{W}^{\sharp})[#1]}}
\providecommand{\UBm}{\ensuremath{(\mathrm{UB}^{-})}}
\providecommand{\NEGinf}{\ensuremath{(\mathrm{NEG}_{\infty})}}
\providecommand{\Zk}{\ensuremath{Z_{k}}}
\providecommand{\Ik}[1]{\ensuremath{I^{(#1)}_{k}}}
\providecommand{\Ihat}[1]{\ensuremath{\hat I^{(#1)}_{k}}}

\section{King's existence theorem on the cube}\label{sec:k01}

The existence hypothesis \EC\ of \cref{hyp:s03:EC}, and the form of \Wone\ that the Wick leg consumes, are proved in this section and the next from estimates on King's renormalisation-group integrals on the cube $\CK$. The argument is the one King gives for his Theorem~2.1 \cite[Sect.~3.1]{K2} and \cite[Sect.~3.2]{K1}: the generating functional is approximated, up to a small-field remainder, by an integral of the effective action after $j$ renormalisation steps, and two spacings whose effective actions live on the same unit lattice are compared. We write it as a Cauchy property of the \emph{normalised} functional, in a ratio form in which every field-independent constant of the effective action cancels; this removes from King's route the convergence of the vacuum energies, which his printed argument uses and which is not available on the cube.

The inputs are of two kinds. Estimates that are proved at length elsewhere, namely Ba{\l}aban's block-factor-$5$ analysis \cite{Ba1,Ba2,Ba3,Ba4} on tori and its transfer to King's cube, are stated here as propositions with the sentence ``proved in Appendix [to be added]''. Four statements remain open; they are named hypotheses (\BALint, \RENorb, \CUTM\ in this section, \Kh\ in the next), stated once with a precise formulation and a one-line description, and every proposition that rests on one of them says so. Everything else in the two sections is proved here.

\subsection{Renormalisation-group integrals on King's cube}\label{ssec:k01:setting}

Fix the couplings, a cube $\CJ$ and King's RG cube $\CK$ of \cref{def:s02:Kcube}. In this section and the next the index $k$ is the spacing index ($\eps_k=5^{-k}$, as throughout), and $j$ counts renormalisation steps, $0\le j\le k-J_*$ for a fixed $J_*$ (King's ``$J$'' of \cite[Thm.~3.1]{K1}); the block scale after $j$ steps is $\eta\defeq5^{j}\eps_k=5^{j-k}\le1$, and the unit lattice after $j$ steps is $\eta^{-1}(\CJ\cap\eta\Z^3)$. Since the side $2MJ$ of $\CJ$ is an integer, the faces of $\CJ$ lie on block planes at every scale.

\paragraph{Sources.} A real source is $s=(u,h)$ with $u=u_x\defeq\sum_ix_if_i$ a real combination of $2$-forms $f_i\in C^\infty_c(\R^3;\Lambda^2)$ sampled at the plaquette centres in the lattice's own frame, and $h\in C^\infty_c(\R^3)$ real; all supports lie at $\ell^\infty$-distance $>1$ from $\R^3\setminus\CJ$ (so $J\ge J_0$). Put $X(s)\defeq F(u)+\Wick{h}$, with the free Wick constant of \cref{def:s02:sources}; $X$ is real and linear in $s$. The unnormalised functional is
\[
\Zk(s)\ \defeq\ \int e^{-S}\,e^{X(s)}\ \prod dA\prod d^2\phi
\]
over the variables of $\CK$ at spacing $\eps_k$ (any fixed, $s$-independent normalisation of the measure; it cancels everywhere below), and $N_k(s)\defeq\Zk(s)/\Zk(0)=\avK{e^{X(s)}}$. At $h=0$ we write $a_k(x)\defeq N_k(u_x,0)$, which is the $a^K_{k,J}(x)$ of \cref{cor:s04:G2}.

\paragraph{King's flow on the cube.} The block-spin transformations of \cite[Sect.~2]{K1} are run on $\CK$ with block factor $5$: the photon field is averaged over the $5^{j}\eps_k$-blocks that tile $\XJ$, with the Neumann condition of \cref{def:s02:Kcube}; the scalar field is averaged with King's Dirichlet weight \cite[(3.8)--(3.9)]{K2} (in the corrected form in which the averaging blocks have their top layers on the Dirichlet planes) at the finest step and with the uniform weight above it. The Dirichlet planes of $\CK$ are $2MJ$ apart, a multiple of every block side, which is what makes this flow possible on $\CK$ and not on the comparison cube (\cref{rem:k01:cube}). Let $S^{(j)}_k(A_j,\phi_j;s)$ be King's effective action after $j$ steps \cite[(3.14)]{K1}, \cite[(3.12)]{K2}, with the source vertices \cite[(3.20)--(3.21)]{K1}; let $\chi_j$ be King's small-field indicator, \cite[(3.2)]{K1} times the face factor \cite[(5.41)]{K2}, a function of $(A_j,\phi_j)$ alone; and
\[
\Ik{j}(s)\ \defeq\ \int dA_j\,d\phi_j\ \chi_j\ e^{-S^{(j)}_k(A_j,\phi_j;s)} .
\]
The \emph{vacuum constant} $v_{k,j}\defeq S^{(j)}_k(0,0;0)$ is finite; it contains King's $E_0$, $E_1$, the normalisation constants $\ln Z_j$, $\ln Z_j(0)$ of the block-spin kernels and every field- and source-free graph. The \emph{normalised action} is $S^{(j)}_{N,k}(\cdot;s)\defeq S^{(j)}_k(\cdot;s)-v_{k,j}$ for every $s$ (the field-free, source-dependent part is kept), and $\Ihat{j}(s)\defeq\int\chi_je^{-S^{(j)}_{N,k}(\cdot;s)}=e^{v_{k,j}}\Ik{j}(s)$. The flow is a device of the proof: the state $\avK{\cdot}$ does not depend on it.

\paragraph{Diagonal pairs.} Fix $\alpha\in(0,\frac12)$ and put $j(k)\defeq\lfloor(1-\alpha)k\rfloor$, $\eta_k\defeq5^{j(k)-k}$; then $\eta_k\le5^{-\alpha k}$ and $k-j(k)\le\alpha k+1$. The \emph{diagonal pair} at $k$ consists of the $k$-model after $j(k)$ steps and the $(k+1)$-model after $j(k)+1$ steps. Both have the block scale $\eta_k$ ($5^{j+1}\eps_{k+1}=5^{j}\eps_k$), hence the same unit lattice and the same indicator $\chi_{j(k)}$; this is the situation of \cite[p.~656]{K1} (``the fields $A_k,\phi_k$ have the same bounds in each integral''). No relation between $j(k+1)$ and $j(k)+1$ is used anywhere.

\paragraph{Ba{\l}aban's decomposition.} In Ba{\l}aban's block-spin representation \cite{Ba1,Ba2}, $\Zk(s)=\int d\mu_{k,s}$ with $d\mu_{k,s}=e^{X(s)}d\mu_{k,0}$ a positive measure on the fields and on the block variables of the scales $\le j$, and $P=P_j\in[0,1]$ is the product of all small-field factors of those scales, a function of the block and fluctuation variables alone. The identity $\Zk(s)=\int P\,d\mu_{k,s}+\int(1-P)\,d\mu_{k,s}$ is exact; the first term is the \emph{small-field part}, the second the \emph{large-field remainder}.

\subsection{The residual open statements}\label{ssec:k01:open}

\begin{hypothesis}[\BALint: Ba{\l}aban's interaction-term bound at a nonzero background]\label{hyp:k01:BALint}
In the setting of \cite[Sect.~2 and Prop.~2.6]{Ba2} at block factor $L=5$, in unit-lattice form at a scale $k$: let $\Omega_A\subset\Omega_C$ be two of the propagator domains of the $k$-th step, let $\tilde B$ be a background gauge field satisfying the regularity bounds imposed there, let $G_k(\Omega,\tilde B)$ be the scalar propagator of the step on $\Omega$ in the background $\tilde B$, and let $E_A=E_C+\delta E$ be the corresponding external fields. For every $\kappa>0$ there are $t_0(\kappa)>0$ and $C_\kappa<\infty$ such that for $t\le t_0(\kappa)$ the terms of the interaction part of the effective action which contain at least one difference $G_k(\Omega_A,\tilde B)-G_k(\Omega_C,\tilde B)$ or $\delta E$, summed over their localisations, are bounded by $C_\kappa\,t^{\kappa}\,|\Lambda^{(k)}_7|$, with $C_\kappa$ and $t_0(\kappa)$ independent of the domains, of the background $\tilde B$ and of the scale.
\end{hypothesis}
This is the interaction-term half of Ba{\l}aban's domain-change estimate \cite[Prop.~2.6, (5.8)]{Ba2} at a nonzero background. Its printed proof applies the renormalisation-cancellation bound to the forest families produced by telescoping, which carry their counterterm partners but not the complete renormalised classes, and at $\tilde B\ne0$ the zero-background regrouping into complete classes does not transfer; a proof through classes complete at every background is written but not independently checked. Every use of the upper stability bound \cite[(1.3)]{Ba2} and of \cite[Thm.~3.1, (3.4)]{K1} at block factor $5$ carries this statement.

\begin{hypothesis}[\RENorb: orbit regrouping of cut renormalised families]\label{hyp:k01:RENorb}
Let $R$ be a finite union of the renormalised classes of \cite[Sects.~2--3]{Ba3}, with every line object decomposed as $o=o'-\Delta o$, and let $S$ be a set of lines. Call $(G,S)\sim(G',S')$ if $G'$ arises from $G$ by finitely many of the three completions of \cite[pp.~420, 429--432]{Ba3} (the counterterm partner; the permutation of the external scalar legs of a divergent $\phi\phi A$ subgraph; the two-vector-leg classes) at divergent subgraphs not meeting $S$, with $S'$ the image of $S$ under the completion. Then the sum, over one $\sim$-class, of the $S$-terms of the telescoping expansion of $\prod(o'-\Delta o)$ over $R$ is a weighted union of renormalised classes of the cut graphs $G\setminus S$, in which every divergent subgraph of a cut graph has the degree and the external-leg types of the corresponding subgraph of $G$; and the sum over all $\sim$-classes is the full telescoping sum.
\end{hypothesis}
This is Ba{\l}aban's regrouping step in the proof of \cite[Prop.~1, p.~432]{Ba3} (``if we remove a line in any graph of the class, then for the obtained graph we can find in this class all graphs necessary to form a new renormalized class''), made explicit for decomposed line objects. It is needed wherever the renormalisation cancellation is applied to a family some of whose lines have been cut, in particular in the comparison of two propagator domains and in the telescoping of the cube's large-field remainders; a proof is written but not independently checked.

\begin{hypothesis}[\CUTM: the mass-mismatch part of the counterterm-cut bound]\label{hyp:k01:CUTM}
In the interpolated model of the companion $O(3)$ paper \cite{RotO3} (interpolation parameter $\kappa$, mass-mismatch vertex $\mathcal M=-\tfrac12l_k^2\sum_x\eta^3\,\Delta\delta(x)\,|\phi'+\phi^{(k)}_\kappa|^2$): for a divergent two-photon-leg class $H$ with photon line slots $s_1,s_2$ and window scale $J=j_1$, containing $r\ge1$ vertices $\mathcal M$ on its window lines, the window-vertex kernel $\Pi^{(0)}_{H,[J]}$ obeys the bound of the counterterm-cut lemma for $\mathcal M$-free classes with the additional factor $(C(\lambda+e^2)l_J)^r$, uniformly in $\kappa$.
\end{hypothesis}
This statement is listed with the open points of the companion paper's chain. No step of this section, of the next, or of the appendix material consumes it: every class containing $\mathcal M$ is bounded there without a window identity. It is stated so that the list of named statements is complete.

\begin{remark}[the cube]\label{rem:k01:cube}
The propositions below are stated for King's RG cube $\CK$, on which King's Dirichlet weight is compatible with the image group of the two Dirichlet planes (they are $2MJ$ apart). On the comparison cube $\NDa$ the planes are $2MJ+\eps_k$ apart, no partition into consecutive blocks is invariant under the image group, and the flow does not exist; this is why \EC\ is posed on King's convention, and why the fixed-cube limits of the comparison family are reached only through \cref{cor:s04:G2}. The cube of the companion paper's Theorem~CA is $\CK$ itself (\cref{rem:s03:cube-identification} is thereby settled).
\end{remark}

\subsection{The inputs at zero scalar source}\label{ssec:k01:inputs}

\begin{proposition}[\Hrel: the small-field representation on King's cube, zero scalar source]\label{prop:k01:H0rel}
Assume \BALint\ and \RENorb. Fix $J\ge J_0$, a family $f_1,\dots,f_N$ and $x\in\R^N$, and let $s\in\{u_x,0\}$. There are $\eta_1>0$, $k_0$ and $C,c,\sigma>0$, depending on the couplings, on $J$, on $f$ and $x$, and on nothing else (not on $k$, $j$), such that for all $k\ge k_0$ and all $j$ with $J_*\le j\le k-J_*$ and $\eta=5^{j-k}\le\eta_1$:
\begin{equation}\label{eq:k01:H0rel}
e^{-C\eta^{\sigma}}\,\Ik{j}(s)\ \le\ \Zk(s)\ \le\ e^{C\eta^{\sigma}}\,\Ik{j}(s)+\rho(\eta)\,\Zk(0),\qquad \rho(\eta)\defeq e^{-c\,p(\eta)^2+C},
\end{equation}
where $p(\eta)=b_0(1+\ln\eta^{-1})^{p}$, $p\ge1$, is Ba{\l}aban's small-field threshold. Moreover, at $s=0$ the two terms of the upper line are the two parts of Ba{\l}aban's decomposition: $\int P\,d\mu_{k,0}\le e^{C\eta^\sigma}\Ik{j}(0)$ and $\int(1-P)\,d\mu_{k,0}\le\rho(\eta)\Zk(0)$ \textup{(\Rzero)}.
\end{proposition}
Proved in Appendix [to be added] \moving{proof appended from the companion work}. This is \cite[Thm.~3.1]{K1} on the cube, which \cite[Sect.~5.2]{K2} asserts (``the proof follows by modifying suitably the material in [Ba~1] and [Ba~2]''), in the form in which the remainder is measured against $\Zk(0)$ rather than against an absolute constant, and with constants uniform in the step $j$ (the form the diagonal pairs need, since $\eta_k\to0$). The lower line is the weaker of the two readings of the sign in \cite[(3.3)]{K1}. No lower bound on $\Zk(0)$ appears. Of the four open statements only \BALint\ (through Ba{\l}aban's torus bounds) and \RENorb\ (through the comparison of propagator domains and the telescoping of the large-field remainders) enter.

\begin{proposition}[\Kpos: one-step bound on the normalised action, zero scalar source]\label{prop:k01:Kpos}
Assume \BALint\ and \RENorb. With $J$, $f$, $x$ and $s\in\{u_x,0\}$ as above there are $\gamma\in(0,1)$, $\beta,\sigma>0$ and $C_3<\infty$, depending on the couplings, $J$, $f$, $x$ only, such that for all $k\ge k_0$, all $J_*\le j\le k-J_*$ with $\eta=5^{j-k}\le\eta_1$, and all $(A_j,\phi_j)$ in the support of $\chi_j$,
\begin{equation}\label{eq:k01:Kpos}
\big|S^{(j)}_{N,k}(A_j,\phi_j;s)-S^{(j+1)}_{N,k+1}(A_j,\phi_j;s)\big|\ \le\ C_3\big(5^{-\gamma j}\eta^{-\beta}+\eta^{\sigma}\big).
\end{equation}
\end{proposition}
Proved in Appendix [to be added] \moving{proof appended from the companion work}. This is \cite[Thm.~3.4]{K1} at $n=1$ on the cube, for the field- and source-dependent part of the effective action: the two members of a diagonal pair live on one unit lattice with one indicator $\chi_j$, and the statement compares them there. It is assembled from six pieces: the photon quadratic form (an even-restriction identity reduces the cube to the doubled torus, where \cite[Prop.~3.10]{K1} applies); the scalar quadratic form at zero background (the Dirichlet symbol) and with background (a two-leg tree graph); the background determinant (one-loop graphs with a remainder, \cite[(3.94)--(3.98)]{K1}, of which the middle line of (3.98) is false as printed and is replaced); the graphs with curvature vertices; the source-free field-dependent graphs away from the faces and near them (the photon face terms and the scalar boundary remainders of \cite[pp.~331--332]{K2}); and the block-vertex classes. The vacuum constant $v_{k,j}$ is not part of the statement.

\begin{lemma}[bounds for free]\label{lem:k01:free}
Put $\beta_f\defeq C_0\sum_i\fournorm{f_i}^2$ with $C_0$ as in \cref{lem:s02:E}. For all $k\ge0$, $J\ge J_0$ and $x\in\R^N$,
\[
e^{-\beta_f|x|^2}\ \le\ a_k(x)\ \le\ e^{\beta_f|x|^2}.
\]
\end{lemma}
\begin{proof}
By \cref{thm:s05:CK}(a), $a_k(x)\le\exp(\norm{u_x}^2_{2,k})$, and $\norm{u_x}_{2,k}\le\sum_i|x_i|\norm{f_i}_{2,k}\le|x|(\sum_i\norm{f_i}^2_{2,k})^{1/2}\le|x|\,\beta_f^{1/2}$ by \cref{lem:s02:E}(E3). The lower bound is $a_k(x)a_k(-x)\ge1$, \cref{lem:s02:alg}(e).
\end{proof}

\subsection{The Cauchy property in the spacing}\label{ssec:k01:CA}

\begin{theorem}[the Cauchy property, ratio form]\label{thm:k01:CA}
Fix $J\ge J_0$, $f$, $x\in\R^N$, and assume \eqref{eq:k01:H0rel} and \eqref{eq:k01:Kpos} for $s\in\{u_x,0\}$ with the constants of \cref{prop:k01:H0rel,prop:k01:Kpos}. Choose $\alpha\in(0,\min\{\frac12,\gamma/(\gamma+\beta)\})$ and the diagonal pairs $j(k)=\lfloor(1-\alpha)k\rfloor$, $\eta_k=5^{j(k)-k}$. Put
\[
b_k\defeq C\eta_k^{\sigma},\qquad \theta_k\defeq C_3\big(5^{-\gamma j(k)}\eta_k^{-\beta}+\eta_k^{\sigma}\big),\qquad \rho_k\defeq\rho(\eta_k),\qquad r_k(0)\defeq2\rho_k,\qquad r_k(u_x)\defeq2\rho_ke^{\beta_f|x|^2}.
\]
Then there is $k_1$ such that for all $k\ge k_1$
\begin{equation}\label{eq:k01:CA}
\big|\ln a_{k+1}(x)-\ln a_k(x)\big|\ \le\ 4b_k+2\theta_k+r_k(u_x)+r_k(0).
\end{equation}
The right side is summable in $k$, and its tail from $k$ on is at most $C'5^{-\kappa k}$ for some $\kappa>0$. Hence $a_k(x)$ converges, and its limit lies in $[e^{-\beta_f|x|^2},e^{\beta_f|x|^2}]$.
\end{theorem}

\begin{proof}
Fix $k\ge k_0$ so large that $j(k)\ge J_*$, $\eta_k\le\eta_1$ and $r_k(u_x)\le1$ (possible by Step~6), and put $j\defeq j(k)$, $b\defeq b_k$, $\theta\defeq\theta_k$. For $s\in\{u_x,0\}$ write $Z\defeq\Zk(s)$, $I\defeq\Ik{j}(s)$, $\hat I\defeq\Ihat{j}(s)$ for the $k$-model after $j$ steps, and $Z'\defeq Z_{k+1}(s)$, $I'\defeq I^{(j+1)}_{k+1}(s)$, $\hat I'\defeq\hat I^{(j+1)}_{k+1}(s)$ for the $(k+1)$-model after $j+1$ steps. Both members have the block scale $\eta_k$ and the indicator $\chi_j$, and \eqref{eq:k01:H0rel} applies to both, since $j+1\le(k+1)-J_*$; their remainders are $R\defeq\rho_k\Zk(0)$ and $R'\defeq\rho_kZ_{k+1}(0)$.

\emph{Step 1 (the remainders relative to $Z$).} For $s=0$, $R/Z=\rho_k=r_k(0)/2$ and $R'/Z'=\rho_k$. For $s=u_x$, $Z=a_k(x)\Zk(0)\ge e^{-\beta_f|x|^2}\Zk(0)$ by \cref{lem:k01:free}, so $R/Z\le\rho_ke^{\beta_f|x|^2}=r_k(u_x)/2$, and likewise $R'/Z'\le r_k(u_x)/2$. In both cases $R/Z,\ R'/Z'\le r_k(s)/2\le\frac12$. No lower bound on $\Zk(0)$ and no bound on $\Zk(u_x)$ beyond \cref{lem:k01:free} is used.

\emph{Step 2 (one model).} From $e^{-b}I\le Z\le e^{b}I+R$: $I\le e^{b}Z$ and $I\ge e^{-b}(Z-R)=e^{-b}Z(1-R/Z)$. So $I\in[e^{-b}Z(1-R/Z),\,e^{b}Z]$, and in the same way $I'\in[e^{-b}Z'(1-R'/Z'),\,e^{b}Z']$. In particular $I,I'>0$.

\emph{Step 3 (across the pair).} On the support of $\chi_j$, \eqref{eq:k01:Kpos} gives $|S^{(j)}_{N,k}(\cdot;s)-S^{(j+1)}_{N,k+1}(\cdot;s)|\le\theta$. The two integrals $\hat I$, $\hat I'$ run over the same variables with the same $\chi_j$ and have non-negative integrands, so $\hat I'/\hat I\in[e^{-\theta},e^{\theta}]$. Since $I=e^{-v_{k,j}}\hat I$ and $I'=e^{-v_{k+1,j+1}}\hat I'$,
\[
I'/I\ \in\ e^{\Delta}\,[e^{-\theta},e^{\theta}],\qquad \Delta\defeq v_{k,j}-v_{k+1,j+1},
\]
and $\Delta$ does not depend on $s$: the vacuum constant is the same number in the numerator and in the denominator of $a_k$.

\emph{Step 4 (ratio of partition functions).} By Steps 2 and 3,
\[
\frac{Z'}{Z}=\frac{Z'}{I'}\cdot\frac{I'}{I}\cdot\frac{I}{Z}\ \in\ e^{\Delta}\Big[e^{-2b-\theta}\big(1-\tfrac RZ\big),\ e^{2b+\theta}\big(1-\tfrac{R'}{Z'}\big)^{-1}\Big].
\]
With $|\ln(1-t)|\le2t$ for $0\le t\le\frac12$ and Step~1, $|\ln(1-R/Z)|,\,|\ln(1-R'/Z')|\le r_k(s)$, hence
\[
\big|\ln(Z'/Z)-\Delta\big|\ \le\ 2b+\theta+r_k(s).
\]

\emph{Step 5 (ratio of ratios).} $\ln\dfrac{a_{k+1}(x)}{a_k(x)}=\Big[\ln\dfrac{Z_{k+1}(u_x)}{\Zk(u_x)}-\Delta\Big]-\Big[\ln\dfrac{Z_{k+1}(0)}{\Zk(0)}-\Delta\Big]$, and Step~4 at $s=u_x$ and at $s=0$ gives \eqref{eq:k01:CA}. (If the constants of \eqref{eq:k01:H0rel}, \eqref{eq:k01:Kpos} differ between the two sources, take the larger ones.)

\emph{Step 6 (summability).} Since $k-j(k)\ge\alpha k$, $\eta_k\le5^{-\alpha k}$, so $b_k$ and $C_3\eta_k^\sigma$ decay geometrically. Since $j(k)\ge(1-\alpha)k-1$ and $k-j(k)\le\alpha k+1$,
\[
5^{-\gamma j(k)}\eta_k^{-\beta}\ \le\ 5^{\gamma+\beta}\,5^{-(\gamma(1-\alpha)-\beta\alpha)k},
\]
which decays geometrically because $\alpha<\gamma/(\gamma+\beta)$. Finally $p(\eta_k)^2\ge b_0^2(1+\alpha k\ln5)^{2}\ge c_1k^2$ for $p\ge1$, so $\rho_k\le e^{-cc_1k^2+C}$; thus $r_k(u_x)\le1$ for large $k$, $\sum_kr_k(s)<\infty$, and the tail of $\sum_k\rho_k$ is eventually below any geometric tail. The series in \eqref{eq:k01:CA} converges with a geometric tail, so $(\ln a_k(x))_k$ is Cauchy, $a_k(x)$ converges, and the limit lies in the closed interval of \cref{lem:k01:free}.
\end{proof}

\begin{remark}[what the argument does not use]\label{rem:k01:notused}
(i) No convergence, and no bound, of the vacuum constants $v_{k,j}$ (King's $E_0$, $E_1$, the kernel normalisations, the face vacuum energies) enters: they cancel in Step~3. King's printed route \cite[p.~329]{K2} goes through the convergence of $\Zk$ itself, which requires exactly those constants to converge. (ii) No upper stability bound, no lower bound on $\Zk(0)$ and no bound on $\Zk(u_x)$ beyond \cref{lem:k01:free} is used; the remainder of \eqref{eq:k01:H0rel} is compared with $\Zk(0)$ only. (iii) No reflection positivity, no Ward identity, and no subsequence. (iv) The all-$n$ form of \cite[(3.13)]{K1} (comparing the $k$-model with the $(k+n)$-model at a fixed stopping scale $\eta$ for every $n$) also gives the Cauchy property, provided the constants of \eqref{eq:k01:H0rel} and \eqref{eq:k01:Kpos} are taken at fixed $\eta$ with no growth in $k$: then $|\ln a_{k+n}-\ln a_k|\le4b(\eta)+2\theta_k(\eta)+4\rho(\eta)(e^{\beta_f|x|^2}+1)$ for every $n$, and one lets $k\to\infty$ and then $\eta\to0$. The diagonal form above uses $n=1$ only, at the price of constants uniform in $j$; the two forms need different readings of the inputs, and neither implies the other.
\end{remark}

\begin{theorem}[King's existence theorem on the cube: the curvature field]\label{thm:k01:EC}
Assume \BALint\ and \RENorb. Then for every finite family $f_1,\dots,f_N$ of real $2$-forms in $C^\infty_c(\R^3;\Lambda^2)$, every $J\ge J_0(f)$ and every $x\in\R^N$ the limit
\[
\ZKJ(x)=\lim_{k\to\infty}\avK{e^{\sum_ix_iF(f_i)}}
\]
exists and lies in $[e^{-\beta_f|x|^2},e^{\beta_f|x|^2}]\subset(0,\infty)$, with $\beta_f$ as in \cref{lem:k01:free}; the convergence is geometric. That is, \EC\ holds. Consequently $\ZKJ$ and $\ZKinf$ exist locally uniformly on $\C^N$ and are entire, with the bounds and the monotonicity of \cref{thm:s05:CK}(b), and every statement of this paper made under \EC\ holds under \BALint\ and \RENorb.
\end{theorem}

\begin{proof}
\Cref{prop:k01:H0rel,prop:k01:Kpos} supply the hypotheses of \cref{thm:k01:CA} at every $x$; the theorem gives the limit and its bounds; \cref{thm:s05:CK}(b) (that is, \cref{cor:s04:G2}(ii)) gives the rest. The constants $k_0$, $\eta_1$ and the rates may depend on $x$ and on $f$; existence needs no uniformity in $x$.
\end{proof}

\begin{remark}[the comparison family]\label{rem:k01:compfamily}
\Cref{thm:k01:EC} says nothing about the fixed-cube limits of the comparison family $\NDa$ (\cref{rem:k01:cube}); for that family \cref{cor:s04:G2}(ii)($\gamma$) squeezes the limits inferior and superior at fixed $J$ between $\ZKJ$ and $Z^K_{J+1}$, and \cref{cor:s04:G2}(iii) identifies every subsequential limit, followed by $J\to\infty$, with $\ZKinf$. At fixed $J$ every inequality available for the comparison family bounds its oscillation in $k$ only by the difference of two distinct limits, which tends to zero only as $J\to\infty$; a Cauchy statement at fixed $J$ for that family would need the small-field analysis at its faces, where no flow of the kind used above exists.
\end{remark}


\part*{The axioms}
% s05 Regularity in the curvature sector
\section{Regularity in the curvature sector}

This section proves the regularity axiom $\Ezz$ for the curvature field, for the comparison convention $\NDa$ (\cref{def:s02:acube}) and for King's convention $\NDK$ (\cref{def:s02:Kcube}), on which King's printed state $\ZKinf$ lives. Gaussian domination (\cref{lem:s02:GD}) and a bound on the Neumann photon form of the cube (\cref{lem:s05:C1}) give a bound on every lattice generating functional, uniform in the spacing, the cube and the scalar sector, with absolute constants (\cref{rem:s05:absolute}); under the existence hypothesis it passes to the limits by Montel's theorem, and Cauchy's estimates turn it into $\Ezz$ with $s=4$. The section closes with the curvature-sector properties of the fixed-cube limits of the comparison family (locally uniform convergence, monotonicity in the cube, positivity, continuity and translations), proved with cone-direction correlation inequalities only (\cref{lem:s05:TWC}).

Throughout, $f_1,\dots,f_n$ are real $2$-forms in $C^\infty_c(\R^3;\Lambda^2)$ with supports at $\ell^\infty$-distance $>1$ from $\R^3\setminus C_{J_0}$, $J\ge J_0$, $z\in\C^n$, and $g(z)\defeq\sum_j(\Re z_j)f_j$, so that $\Re\sum_jz_jF(f_j)=F(g(z))$ (\cref{lem:s02:alg}(d)). For a box $X$, $\ip AB_k$ and $\dstar_kg$ are as in \eqref{eq:s02:riemann}: $\dstar_kg$ is the lattice $1$-form with $F_{\Lk}(g)=\ip{\dstar_kg}{A}_k$, defined whenever every plaquette at which $g\ne0$ has its four corners in $X$. The support condition implies this for $X=\XJ$: the plaquette centre is at $\ell^\infty$-distance $>1$ from $\R^3\setminus\CJ$ and the corners are within $\eps/2\le\frac12$ of it.

\subsection{The Neumann cube photon form}

\begin{lemma}[the Neumann cube photon form]\label{lem:s05:C1}
Let $X=\CJ\cap\Lk$, let $C$ be the photon covariance of the cube states on $X$ (the same for $\NDa_{k,J}$ and $\NDK_{k,J}$) with respect to $\ip\cdot\cdot_k$, let $g$ be a real $2$-form on the plaquettes $p=(x;\alpha<\beta)$ all four of whose corners lie in $X$, and $v\defeq\dstar_kg$. Then
\[
\ip v{Cv}_k\ \le\ \min\big\{\,2\norm{g}_{2,k}^2,\ \lmu^{-2}\norm{v}_{2,k}^2\,\big\}.
\]
\end{lemma}

\begin{proof}
The Gaussian density of the photon field is $e^{-Q(A)/2}$ with
\[
Q(A)=\ip A{C^{-1}A}_k=\sum_{\alpha,\beta}\ \sum_{x,\,x+\eps e_\beta\in X}\eps\,(\partial_\beta A_\alpha(x))^2+\lmu^2\norm A_{2,k}^2,\qquad \partial_\beta A_\alpha(x)\defeq A_\alpha(x+\eps e_\beta)-A_\alpha(x):
\]
the Neumann rule keeps exactly the gradient terms of pairs inside $X$. By Legendre duality, $\ip v{Cv}=\sup_A\big[2\ip vA-Q(A)\big]$, attained at $A=Cv$. \emph{Second bound.} $Q(A)\ge\lmu^2\norm A^2$, so $C\le\lmu^{-2}$ as forms. \emph{First bound.} Let $P_g$ be the set of plaquettes on which $g\ne0$. For $p=(x;\alpha<\beta)\in P_g$ both gradient terms $(\partial_\alpha A_\beta(x))^2$ (pair $(x,x+\eps e_\alpha)$, component $\beta$) and $(\partial_\beta A_\alpha(x))^2$ (pair $(x,x+\eps e_\beta)$, component $\alpha$) are present in $Q$, because the pairs lie in $X$. With $F(p)=\eps^{-1}(\partial_\alpha A_\beta(x)-\partial_\beta A_\alpha(x))$,
\[
\norm{dA}^2_{P_g}\defeq\sum_{p\in P_g}\eps^3F(p)^2=\sum_{p\in P_g}\eps\,(\partial_\alpha A_\beta(x)-\partial_\beta A_\alpha(x))^2\le2\sum_{p\in P_g}\eps\big[(\partial_\alpha A_\beta(x))^2+(\partial_\beta A_\alpha(x))^2\big]\le2Q(A),
\]
where the last step is a counting statement: the ordered triple (site $x$, component $\beta$, direction $\alpha\ne\beta$) determines the plaquette $(x;\{\alpha,\beta\})$, so distinct plaquettes use distinct gradient terms of $Q$ and no term is counted twice. Also $\ip vA_k=F_{\Lk}(g)=\sum_{p\in P_g}\eps^3g_pF(p)=:\ip g{dA}_{P_g}$. Hence
\[
\ip v{Cv}\le\sup_A\big[2\ip g{dA}_{P_g}-\tfrac12\norm{dA}^2_{P_g}\big]\le\sup_{B\in\R^{P_g}}\big[2\ip gB_{P_g}-\tfrac12\norm B^2_{P_g}\big]=2\norm g_{2,k}^2,
\]
the last supremum being attained at $B=2g$. The inequality $Q(A)-\frac12\norm{dA}^2_{P_g}\ge0$ used in the middle is in fact the identity
\[
Q(A)-\tfrac12\norm{dA}^2_{P_g}=\tfrac12\sum_{p\in P_g}\eps\,(\partial_\alpha A_\beta(x)+\partial_\beta A_\alpha(x))^2+\big(\text{gradient terms of $Q$ used by no $p\in P_g$}\big)+\lmu^2\norm A^2_{2,k}.\qedhere
\]
\end{proof}

\begin{remark}[on the constant $2$]\label{rem:s05:sharp}
On a torus the Hodge decomposition gives the constant $1$ in place of $2$ (\cite[lem:C-Fsector]{SymPhase}); on the Neumann cube no Hodge argument is available, and the constant $2$ appears to be nearly sharp: the exact quantity $\rho\defeq\sup_g\Var F(g)/\norm g^2_{2,k}$, which depends on $\eps$ and $\lmu$ only through $\lmu\eps$, exceeds $1.94$ on a cube of $28^3$ sites at $\lmu\eps=10^{-4}$ and increases with the cube at small $\lmu\eps$. Whether it tends to $2$ is not proved; nothing below depends on it. The second bound is the cube form of the two-point bound $0\le S_2(f)\le\lmu^{-2}\norm{\dstar f}^2_{L^2}$; it is the worse of the two at small $\lmu$.
\end{remark}

\subsection{The comparison convention}

\begin{lemma}[standalone local bound]\label{lem:s05:C2}
For either convention and every compact $K\subset\C^n$,
\[
\sup_{k\ge0}\ \sup_{J\ge J_0}\ \sup_{z\in K}\ \big|\langle e^{\sum_jz_jF_{\Lk}(f_j)}\rangle_{k,J}\big|\ \le\ \exp\Big(3C_0\sup_{z\in K}\fournorm{g(z)}^2\Big)<\infty,
\]
with $C_0$ the constant of \cref{lem:s02:E}\textup{(E1)}. No hypothesis and no correlation inequality is used.
\end{lemma}

\begin{proof}
$|\langle e^{\sum z_jF(f_j)}\rangle|\le\langle e^{F(g(z))}\rangle$. Both cube states are BFS measures whose Gaussian is the Neumann photon Gaussian of $\XJ$ with covariance $C$ (the matter on $\TJ$ does not touch the photon variables), so \cref{lem:s02:GD} with $s=1$ and $v=\dstar_kg(z)$ gives $\langle e^{F(g(z))}\rangle\le e^{\ip v{Cv}_k/2}\le e^{\norm{g(z)}^2_{2,k}}$ by \cref{lem:s05:C1}. The plaquette centres of each of the three orientations form a translate of $\eps_k\Z^3$, and $|g_{c(p)}(e_\alpha,e_\beta)|\le(1+|c(p)|^2)^{-2}\fournorm g$, so by (E1) each orientation contributes at most $C_0\fournorm g^2$ to $\norm g^2_{2,k}$ and $\norm{g(z)}^2_{2,k}\le3C_0\fournorm{g(z)}^2$ for every $k$ (this is the count behind (E3) of \cref{lem:s02:E}). Finally $\fournorm{g(z)}\le\sum_j|\Re z_j|\fournorm{f_j}$ is bounded on $K$.
\end{proof}

The existence theorem \cref{hyp:s03:EC} is posed on King's convention. For the comparison family the analogous clause is its convergence at fixed cubes (\cref{def:s03:conv}): for every such family, every $J\ge J_0$ and every real $z$, $\lim_k\ava{e^{\sum_jz_jF(f_j)}}$ exists in $(0,\infty)$. It is not needed for the main theorem (the two clauses do not imply one another by the inputs of this paper, \cref{cor:s04:G2}); \cref{thm:s05:C} is stated under it because the rotation part and the tube estimates are written for the comparison family, and \cref{thm:s05:CK} is its form for King's printed state, where the limits exist by \cref{hyp:s03:EC}.

\begin{theorem}[regularity, comparison convention]\label{thm:s05:C}
Let $f_1,\dots,f_n$, $J_0$, $g(z)$ be as above.
\begin{enumerate}[label=\textup{(\alph*)},leftmargin=*]
\item \textup{(Finite volume; no hypothesis.)} For all $k\ge0$, $J\ge J_0$, $\lmu>0$, $z\in\C^n$, and in the presence of any real scalar source $h\in C^\infty_c$,
\[
\Bigg|\frac{\ava{e^{\sum_jz_jF_{\Lk}(f_j)}e^{\Wick h}}}{\ava{e^{\Wick h}}}\Bigg|
\ \le\ \frac{\ava{e^{F_{\Lk}(g(z))}e^{\Wick h}}}{\ava{e^{\Wick h}}}
\ \le\ \exp\Big(\min\big\{\norm{g(z)}^2_{2,k},\ \tfrac12\lmu^{-2}\norm{\dstar_kg(z)}^2_{2,k}\big\}\Big).
\]
\item \textup{(Limits; if the comparison family is convergent at fixed cubes.)} $Z_J(\sum_jz_jf_j,0)\defeq\lim_k\ava{e^{\sum_jz_jF(f_j)}}$ exists locally uniformly on $\C^n$ and is entire; $Z_\infty(\sum_jz_jf_j,0)\defeq\lim_JZ_J(\sum_jz_jf_j,0)$ exists locally uniformly on $\C^n$ and is entire; and for $J\ge J_0$
\[
|Z_J(\textstyle\sum_jz_jf_j,0)|,\ |Z_\infty(\textstyle\sum_jz_jf_j,0)|\ \le\ \exp\Big(\min\big\{\norm{g(z)}_2^2,\ \tfrac12\lmu^{-2}\norm{\dstar g(z)}^2_{L^2}\big\}\Big).
\]
\item \textup{(Schwinger functions; if the comparison family is convergent at fixed cubes.)} $S_n(f_1,\dots,f_n)\defeq\partial_{z_1}\cdots\partial_{z_n}Z_\infty(\sum_jz_jf_j,0)\big|_{z=0}$ satisfies, for real $f_i$,
\[
|S_n(f_1,\dots,f_n)|\ \le\ (\sqrt2\,e)^n(n!)^{1/2}\prod_{i=1}^n\norm{f_i}_2\ \le\ (e\pi/2)^n(n!)^{1/2}\prod_{i=1}^n\fournorm{f_i}.
\]
$S_n$ extends to a continuous symmetric multilinear form on $\cS(\R^3;\Lambda^2_\C)^n$ and defines a tempered distribution; for complex test functions the bound carries the factor $2^n$. This is the axiom $\Ezz$ of \cite{OS75} with $s=4$ and $\sigma_n\defeq(e\pi)^n(n!)^{1/2}\le\alpha_\beta(n!)^\beta$ for every $\beta>\frac12$ (and for no $\alpha$ at $\beta=\frac12$).
\end{enumerate}
\end{theorem}

\begin{proof}
(a) The first inequality is $\Re\sum_jz_jF(f_j)=F(g(z))$. The measure $e^{\Wick h}d\langle\cdot\rangle^{(a)}_{k,J}$, normalised, is a BFS measure whose single-site weights carry the extra factor $e^{\eps^3h(x)r_x^2}$ and whose photon covariance is the same $C$; \cref{lem:s02:GD} with $s=1$ and $v=\dstar_kg(z)$ bounds the middle expression by $e^{\ip v{Cv}_k/2}$, and \cref{lem:s05:C1} gives both terms of the minimum.

(b) \emph{Step 1: $k\to\infty$ at fixed $J\ge J_0$.} The functions $\Phi_k(z)\defeq\ava{e^{\sum_jz_jF(f_j)}}$ are entire on $\C^n$: finite-dimensional integrals with Gaussian decay in $A$ ($\lmu>0$) and $e^{-\lambda\eps^3r^4}$ decay in $r$, converging absolutely and locally uniformly in $z$. By \cref{lem:s05:C2} they are locally uniformly bounded; by the convergence at fixed cubes they converge at every real point. $\R^n$ is a set of uniqueness for entire functions on $\C^n$ (a holomorphic function vanishing on $\R^n$ has all its Taylor coefficients at $0$ equal to $0$). By Montel's theorem every subsequence has a locally uniformly convergent sub-subsequence; all subsequential limits are entire and agree on $\R^n$, hence everywhere \cite{Hor90}; so $\Phi_k\to Z_J(\cdot,0)$ locally uniformly on $\C^n$ and $Z_J(\cdot,0)$ is entire. The bound of (a) at $h=0$ passes to the limit: $|Z_J(\sum z_jf_j,0)|\le\lim_ke^{\norm{g(z)}^2_{2,k}}=e^{\norm{g(z)}_2^2}$ by (E3) of \cref{lem:s02:E}, and likewise $|Z_J(\sum z_jf_j,0)|\le\lim_k\exp(\frac12\lmu^{-2}\norm{\dstar_kg(z)}^2_{2,k})=\exp(\frac12\lmu^{-2}\norm{\dstar g(z)}^2_{L^2})$ by \eqref{eq:s02:riemann}. Both bounds are free of $J$.

\emph{Step 2: monotonicity in $J$ on the imaginary axes.} For $t\in\R^n$ and $f\defeq\sum_jt_jf_j$, $\ava{e^{\mathrm iF(f)}}$ is real (the state is invariant under $(A,\vartheta)\mapsto(-A,-\vartheta)$, which maps $F(f)$ to $-F(f)$) and of modulus $\le1$, and $J\mapsto\ava{e^{\mathrm iF(f)}}$ is nondecreasing on $J\ge J_0$ at fixed $k$: this is the first inequality of \cref{lem:s02:cut} with $Y=C_{J+1}\cap\Lk$, $X=\CJ\cap\Lk$ and the class-$\cI$ observable $\Theta(\mathrm if,0)$, whose plaquettes lie in $X$ by the support condition (its proof uses \cref{lem:s02:BFS}(a),(b) only, and the site-independence of the counterterm). Letting $k\to\infty$ (Step 1), $J\mapsto Z_J(\mathrm if,0)$ is nondecreasing and bounded by $1$, hence convergent, for every $t\in\R^n$.

\emph{Step 3: $J\to\infty$.} $\{Z_J(\cdot,0)\}_{J\ge J_0}$ is locally uniformly bounded on $\C^n$ by the $J$-free bound of Step 1, and converges on $(\mathrm i\R)^n$ (Step 2), a totally real subspace of maximal dimension, hence a set of uniqueness for entire functions (at a point $\mathrm it_0$ the restriction to $\mathrm it_0+\mathrm i\R^n$ determines the Taylor series). Montel's theorem and the identity theorem as in Step 1 give $Z_J(\cdot,0)\to Z_\infty(\cdot,0)$ locally uniformly on $\C^n$, $Z_\infty(\cdot,0)$ entire, and both bounds of Step 1 pass to the limit.

(c) By multilinearity normalise $\norm{f_i}_2=1$. On the polydisc $|z_i|\le r$, $\norm{g(z)}_2\le\sum_i|\Re z_i|\le nr$, so $|Z_\infty(\sum z_jf_j,0)|\le e^{n^2r^2}$ by (b). Cauchy's estimate for the mixed first derivative in the $n$ distinct variables gives $|S_n|\le r^{-n}\sup_{|z_i|=r}|Z_\infty|\le r^{-n}e^{n^2r^2}$; at $r^2=1/(2n)$ this is $(2en)^{n/2}=(2e)^{n/2}n^{n/2}\le(\sqrt2\,e)^n(n!)^{1/2}$ by (E4) of \cref{lem:s02:E}. Then (E2), $\norm{f_i}_2\le\frac\pi{2\sqrt2}\fournorm{f_i}$, gives the second bound. $S_n$ is symmetric and multilinear (derivatives of an entire function of a linear combination) and continuous for the $L^2$ norms, so it extends by density from $C^\infty_c$ to $\cS$; for complex $f_i=f_i'+\mathrm if_i''$ with real parts, multilinearity writes $S_n(f_1,\dots,f_n)$ as a sum of $2^n$ terms $\mathrm i^mS_n(f^\sharp_1,\dots,f^\sharp_n)$, each bounded by $(e\pi/2)^n(n!)^{1/2}\prod\fournorm{f^\sharp_i}\le(e\pi/2)^n(n!)^{1/2}\prod\fournorm{f_i}$, since $\fournorm{\Re f},\fournorm{\Im f}\le\fournorm f$; the kernel theorem gives the tempered distribution. The axiom $\Ezz$ of \cite[p.~287]{OS75} asks for $|S_n(f_1,\dots,f_n)|\le\sigma_n\prod_i|f_i|_s$ on $\cS$ with $\sigma_n\le\alpha(n!)^\beta$; here $s=4$ ($\fournorm\cdot$ is the norm that \cite{SymPhase} identifies with that of \cite{OS75} at $s=4$), and $\sigma_n=(e\pi)^n(n!)^{1/2}\le\alpha_\beta(n!)^\beta$ exactly when $(e\pi)^n/(n!)^{\beta-1/2}$ is bounded, that is for every $\beta>\frac12$.
\end{proof}

\begin{remark}[the constants are absolute]\label{rem:s05:absolute}
The bound $e^{\norm{g}_2^2}$ of \cref{thm:s05:C} and the constants $(\sqrt2\,e)^n(n!)^{1/2}$, $(e\pi/2)^n(n!)^{1/2}$, $\sigma_n$ of $\Ezz$ depend on none of $e$, $\lambda$, $m$, $\lmu$ and $M$: the only inputs are Gaussian domination, the Neumann form's constant $2$ and the elementary facts (E1)--(E4). The second bound of (a) and (b) is the one that carries $\lmu^{-2}$. The proof of (b) uses, besides the existence clause, only \cite[Thm.~6.2(1),(2)]{BFS1} through \cref{lem:s02:cut}, Gaussian domination and \cref{lem:s05:C1}; it uses neither \cite[Thm.~6.2(3),(4)]{BFS1} (\cref{lem:s02:BFS}(d)) nor any local bound resting on that reversed inequality.
\end{remark}

\subsection{King's convention}

For the K family write $a^K_{k,J}(z)\defeq\avK{e^{\sum_jz_jF(f_j)}}$ as in \cref{cor:s04:G2}, and $\ZKJ$, $\ZKinf$ for its limits.

\begin{theorem}[regularity of King's printed state]\label{thm:s05:CK}
Let $f=(f_1,\dots,f_n)$, $J_0$, $g(z)$ be as above.
\begin{enumerate}[label=\textup{(\alph*)},leftmargin=*]
\item \textup{(No hypothesis.)} For all $k\ge0$, $J\ge J_0$, $z\in\C^n$,
\[
|a^K_{k,J}(z)|\ \le\ \avK{e^{F(g(z))}}\ \le\ \exp\Big(\min\big\{\norm{g(z)}^2_{2,k},\ \tfrac12\lmu^{-2}\norm{\dstar_kg(z)}^2_{2,k}\big\}\Big),
\]
and for every real $2$-form $u$ whose plaquettes have their corners in $\XJ$: $\avK{F(u)}=0$ and $\avK{F(u)^2}\le\ip v{Cv}_k\le2\norm u^2_{2,k}$, $v=\dstar_ku$.
\item \textup{(Limits.)} $\ZKJ(z)\defeq\lim_ka^K_{k,J}(z)$ exists locally uniformly on $\C^n$ and is entire; $\ZKinf(z)\defeq\lim_J\ZKJ(z)$ exists locally uniformly on $\C^n$ and is entire; $|\ZKJ(z)|,|\ZKinf(z)|\le e^{\norm{g(z)}_2^2}$ and also $\le\exp(\frac12\lmu^{-2}\norm{\dstar g(z)}^2_{L^2})$; $J\mapsto\ZKJ(\mathrm it)$ is nondecreasing on $J\ge J_0$ with values in $[0,1]$ for $t\in\R^n$; and $\ZKinf(x)\ge1$ at every real $x\in\R^n$.
\item \textup{(Schwinger functions.)} The Schwinger functions $S^K_n(f_1,\dots,f_n)\defeq\partial_{z_1}\cdots\partial_{z_n}\ZKinf(0)$ satisfy $|S^K_n|\le(\sqrt2\,e)^n(n!)^{1/2}\prod_i\norm{f_i}_2\le(e\pi/2)^n(n!)^{1/2}\prod_i\fournorm{f_i}$ for real $f_i$; the lattice Schwinger functions $S^{K,(k,J)}_n\defeq\partial_{z_1}\cdots\partial_{z_n}a^K_{k,J}(0)$ satisfy the same bound at every $(k,J)$ with $\norm{f_i}_{2,k}$ in place of $\norm{f_i}_2$, and $S^{K,(k,J)}_n\to S^K_n$ as $k\to\infty$ and then $J\to\infty$. Each $S^K_n$ extends to a continuous symmetric multilinear form on $\cS(\R^3;\Lambda^2_\C)^n$ with the same bound, carrying the factor $2^n$ for complex test functions: $\Ezz$ holds with $s=4$ and $\sigma_n=(e\pi)^n(n!)^{1/2}\le\alpha_\beta(n!)^\beta$ for every $\beta>\frac12$.
\end{enumerate}
\end{theorem}

\begin{proof}
(a) The two inequalities are the bound of \cref{lem:s05:C2} for the K family, with \cref{lem:s05:C1} kept in both forms: $\Re\sum z_jF(f_j)=F(g(z))$; $\langle\cdot\rangle^K_{k,J}$ is a BFS measure (the killed sites carry point masses, which \cref{lem:s02:GD} allows: at a killed site $I_q(0)=\delta_{q0}$ and the currents vanish on its bonds), and its Gaussian is the Neumann photon Gaussian of $\XJ$, the same as the comparison cube's, so \cref{lem:s02:GD} with $s=1$, $v=\dstar_kg(z)$ and \cref{lem:s05:C1} give $\avK{e^{F(g(z))}}\le e^{\ip v{Cv}_k/2}\le\exp(\frac12\min\{2\norm g^2_{2,k},\lmu^{-2}\norm v^2_{2,k}\})$. \emph{Second-moment clause.} $\avK{F(u)}=0$ because the state is invariant under $(A,\vartheta)\mapsto(-A,-\vartheta)$ and $F(u)$ is odd. For real $s$, $e^x+e^{-x}\ge2+x^2$ pointwise and $e^{\pm sF(u)}$ are integrable, so $2+s^2\avK{F(u)^2}\le\avK{e^{sF(u)}}+\avK{e^{-sF(u)}}\le2e^{s^2\ip v{Cv}_k/2}$ (\cref{lem:s02:GD} at $\pm s$; $\pm v$ give the same quadratic form). Hence $\avK{F(u)^2}\le2(e^{s^2\ip v{Cv}_k/2}-1)/s^2\to\ip v{Cv}_k$ as $s\downarrow0$, and \cref{lem:s05:C1} applies since the corners lie in $\XJ$.

(b) The two limits, their entireness and the $J$-free bound $e^{\norm{g(z)}_2^2}$ are \cref{cor:s04:G2}(ii)($\alpha$),($\beta$): the proof of \cref{thm:s05:C}(b) run on the K family, with the local bound of (a) (uniform in $k$, $J$ by $\norm g^2_{2,k}\le3C_0\fournorm g^2$), \cref{hyp:s03:EC} for the convergence at real points, Montel's theorem with the uniqueness set $\R^n$, and then, for $J\to\infty$, the monotonicity of \cref{cor:s04:G1} on the cone in place of \cref{lem:s02:cut} and the uniqueness set $(\mathrm i\R)^n$. Three additions. (1) \emph{Values in $[0,1]$.} $a^K_{k,J}(\mathrm it)$ is real by the symmetry used in (a) and of modulus $\le1$; it is $\ge0$ by the current representation of \cref{lem:s02:GD}: with $\Psi(A)=\sum_qw(q)\cos\ip{\kappa(q)}A$, $w(q)\ge0$, and the Gaussian identity at $a=\mathrm iv\pm\mathrm i\kappa$,
\[
\avK{e^{\mathrm i\ip vA}}=e^{-\ip v{Cv}_k/2}\,\E^\sharp\cosh\ip{Cv}{\kappa}\ \ge\ 0,\qquad v=\dstar_k\textstyle\sum_jt_jf_j,
\]
the killed sites being covered as in (a). Monotonicity in $J$ is \cref{cor:s04:G1} at every $k$ (condition \eqref{eq:s04:S} holds for $X_{J_0}$ by the support condition), passed to the limit $k\to\infty$. (2) \emph{The $\lmu^{-2}$ bound.} $\norm{\dstar_kg}_{2,k}\to\norm{\dstar g}_{L^2}$ by \eqref{eq:s02:riemann}, so the second bound of (a) passes to $\ZKJ$ and then to $\ZKinf$ exactly as the first. (3) \emph{$\ZKinf(x)\ge1$.} For real $x$, $a^K_{k,J}(x)=\avK{e^{F(g(x))}}\ge e^{\avK{F(g(x))}}=1$ by Jensen's inequality and (a); the inequality is closed under both limits.

(c) Let $\norm{f_j}_{2,k}>0$ for all $j$ (otherwise $F_{\Lk}(f_j)=0$ and $S^{K,(k,J)}_n=0$). Cauchy's estimates on the polydisc $|z_j|\le\rho_j\defeq\rho/(n\norm{f_j}_{2,k})$ with the bound (a): there $\norm{g(z)}_{2,k}\le\sum_j|z_j|\norm{f_j}_{2,k}\le\rho$, so $|a^K_{k,J}|\le e^{\rho^2}$ and $|S^{K,(k,J)}_n|\le e^{\rho^2}\prod_j\rho_j^{-1}=e^{\rho^2}n^n\rho^{-n}\prod_j\norm{f_j}_{2,k}$; at $\rho^2=n/2$ the prefactor is $e^{n/2}2^{n/2}n^{n/2}\le e^{n/2}2^{n/2}e^{n/2}(n!)^{1/2}=(\sqrt2\,e)^n(n!)^{1/2}$ by (E4). The Taylor coefficients at $0$ converge under the locally uniform convergences of (b) (Weierstrass), and $\norm{f_j}_{2,k}\to\norm{f_j}_2$ by (E3), which gives the bound for $S^K_n$; (E2) gives the $\fournorm\cdot$ form. The $2^n$ factor for complex test functions, the extension to $\cS$ and the growth class $\sigma_n$ are as in \cref{thm:s05:C}(c).
\end{proof}

\subsection{The curvature sector of the comparison family's limits}

For the comparison family and an admissible scheme $a$, the general-source functionals $Z^a_J(\zeta)=\lim_kZ^a_{k,J}(\zeta)$, $\zeta\in\C^{2N}$, have, under \Wtwo\ (the optional extension), the following properties: (i) locally uniform convergence to an entire $Z^a_J$; (ii) on the cone, $J\mapsto Z^a_J(\mathrm if,g)$ nondecreasing and $J\mapsto Z^a_J(f,-g)$ nonincreasing; (iii) the limit $Z^a_\infty=\lim_JZ^a_J$, locally uniform and entire, positive at real sources, with $Z^a_J\le Z^a_\infty$ on the cone; (iv) $\Phi^a=\ln Z^a_\infty$ finite, convex and continuous on the Fr\'echet space $E_\rho=C^\infty_c(\bar B_\rho;\Lambda^2)\times C^\infty_c(\bar B_\rho)$; (v) translation invariance; (vi) the change of scheme. At $h=0$ these conclusions hold under the existence clause alone; the natural local bound, $|Z_J(\zeta)|\le Z_J(\Re\zeta)\le Z_{J_0}(\Re\zeta)$, would use the nonincreasing clause of (ii), that is \cref{lem:s02:BFS}(d). The following lemma states the curvature sector with that use removed; the Wick sector on the tube is \cref{lem:s11:TWtube}.

\begin{lemma}[the comparison family's limits, curvature sector, with cone-direction inequalities only]\label{lem:s05:TWC}
Assume that the comparison family is convergent at fixed cubes (\cref{def:s03:conv}). For every finite family $f_1,\dots,f_N$ of real $2$-forms in $C^\infty_c$ supported at $\ell^\infty$-distance $>1$ from $\R^3\setminus C_{J_0}$, with $Z_{k,J}(\sum_jz_jf_j,0)\defeq\ava{e^{\sum_jz_jF(f_j)}}$ and $g(z)=\sum_j(\Re z_j)f_j$:
\begin{enumerate}[label=\textup{(\roman*)},leftmargin=*]
\item $Z_{k,J}(\cdot,0)\to Z_J(\cdot,0)$ locally uniformly on $\C^N$ as $k\to\infty$, and $Z_J(\cdot,0)$ is entire, for every $J\ge J_0$;
\item $J\mapsto Z_J(\mathrm if,0)$ is nondecreasing on $J\ge J_0$ for every real $f=\sum_jt_jf_j$;
\item for all $k\ge0$, $J\ge J_0$ and $z\in\C^N$,
\[
|Z_{k,J}(\textstyle\sum_jz_jf_j,0)|\le\ava{e^{F(g(z))}}\le e^{\frac12\ip v{Cv}_k}\le\exp\Big(\tfrac12\lmu^{-2}\norm{\dstar_kg(z)}^2_{2,k}\Big)\le\exp\Big(\tfrac12\lmu^{-2}\sup_{k'}\norm{\dstar_{k'}g(z)}^2_{2,k'}\Big),\quad v=\dstar_kg(z),
\]
a bound uniform in $k$ and $J$ and locally uniform in $z$; consequently $Z_\infty(\cdot,0)\defeq\lim_JZ_J(\cdot,0)$ exists locally uniformly on $\C^N$, is entire, satisfies $Z_\infty(f,0)>0$ for real $f$, and $Z_J(\mathrm if,0)\le Z_\infty(\mathrm if,0)$ for real $f$;
\item for every $\rho>0$, $\Phi(f)\defeq\ln Z_\infty(f,0)$ is finite, convex and continuous on the Fr\'echet space $C^\infty_c(\bar B_\rho;\Lambda^2)$;
\item $Z_\infty(\tau_bf,0)=Z_\infty(f,0)$ for all $b\in\R^3$ and all complex coefficient vectors, $\tau_b$ acting on every $f_j$.
\end{enumerate}
Clauses \textup{(i)}, \textup{(iii)}--\textup{(v)} use, besides the convergence at fixed cubes, only \cite[Thm.~6.2(1),(2)]{BFS1} through \cref{lem:s02:cut}, \cref{lem:s02:GD}, \cref{lem:s02:alg}\textup{(d)}, \eqref{eq:s02:riemann}, \cref{lem:s02:baire}, Montel's theorem and the identity theorem. The nonincreasing clause \textup{(ii)} at $h=0$, $J\mapsto Z_J(f,0)$ nonincreasing for real $f$, holds as well and uses \cref{lem:s02:BFS}\textup{(d)}; it is not consumed by the curvature sector of this paper. Clause \textup{(vi)} is void at $h=0$, no scheme entering.
\end{lemma}

\begin{proof}
(i) is Step 1 of the proof of \cref{thm:s05:C}(b), which uses \cref{lem:s05:C2}, the convergence at fixed cubes and Montel's theorem. (ii) is its Step 2. (iii) The first inequality is \cref{lem:s02:alg}(d), the second is \cref{lem:s02:GD} with $s=1$ (the cube state is a BFS measure), the third is $C\le\lmu^{-2}$ as forms (the photon form is bounded below by $\lmu^2\norm A^2_{2,k}$, the second bound of \cref{lem:s05:C1}), and the supremum over $k'$ is finite by \eqref{eq:s02:riemann}. The bound passes to $Z_J$ by (i) and makes $\{Z_J(\cdot,0)\}_{J\ge J_0}$ locally bounded on $\C^N$; with the convergence on $(\mathrm i\R)^N$ from (ii) and the bound $|Z_J(\mathrm if,0)|\le1$, Montel's theorem and the identity theorem give the locally uniform convergence to an entire $Z_\infty(\cdot,0)$ as in Step 3 of \cref{thm:s05:C}(b). Positivity: $Z_{k,J}(f,0)\ge e^{\ava{F(f)}}=1$ for real $f$ by Jensen's inequality and $\ava{F(f)}=0$ (the symmetry $(A,\vartheta)\mapsto(-A,-\vartheta)$), closed under both limits. The last clause is (ii) in the limit.

(iv) $\Phi$ is finite by (iii). Convexity: for $f,f'\in C^\infty_c(\bar B_\rho;\Lambda^2)$ and $t\in[0,1]$ apply (i) and (iii) to the family $(f,f')$ at the real coefficients $(t,1-t)$; at fixed $k,J$, $\ln Z_{k,J}(\cdot,0)$ is convex in the coefficients by H\"older's inequality, and convexity survives the two pointwise limits. Borel measurability: fix $J$ with $MJ>\rho+1$, so that every $f\in C^\infty_c(\bar B_\rho;\Lambda^2)$ is supported at $\ell^\infty$-distance $>1$ from $\R^3\setminus\CJ$; at each $k$, $\ln Z_{k,J}(f,0)$ depends on $f$ only through the finitely many sampled values $f_{c(p)}(e_\alpha,e_\beta)$, continuous functionals of $f$ in the $C^\infty$ topology, and the finite-dimensional integral is continuous in them (dominated convergence); so $\Phi=\lim_J\lim_k$ of continuous functions is Borel (of Baire class $\le2$). $C^\infty_c(\bar B_\rho;\Lambda^2)$ is a closed subspace of $C^\infty(\bar B_{\rho+1};\Lambda^2)$, hence Fr\'echet. \cref{lem:s02:baire} gives the continuity.

(v) is \cref{prop:s06:translations}, whose proof uses the lattice sandwich between $C_{J-1}$ and $C_{J+1}$ (first inequality of \cref{lem:s02:cut}), $k\to\infty$ by (i) (the convergence at fixed cubes covers the translated family), $J\to\infty$ by (iii), the identity theorem with the uniqueness set $(\mathrm i\R)^N$, the continuity of (iv) and the density of $\bigcup_k5^{-k}\Z^3$.
\end{proof}

For King's convention the same clauses hold by \cref{hyp:s03:EC}: (i) and (iii) are \cref{thm:s05:CK}(b), with the local bound of \cref{thm:s05:CK}(a); (ii) is \cref{cor:s04:G1}; (iv) follows from \cref{thm:s05:CK}(b) by the argument above, word for word; (v) is \cref{cor:s06:TK}. In every clause the inequalities used are the cone directions of \cref{lem:s02:BFS}(a),(b) and of \cref{lem:s04:G}; the reversed inequalities of \cref{lem:s02:BFS}(d) enter the curvature sector of this paper nowhere.

% s06 Translations, the cubic group, and rotations
\section{Translations, the cubic group, and rotations}

This section proves the clause \OSo\ of King's printed state as far as the lattice allows: invariance under all translations and under the cubic group $\Oh$ about every point; full $O(3)$ invariance comes from the rotation theorem (\cref{hyp:s03:RC}) and is stated at the end. No cube is invariant under a non-trivial translation or a reflection; what replaces exactness is a sandwich between a smaller and a larger cube of the same family, in the direction of the cone inequalities of \cref{lem:s02:cut,lem:s04:G3}, followed by the limits $k\to\infty$, $J\to\infty$. Coordinate permutations alone are exact (\cref{lem:s06:PK}).

\emph{Conventions.} We use the notation of \cref{def:s02:sources,cor:s04:G2}: for a finite family $f=(f_1,\dots,f_N)$ of real $2$-forms in $C^\infty_c(\R^3;\Lambda^2)$ and $z\in\C^N$, $a^K_{k,J}(z)=\avK{e^{\sum_jz_jF(f_j)}}$ and $a^{(a)}_{k,J}(z)=\ava{e^{\sum_jz_jF(f_j)}}$ are the lattice functionals, $\ZKJ$, $\ZKinf$ and $Z_J$, $Z_\infty$ their limits where these exist, and $g(z)=\sum_j(\Re z_j)f_j$. Since $F_{\Lk}$ is linear in the source (extended $\C$-linearly), a lattice functional depends on a family only through its total source, and we use the same symbol for a coefficient vector and for its total source: $a^K_{k,J}(z)=Z^K_{k,J}(\sum_jz_jf_j)$ with $Z^K_{k,J}(u)\defeq\avK{e^{F(u)}}$ for a complex-valued $2$-form $u$; in particular, for a real $2$-form $u$, $Z^K_{k,J}(\mathrm iu)$ is the one-member functional at the coefficient $\mathrm i$, and the same convention is carried to $\ZKJ$, $\ZKinf$ (so $\ZKinf(\mathrm iu_t)=Z^{K,f}_\infty(\mathrm it)$ for $u_t=\sum_jt_jf_j$). When several families are in play we write $a^{K,f}_{k,J}$, $Z^{K,f}_J$, $Z^{K,f}_\infty$. $J_0(f)$ is the least $J$ such that every $\supp f_j$ lies at $\ell^\infty$-distance $>1$ from $\R^3\setminus\CJ$; for $J\ge J_0(f)$ the support condition \eqref{eq:s04:S} holds for $\XJ$. The transport $(O\cdot f)_x(u,v)=f_{O^{-1}x}(O^{-1}u,O^{-1}v)$ and the translation $(\tau_bf)_x=f_{x-b}$ are those of \cref{def:s02:sources}; $O\cdot(O'\cdot f)=(OO')\cdot f$, and both are linear in $f$. Translations and the cubic group act on the Schwinger functions by this signed push-forward of $2$-forms, as in the torus paper \cite{SymPhase}: for an improper $O$ no extra sign appears.

\subsection{Translations on the comparison convention}

\begin{proposition}[translations, comparison convention]\label{prop:s06:translations}
Assume that the comparison family is convergent at fixed cubes (\cref{def:s03:conv}), so that by \cref{lem:s05:TWC}(i),(iii) the limits $Z_J(z)=\lim_ka^{(a)}_{k,J}(z)$ and $Z_\infty(z)=\lim_JZ_J(z)$ exist locally uniformly on $\C^N$ for every family and are entire. Then for every finite family $f$ of real $2$-forms, every $b\in\R^3$ and every $z\in\C^N$,
\[
Z_\infty\Big(\sum_jz_j\tau_bf_j\Big)=Z_\infty\Big(\sum_jz_jf_j\Big).
\]
\end{proposition}

\begin{proof}
\emph{Step 1 (lattice covariance).} Let $k_0\ge0$, $b\in\Lambda_{k_0}$ and $k\ge k_0$. Then $\tau_b$ maps $\Lk$ onto itself, and the substitution $A'_\alpha(x)\defeq A_\alpha(x+b)$, $\phi'(x)\defeq\phi(x+b)$ is a bijection of the configurations of the comparison state on $\XJ$ onto those of the comparison state on $\XJ-b$ (\cref{def:s02:acube} on the translated box) which carries the Lebesgue measure and the action to themselves: the action is a sum of local terms with site-independent coefficients ($\dm_\eps$, $\sW_\eps$ depend on $\eps$ only), and the boundary conditions are transported. The plaquette $p$ corresponds to $p-b$ with the same orientation and $F(p)(A)=F(p-b)(A')$, so $F_{\Lk}(\tau_bf)(A)=F_{\Lk}(f)(A')$. Hence, for all coefficients,
\begin{equation}\label{eq:s06:cov}
Z_{k,\XJ}\Big(\sum_jz_j\tau_bf_j\Big)=Z_{k,\XJ-b}\Big(\sum_jz_jf_j\Big)\qquad(b\in\Lambda_{k_0},\ k\ge k_0),
\end{equation}
the right side being the functional of the comparison state on $\XJ-b$.

\emph{Step 2 (interlacing of the translated box).} For $|b|_\infty\le M$: $X_{J-1}\subset\XJ-b\subset X_{J+1}$, since $-MJ-b_\mu\le-M(J-1)$, $MJ-b_\mu\ge M(J-1)$, $-MJ-b_\mu\ge-M(J+1)$ and $MJ-b_\mu\le M(J+1)$ for every $\mu$.

\emph{Step 3 (the sandwich on the imaginary axes).} Let $J\ge J_0(f)+1$, $t\in\R^N$ and $u_t=\sum_jt_jf_j$, so that $(u_t,0)\in\cK$ is supported at distance $>1$ from $\R^3\setminus C_{J-1}$. Apply the first inequality of \cref{lem:s02:cut} twice: with $Y=\XJ-b$ and $X=X_{J-1}$, and with $Y=X_{J+1}$ and $X=\XJ-b$. The support condition of the lemma holds because the plaquettes carrying $u_t$ have their corners in $X_{J-1}$; no site is killed, so each cut-out factor is the comparison state on the smaller box (the Dirichlet terms of the smaller state are the diagonal parts of the cut bonds, and where the boundaries of two of the boxes touch they coincide). With \eqref{eq:s06:cov},
\[
a^{(a)}_{k,J-1}(\mathrm it)\ \le\ a^{(a),\tau_bf}_{k,J}(\mathrm it)\ \le\ a^{(a)}_{k,J+1}(\mathrm it)\qquad(k\ge k_0,\ J\ge J_0(f)+1).
\]
Let $k\to\infty$: the two ends converge by \cref{lem:s05:TWC}(i) for the family $f$ at $J\mp1\ge J_0(f)$, and the middle term by \cref{lem:s05:TWC}(i) for the family $\tau_bf$, which is admissible at $J$ because $\supp f_j+b$ lies at distance $>1$ from $\R^3\setminus\CJ$ (as $C_{J-1}+b\subset\CJ$). Then let $J\to\infty$ by \cref{lem:s05:TWC}(iii): $Z^{\tau_bf}_\infty(\mathrm it)=Z^f_\infty(\mathrm it)$ for all $t\in\R^N$.

\emph{Step 4 (all coefficients).} $D(z)\defeq Z^{\tau_bf}_\infty(z)-Z^f_\infty(z)$ is entire on $\C^N$ and vanishes on $(\mathrm i\R)^N$, a totally real subspace of maximal real dimension $N$; every partial derivative of $D$ at $0$ vanishes ($\partial_{y_j}D(\mathrm iy)=\mathrm i\,\partial_{z_j}D$), so $D\equiv0$. Thus the identity holds for $b\in\Lambda_{k_0}$, $|b|_\infty\le M$, and all $z\in\C^N$.

\emph{Step 5 (all lattice translations).} A vector $b\in\bigcup_k\Lk$ lies in some $\Lambda_{k_0}$ and is a finite sum $b_1+\dots+b_n$ of vectors $b_i\in\Lambda_{k_0}$ with $|b_i|_\infty\le M$ (each coordinate splits into steps $\pm M$ and a remainder in $\eps_{k_0}\Z$ of modulus $<M$). Steps 3--4 hold for every family, in particular for the translated families $\tau_{b_1+\dots+b_{i-1}}f$; composing gives the identity for $b$.

\emph{Step 6 (all $b\in\R^3$ at real coefficients).} Fix $\rho>0$ and let $E_\rho\defeq C^\infty_c(\bar B_\rho;\Lambda^2)$ with its $C^\infty$ topology, a Fr\'echet space. On $E_\rho$ the function $\Phi(u)\defeq\ln Z_\infty(u)$ at real $u$ is finite ($Z_\infty(u)\ge1$ by Jensen's inequality, as $\av{F(u)}=0$ by charge conjugation), convex (a pointwise limit of the convex functions $u\mapsto\ln\ava{e^{F(u)}}$) and Borel (each lattice functional depends continuously on finitely many values of $u$), hence continuous by \cref{lem:s02:baire}; this is \cref{lem:s05:TWC}(iv). For $|b|\le\rho'-\rho$ the map $b\mapsto\tau_bu\in E_{\rho'}$ is continuous, $\bigcup_k\Lk$ is dense, and $\Phi$ on $E_{\rho'}$ restricted to $E_\rho$ is the same pointwise limit; so Step 5 gives $Z_\infty(\tau_bu)=Z_\infty(u)$ for all $b\in\R^3$ and real $u\in E_\rho$, and exhausting over $\rho$ for all real $u\in C^\infty_c$.

\emph{Step 7 (complex coefficients).} For a family $f$ and $b\in\R^3$, both sides of the asserted identity are entire in $z$ and agree on $\R^N$ by Step 6 (applied to $u=\sum_jx_jf_j$), a set of uniqueness; hence they agree on $\C^N$.
\end{proof}

Only the first inequality of \cref{lem:s02:cut} enters: the sandwich is taken on the cone, and the passage off it is the identity theorem. The translation clause for the joint field on the tube, by the same sandwich at the cone points of the tube, is \cref{lem:s11:TWtube}(v).

\subsection{The translation sandwich on King's convention}

On King's convention the translated cube has its killed sites on $\TJ-b$, which neither contains nor is contained in $T_{J\pm1}$, so \cref{lem:s02:cut} does not compare it with $C^K_{J\pm1}$ directly; one kill step (\cref{lem:s02:BFS}(c) through \cref{lem:s02:Kprime}) and one cut-out with killed sites (\cref{lem:s04:G3}) at each end do.

\emph{Lattice covariance on King's convention.} For $b\in\Lambda_{k_0}$ and $k\ge k_0$ let $\av{\cdot}^{K,-b}_{k,J}$ denote the comparison state on $\XJ-b$ with the matter killed on $\TJ-b$ (King's cube translated by $-b$). The substitution of Step 1 above is a bijection of the configurations on $\XJ$ killed on $\TJ$ onto those on $\XJ-b$ killed on $\TJ-b$, carrying the point masses to the point masses and the action to the action (the kill set translates with the box). Hence, for every finite family and all coefficients,
\begin{equation}\label{eq:s06:covK}
Z^{a,K}_{k,J}\Big(\sum_jz_j\tau_bf_j,\sum_jw_j\tau_bh_j\Big)=Z^{a,K,-b}_{k,J}\Big(\sum_jz_jf_j,\sum_jw_jh_j\Big)\qquad(b\in\Lambda_{k_0},\ k\ge k_0),
\end{equation}
the right side being the scheme functional of $\av{\cdot}^{K,-b}_{k,J}$ with the same multiplier ($\sS_{\Lk}(\tau_bh)=\sS_{\Lk}(h)$).

\begin{lemma}[translation sandwich, King's convention]\label{lem:s06:TK}
Let $(f,g)\in\cK$ be supported at $\ell^\infty$-distance $>1$ from $\R^3\setminus C_{J_0}$, and let $b\in\Lambda_{k_0}$ with $|b|_\infty\le M$. Then for every $k\ge k_0$, every $M\ge1$ and every $J\ge J_0+1$,
\[
Z^{a,K}_{k,J-1}(\mathrm if,g)\ \le\ Z^{a,K}_{k,J}(\mathrm i\,\tau_bf,\tau_bg)\ \le\ Z^{a,K}_{k,J+1}(\mathrm if,g),
\]
and the reversed inequalities hold for $(f,-g)$.
\end{lemma}

\begin{proof}
By \eqref{eq:s06:covK} the middle term is $Z^{a,K,-b}_{k,J}(\mathrm if,g)$; the multipliers of the three terms are the same positive constant (\cref{lem:s02:alg}(a); $\sS$ is translation invariant), so it suffices to prove the inequalities for $\av{\Theta(\mathrm if,g)}$, and the reversed ones for $\Theta(f,-g)$; all these expectations are real.

\emph{Geometry.} Here $b_\mu\in\eps_{k_0}\Z$, $|b_\mu|\le M$, and all sets are subsets of $\Lk$. (G1) $\alive^K_J-b=\prod_\mu\{-MJ-b_\mu,\dots,MJ-2\eps-b_\mu\}\subset\alive^K_{J+1}=\{-MJ-M,\dots,MJ+M-2\eps\}^3$, since $-MJ-b_\mu\ge-MJ-M$ and $MJ-2\eps-b_\mu\le MJ+M-2\eps$. (G2) $\XJ-b\subset X_{J+1}$: the same arithmetic with $\eps$ in place of $2\eps$. (G3) $\alive^K_{J-1}=\{-MJ+M,\dots,MJ-M-2\eps\}^3\subset\alive^K_J-b$, since $-MJ+M\ge-MJ-b_\mu$ and $MJ-M-2\eps\le MJ-2\eps-b_\mu$. (G4) $X_{J-1}\subset\XJ-b$, likewise. All four inclusions hold with margin $M-|b_\mu|\ge0$ on every face, at every spacing, including $k=0$, $M=1$, $|b_\mu|=M$; non-strict inclusions are all that \cref{lem:s04:G3} needs.

\emph{Sources.} By \eqref{eq:s04:S} for $X_{J-1}$, every plaquette carrying $f$ has its corners in $X_{J-1}$ and every site carrying $g$ lies in $\alive^K_{J-1}$; by (G1)--(G4) these lie in the site sets and alive sets of all three states. So the observable is one function of variables present and alive in all three states, and its scalar factor equals $1$ at every site killed below.

\emph{Upper end.} Start from $\av{\cdot}^K_{k,J+1}$, the comparison state on $X_{J+1}$ killed on $T_{J+1}$, a BFS measure of King type (\cref{def:s02:Kcube}). (i) Kill $U_+\defeq\alive^K_{J+1}\setminus(\alive^K_J-b)$. These are alive sites of the state, $g=0$ on $U_+$ (by (G3), $U_+\cap\alive^K_{J-1}=\emptyset$), and $\Theta(\mathrm if,g)$ is continuous with $|\Theta|=e^{\sum_xg_xr_x^2}$; so by \cref{lem:s02:BFS}(c), through \cref{lem:s02:Kprime} (which allows the already killed sites $T_{J+1}$), killing $U_+$ does not increase $\av{\Theta(\mathrm if,g)}$, and by \cref{lem:s02:BFS}(d) with \cref{lem:s02:Kprime} it does not decrease $\av{\Theta(f,-g)}$. The result is the comparison state on $X_{J+1}$ killed on $T'\defeq T_{J+1}\cup U_+=X_{J+1}\setminus(\alive^K_J-b)$ (by (G1); nothing is revived). (ii) Apply \cref{lem:s04:G3} with $V=X_{J+1}$, $T=T'$ and $X=\XJ-b\subset V$ (by (G2)); the lemma allows any $T\subset V$ and any $X\subset V$. The cut-out factor is the comparison state on $\XJ-b$ killed on $(\XJ-b)\cap T'=(\XJ-b)\setminus(\alive^K_J-b)=\TJ-b$, that is $\av{\cdot}^{K,-b}_{k,J}$; the support condition (S$_X$) holds for $X=\XJ-b$ by \emph{Sources}. Hence $\av{\Theta(\mathrm if,g)}^{K,-b}_{k,J}\le\av{\Theta(\mathrm if,g)}^{(a),T'}_{X_{J+1}}\le\av{\Theta(\mathrm if,g)}^K_{k,J+1}$, and $\ge$ for class $\cD$.

\emph{Lower end.} Start from $\av{\cdot}^{K,-b}_{k,J}$, the comparison state on $\XJ-b$ killed on $\TJ-b$ (of King type, a translate of a state of \cref{def:s02:Kcube}). (i) Kill $U_-\defeq(\alive^K_J-b)\setminus\alive^K_{J-1}$, alive sites at which $g=0$ by \emph{Sources}; as above, $\av{\Theta(\mathrm if,g)}$ does not increase and $\av{\Theta(f,-g)}$ does not decrease, and the result is the comparison state on $\XJ-b$ killed on $T''\defeq(\XJ-b)\setminus\alive^K_{J-1}$ (by (G3)). (ii) \cref{lem:s04:G3} with $V=\XJ-b$, $T=T''$ and $X=X_{J-1}\subset V$ (by (G4)): the factor is the comparison state on $X_{J-1}$ killed on $X_{J-1}\cap T''=X_{J-1}\setminus\alive^K_{J-1}=T_{J-1}$, that is $\av{\cdot}^K_{k,J-1}$; (S$_X$) for $X_{J-1}$ is the hypothesis. Hence $\av{\Theta(\mathrm if,g)}^K_{k,J-1}\le\av{\Theta(\mathrm if,g)}^{K,-b}_{k,J}$, and $\ge$ for class $\cD$.

There is no edge case because a translation by $b\in\Lambda_{k_0}$, $|b|_\infty\le M$, moves King's cube rigidly, with the alive sets nested at margins $M-|b_\mu|\ge0$ at every spacing; the reflection of \cref{lem:s04:C} shifts the cube by one lattice layer instead, which is why its one-step form needs $M\ge2\eps_k$.
\end{proof}

The class-$\cI$ clause uses \cref{lem:s02:BFS}(a),(b),(c),(e) and \cref{lem:s02:Kprime} only; the class-$\cD$ clause uses \cref{lem:s02:BFS}(d) in addition. Only the class-$\cI$ clause is used in this section.

\begin{corollary}[translations of King's printed state]\label{cor:s06:TK}
For every finite family $f$ of real $2$-forms, every $b\in\R^3$ and every $z\in\C^N$, $\ZKinf(\sum_jz_j\tau_bf_j)=\ZKinf(\sum_jz_jf_j)$. Moreover, for all real $2$-forms $u,u'\in C^\infty_c$,
\[
|\ZKinf(\mathrm iu)-\ZKinf(\mathrm iu')|\le\sqrt2\,\norm{u-u'}_2,
\]
the $L^2$ norm of a $2$-form being computed in a basis $(e_c)$ of $\Lambda^2$ orthonormal for the norm formula.
\end{corollary}

\begin{proof}
Write $u_t=\sum_jt_jf_j$ for $t\in\R^N$.
(1) \emph{$b\in\Lambda_{k_0}$, $|b|_\infty\le M$, on the imaginary axes.} For $J\ge J_0(f)+1$ and $k\ge k_0$, \cref{lem:s06:TK} at $(u_t,0)\in\cK$ gives $a^K_{k,J-1}(\mathrm it)\le a^{K,\tau_bf}_{k,J}(\mathrm it)\le a^K_{k,J+1}(\mathrm it)$. Let $k\to\infty$: the ends converge by \cref{thm:s05:CK}(b) for $f$ at $J\mp1\ge J_0(f)$; the middle term by \cref{thm:s05:CK}(b) for the family $\tau_bf$, which is admissible at $J$ since $\supp f_j+b$ lies at distance $>1$ from $\R^3\setminus\CJ$ ($C_{J-1}+b\subset\CJ$ for $|b|_\infty\le M$), so that \cref{hyp:s03:EC}, posed for every finite family, applies to it. Hence $Z^K_{J-1}(\mathrm it)\le Z^{K,\tau_bf}_J(\mathrm it)\le Z^K_{J+1}(\mathrm it)$, and $J\to\infty$ (\cref{cor:s04:G2}(ii)($\beta$) at both ends) gives $Z^{K,\tau_bf}_\infty(\mathrm it)=Z^{K,f}_\infty(\mathrm it)$ for all $t\in\R^N$.
(2) \emph{All coefficients, same $b$.} $D(z)\defeq Z^{K,\tau_bf}_\infty(z)-Z^{K,f}_\infty(z)$ is entire on $\C^N$ (\cref{thm:s05:CK}(b) for both families) and vanishes on $(\mathrm i\R)^N$, a totally real subspace of maximal real dimension, so every partial derivative of $D$ at $0$ vanishes and $D\equiv0$.
(3) \emph{All $b\in\bigcup_k\Lk$}: as in Step 5 of \cref{prop:s06:translations}, by composition over the translated families.
(4) \emph{The Lipschitz bound.} For real $2$-forms $u,u'$ whose plaquettes have corners in $\XJ$, at every $(k,J)$, $|\avK{e^{\mathrm iF(u)}}-\avK{e^{\mathrm iF(u')}}|\le\avK{|F(u-u')|}\le{\avK{F(u-u')^2}}^{1/2}\le\sqrt2\norm{u-u'}_{2,k}$, by $|e^{\mathrm ix}-e^{\mathrm iy}|\le|x-y|$, the linearity of $F_{\Lk}$ and the second-moment clause of \cref{thm:s05:CK}(a). Both sides have limits as $k\to\infty$ and $J\to\infty$ (\cref{thm:s05:CK}(b) for the one-member families at the coefficient $\mathrm i$; $\norm{u-u'}_{2,k}\to\norm{u-u'}_2$ by \eqref{eq:s02:riemann}).
(5) \emph{All $b_0\in\R^3$.} Pick $b_n\to b_0$ with $b_n\in\bigcup_k\Lk$. For real $u\in C^\infty_c$, $\tau_{b_n}u\to\tau_{b_0}u$ in $L^2$ (uniform continuity and compact support), so by (4) and (3), $\ZKinf(\mathrm i\tau_{b_0}u)=\lim_n\ZKinf(\mathrm i\tau_{b_n}u)=\ZKinf(\mathrm iu)$. With $u=u_t$ this is $Z^{K,\tau_{b_0}f}_\infty(\mathrm it)=Z^{K,f}_\infty(\mathrm it)$ for all $t$, and (2) extends it to all $z\in\C^N$.
\end{proof}

Nothing here uses a reflection or any property of the limit beyond \cref{thm:s05:CK}; King's derivation of translation invariance from the uniqueness of his limit \cite[Thm.~2.3(i)]{K2} is not used.

\subsection{The cubic group}

\begin{lemma}[passive covariance with killed sites]\label{lem:s06:PK}
Let $O$, $\Lambda$, $(f_\alpha)$ and $X$ be as in \cref{lem:s02:pass}, and let $U\subset X$ be a set of killed sites. Then the identities of \cref{lem:s02:pass} hold between the state on $X$ with the matter killed on $U$ and the state on $OX$ (frame $(Of_\alpha)$, transported boundary conditions and wall) with the matter killed on $OU$:
\[
\big\langle\Theta(z\,O\cdot f,w\,O\cdot h)\big\rangle_{O\Lambda,OX;\,OU}=\big\langle\Theta(zf,wh)\big\rangle_{\Lambda,X;\,U},\qquad\sS_{O\Lambda}(O\cdot h)=\sS_\Lambda(h).
\]
\end{lemma}

\begin{proof}
The configuration-space bijection $\phi'(Ox)=\phi(x)$, $A'_\alpha(Ox)=A_\alpha(x)$ of \cref{lem:s02:pass} carries $\prod_{x\in U}\delta_0(d\phi(x))\,d\nu$ to $\prod_{x\in OU}\delta_0(d\phi'(Ox))\,d\nu$ (the point mass at $\phi(x)=0$ goes to the point mass at $\phi'(Ox)=0$), and the diagonal parts $\frac\eps2r^2$ at the alive ends of the bonds into $U$ are frame-free; the rest of that proof is unchanged.
\end{proof}

\begin{theorem}[cubic symmetry of King's printed state]\label{thm:s06:S1K}
For all real $2$-forms $f_1,\dots,f_N\in C^\infty_c(\R^3;\Lambda^2)$, all $z\in\C^N$, all $O\in\Oh$ (the signed permutation matrices) and all $b\in\R^3$,
\[
\ZKinf\Big(\sum_jz_j\,\tau_bO\!\cdot\!f_j\Big)=\ZKinf\Big(\sum_jz_jf_j\Big).
\]
Hence the Schwinger functions $S^K_n$ of \cref{thm:s05:CK}(c) are $\Cub$-covariant, $F$ transforming as a $2$-form.
\end{theorem}

\begin{proof}
``$(k,J)$ exact'' means an identity of the lattice functionals for all complex coefficients and all $k$, $J$ for which both families have their plaquettes in $\XJ$.

\emph{Step 1 (permutations are exact).} Let $P$ be a permutation matrix, $Pe_\alpha=e_{\pi(\alpha)}$. Then $P\Lk=\Lk$, $P\XJ=\XJ$ and $P\TJ=\TJ$ (the condition ``$x_\mu=MJ-\eps$ for some $\mu$'' is permutation invariant), hence $P\,\alive^K_J=\alive^K_J$. By \cref{lem:s06:PK} with $O=P$, $X=\XJ$, $U=\TJ$, $\langle\Theta(zP\cdot f)\rangle_{P\Lk,P\XJ;\,P\TJ}$ in the frame $(e_{\pi(\alpha)})$ equals $\langle\Theta(zf)\rangle^K_{k,J}$, and the state on the left is $\av{\cdot}^K_{k,J}$ itself: the Dirichlet rule is a condition on the site set $\XJ$, the kill set is $\TJ$, the Neumann rule keeps the same set of (bond, direction) pairs relabelled, the mass term runs over the same bonds, and the hopping and single-site terms are frame-free. The plaquette set and the centres are the same; for a transposed pair $(x;\beta,\alpha)$ both $F(p)$ and $f_{c(p)}(e_\beta,e_\alpha)$ change sign by antisymmetry, so their product is unchanged. Thus $a^{K,P\cdot f}_{k,J}(z)=a^{K,f}_{k,J}(z)$ is $(k,J)$ exact. The distance-$>1$ interior $(-MJ+1,MJ-1)^3$ of $\CJ$ is $P$-invariant, so $J_0(P\cdot f)=J_0(f)$, and \cref{thm:s05:CK}(b) for both families gives $Z^{K,P\cdot f}_\infty=Z^{K,f}_\infty$ on $\C^N$.

\emph{Step 2 (the reflection $\sigma_3$ on the imaginary axes).} Let $\sigma_3=\diag(1,1,-1)$, $t\in\R^N$, $u_t=\sum_jt_jf_j$, $J_0=J_0(f)$. \cref{lem:s04:C}(b), for every $k\ge0$, every $M\ge1$ and every $J\ge J_0+2$, gives
\[
a^{K,f}_{k,J-2}(\mathrm it)\ \le\ a^{K,\sigma_3\cdot f}_{k,J}(\mathrm it)\ \le\ a^{K,f}_{k,J+2}(\mathrm it),
\]
since $\sigma_3\cdot u_t=\sum_jt_j\sigma_3\cdot f_j$. The family $\sigma_3\cdot f$ is admissible at $J$, the interior being $\sigma_3$-symmetric, so $J_0(\sigma_3\cdot f)=J_0(f)$. Let $k\to\infty$ (\cref{thm:s05:CK}(b) for $f$ at $J\mp2$ and for $\sigma_3\cdot f$ at $J$), then $J\to\infty$ (\cref{cor:s04:G2}(ii)($\beta$) at both ends): $Z^{K,\sigma_3\cdot f}_\infty(\mathrm it)=Z^{K,f}_\infty(\mathrm it)$ for all $t\in\R^N$. The one-step sandwich of \cref{lem:s04:C}(a) would serve whenever $M\ge2\eps_k$, that is for all but finitely many $k$, which the $k$-limit ignores; the two-step form is used because it needs no case distinction, and the $J$-limit absorbs the extra step. This is the use King makes of the correlation inequalities in the proof of his Theorem~3.2 \cite{K2}, with the kill steps added; no reversed inequality enters.

\emph{Step 3 ($\sigma_1,\sigma_2$).} $\sigma_\nu=P\sigma_3P^{-1}$ with $P$ the transposition of the axes $\nu$ and $3$. Since the transport is a left action, for every real $2$-form $u$,
\[
\ZKinf(\mathrm i\sigma_\nu\!\cdot\!u)=\ZKinf(\mathrm iP\!\cdot\!(\sigma_3P^{-1}\!\cdot\!u))=\ZKinf(\mathrm i\sigma_3\!\cdot\!(P^{-1}\!\cdot\!u))=\ZKinf(\mathrm iP^{-1}\!\cdot\!u)=\ZKinf(\mathrm iu),
\]
by Step 1 for the family $\sigma_3P^{-1}\cdot u$, Step 2 for the family $P^{-1}\cdot u$, and Step 1 again; each step is a statement about every family.

\emph{Step 4 (the group, on the imaginary axes).} Every $O\in\Oh$ is a product of a permutation matrix and some of $\sigma_1,\sigma_2,\sigma_3$, and if $\ZKinf(\mathrm iO\cdot u)=\ZKinf(\mathrm iu)$ for all real $u$ holds for $O_1$ and $O_2$, it holds for $O_1O_2$. By Steps 1--3, $Z^{K,O\cdot f}_\infty(\mathrm it)=Z^{K,f}_\infty(\mathrm it)$ for every $O\in\Oh$, every family $f$ and every $t\in\R^N$.

\emph{Step 5 (all coefficients).} $D(z)\defeq Z^{K,O\cdot f}_\infty(z)-Z^{K,f}_\infty(z)$ is entire on $\C^N$ (\cref{thm:s05:CK}(b); $J_0(O\cdot f)=J_0(f)$) and vanishes on $(\mathrm i\R)^N$, a set of uniqueness; so $D\equiv0$. No splitting of sources is needed, there being no scalar source.

\emph{Step 6 (about every point).} $\Oh$ about $b$ is $\tau_bO\tau_{-b}$, and translations are \cref{cor:s06:TK} (every family, every $b\in\R^3$, all coefficients): $\ZKinf(\sum_jz_j\tau_bO\cdot f_j)=\ZKinf(\sum_jz_jO\cdot f_j)=\ZKinf(\sum_jz_jf_j)$.

\emph{Step 7 (Schwinger functions).} Differentiating the identity of entire functions at $z=0$ gives $S^K_n(\tau_bO\cdot f_1,\dots,\tau_bO\cdot f_n)=S^K_n(f_1,\dots,f_n)$ for all $n$ and $(b,O)\in\Cub$.
\end{proof}

\begin{remark}[what cubic symmetry needs]\label{rem:s06:K4}
King's cube is exactly symmetric under the coordinate permutations (Step 1) but not under the reflections; there the sandwich of \cref{lem:s04:C} replaces exactness, as that of \cref{lem:s02:refl} does on the comparison convention. So $\Oh$-invariance of King's printed state follows from the uniqueness of the monotone cube limit (\cref{cor:s04:G1} with \cref{thm:s05:CK}(b)) together with the cut sandwich, and does not need the uniqueness of the infinite-volume state. The remark of the torus paper \cite{SymPhase} that $\Oh$-invariance ``would follow from the uniqueness of the infinite-volume limit, which is not known'' is superseded by \cref{thm:s06:S1K}. Only class-$\cI$ clauses enter, and no scheme or Wick constant.
\end{remark}

\subsection{Rotations}

\begin{theorem}[rotations]\label{thm:s06:O3}
$\ZKinf$ is $\Euc$-invariant: for every finite family $f$, every $z\in\C^N$ and every $(b,R)\in\Euc$, $\ZKinf(\sum_jz_j\tau_bR\cdot f_j)=\ZKinf(\sum_jz_jf_j)$, and the Schwinger functions $S^K_n$ are Euclidean covariant, $F$ transforming as a $2$-form.
\end{theorem}

\begin{proof}
By the rotation theorem (\cref{hyp:s03:RC}) the limit state $Z^{\mathrm{LIM}}_\infty(\cdot,0)$ of the comparison family is $\Euc$-invariant for all complex coefficients, and by \cref{cor:s04:G2}(iii), with the fixed-cube limits of \cref{hyp:s03:EC}, it equals $\ZKinf(\cdot,0)$; so $\ZKinf(\cdot,0)$ is $O(3)$-invariant for all complex coefficients; translations are \cref{cor:s06:TK}; the Schwinger functions follow by differentiation at $z=0$.
\end{proof}

The rotation theorem is proved in the rotation part \moving{rotation part}. Nothing in \cref{cor:s06:TK,thm:s06:S1K} depends on it. Its consequences for reflection positivity about every plane and for the components of the measure are drawn in \cref{thm:s09:S2K}(c).

% s07 Reflection positivity I: the modified box
\section{Reflection positivity I: the modified box}

King's Neumann--Dirichlet cube is not reflection symmetric about any plane, and the natural symmetrisation of its photon form is not reflection positive on gauge-invariant functionals (\cref{rem:s07:R1}). Reflection positivity is therefore proved first for the modified box state of \cite[(3.38)]{K2} and then carried to the cube states by a two-sided comparison on the cone: the cube $C_J$ sits inside the box (\cref{lem:s07:L}), and the box sits inside a larger cube once a margin proportional to $\lmu^{-1}$ is allowed (\cref{lem:s07:U}). The passage to the limits is in \cref{thm:s08:OS2C}. Throughout, $\eps=\eps_k$, the lattice is $\Lk\subset\R^d$ with $d=3$ (the certificates of \cref{lem:s07:R} are also given for $d=2$), a \emph{lattice box} is a product $\Omega=\prod_\mu\{a_\mu,a_\mu+\eps,\dots,b_\mu\}\subset\Lk$, $\Omega^*$ is the set of its internal bonds (both ends in $\Omega$) and $\Omega^{**}$ the set of its plaquettes (all four corners in $\Omega$).

\subsection{The lattice reflection and the modified box}

\begin{definition}[mid-planes and the lattice reflection]\label{def:s07:reflection}
The \emph{mid-planes} are the planes $\{x_\nu=c\}$ with $c\in D\defeq\bigcup_{k\ge0}\eps_k(\Z+\frac12)$. If $c\in\eps_{k_0}(\Z+\frac12)$ then $c\in\eps_k(\Z+\frac12)$ for every $k\ge k_0$, because $5^j(m+\frac12)=5^jm+\frac{5^j-1}2+\frac12$ and $5^j$ is odd; the same set $D$ is obtained from any tail $k\ge k_1$ of the spacings. $D$ is dense in $\R$ but not all of $\R$.

Fix $\nu$, a mid-plane $c\in\eps_k(\Z+\frac12)$ and a lattice box $\Omega\subset\Lk$ with $\rho\Omega=\Omega$, $\rho=\rho_{\nu,c}$ the reflection of \cref{def:s02:sources}. On the fields put $(\hat\rho\phi)(x)\defeq\phi(\rho x)$ and $(\hat\rho A)(b)\defeq A(\rho b)$, the image bond taken with its orientation: $(\hat\rho A)_\alpha(x)=A_\alpha(\rho x)$ for $\alpha\ne\nu$ and $(\hat\rho A)_\nu(x)=-A_\nu(\rho x-\eps e_\nu)$. Then $\hat\rho$ is a signed permutation of the variables attached to $\Omega$ and $\Omega^*$ (it preserves Lebesgue measure and maps $\Omega^*$, $\Omega^{**}$ to themselves); the \emph{crossing bonds} $B_0\defeq\{\langle x^-,x^-+\eps e_\nu\rangle:\ x^-_\nu=c-\eps/2\}$ are mapped to their reversals, $(\hat\rho A)(b)=-A(b)$ for $b\in B_0$. Let $\Omega_\pm\defeq\{x\in\Omega:\ \pm(x_\nu-c)>0\}$ (no site lies on the plane), $B_+$ the bonds with both ends in $\Omega_+$, $\mathfrak B_+$ the Borel functions of $(\phi|_{\Omega_+},A|_{B_+})$, and
\[
\theta G\defeq\overline{G\circ\hat\rho}\qquad(G\in\mathfrak B_+).
\]
\end{definition}

On the smeared fields, $F_{\Lk}(f)\circ\hat\rho=F_{\Lk}(\rho^*f)$ and $W^a_{\Lk}(h)\circ\hat\rho=W^a_{\Lk}(h\circ\rho)$. Indeed, $F(p)\circ\hat\rho=\sigma_pF(\rho p)$ with $\sigma_p=-1$ exactly when $\nu\in\{\alpha,\beta\}$ (the orientation of the image plaquette, whose base point is shifted back one layer), the centres satisfy $c(\rho p)=\rho c(p)$, and $(\rho^*f)_{c(\rho p)}(e_\alpha,e_\beta)=f_{c(p)}(\rho e_\alpha,\rho e_\beta)=\sigma_pf_{c(p)}(e_\alpha,e_\beta)$; reindexing $p\mapsto\rho p$ gives the first identity. The scalar case is a relabelling of sites, and the $c$-numbers $2\sW_\eps\sS(h)$ and $a_k\sS(h)$ are invariant because $\rho$ permutes $\Lk$ ($2(x_\nu-c)\in\eps\Z$). Consequently, for families $\zeta_i$ whose lattice observables $\Theta(\zeta_i)$ lie in $\mathfrak B_+$,
\begin{equation}\label{eq:s07:thetaTheta}
\theta(\Theta(\zeta_i))\,\Theta(\zeta_j)=\Theta(\theta\zeta_i^*\sqcup\zeta_j),\qquad m(\theta\zeta_i^*\sqcup\zeta_j)=\overline{m(\zeta_i)}\,m(\zeta_j)
\end{equation}
for the scheme multiplier $m(\zeta)\defeq e^{-\sum_jw_ja_k\sS(h_j)}$ ($a_k$ real), so that for any state $\av{\cdot}$ on the box
\begin{equation}\label{eq:s07:matrix}
\sum_{i,j}\bar\gamma_i\gamma_j\,m(\theta\zeta_i^*\sqcup\zeta_j)\av{\Theta(\theta\zeta_i^*\sqcup\zeta_j)}=\sum_{i,j}\overline{(\gamma_im(\zeta_i))}\,(\gamma_jm(\zeta_j))\,\big\langle\theta(\Theta(\zeta_i))\,\Theta(\zeta_j)\big\rangle .
\end{equation}
\emph{Support.} If the sources of $\zeta_i$ are supported in $\{x_\nu\ge c+\delta\}$ with $\delta>0$, then for $\eps_k<2\delta$ every site carrying a sampled $h$ and every corner of every plaquette carrying a sampled $f$ (centre in the support, corners at $x_\nu\ge c+\delta-\eps_k/2>c$) lies in $\Omega_+$, so $\Theta(\zeta_i)\in\mathfrak B_+$. The splitting $h=h^+-h^-$ of \cref{def:s08:enlarged} is taken with cutoffs supported in $\{x_\nu\ge c+\delta/2\}$, so that $h^\pm$ are supported in $\{x_\nu>c\}$ as well.

\begin{definition}[the modified box state]\label{def:s07:box}
Let $\Omega\subset\Lk$ be a lattice box, $F(p)=\eps^{-1}\sum_{b\in\partial p}A(b)$, $\Omega(y)$ the indicator of $\Omega$, and $x^*$ the nearest neighbours of $x$ in $\Lk$. King's modified action \cite[(3.38)]{K2} on $\Omega$ is
\begin{align*}
S^\eps_\Omega(A,\phi)&=\tfrac12\sum_{p\in\Omega^{**}}\eps^d|F(p)|^2+\tfrac12\lmu^2\sum_{b\in\Omega^*}\eps^d|A(b)|^2+\tfrac12\sum_{x\in\Omega}\eps^d\Big|\sum_{y\in x^*}\eps^{-1}\Omega(y)A(\langle x,y\rangle)\Big|^2\\
&\quad+\tfrac12\sum_{\langle x,y\rangle}\eps^{d-2}\big[(\Omega r)(x)^2+(\Omega r)(y)^2-2(\Omega r)(x)(\Omega r)(y)\cos(\vartheta_x-e\eps A(\langle x,y\rangle)-\vartheta_y)\big]\\
&\quad+\sum_{x\in\Omega}\eps^d\big[\tfrac12m^2r_x^2+\lambda r_x^4+\tfrac12\dm_\eps\,r_x^2\big],
\end{align*}
with the \emph{constant} free-boundary counterterm $\dm_\eps$ of \cref{def:s02:Kcube} at every site, and with the free Wick constant $\sW_\eps$ in the Wick observables. (King prints a site-dependent counterterm in the last term; it is not used here, and the statements below are made for the constant one.) The photon variables are the internal bonds $\Omega^*$ only. The matter part is exactly the Dirichlet matter of the Neumann--Dirichlet state on $X=\Omega$: cosines on the internal bonds, and the diagonal $\frac12\eps^{d-2}r_x^2$ for every bond leaving $\Omega$. The \emph{modified box state} $\langle\cdot\rangle^{(3.38)}_{k,\Omega}$ is the normalised measure $e^{-S^\eps_\Omega}\prod_{x\in\Omega}d^2\phi_x\prod_{b\in\Omega^*}dA_b$; it is a BFS measure (the Gaussian is non-degenerate since $\lmu>0$, the single-site measures have all exponential moments since $\lambda>0$). We write $Z^{a,(3.38)}_{k,\Omega}(\zeta)\defeq m(\zeta)\langle\Theta(\zeta)\rangle^{(3.38)}_{k,\Omega}$ with the same observables and multiplier as in \cref{def:s02:sources}.

The third term is the graph divergence of $(\Omega,\Omega^*)$: with $d_0:\R^\Omega\to\R^{\Omega^*}$, $(d_0\chi)(\langle x,x+\eps e_\alpha\rangle)=\eps^{-1}(\chi(x+\eps e_\alpha)-\chi(x))$, one has $\sum_{y\in x^*}\Omega(y)A(\langle x,y\rangle)=-\eps(d_0^TA)(x)$ (only internal bonds contribute). With $d_1A\defeq(\eps F(p))_{p\in\Omega^{**}}$, the \emph{unitary form} on the box and the modified form are
\[
Q^u_\Omega(A)\defeq\tfrac12\sum_{p\in\Omega^{**}}\eps^d|F(p)|^2+\tfrac12\lmu^2\sum_{b\in\Omega^*}\eps^d|A(b)|^2,\qquad Q^{(3.38)}_\Omega=Q^u_\Omega+\tfrac12\eps^d\|d_0^TA\|^2 ,
\]
\ie the modified form is the unitary form plus a gauge-fixing term depending on $A$ only through $d_0^TA$. The modified action is invariant under $\hat\rho$ for every mid-plane $c$ with $\rho\Omega=\Omega$: every term is written through $A(\langle x,y\rangle)$ and nearest-neighbour pairs, which $\rho$ permutes, $(d_0^T\hat\rho A)(x)=(d_0^TA)(\rho x)$, and $\dm_\eps$ is site independent.
\end{definition}

\emph{The boxes.} Fix $c\in D$, $k\ge k_0(c)$, $J\ge J_0$ such that every source of the enlarged family of \cref{def:s08:enlarged} and its reflection is supported at $\ell^\infty$-distance $>1$ from $\R^3\setminus C_{J_0}$, an integer $n\ge1$, and put $a_k\defeq c+\eps_k/2\in\eps_k\Z$ and
\begin{equation}\label{eq:s07:margin}
\Omega_{J,k}\defeq(C_{J+n}+a_ke_\nu)\cap\Lk,\qquad Mn\ \ge\ |c|+2+R(\lmu),\qquad R(\lmu)\defeq\frac{\pi\sqrt3}{2\lmu}.
\end{equation}
$C_{J+n}\cap\Lk$ is symmetric in direction $\nu$ about $-\eps_k/2$ (sites $x_\nu\in\{-M(J+n),\dots,M(J+n)-\eps_k\}$, and $M(J+n)\in\eps_k\Z$), so $\Omega_{J,k}$ is a lattice box symmetric about $\{x_\nu=c\}$ with an even number $2M(J+n)5^k$ of layers. Under the margin condition of \eqref{eq:s07:margin}, for every $k$: (i) $C_J$ plus one layer is contained in $\Omega_{J,k}$ (this needs $Mn\ge|c|+\frac32\eps_k$); (ii) $C_{J+2n}$ contains every site at $\ell^\infty$-distance $\le R(\lmu)+\eps_k$ from $\Omega_{J,k}$ (this needs $Mn\ge|c|+R(\lmu)+\frac32\eps_k$). We fix such an $n=n(\lmu,c)$, independent of $k$ and $J$. The two margins are exactly the hypotheses of \cref{lem:s07:L,lem:s07:U} below; by \cref{lem:s04:D} they serve King's convention as they stand.

\subsection{Reflection positivity of the modified box}

\begin{lemma}[lattice reflection positivity of the modified box]\label{lem:s07:RPbox}
Let $\lmu>0$, $c$ a mid-plane and $\Omega$ a lattice box with $\rho\Omega=\Omega$, $\rho=\rho_{\nu,c}$.
\begin{enumerate}[label=\textup{(\alph*)},nosep]
\item For every gauge-invariant $G\in\mathfrak B_+$ that is bounded, or of the form $\Theta(\zeta)$ with sources supported in $\Omega_+$, or more generally with $\langle|G|^2\rangle^{(3.38)}_{k,\Omega}<\infty$: $\langle\theta(G)G\rangle^{(3.38)}_{k,\Omega}\ge0$.
\item $(G,G')\mapsto\langle\theta(G')G\rangle^{(3.38)}_{k,\Omega}$ is a Hermitian positive semidefinite form on such $G$; hence $\sum_{i,j}\bar\gamma_i\gamma_j\langle\theta(G_i)G_j\rangle^{(3.38)}_{k,\Omega}\ge0$, and by \eqref{eq:s07:matrix} and the support clause of \cref{def:s07:reflection}, $\sum_{i,j}\bar\gamma_i\gamma_jZ^{a,(3.38)}_{k,\Omega}(\theta\zeta_i^*\sqcup\zeta_j)\ge0$ for families $\zeta_i$ with sources supported in $\{x_\nu\ge c+\delta\}$, $\eps_k<2\delta$.
\end{enumerate}
\end{lemma}

\begin{proof}
\emph{Step 1: gauge-invariant expectations of the modified form equal those of the unitary form.} The graph $(\Omega,\Omega^*)$ is connected, $\ker d_0$ consists of the constants, $\R^{\Omega^*}=\operatorname{ran}d_0\oplus\ker d_0^T$ orthogonally, and $d_1d_0=0$. Write $A=A^T+d_0\chi$ with $d_0^TA^T=0$ and $\chi\perp1$; $(A^T,\chi)\mapsto A$ is a linear bijection with constant Jacobian. Then $Q^u_\Omega(A)=Q^u_\Omega(A^T)+\frac{\lmu^2\eps^d}2\|d_0\chi\|^2$ (orthogonality and $d_1d_0=0$) and $Q^{(3.38)}_\Omega(A)=Q^u_\Omega(A^T)+q(\chi)$ with $q(\chi)=\frac{\lmu^2\eps^d}2\|d_0\chi\|^2+\frac{\eps^d}2\|d_0^Td_0\chi\|^2$; both $\chi$-forms are positive definite on $1^\perp$ because $\lmu>0$. A gauge transformation is $A\mapsto A+d_0\chi$ together with $\vartheta\mapsto\vartheta-e\chi$: the hopping phase $\vartheta_x-e\eps A_b-e(\chi(y)-\chi(x))-\vartheta_y$ equals $(\vartheta_x+e\chi_x)-e\eps A_b-(\vartheta_y+e\chi_y)$ and the Dirichlet diagonals carry no $A$, so for gauge-invariant $G$ the matter integral $\Psi(A)\defeq\int\prod_xd^2\phi_x\,(Ge^{-S_m})(A,\phi)$ satisfies $\Psi(A^T+d_0\chi)=\Psi(A^T)$ (rotate $\phi$ site by site; $\prod d^2\phi$ is invariant). Hence $\int dA\,e^{-Q^{(3.38)}_\Omega}\Psi=\int dA^Te^{-Q^u_\Omega(A^T)}\Psi(A^T)\int d\chi\,e^{-q(\chi)}$, and the same with $Q^u_\Omega$ in place of $Q^{(3.38)}_\Omega$ and a different finite $\chi$-factor; the factors cancel in the normalised expectation. Fubini applies, since $|G|\le Ce^{a\sum_b|A(b)|}e^{C\sum_xr_x^2}$ for $G=\Theta(\zeta)$ (trivially for bounded $G$) is dominated by the Gaussian and by $e^{-\lambda\eps^d r^4}$. This is the mechanism of \cite[Thm.~5.1]{BFS1} and of \cite[lem:C-FP]{SymPhase}, carried to the box graph.

\emph{Step 2: the unitary box measure is reflection positive on all of $\mathfrak B_+$.} Let $S^u_\Omega\defeq Q^u_\Omega+S_m$, write $\omega=(\omega_+,\omega_-,a)$ with $a=(a_b)_{b\in B_0}$ the crossing-bond variables, and put $\mathsf M\defeq\lmu^2\eps^d$. The terms of $S^u_\Omega$ are of four kinds. Single-site terms (local matter and Dirichlet diagonals), hopping cosines on bonds of $B_\pm$, masses of bonds of $B_\pm$, and plaquettes with all corners in $\Omega_+$ or all in $\Omega_-$ involve one side only; by $\hat\rho$-invariance the minus-side terms are the $\hat\rho$-transport of the plus-side terms. The crossing masses are $\frac{\mathsf M}2a_b^2$, $b\in B_0$. A crossing plaquette $p=(x^-;\alpha,\nu)$, $\alpha\ne\nu$, is present iff $x^-,x^-+\eps e_\alpha\in\Omega$ (then all four corners are in $\Omega$ by symmetry); with $x^+\defeq x^-+\eps e_\nu$, $b=\langle x^-,x^+\rangle$ and $b'=\langle x^-+\eps e_\alpha,x^++\eps e_\alpha\rangle$ it reads $\eps F(p)=\pm[A_\alpha(x^-)+a_{b'}-A_\alpha(x^+)-a_b]$, the overall sign immaterial after squaring. A crossing hopping term is $-\eps^{d-2}r_{x^-}r_{x^+}\cos(\vartheta_{x^-}-e\eps a_b-\vartheta_{x^+})$, $b\in B_0$.

Split each crossing variable: $\int_\R g(a_b)e^{-\frac{\mathsf M}2a_b^2}da_b=\sqrt{2\pi/\mathsf M}\ \E\,g(\alpha_b-\alpha'_b)$ with $\alpha_b,\alpha'_b$ independent $N(0,\frac1{2\mathsf M})$, whose difference is $N(0,\frac1{\mathsf M})$. Put $\xi\defeq(\omega_+,\alpha)$ and $\xi'\defeq(\omega_-^\rho,\alpha')$, where $\omega_-^\rho$ is the $\hat\rho$-image of the minus data transported to the plus side ($(\omega^\rho_-)_\alpha(x^+)=A_\alpha(x^-)$ for the tangential bonds, $\phi^\rho(x^+)=\phi(x^-)$, the $\nu$-bonds of $B_-$ carried with the sign of \cref{def:s07:reflection}); by $\hat\rho$-invariance the minus-side one-sided terms are the plus-side terms evaluated on $\omega_-^\rho$, and the a priori measures of $\xi$ and $\xi'$ coincide. For a crossing plaquette put $X_p(\xi)\defeq A_\alpha(x^+)-\alpha_{b'}+\alpha_b$, a function of plus variables and $\alpha$ (note $A_\alpha(x^+)\in B_+$ since $x^+,x^++\eps e_\alpha\in\Omega_+$), and the same function on $\xi'$, $X_p(\xi')=A_\alpha(x^-)-\alpha'_{b'}+\alpha'_b$; then $\eps F(p)=\pm[X_p(\xi')-X_p(\xi)]$ and
\[
\tfrac{\eps^d}2F(p)^2=\tfrac{\eps^{d-2}}2X_p(\xi)^2+\tfrac{\eps^{d-2}}2X_p(\xi')^2-\eps^{d-2}X_p(\xi)X_p(\xi').
\]
For a crossing hopping term put $\psi_b(\xi)\defeq\phi(x^+)e^{\mathrm ie\eps\alpha_b}$, $\psi_b(\xi')=\phi(x^-)e^{\mathrm ie\eps\alpha'_b}$; then $r_{x^-}r_{x^+}\cos(\vartheta_{x^-}-e\eps(\alpha_b-\alpha'_b)-\vartheta_{x^+})=\Re[\psi_b(\xi')\overline{\psi_b(\xi)}]=\sum_{i=1,2}\psi^i_b(\xi)\psi^i_b(\xi')$. Hence, with $S'$ the action minus the crossing masses,
\[
e^{-S'}=\cG(\xi)\,\cG(\xi')\,\exp\Big(\sum_\iota\gamma_\iota Y_\iota(\xi)Y_\iota(\xi')\Big),\qquad\cG(\xi)\defeq e^{-S_+(\xi)-\frac{\eps^{d-2}}2\sum_pX_p(\xi)^2},
\]
where $S_+$ collects the one-sided plus terms and $\{(\gamma_\iota,Y_\iota)\}=\{(\eps^{d-2},X_p)\}_p\cup\{(\eps^{d-2},\psi^i_b)\}_{b,i}$: all $\gamma_\iota>0$, all $Y_\iota$ real. For $G\in\mathfrak B_+$ we have $G(\omega)=G(\xi)$ and $\theta G(\omega)=\overline{G(\xi')}$, the same function on the transported data ($G$ depends neither on $a$ nor on $\alpha,\alpha'$). Expanding the exponential,
\[
Z\,\langle\theta(G)G\rangle^u_{k,\Omega}=K\sum_{n\in\N^I}\prod_\iota\frac{\gamma_\iota^{n_\iota}}{n_\iota!}\Big|\int d\xi\,G(\xi)\cG(\xi)\prod_\iota Y_\iota(\xi)^{n_\iota}\Big|^2\ \ge0,\qquad K>0,
\]
the two integrals of each term being complex conjugates because measures and functions coincide on the two copies. The interchange of sum and integral is justified by $\sum_n\prod_\iota\gamma_\iota^{n_\iota}|Y_\iota Y'_\iota|^{n_\iota}/n_\iota!\le\exp(\frac12\sum_\iota\gamma_\iota(Y_\iota^2+Y_\iota'^2))$: the plaquette parts cancel against $\cG$ exactly, $|\psi_b|=r_{x^+}$, and $\int|G|e^{-S_+(\xi)+\frac{\eps^{d-2}}2\sum_br_{x^+}^2}d\xi<\infty$ by $\lambda>0$ (a quartic term at every plus site), $\lmu>0$ (a mass on every $B_+$ bond and the $\alpha$-Gaussians) and the growth bound on $G$. This is the mechanism of \cite[lem:C-RPP(ii)]{SymPhase} with one plane in place of a pair and the Dirichlet diagonals as single-site terms; the classical form is the positive cross-term criterion of \cite{GJ,FILS78}. The Hermitian property $\overline{\langle\theta(G')G\rangle}=\langle(G'\circ\hat\rho)\bar G\rangle=\langle G'\,\overline{G\circ\hat\rho}\rangle=\langle\theta(G)G'\rangle$ follows from the $\hat\rho$-invariance of the measure.

\emph{Step 3: combination and $L^2$ extension.} For gauge-invariant $G\in\mathfrak B_+$, bounded or equal to $\Theta(\zeta)$, the product $\theta(G)G$ is gauge invariant ($\hat\rho$ maps gauge transformations to gauge transformations: $\hat\rho\,d_0\chi=d_0(\chi\circ\rho)$), so $\langle\theta(G)G\rangle^{(3.38)}_{k,\Omega}=\langle\theta(G)G\rangle^u_{k,\Omega}\ge0$ by Steps 1--2. For gauge-invariant $G\in\mathfrak B_+$ with $\langle|G|^2\rangle<\infty$ (for instance polynomials in the plus-side smeared fields, \cref{prop:s08:poly}), $G_n\defeq G\,\mathbf 1\{|G|\le n\}$ is bounded, gauge invariant and in $\mathfrak B_+$, so $\langle\theta(G_n)G_n\rangle\ge0$; moreover $|\theta(G_n)G_n|\le|\theta G||G|$ and $\langle|\theta G||G|\rangle\le\langle|G|^2\rangle^{1/2}\langle|\theta G|^2\rangle^{1/2}=\langle|G|^2\rangle$ by Cauchy--Schwarz and $\hat\rho$-invariance, so dominated convergence gives (a). The modified measure is itself $\hat\rho$-invariant (\cref{def:s07:box}), so the computation of the Hermitian property in Step 2 holds verbatim for it on $L^2$; the form is sesquilinear, Hermitian and nonnegative on the diagonal, hence positive semidefinite (apply (a) to $G=\sum_i\gamma_iG_i$). With \eqref{eq:s07:matrix} and the support clause this is (b).
\end{proof}

\begin{remark}[on King's Proposition 3.3]
\Cref{lem:s07:RPbox} is the statement of \cite[Prop.~3.3]{K2}, which King derives from \cite[Cor.~5.4]{BFS1}. That corollary (with \cite[Lemma~5.2, Thm.~5.3]{BFS1}) concerns the lattice actions of \cite[\S5]{BFS1}, whose photon Gaussian is translation invariant with a reflection-positive boundary term; King's finite-box Gaussian with its masked divergence is not one of them, the obstruction being the box geometry of the gauge-fixing term. So his Proposition 3.3 is true, and it is proved here rather than derived from \cite[Cor.~5.4]{BFS1}. The gauge-fixing term $\|d_0^TA\|^2$ is not itself compatible with the mechanism of Step 2 (at a site next to the plane it produces cross terms $Y(\xi)Z(\xi')+Z(\xi)Y(\xi')$ of indefinite sign); on the torus it drops out for gauge-invariant observables, on King's cube it does not reduce to a divergence at the boundary (\cref{lem:s07:B}), and this is why the photon form has to be changed before reflecting (\cref{rem:s07:R1}).
\end{remark}

\subsection{The same box: the boundary identity and the reversed inequality}

In this subsection we use lattice units, $\eps=1$ (every term of the three photon forms carries the same weight $\frac12\eps^{d-2}$), and write $\mu\defeq\lmu\eps$. For a lattice box $\Omega$ and $A\in\R^{\Omega^*}$ let $K^{\rm int}_\Omega(A)\defeq\sum_\mu\sum_\nu\sum|A_\mu(x+e_\nu)-A_\mu(x)|^2$, the sum over pairs with both bonds $(x,\mu),(x+e_\nu,\mu)$ in $\Omega^*$ (the componentwise Neumann gradient form on the internal bonds), and $N_\partial(A)\defeq\sum|A_b|^2$, the sum over the normal crossing pairs: for $b=(x,\mu)\in\Omega^*$ one term for each of $(x\pm e_\mu,\mu)$ that is not in $\Omega^*$. Split $N_\partial=N^{\rm up}_\partial+N^{\rm low}_\partial$ according to whether the missing partner is $(x+e_\mu,\mu)$ (\ie $x+2e_\mu\notin\Omega$) or $(x-e_\mu,\mu)$ (\ie $x-e_\mu\notin\Omega$).

\begin{lemma}[the boundary identity]\label{lem:s07:B}
For every lattice box $\Omega$ and every $A\in\R^{\Omega^*}$,
\[
\|d_1A\|^2_{\Omega^{**}}+\|d_0^TA\|^2=K^{\rm int}_\Omega(A)+N_\partial(A).
\]
Consequently $Q^{(3.38)}_\Omega=\frac12\big[K^{\rm int}_\Omega+N_\partial+\mu^2\|\cdot\|^2_{\Omega^*}\big]$: the modified photon form is the internal Neumann--Feynman form plus the positive boundary term $N_\partial$. (By the proof of \cref{lem:s07:R}(a), the marginal of the Neumann--Dirichlet Gaussian of $\Omega$ on the internal bonds has the precision $M^{(3.38)}-N^{\rm low}_\partial-P_{12}P_{22}^{-1}P_{21}$: the modified form minus the squares of the near-normal bonds and minus the Schur term of the far-face variables.)
\end{lemma}

\begin{proof}
Extend $A$ by $0$ to a $1$-form $a$ on $\Z^d$. For finitely supported $1$-forms, $\sum_p|da(p)|^2+\sum_x|d^*a(x)|^2=\sum_\mu\sum_\nu\sum_x|a_\mu(x+e_\nu)-a_\mu(x)|^2$ (the Hodge Laplacian on $\Z^d$ is the componentwise Laplacian; summation by parts). On the left, $d^*a(x)$ is the masked divergence for $x\in\Omega$ and $0$ outside, since no internal bond has an end outside $\Omega$; a plaquette $p\notin\Omega^{**}$ meeting $\Omega^*$ has exactly one internal bond $b$ (two parallel internal edges put all four corners in $\Omega$; two adjacent ones put three corners in $\Omega$, and a product of intervals then contains the fourth), so $|da(p)|^2=|A_b|^2$. On the right, pairs with both bonds internal give $K^{\rm int}_\Omega$; a pair $(x,\mu),(x+e_\nu,\mu)$ with exactly one internal bond contributes $|A_b|^2$; for $\nu\ne\mu$ the partner is non-internal iff $x+e_\nu\notin\Omega$ (as $x,x+e_\mu\in\Omega$, and in a box $x+e_\nu\in\Omega\iff x+e_\mu+e_\nu\in\Omega$), \ie iff the plaquette spanned by $b$ and $e_\nu$ lies outside $\Omega^{**}$; these pairs are in bijection with the plaquettes counted on the left and cancel; for $\nu=\mu$ they are the terms of $N_\partial$.
\end{proof}

\begin{lemma}[the same box: the second inequality of King's (3.42) is reversed]\label{lem:s07:R}
Let $\Omega$ be a lattice box, $\langle\cdot\rangle^{\ND}_\Omega$ the Neumann--Dirichlet state on $X=\Omega$ (variables $A_\mu(x)$, $x\in\Omega$, all outgoing bonds; photon form $Q^{\ND}_\Omega=\frac12[K_\Omega+\mu^2\|\cdot\|^2]$ of \cref{def:s02:acube}) and $\langle\cdot\rangle^{(3.38)}_\Omega$ the modified box state of \cref{def:s07:box} (variables $\Omega^*$); both carry the same Dirichlet matter on $\Omega$ (constant $\dm_\eps$, free Wick constant), coupled to $A$ through the internal bonds only. Let $\Sigma^{\ND}_\Omega$ and $\Sigma^u_\Omega$ be the covariances of $F=(d_1A(p))_{p\in\Omega^{**}}$ under the two photon Gaussians. Then:
\begin{enumerate}[label=\textup{(\alph*)},nosep]
\item $\Sigma^u_\Omega=\Sigma^{(3.38)}_\Omega\le\Sigma^{\ND}_\Omega$ on $\R^{\Omega^{**}}$;
\item for every gauge-invariant $Y\in\cI$ with sources in $\Omega$ and every $e\ge0$: $\langle Y\rangle^{\ND}_\Omega\le\langle Y\rangle^{(3.38)}_\Omega$, and $\langle X\rangle^{\ND}_\Omega\ge\langle X\rangle^{(3.38)}_\Omega$ for $X\in\cD$. This is the reverse of the second inequality of \cite[(3.42)]{K2}, whose middle term is the modified box state on $C_{J+1}$; it holds weakly, for the modified state with the constant $\dm_\eps$ (nothing is asserted for King's printed site-dependent counterterm);
\item King's inequality fails strictly in the Gaussian sector: at $e=0$ and $Y=e^{\mathrm it\langle v,F\rangle}$ it reads $v^T\Sigma^{\ND}_\Omega v\le v^T\Sigma^u_\Omega v$, while $\Delta_\Omega\defeq\Sigma^{\ND}_\Omega-\Sigma^u_\Omega\ge0$ is nonzero. In exact rational arithmetic the inertia (negative, zero, positive) of $\Delta_\Omega$ is $(0,1,8)$ on the $4\times4$ box ($d=2$) for $\mu^2\in\{\frac1{20},\frac14,1\}$ and $(0,3,13)$ on the $4\times2\times2$ box ($d=3$) for $\mu^2\in\{\frac1{20},\frac14\}$; single-plaquette certificates are $\Delta_{pp}=\Var^{\ND}F(p)-\Var^{(3.38)}F(p)=\frac{14285248}{44314767}$ at $p=((0,0);0,1)$ on $4\times4$, $\mu^2=\frac14$, and $\frac{121897601600}{395580222961}$ at $p=((0,0,0);1,2)$ on $4\times2\times2$, $\mu^2=\frac1{20}$.
\end{enumerate}
\end{lemma}

\begin{proof}
(a) Write the Neumann--Dirichlet precision $P\defeq K_\Omega+\mu^2$ in blocks, the internal variables $\Omega^*$ first and the outgoing boundary variables second. Its internal block is $P_{11}=K^{\rm int}_\Omega+N^{\rm up}_\partial+\mu^2$: pairs of $K_\Omega$ with both bonds internal give $K^{\rm int}_\Omega$; a pair $(x,\mu),(x+e_\nu,\mu)$ with exactly one internal bond is impossible for $\nu\ne\mu$ (box convexity, as in \cref{lem:s07:B}) and for $\nu=\mu$ is $(x,\mu)$ internal with $(x+e_\mu,\mu)$ outgoing, \ie $x+2e_\mu\notin\Omega$ (the other assignment would need $x+e_\mu\notin\Omega$ and $x+e_\mu\in\Omega$ at once); its diagonal contribution to the internal block is $|A_\mu(x)|^2$, which is $N^{\rm up}_\partial$. Hence, by \cref{lem:s07:B}, $M^{(3.38)}\defeq K^{\rm int}_\Omega+N_\partial+\mu^2=P_{11}+N^{\rm low}_\partial\ge P_{11}$. The marginal covariance of the internal variables under the Neumann--Dirichlet Gaussian is the inverse of the Schur complement, $(P^{-1})_{11}=(P_{11}-P_{12}P_{22}^{-1}P_{21})^{-1}\ge P_{11}^{-1}\ge(M^{(3.38)})^{-1}$, by operator monotonicity of the inverse on positive definite matrices and $P_{12}P_{22}^{-1}P_{21}\ge0$. Since $F=d_1A$ on $\Omega^{**}$ uses internal bonds only, $\Sigma^{\ND}_\Omega=D_1(P^{-1})_{11}D_1^T\ge D_1(M^{(3.38)})^{-1}D_1^T=\Sigma^{(3.38)}_\Omega$, and $\Sigma^{(3.38)}_\Omega=\Sigma^u_\Omega$ by Step 0 in the proof of \cref{lem:s07:U} (the field strength depends only on the $\ker d_0^T$-component, whose law is the same under both Gaussians, Step 1 of \cref{lem:s07:RPbox}). (b) The mechanism of the second part of \cref{lem:s07:U} with $B=\Omega$ and no hopping removal: both states have the Dirichlet matter of $\Omega$ coupled to $\Omega^*$, and for gauge-invariant $Y$ with sources in $\Omega$ the expectation depends on the photon law only through the centred Gaussian law of $F$ on $\operatorname{ran}d_1$; the two $F$-laws are ordered by (a); realise them, regularise, apply \cref{lem:s02:BFS}(a) (resp.\ (d)) by interpolation and let the regularisation go to zero exactly as there. The larger covariance has the smaller $\cI$-expectation, so $\langle Y\rangle^{\ND}_\Omega\le\langle Y\rangle^{(3.38)}_\Omega$. (c) At $e=0$ the matter decouples from $A$ and is identical in the two states, so for $Y=e^{\mathrm it\langle v,F\rangle}$ both sides equal $\exp(-\frac{t^2}2v^T\Sigma v)$ with the respective $\Sigma$; King's $\le$ becomes $v^T\Sigma^{\ND}_\Omega v\le v^T\Sigma^u_\Omega v$, which is false for every $v$ with $v^T\Delta_\Omega v>0$, for instance $v=e_p$ at the plaquettes listed. The inertias and the certificates are obtained from the two covariance matrices in exact rational arithmetic (symmetric elimination with diagonal pivoting over $\mathbb Q$).
\end{proof}

\emph{Scope.} (i) King's passage from the box states to the cube limit on \cite[p.~336]{K2} needs the sandwich for the box functionals with purely imaginary curvature coefficients, \ie for the whole class $\cI$; (c) refutes the second inequality on that class exactly. (ii) For the scalar observable $\Theta(0,2h_+)$ of the printed display, (b) shows that the inequality can hold only with equality; at $e=0$ it does (identical decoupled matter), and at $e>0$ a strict failure is to be expected (the Wick square couples to $F$ at order $e^2$) but is not certified here. Both clauses are for the modified state with the constant $\dm_\eps$: if King's site-dependent counterterm differs from the free-boundary number, the two matter measures differ, the mechanism of (b) does not apply, and neither clause is proved for his printed objects; (c) alone is independent of the counterterm, the matter decoupling from $A$ at $e=0$.

\begin{remark}[what is false in the printed route]\label{rem:s07:R1}
(i) King's Neumann--Dirichlet cube state is not reflection symmetric about any plane: the site set $C_J\cap\Lk$ is symmetric in direction $\nu$ about $x_\nu=-\eps_k/2$, the set of $\nu$-bonds about $x_\nu=0$, and no reflection maps the variable set to itself (\cref{lem:s02:refl} accordingly obtains the reflection symmetry of the limit from a cut sandwich, not from a lattice symmetry). (ii) The natural symmetrisation, the Feynman gradient form on the internal bonds $K^{\rm int}_\Omega+\mu^2$ (the modified form without $N_\partial$; reflection symmetric), is \emph{not} reflection positive on gauge-invariant functionals. This is exact at $e=0$: on the $4\times4$ box at $\mu^2=\frac1{20}$ the reflection kernel $R_{pq}=\sigma_p\Sigma[\rho p,q]$ of the plus-side plaquettes is negative definite (inertia $(3,0,0)$), and on $4\times2\times2$ at $\mu^2=\frac1{20}$ it has inertia $(2,2,2)$ with $R_{pp}=\Cov(\theta F(p),F(p))<0$ for $p=((2,0,0);0,1)$. Indeed, at $e=0$ the matter decouples, and the gauge-invariant plus-side functionals $G_1=e^{\mathrm it\langle v,F_+\rangle}$, $G_2=1$ give the matrix $(\langle\theta(G_i)G_j\rangle)$ with determinant $e^{-t^2v^T\Sigma_{++}v}(e^{t^2v^TRv}-1)<0$ for every $t>0$ when $v^TRv<0$, since $\langle\theta(G_1)G_1\rangle=\exp(-t^2v^T\Sigma_{++}v+t^2v^TRv)$ and $\langle\theta(G_1)\rangle=\langle G_1\rangle=\exp(-\frac12t^2v^T\Sigma_{++}v)$ by the reflection symmetry $\Sigma_{--}=\Sigma_{++}$. By continuity at fixed spacing the symmetrised cube state is not reflection positive for small $e>0$ either. A proof must work at every coupling, so reflection positivity of King's own cube form, or of this symmetrisation, at fixed spacing is not available as a route: any route through cube states at fixed spacing has to change the photon form, as the modified box does (a reflection-symmetric, reflection-positive modification whose values converge to the same limit, \cref{thm:s08:OS2C}). Nothing is claimed about the symmetrised form at a given larger $e$ or about other proofs; the route taken here needs the comparison of \cref{lem:s07:L,lem:s07:U}, and its reflecting plane must be a lattice mid-plane for the half-bond mechanism of \cref{lem:s07:RPbox} (King's cubes have an even number of layers). (iii) Accordingly, the second inequality of \cite[(3.42)]{K2} is false as printed in the following sense: false for the class of observables that the argument on \cite[p.~336]{K2} needs (exactly, \cref{lem:s07:R}(c)), and never strictly true for the modified state with the constant counterterm (\cref{lem:s07:R}(b)). Its valid first half is \cref{lem:s07:L}; the invalid half is replaced by \cref{lem:s07:U} with the margin $R(\lmu)$ of \eqref{eq:s07:margin}, $n(\lmu,c)$ cube steps instead of one.
\end{remark}

\begin{remark}[the box form is the larger one]\label{rem:s07:K8}
King describes the modified box action as obtained from the cube action by omitting some terms. By \cref{lem:s07:B} the relation is the reverse: on one box the modified photon form is the internal Neumann--Dirichlet form $P_{11}$ plus the positive term $N^{\rm low}_\partial$ and without the outgoing boundary variables, so the modified state is the one with the \emph{larger} form, the smaller field-strength covariance, and the larger $\cI$-expectations: $\langle Y\rangle^{\ND}_{C}\le\langle Y\rangle^{(3.38)}_{C}$ on the cone for every box $C$ (\cref{lem:s07:R}(b)), the direction that \cref{lem:s02:cut} requires for consistency with \cref{lem:s07:L,lem:s07:U}. For King's convention the same comparison holds with Step~1 of \cref{lem:s04:G} prepended (\cref{rem:s04:F}): $\avK{Y}\le\ava{Y}\le\langle Y\rangle^{(3.38)}_{k,\XJ}$ on $\cI$.
\end{remark}

\subsection{The sandwich}

\begin{lemma}[the lower half]\label{lem:s07:L}
Let $C\subset\Omega$ be lattice boxes with $C$ plus one layer contained in $\Omega$, so that the variables $A_\mu(x)$, $x\in C$, of the Neumann--Dirichlet state on $C$ (the outgoing bonds) are internal bonds of $\Omega$. Then
\begin{enumerate}[label=\textup{(\alph*)},nosep]
\item $Q^{(3.38)}_\Omega\ge Q^{\ND}_C\oplus\frac{\mu^2}2I_{\rm rest}$ as quadratic forms on $\R^{\Omega^*}=\R^{C\text{-variables}}\oplus\R^{\rm rest}$;
\item for sources such that every plaquette and every site at which they do not vanish lies in $C$,
\[
\langle Y\rangle^{\ND}_C\le\langle Y\rangle^{(3.38)}_\Omega\ \ (Y\in\cI),\qquad\langle X\rangle^{\ND}_C\ge\langle X\rangle^{(3.38)}_\Omega\ \ (X\in\cD),
\]
and the same for $Z^a$ (equal positive multipliers, \cref{lem:s02:alg}(a)). The $\cD$ half uses \cref{lem:s02:BFS}(d); the $\cI$ half only (a), (b), (e).
\end{enumerate}
\end{lemma}

\begin{proof}
(a) $K^{\rm int}_\Omega$ contains every gradient pair of $K_C$ (both base points in $C$, both bonds internal to $\Omega$); the remaining pairs and $N_\partial$ are nonnegative; the mass term splits. (b) Start from $\langle\cdot\rangle^{(3.38)}_\Omega$ and remove the hopping cosines across $\partial C$ and all cosines not internal to $C$; by \cref{lem:s02:BFS}(b),(d) $\langle Y\rangle$ does not increase and $\langle X\rangle$ does not decrease, every intermediate measure being a BFS measure. The matter is now the Dirichlet matter of $C$: the diagonal parts $\frac12r^2$ of the cut bonds stay and are exactly the Dirichlet terms, as in the proof of \cref{lem:s02:cut}; it is coupled to $A$ only through bonds internal to $C$, which are $C$-variables, and the matter of $\Omega\setminus C$ is decoupled and cancels (\cref{lem:s02:BFS}(e)). Then pass from the form $Q^{(3.38)}_\Omega$ to the smaller form $Q^{\ND}_C\oplus\frac{\mu^2}2I_{\rm rest}$: by (a) the difference is positive semidefinite, a finite sum $\sum_is_i|\ell_i(A)|^2$ with $s_i\ge0$, and adding each term does not decrease $\langle Y\rangle$ and does not increase $\langle X\rangle$ (\cref{lem:s02:BFS}(a),(d)). So the state with the larger form, the modified one, has the larger $\langle Y\rangle$ and the smaller $\langle X\rangle$. The state with the smaller form factorises as $\langle\cdot\rangle^{\ND}_C$ times an independent Gaussian on the rest, which cancels.
\end{proof}

For King's convention the lower end is used in the form $\avK{Y}\le\ava{Y}\le\langle Y\rangle^{(3.38)}_{\Omega}$ on $\cI$ (Step~1 of \cref{lem:s04:G}, then \cref{lem:s07:L}; \cref{rem:s04:F}), and the upper end needs no change (\cref{lem:s04:D}).

\begin{lemma}[the upper half, with a margin]\label{lem:s07:U}
Let $\Omega\subset B$ be lattice boxes at spacing $\eps$ such that $B$ contains every site within $\ell^\infty$-distance $N$ lattice steps of $\Omega$, and let $\mu=\lmu\eps$ satisfy $d(\pi/2N)^2\le\mu^2$, \ie
\[
N\eps\ \ge\ \frac{\pi\sqrt d}{2\lmu}\qquad\big(=R(\lmu)=\tfrac{\pi\sqrt3}{2\lmu}\text{ in }d=3\big).
\]
Let $\Sigma_X$ denote the covariance of $(\eps F(p))_{p\in\Omega^{**}}$ under the photon Gaussian of the state on $X$. Then
\begin{enumerate}[label=\textup{(\alph*)},nosep]
\item $\Sigma^{\ND}_B\le\Sigma^{(3.38)}_\Omega=\Sigma^u_\Omega$ on $\R^{\Omega^{**}}$;
\item for sources supported in $\Omega$ as in \cref{lem:s07:L},
\[
\langle Y\rangle^{(3.38)}_\Omega\le\langle Y\rangle^{\ND}_B\ \ (Y\in\cI),\qquad\langle X\rangle^{(3.38)}_\Omega\ge\langle X\rangle^{\ND}_B\ \ (X\in\cD),
\]
and the same for $Z^a$. Only the $\cI$ half is used in the sequel.
\end{enumerate}
\end{lemma}

\begin{proof}
\emph{(a), in lattice units.} Step 0: $\Sigma^{(3.38)}_\Omega=\Sigma^u_\Omega$, because $d_1A=d_1A^T$ depends only on the $\ker d_0^T$-component, whose law $\propto e^{-Q^u_\Omega(A^T)}$ is the same under both Gaussians (Step 1 of \cref{lem:s07:RPbox}). Step 1 (variational identity): for a Gaussian with density $\propto e^{-\frac12A^TM_XA}$, $\langle g,\Sigma_Xg\rangle=\Var\langle g,d_1A\rangle=\sup_A[2\langle g,d_1A\rangle-A^TM_XA]$, the supremum attained at $A=M_X^{-1}d_1^Tg$. So it suffices to show that for every $A$ on the $B$-variables there is $A'\in\R^{\Omega^*}$ with $d_1A'=d_1A$ on $\Omega^{**}$ and $\|d_1A'\|^2_{\Omega^{**}}+\mu^2\|A'\|^2_{\Omega^*}\le K_B(A)+\mu^2\|A\|^2_B$; take $A'\defeq A|_{\Omega^*}$ (the plaquettes of $\Omega^{**}$ use only $\Omega^*$-bonds). Step 2 (cutoff): let $\psi(x)=\prod_\nu\phi_\nu(x_\nu)$ with $\phi_\nu=1$ on the $\nu$-range of $\Omega$ and on the first layer outside, $\phi_\nu(t)=\cos(\pi(t-1)/2N)$ at the $t$-th layer outside for $1\le t\le N+1$ (so $\phi_\nu(N+1)=0$), and $0$ beyond; then $0\le\psi\le1$, $\psi=1$ on $\Omega$ and its neighbours, $\supp\psi\subset B$. Put $a\defeq\psi A$ on $\Z^d$ (the bond $\langle x,x+e_\mu\rangle$ weighted by $\psi(x)$; $A\defeq0$ off the $B$-variables). Step 3 (Hodge on $\Z^d$): $\|d_1A\|^2_{\Omega^{**}}=\|d_1A'\|^2_{\Omega^{**}}\le\sum_p|da(p)|^2\le\sum_p|da|^2+\sum_x|d^*a|^2=\sum_\mu\sum_\nu\sum_x|a_\mu(x+e_\nu)-a_\mu(x)|^2$, since the bonds of plaquettes in $\Omega^{**}$ have base points in $\Omega$, where $\psi=1$. Step 4 (discrete IMS identity): for finitely supported $\psi f$ and a direction $e$,
\[
\sum_x|\psi(x+e)f(x+e)-\psi(x)f(x)|^2=\sum_x\psi(x)\psi(x+e)|f(x+e)-f(x)|^2+\sum_x\psi(x)\big(2\psi(x)-\psi(x+e)-\psi(x-e)\big)f(x)^2 ,
\]
exactly (per bond the difference of the two sides is $\psi'(\psi'-\psi)f'^2+\psi(\psi-\psi')f^2$; collect by site). With $f=A_\mu$ and $e=e_\nu$: $\psi\psi'\le1$, and $\psi\psi'>0$ only for pairs with both base points in $B$, which are gradient pairs of $K_B$; so the first sums, over $\mu,\nu$, total at most $K_B(A)$. The potential is $V(x)=\sum_\nu\big(\prod_{\kappa\ne\nu}\phi_\kappa^2\big)\phi_\nu(-\Delta\phi_\nu)(x_\nu)$. Step 5 (the profile): with $\vartheta\defeq\pi/2N$, $\phi(-\Delta\phi)=0$ on the range of $\Omega$ (all neighbours equal $1$), $=1-\cos\vartheta\le\vartheta^2/2$ at $t=1$, $=2\cos^2(\vartheta(t-1))(1-\cos\vartheta)\le\vartheta^2$ for $2\le t\le N$ (using $\phi(N+1)=\cos(\vartheta N)=0$ at $t=N$), and $=0$ for $t\ge N+1$; hence $V=0$ on $\Omega$ and $V\le d\vartheta^2$ everywhere. Step 6: $\|d_1A'\|^2_{\Omega^{**}}\le K_B(A)+d\vartheta^2\sum_{x\in B\setminus\Omega}|A(x)|^2$; the variables with base point in $B\setminus\Omega$ and those of $\Omega^*$ (base point in $\Omega$) are disjoint subsets of the $B$-variables, so $\|d_1A'\|^2_{\Omega^{**}}+\mu^2\|A'\|^2_{\Omega^*}\le K_B(A)+\mu^2\|A\|^2_B$ as soon as $d\vartheta^2\le\mu^2$.

\emph{(b).} Start from $\langle\cdot\rangle^{\ND}_B$ and remove all hopping cosines touching $B\setminus\Omega$ (\cref{lem:s02:BFS}(b): $\cI$ does not increase; (d): $\cD$ does not decrease). The matter is now the Dirichlet matter of $\Omega$ (the diagonals of the bonds leaving $\Omega$ kept; every such bond ends in $B$ since $N\ge1$), coupled to $A$ on $\Omega^*$, and the sites of $B\setminus\Omega$ decouple. This matter is identical to that of the modified box state: the same cosines on $\Omega^*$, the same diagonals $\frac12r_x^2$ for every bond leaving $\Omega$, the same single-site weights with the constant $\dm_\eps$ and the same Wick constant. For a gauge-invariant observable with sources in $\Omega$ the matter integral at fixed $A$, $\Psi_Y(A)=\int\prod_{x\in\Omega}d^2\phi_x\,Y\,e^{-S_m}$, is invariant under $A\mapsto A+d_0\chi$, $\chi\in\R^\Omega$ (as in Step 1 of \cref{lem:s07:RPbox}), hence a function of the class of $A|_{\Omega^*}$ modulo $\operatorname{ran}d_0=\ker d_1$ (a box has $H^1=0$: the cycle space of a grid graph is generated by its $4$-cycles), hence a function $\tilde\Psi_Y(d_1A|_{\Omega^{**}})$ of the field strengths on $\Omega^{**}$. So $\langle Y\rangle=\E[\tilde\Psi_Y(F)]/\E[\tilde\Psi_1(F)]$ depends on the photon law only through the centred Gaussian law of $F=(d_1A(p))_{p\in\Omega^{**}}$, with covariance $\Sigma^{\ND}_B$ on the Neumann--Dirichlet side after the removal and $\Sigma^u_\Omega$ on the modified side (Step 0). Both laws live on $\operatorname{ran}d_1$. Let $T:\R^{\Omega^{**}}\to\R^{\Omega^*}$ be a linear map with $d_1T=\mathrm{id}$ on $\operatorname{ran}d_1$, extended linearly to all of $\R^{\Omega^{**}}$ (any extension serves, both laws being supported on $\operatorname{ran}d_1$), and realise each law as the photon Gaussian of $A''\defeq TF\in\R^{\Omega^*}$, covariance $T\Sigma T^T$: then $d_1A''=F$ almost surely, the gauge-invariant expectations are reproduced, and $T\Sigma^{\ND}_BT^T\le T\Sigma^u_\Omega T^T$ by (a). These covariances are degenerate; regularise to $T\Sigma T^T+\delta I$, $\delta>0$, which are non-degenerate, so the regularised states are BFS measures, the ordering of the covariances is preserved, and the law of $F''=d_1A''$ is $N(0,\Sigma+\delta d_1d_1^T)$, still ordered. \Cref{lem:s02:BFS}(a) by interpolation (the difference of the precision matrices is positive semidefinite by operator monotonicity of the inverse, hence a finite sum of terms $s|\ell(A)|^2$, $s\ge0$) gives $\langle Y\rangle_{\Sigma^u+\delta}\le\langle Y\rangle_{\Sigma^{\ND}+\delta}$ for $Y\in\cI$, and the reverse for $\cD$ by \cref{lem:s02:BFS}(d). \emph{The limit $\delta\downarrow0$.} For the $\cI$ half, $\tilde\Psi_Y$ is bounded and continuous: $|\Psi_Y(A)|\le\int e^{r^2(g)}e^{\sum\eps^{d-2}r_xr_y}\prod_xd\rho_x<\infty$ independently of $A$ ($|e^{\mathrm iF}|=1$; the hopping bound $e^{\eps^{d-2}r_xr_y}\le e^{\frac12\eps^{d-2}(r_x^2+r_y^2)}$ cancels the diagonals and $\lambda>0$ does the rest), and so is $\tilde\Psi_1$, with $\tilde\Psi_1>0$; the Gaussians $N(0,\Sigma+\delta d_1d_1^T)$ converge weakly to $N(0,\Sigma)$, so numerator and denominator converge and the inequality passes to the limit. For the $\cD$ half (not used below), $X=\Theta(f,-g)$ gives $\tilde\Psi_X(F)=e^{\langle v,F\rangle}\tilde\Psi^0_X(F)$ with $v$ the plaquette vector of $f$ and $0<\tilde\Psi^0_X\le e^{2\sW_\eps\sS(g)}\tilde\Psi_1\le C$ bounded and continuous, so $\tilde\Psi_X$ is continuous but unbounded; under $N(0,\Sigma+\delta d_1d_1^T)$ one has $\E[e^{2\langle v,F''\rangle}]=e^{2v^T(\Sigma+\delta d_1d_1^T)v}\le e^{2v^T\Sigma v+2\|d_1^Tv\|^2}$ for $\delta\le1$, so $\{\tilde\Psi_X(F'')\}_{0<\delta\le1}$ is bounded in $L^2$, hence uniformly integrable, and weak convergence with uniform integrability and continuity gives $\E[\tilde\Psi_X(F'')]\to\E[\tilde\Psi_X(F)]$. Altogether $\langle Y\rangle^{(3.38)}_\Omega\le\langle Y\rangle^{\text{after removal}}\le\langle Y\rangle^{\ND}_B$, and the reverse for $\cD$.
\end{proof}

\begin{remark}[why a margin, and why a fixed photon mass]
Part (a) of \cref{lem:s07:U} fails without a margin: for $B=\Omega$ the comparison is reversed (\cref{lem:s07:R}), and with a single layer ($N=1$) the difference $\Sigma^u_\Omega-\Sigma^{\ND}_B$ has a negative eigenvalue already at moderate $\mu$. No fixed margin works at $\lmu=0$: there $\Sigma^u_\Omega$ is the projector onto exact $2$-forms, while $\Var_B\langle g,dA\rangle>\|g\|^2$ for a Coulomb potential whose gradient leaks out of $B$. So \cref{lem:s07:U} is a statement at fixed $\lmu>0$ with $R(\lmu)\to\infty$ as $\lmu\to0$; a limit of vanishing photon mass inherits reflection positivity, if it exists, as a pointwise limit of positive semidefinite matrices. On one box the three states are ordered on $\cI$ as $\langle Y\rangle^{\ND}_\Omega\le\langle Y\rangle^{(3.38)}_\Omega\le\langle Y\rangle^{\ND}_B$ (\cref{lem:s07:R,lem:s07:U}), as the monotonicity of \cref{lem:s02:cut} requires. The sandwich with the boxes \eqref{eq:s07:margin} is assembled in the proof of \cref{thm:s08:OS2C}; for King's convention its lower end carries the kill step of \cref{rem:s04:F} and its upper end is unchanged (\cref{lem:s04:D}).
\end{remark}

% s08 Reflection positivity II: the limit, OS3 and the polynomial form
\section{Reflection positivity II: the limit, \texorpdfstring{$\OSth$}{OS3} and the polynomial form}

The lattice statements of the previous section are now carried to the limits $k\to\infty$ and $J\to\infty$. The scheme of proof is the same for the curvature sector and for the Wick tube: uniform local bounds for the modified-box functionals at complex coefficients (\cref{lem:s08:claim4}), Montel's theorem in $k$ at fixed $J$, the two-sided sandwich of \cref{lem:s07:L,lem:s07:U} on the cone as $J\to\infty$ with the identity theorem off it, and a translation for the planes that are not mid-planes. The sector-specific input is the convergence of the cube functionals themselves: the convergence at fixed cubes (\cref{def:s03:conv}; on King's convention \cref{hyp:s03:EC}) in the curvature sector, and on the Wick tube the positive-leg bound of \cref{cor:s11:W2prime} together with \cref{hyp:s03:W1} (\cref{prop:s08:T}). The arbitrary-amplitude form of the Wick sector, which needs \Wtwo, is the optional extension of \cref{thm:s23:OS2W}; the one lemma written for it here, \cref{lem:s08:W1}, is marked as such and is not used by the main theorem.

\subsection{The enlarged family}

\begin{definition}[the enlarged family, its tube and its bounded domain]\label{def:s08:enlarged}
Fix $\nu\in\{1,2,3\}$, a mid-plane $c\in D$ (\cref{def:s07:reflection}), and finite families $\zeta_1,\dots,\zeta_N$ of real sources $f_{il}\in C^\infty_c(\R^3;\Lambda^2)$, $h_{il}\in C^\infty_c(\R^3)$ supported in $\{x_\nu\ge c+\delta\}$ for some $\delta>0$, with complex coefficients $z_{il},w_{il}$. Split each scalar source as $h_{il}=h^+_{il}-h^-_{il}$ with $h^-_{il}\defeq\norm{h_{il}}_\infty\chi_{il}$, where $\chi_{il}\in C^\infty_c$, $0\le\chi_{il}\le1$, $\chi_{il}=1$ on $\supp h_{il}$ and $\supp\chi_{il}\subset\{x_\nu\ge c+\delta/2\}$, and $h^+_{il}\defeq h_{il}+h^-_{il}\ge0$. The \emph{enlarged family} is the family of real sources
\[
\cF\defeq\{\rho^*f_{il}\}\cup\{f_{il}\}\cup\{h^\pm_{il}\circ\rho\}\cup\{h^\pm_{il}\},\qquad\text{coefficient space }\C^m,\ m=m_F+m_h ,
\]
whose scalar members $\tilde h_j\ge0$ are all nonnegative. For a real coefficient vector $\xi\in\R^m$ the total source is $(f_\xi,k_1(\xi)-k_2(\xi))$ with $k_1(\xi)\defeq\sum_j(\xi_j)_+\tilde h_j$, $k_2(\xi)\defeq\sum_j(\xi_j)_-\tilde h_j$, both smooth and nonnegative ($x_\pm\defeq\max(\pm x,0)$). Let $J_0$ be the least $J$ such that every member of $\cF$ is supported at $\ell^\infty$-distance $>1$ from $\R^3\setminus C_J$; for $J\ge J_0$, $n=n(\lmu,c)$ and the boxes $\Omega_{J,k}$ of \eqref{eq:s07:margin} put
\[
\Phi_k(\zeta)\defeq Z^{a,(3.38)}_{k,\Omega_{J,k}}(\zeta),\qquad\zeta\in\C^m ,
\]
an entire function (a finite-dimensional integral with Gaussian decay in $A$ and $e^{-\lambda\eps^3r^4}$ decay in $r$). The \emph{matrix points} $\zeta^*_{ij}\in\C^m$ carry $\bar z_{il}$ on $\rho^*f_{il}$, $z_{jl}$ on $f_{jl}$, $\pm\bar w_{il}$ on $h^\pm_{il}\circ\rho$, $\pm w_{jl}$ on $h^\pm_{jl}$ and $0$ elsewhere; their total source is that of $\theta\zeta_i^*\sqcup\zeta_j$, so $\Phi_k(\zeta^*_{ij})$ is the $(i,j)$ entry of \cref{lem:s07:RPbox}(b). Since reflected and unreflected members have disjoint supports, $\norm{h^+}_\infty\le2\norm h_\infty$ and $\norm{h^-}_\infty=\norm h_\infty$,
\begin{equation}\label{eq:s08:ampbound}
\norm{k_1(\Re\zeta^*_{ij})}_\infty\ \le\ 2\max\big(\amp(\zeta_i),\amp(\zeta_j)\big).
\end{equation}
\begin{enumerate}[label=\textup{(\roman*)},leftmargin=*,nosep]
\item The \emph{Wick tube} of $\cF$ is $T_\cF\defeq\{\zeta\in\C^m:\ 2\norm{k_1(\Re\zeta)}_\infty<\sigma_+/2\}$, with $\sigma_+=s_1$ the amplitude of \Itwo\ (\cref{def:s10:I2}, \eqref{eq:s11:constants}). It is open, convex ($\xi\mapsto\norm{k_1(\xi)}_\infty$ is a supremum of convex functions), connected, invariant under $z\mapsto-z$ on the curvature coordinates, contains $0$, and by \eqref{eq:s08:ampbound} contains every $\zeta^*_{ij}$ as soon as $\amp(\zeta_i)<\sigma_+/8=:\sigma_W$ for every $i$.
\item The \emph{bounded domain}: let $B_1\subset\R^m$ be the open convex hull of the unit $\ell^\infty$-balls about $0$ and about the real parts $\Re\zeta^*_{ij}$; let $\beta_j\defeq2\sup_{\xi\in B_1}|\xi_j|$ for the scalar coordinates (so $\beta_j=2(|\Re w_{il}|+1)$ for the coordinate of $h^\pm_{il}$ or of $h^\pm_{il}\circ\rho$), $B''\defeq\prod_j[0,\beta_j]$ over the scalar coordinates, and $q_v\defeq\sum_jv_j\tilde h_j$ for $v\in\vertex(B'')$, the \emph{vertex sources}. All $q_v$ are pointwise $\le q_\zeta\defeq\sum_{i,l}2(|\Re w_{il}|+1)(h^+_{il}+h^-_{il}+h^+_{il}\circ\rho+h^-_{il}\circ\rho)$. Put $\cT_1\defeq\{\zeta\in\C^m:\Re\zeta\in B_1\}$: open, convex, connected, containing every $\zeta^*_{ij}$ and the unit polydisc about $0$, so that $\cT_1\cap T_\cF$ is a non-empty open set. For $\xi\in B_1$ both $2\xi_+$ and $\xi_+$ lie in $B''$.
\item The \emph{uniqueness set} $U\subset\C^m$ is the set of coefficient vectors whose curvature coefficients are purely imaginary and whose scalar coefficients are positive reals. It is a non-empty, relatively open subset of the totally real subspace $(\mathrm i\R)^{m_F}\times\R^{m_h}$ of maximal real dimension $m$ (it is not open in $\C^m$), hence a set of uniqueness for holomorphic functions on any connected open set meeting it: a holomorphic function vanishing on a relatively open piece of a maximal totally real subspace has all derivatives zero there \cite{Hor90}. $U\cap\cT_1\ne\emptyset$ and $U\cap T_\cF\ne\emptyset$ (small positive scalar coefficients). On $U$ the total source is $(\mathrm if_\zeta,g_\zeta)\in\cK$, so $\Theta(\zeta)\in\cI$; the scheme multipliers are positive reals, equal for the Neumann--Dirichlet and the modified states; and every lattice value is real (charge conjugation, \cref{def:s02:sources}: $\overline{Z^a(\mathrm if,g)}=Z^a(-\mathrm if,g)=Z^a(\mathrm if,g)$).
\end{enumerate}
Every lattice functional, and hence every limit below, depends on a family only through its total sources; the splitting is therefore immaterial at the matrix points and at the Taylor coefficients at $0$.
\end{definition}

\subsection{Uniform local bounds}

The following local bound is used only in the optional extension (\cref{thm:s23:OS2W}); it is stated here because its domain $\cT_1$ is what makes ``finitely many instances of \Wtwo'' literally true there. The fixed-cube limits $Z^a_J\defeq\lim_kZ^a_{k,J}$ exist for all complex coefficients and are entire under \Wtwo(a) alone: it gives convergence at the real points, the vertex argument of \cref{lem:s08:claim4} with \cref{lem:s02:alg}(d) gives local boundedness in $k$, and Montel's theorem with the uniqueness set $\R^m$ does the rest.

\begin{lemma}[uniform local bound; optional extension]\label{lem:s08:W1}
Assume \Wtwo(a) and \Wtwo(b) at the finitely many sources $q_v$, $v\in\vertex(B'')$. Then $\{Z^a_J\}_{J\ge J_0}$ is uniformly bounded on $\cT_1$, $Z^a_J\to Z^a_\infty$ locally uniformly on $\cT_1$, $Z^a_\infty$ is holomorphic on $\cT_1$, and $Z^a_J\le Z^a_\infty$ on $U\cap\cT_1$.
\end{lemma}

\begin{proof}
For real $\xi\in B_1$, \cref{lem:s02:alg}(f) gives $Z^a_J(\xi)\le Z^a_J(2f_\xi,-2k_2(\xi))^{1/2}Z^a_J(0,2k_1(\xi))^{1/2}$. The first factor is nonincreasing in $J$: at fixed $k$ this is the second inequality of \cref{lem:s02:cut} with \cref{lem:s02:alg}(a), and it passes to $k\to\infty$ by \Wtwo(a); \Wtwo(b) is not used there. Hence it is $\le Z^a_{J_0}(2f_\xi,-2k_2(\xi))$, and $\xi\mapsto\ln Z^a_{J_0}(2f_\xi,-2k_2(\xi))$ is bounded on the compact $\bar B_1$: the coefficient map is piecewise linear, so the image of $\bar B_1$ lies in a real box $B'$, on which $\ln Z^a_{J_0}$ is convex (\cref{lem:s02:alg}(d)) and so bounded by its maximum over the vertices of $B'$, each finite by \Wtwo(a). For the second factor, $2\xi_+\in B''$ and $\ln Z^a_J(0,\cdot)$ is convex in the coefficients, so $Z^a_J(0,2k_1(\xi))\le\max_{v\in\vertex(B'')}Z^a_J(0,q_v)\le\max_v\sup_{J'\ge J_{q_v}}Z^a_{J'}(0,q_v)=: C_{B''}<\infty$ by \Wtwo\ at the $q_v$ (note $J\ge J_0\ge J_{q_v}$). By \cref{lem:s02:alg}(d), $|Z^a_J(\zeta)|\le Z^a_J(\Re\zeta)$, so $\sup_{J\ge J_0}\sup_{\cT_1}|Z^a_J|<\infty$. On $U\cap\cT_1$ the sources lie in $\cK$, so $J\mapsto Z^a_J(\zeta)$ is nondecreasing (first inequality of \cref{lem:s02:cut} at fixed $k$, then $k\to\infty$) and bounded, hence convergent. $U\cap\cT_1$ is a set of uniqueness for holomorphic functions on the connected $\cT_1$ (\cref{def:s08:enlarged}(iii)); Montel's theorem gives locally uniform convergence on $\cT_1$ to a holomorphic limit, which is $Z^a_\infty$ wherever the latter is defined (the same $J$-limit), and the monotone limit on $U\cap\cT_1$ gives $Z^a_J\le Z^a_\infty$ there. Under \Wtwo(b) for every $0\le q\in C^\infty_c$ the same holds on all of $\C^m$.
\end{proof}

\begin{lemma}[uniform local bounds for the modified-box functionals]\label{lem:s08:claim4}
Let $\cF$, $\Phi_k$, $\cT_1$, $T_\cF$ be as in \cref{def:s08:enlarged}.
\begin{itemize}[leftmargin=*,nosep]
\item[(A)] \textup{(Optional extension.)} Under \Wtwo(a) and \Wtwo(b) at the $q_v$, $v\in\vertex(B'')$: there is $C<\infty$, independent of $k$ and of $J\ge J_0$, with $\limsup_k\sup_{\cT_1}|\Phi_k|\le C$; and at fixed $J$, $\sup_k\sup_Q|\Phi_k|<\infty$ for every compact $Q\subset\C^m$ (this part uses \Wtwo(a) only).
\item[(A$_\infty$)] \textup{(Optional extension.)} Under \Wtwo\ for every $0\le q\in C^\infty_c$: the same with $\cT_1$ replaced by any compact $Q\subset\C^m$, with $C=C_Q$.
\item[(T)] \textup{(The tube.)} Under \cref{hyp:s03:W1}, \Itwo, \Ithree\ and the correlation inequalities of \cref{lem:s02:BFS}(a)--(e) only (no \Wtwo): the same for every compact $Q\subset T_\cF$ and all $J\ge J_0$, with no further threshold on $J$.
\item[(C)] \textup{(Curvature sector, $h\equiv0$; no hypothesis.)} $|\Phi_k(z)|\le\langle e^{F(f_{\Re z})}\rangle^{(3.38)}_{k,\Omega_{J,k}}\le\exp(\tfrac12\Var^{(3.38)}F(f_{\Re z}))$ and $\Var^{(3.38)}F(f)\le\sum_p\eps^3f_{c(p)}^2\to\norm f_{L^2}^2$, uniformly in $k$ and $J$.
\end{itemize}
\end{lemma}

\begin{proof}
$|\Phi_k(\zeta)|\le\Phi_k(\Re\zeta)$ by the proof of \cref{lem:s02:alg}(d) ($|\Theta(\zeta)|=\Theta(\Re\zeta)$ and the modulus of the multiplier), which is valid for any probability measure with the same real multiplier, in particular for the modified box state, which is not one of the states of \cref{def:s02:acube,def:s02:Kcube}. At a real $\xi$ the proof of \cref{lem:s02:alg}(f) (Cauchy--Schwarz with the multiplier identity, again valid for any such measure) gives
\begin{equation}\label{eq:s08:split}
\Phi_k(\xi)\ \le\ Z^{a,(3.38)}_{k,\Omega_{J,k}}(2f_\xi,-2k_2(\xi))^{1/2}\;Z^{a,(3.38)}_{k,\Omega_{J,k}}(0,2k_1(\xi))^{1/2}.
\end{equation}
\emph{First factor} (class $\cD$). By the $\cD$ half of \cref{lem:s07:L} with $C=C_J\cap\Lk$ (its hypothesis is \eqref{eq:s07:margin}(i)) it is $\le Z^a_{k,J}(2f_\xi,-2k_2)\le Z^a_{k,J_0}(2f_\xi,-2k_2)$, the second step being the second inequality of \cref{lem:s02:cut} at fixed $k$. The map $\xi\mapsto\eta(\xi)\defeq$ (coefficient vector of $(2f_\xi,-2k_2(\xi))$) is continuous, so $\eta(\Re Q)$, resp.\ $\eta(\bar B_1)$, lies in a real box $B'$; $\eta\mapsto\ln Z^a_{k,J_0}(\eta)$ is convex (\cref{lem:s02:alg}(d)), so $\sup_\xi Z^a_{k,J_0}(\eta(\xi))\le\max_{v\in\vertex(B')}Z^a_{k,J_0}(v)$, finitely many values, each convergent in $k$ by \cref{hyp:s03:W1} (the sources here have negative Wick coefficients, on which the theorem has no amplitude restriction; in versions (A), (A$_\infty$) by \Wtwo(a)), and convergence in $k$ implies boundedness in $k$. So the first factor is bounded uniformly in $k$, $\xi$ and $J\ge J_0$.
\emph{Second factor} (class $\cI$). By the $\cI$ half of \cref{lem:s07:U} with $B=C_{J+2n}\cap\Lk$ and $N\defeq\lceil R(\lmu)/\eps_k\rceil$ (hypothesis \eqref{eq:s07:margin}(ii)) it is $\le Z^a_{k,J+2n}(0,2k_1(\xi))$. (A): for $\xi\in B_1$, $2\xi_+\in B''$, so by convexity it is $\le\max_{v\in\vertex(B'')}Z^a_{k,J+2n}(0,q_v)$; each of these finitely many values converges in $k$ by \Wtwo(a) (hence is bounded in $k$, which is the fixed-$J$ bound Montel's theorem needs, and the $\limsup_k$ bound passes to any subsequential limit) with limit $Z^a_{J+2n}(0,q_v)\le\sup_{J'\ge J_{q_v}}Z^a_{J'}(0,q_v)<\infty$: this is \Wtwo\ at $q_v$, used directly ($J+2n\ge J_0\ge J_{q_v}$), uniformly in $J$; the finitely many instances $q_v$ are all of \Wtwo\ that version (A) uses. (A$_\infty$): the same with a box $B''\supset2(\Re Q)_+$ depending on $Q$. (T): for $\xi\in\Re Q$, $Q\subset T_\cF$, one has $0\le2k_1(\xi)\le\sigma_+/2\le s_1$ on the lattice (the sampled supremum is at most the continuum one), so \cref{cor:s11:W2prime}, $\ln N^a_{k,C_J}(0,q)\le\kappa\norm q_{1,k}$ for $0\le q\le s_1$ on $C_J\cap\Lk$, uniformly in $k$ and $J$, applied with $N^a_{k,X}(0,q)=Z^a_{k,X}(0,q)$ (the same functional with the same scheme multiplier, so no $J$-dependent factor intervenes), gives $Z^a_{k,J+2n}(0,2k_1(\xi))\le\exp(\kappa\norm{2k_1(\xi)}_{1,k})\le\exp\big(2\kappa\sum_j|\xi_j|\sup_{k'}\norm{\tilde h_j}_{1,k'}\big)$, bounded uniformly in $k$, $J$ and $\xi\in\Re Q$ because $\norm{\tilde h_j}_{1,k'}\to\norm{\tilde h_j}_{L^1}$ (Riemann sums, \eqref{eq:s02:riemann}). No lower bound on $J$ beyond $J_0$ enters: \cref{cor:s11:W2prime} holds for every $k$ and $J$.
\emph{(C).} At $h\equiv0$ the first inequality is \cref{lem:s02:alg}(d) as above and the second is \cref{lem:s02:GD} with $s=1$, the modified box state being a BFS measure (\cref{def:s07:box}). Its photon Gaussian has the unitary form for gauge-invariant observables (Step 1 of \cref{lem:s07:RPbox}), and $\Var^u\langle v,d_1A\rangle=\eps^{-1}\langle d_1^Tv,(d_1^Td_1+\lmu^2\eps^2)^{-1}d_1^Tv\rangle\le\eps^{-1}\norm v^2$; with $v_p=\eps^2f_{c(p)}$ this is $\Var F(f)\le\sum_p\eps^3f_{c(p)}^2$.
\end{proof}

\begin{remark}[one instance of \Wtwo\ in place of the vertex set]\label{rem:s08:oneinstance}
For $v\in\vertex(B'')$ one has $0\le q_v\le q_\zeta$ pointwise, and for $p=1+\eta>1$, $p'=p/(p-1)$, $q_v=\frac1p(pq_\zeta)-\frac1{p'}\big(p'(q_\zeta-q_v)\big)$. H\"older's inequality with the multiplier bookkeeping $e^{-a_k\sS(q_v)}=(e^{-a_k\sS(pq_\zeta)})^{1/p}(e^{a_k\sS(p'(q_\zeta-q_v))})^{1/p'}$ gives, for every state, $Z^a(0,q_v)\le Z^a(0,pq_\zeta)^{1/p}Z^a(0,-p'(q_\zeta-q_v))^{1/p'}$. The second factor is of class $\cD$, nonincreasing in $J$ (\cref{lem:s02:cut}, then $k\to\infty$ by \Wtwo(a)), hence $\le Z^a_{J_0}(0,-p'(q_\zeta-q_v))<\infty$, one real point of \cref{hyp:s03:W1}; the first is bounded uniformly in $J$ by \Wtwo(b) at the single source $(1+\eta)q_\zeta$. So \Wtwo\ at $(1+\eta)q_\zeta$, for any one $\eta>0$, implies the finitely many instances used in \cref{lem:s08:W1,lem:s08:claim4}(A). Only this direction is proved; no converse is claimed. Pointwise domination $q_v\le q_\zeta$ alone gives nothing, since $Z^a(0,\cdot)$ is not monotone in $q$: the Wick square has no sign, and the scheme multiplier $e^{-a_k\sS(q)}$ has no sign control.
\end{remark}

\subsection{King's state on the Wick tube}

\begin{proposition}[the Wick-tube object: existence and translations]\label{prop:s08:T}
Let $a$ be an admissible scheme (\cref{def:s03:adm}), and assume \Itwo\ and \Ithree. For a finite family $\cF'$ of real sources with nonnegative scalar members $\tilde h_j\ge0$ (any family is brought to this form by the splitting of \cref{def:s08:enlarged}) let $T_{\cF'}$ be its Wick tube, as in \cref{def:s08:enlarged}(i).
\begin{enumerate}[label=\textup{(\alph*)},leftmargin=*,nosep]
\item $Z^a_J\to Z^a_{\infty,T}$ locally uniformly on $T_{\cF'}$ as $J\to\infty$; $Z^a_{\infty,T}$ is holomorphic on $T_{\cF'}$; $Z^a_{\infty,T}=\Ginf$ on $T_{\cF'}\cap\Dtube$ and $Z^a_{\infty,T}=Z^a_\infty$ wherever the latter exists (the same $J$-limit of the same $Z^a_J$); $|Z^a_{\infty,T}(\zeta)|\le Z^a_{\infty,T}(\Re\zeta)$; $Z^a_{\infty,T}(\xi)>0$ for real $\xi\in T_{\cF'}$ with $-\xi\in T_{\cF'}$; and $Z^a_{\infty,T}(-z,w)=Z^a_{\infty,T}(z,w)$.
\item \textup{(Translations.)} $Z^a_{\infty,T}(\tau_b\cdot\zeta)=Z^a_{\infty,T}(\zeta)$ for all $b\in\R^3$ and $\zeta\in T_{\cF'}$, where $\tau_b$ acts on all sources; the translated family has the same tube.
\end{enumerate}
The proof uses, besides \cref{hyp:s03:W1}, only \cref{cor:s11:W2prime} (hence \Itwo, \Ithree\ and \cite[Thm.~6.2(1),(2)]{BFS1}), the second inequality of \cref{lem:s02:cut} (hence \cref{lem:s02:BFS}(d)), \cref{lem:s02:alg,lem:s02:baire}, Montel's theorem and the identity theorem; \Wtwo\ is not used.
\end{proposition}

\begin{proof}
(a) \emph{Local boundedness} of $\{Z^a_J\}_{J\ge J_0}$ on $T_{\cF'}$: $|Z^a_J(\zeta)|\le Z^a_J(\Re\zeta)\le Z^a_J(2f,-2k_2)^{1/2}Z^a_J(0,2k_1)^{1/2}\le Z^a_{J_0}(2f,-2k_2)^{1/2}\exp(\kappa\norm{k_1}_{L^1})$, by \cref{lem:s02:alg}(d),(f), the nonincreasing $\cD$-factor (second inequality of \cref{lem:s02:cut} at fixed $k$, then $k\to\infty$), and \cref{cor:s11:W2prime} in the limit $k\to\infty$ for the $\cI$-factor, valid since $2\norm{k_1}_\infty<\sigma_+/2\le s_1$; the first factor is bounded on compacts by the vertex argument of \cref{lem:s08:claim4} with \cref{hyp:s03:W1} at $C_{J_0}$. \emph{Convergence on $U\cap T_{\cF'}$}: there the values are real, nondecreasing in $J$ (first inequality of \cref{lem:s02:cut}, then $k\to\infty$) and bounded, since $Z^a_J(\mathrm if_\zeta,g_\zeta)\le Z^a_J(0,g_\zeta)$ by \cref{lem:s02:alg}(d) and $Z^a_J(0,g_\zeta)\le\exp(\kappa\norm{g_\zeta}_{L^1})$ by \cref{cor:s11:W2prime} in the limit $k\to\infty$ ($0\le g_\zeta\le\sigma_+/4\le s_1$ there). Montel's theorem on the connected open $T_{\cF'}$ with the uniqueness set $U\cap T_{\cF'}$ gives locally uniform convergence to a holomorphic $Z^a_{\infty,T}$; on $T_{\cF'}\cap\Dtube$ this is the limit $\Ginf$ of \cref{thm:s11:W}(i), the same argument on the other tube. The bound $|Z^a_{\infty,T}(\zeta)|\le Z^a_{\infty,T}(\Re\zeta)$ and the evenness in $z$ pass from the lattice (\cref{lem:s02:alg}(d); the invariance $(A,\vartheta)\mapsto(-A,-\vartheta)$ of every state (charge conjugation, \cref{def:s02:sources}), which maps $F(f)\mapsto-F(f)$ and fixes $r$; $T_{\cF'}$ is invariant under $z\mapsto-z$). Positivity at symmetric real points is \cref{lem:s02:alg}(e), $Z^a(\xi)Z^a(-\xi)\ge1$, in the limit.

(b) \emph{Lattice translations.} Let $b\in5^{-k_1}\Z^3$, $|b|_\infty\le M$, $k\ge k_1$. $\tau_b$ preserves $\Lk$ and $\sS_{\Lk}$, so $Z^a_{k,J}(\tau_b\cdot\zeta)=Z^a_{k,(C_J-b)\cap\Lk}(\zeta)$ (the action is translation covariant on $\Lk$). As sets of sites, $C_{J-1}\cap\Lk\subset(C_J-b)\cap\Lk\subset C_{J+1}\cap\Lk$, so the first inequality of \cref{lem:s02:cut} (stated for any finite site set) gives, for $\zeta\in U$ and $J\ge J_0+1$ (so that every plaquette and site at which the sources do not vanish lies in $C_{J-1}$), $Z^a_{k,J-1}(\zeta)\le Z^a_{k,J}(\tau_b\cdot\zeta)\le Z^a_{k,J+1}(\zeta)$. Let $k\to\infty$ and then $J\to\infty$ on $U\cap T_{\cF'}$, where the outer limits converge by (a): $Z^a_{\infty,T}(\tau_b\cdot\zeta)=Z^a_{\infty,T}(\zeta)$ there, hence on all of $T_{\cF'}$ by the identity theorem (both sides are holomorphic on the connected $T_{\cF'}$, the translated family having the same supremum norms). Iterating, the identity holds for all $b\in\bigcup_k5^{-k}\Z^3$.
\emph{All $b\in\R^3$.} For $\rho>0$ let $E_\rho\defeq C^\infty_c(\bar B_\rho;\Lambda^2)\times C^\infty_c(\bar B_\rho)$ (a Fr\'echet space) and $O_\rho\defeq\{(f,h)\in E_\rho:\ 4\norm h_\infty<\sigma_+/2\}$, open, convex and symmetric under $s\mapsto-s$. For $s=(f,h)\in O_\rho$ the one-member family with its splitting has the coefficient vector $(1,1,-1)$ in its tube ($2\norm{k_1}_\infty=2\norm{h^+}_\infty\le4\norm h_\infty<\sigma_+/2$), so $\Phi_T(s)\defeq\ln Z^a_{\infty,T}(f,h)$ is defined and finite on $O_\rho$ by (a). $\Phi_T$ is convex (a pointwise limit of the functions $\ln Z^a_{k,J}$, convex by H\"older's inequality) and Borel (a pointwise limit of continuous functions of finitely many sampled values). By \cref{lem:s02:baire}, $\Phi_T$ is continuous on $O_\rho$. For $s\in O_\rho$ and $b\in\R^3$ pick $b_n\to b$ in $\bigcup_k5^{-k}\Z^3$; then $\tau_{b_n}s,\tau_bs\in O_{\rho+|b|+1}$, $\tau_{b_n}s\to\tau_bs$ in $E_{\rho+|b|+1}$ (translations act continuously on $C^\infty_c$), and $\Phi_T(\tau_{b_n}s)=\Phi_T(s)$; so $\Phi_T(\tau_bs)=\Phi_T(s)$.
\emph{Complex coefficients.} For a family $\cF'$ the two holomorphic functions $\zeta\mapsto Z^a_{\infty,T}(\tau_b\cdot\zeta)$ and $\zeta\mapsto Z^a_{\infty,T}(\zeta)$ on $T_{\cF'}$ agree at the real points of $\{\xi\in\R^m:4\norm{h_\xi}_\infty<\sigma_+/2\}\cap T_{\cF'}$, where both limits exist and both equal the one-member value $\exp\Phi_T(f_\xi,h_\xi)$ (the same $J$-limit of the same $Z^a_J$). This intersection is open in $\R^m$ and contains $0$ (the first set alone need not lie in the tube: with $\tilde h_1=\tilde h_2$ and $\xi=(1,-1)$, $h_\xi=0$ while $k_1=\tilde h_1$), hence is a set of uniqueness in the connected $T_{\cF'}$; so the two functions agree on $T_{\cF'}$.
\end{proof}

\subsection{The two limits and the passage to every plane}

Throughout this subsection $\cF$, $\Phi_k$, $\zeta^*_{ij}$, $\cT_1$, $T_\cF$, $U$ are those of \cref{def:s08:enlarged} for a mid-plane $c\in D$, and ``version'' refers to \cref{lem:s08:claim4}: (C) is the curvature sector when the comparison family is convergent at fixed cubes (the uniform local bound being \cref{thm:s05:C}(b)), (T) the tube under \cref{hyp:s03:W1}, \Itwo, \Ithree, and (A), (A$_\infty$) the optional extension under \Wtwo.

\emph{Step 5 ($k\to\infty$ at fixed $J$).} By the fixed-$J$ bound of \cref{lem:s08:claim4} and Montel's theorem there is a subsequence $k_m\to\infty$ along which $\Phi_{k_m}\to\Phi^J$ locally uniformly on $\C^m$, with $\Phi^J$ entire. No convergence of the full sequence is claimed or needed. By \cref{lem:s07:RPbox}(b), $\sum_{ij}\bar\gamma_i\gamma_j\Phi_k(\zeta^*_{ij})\ge0$ for every $k\ge k_0(c,\delta)$ and $\gamma\in\C^N$; positive semidefiniteness of a Hermitian matrix is closed under entrywise limits, so $\sum_{ij}\bar\gamma_i\gamma_j\Phi^J(\zeta^*_{ij})\ge0$. The uniform bounds of \cref{lem:s08:claim4} pass to $\Phi^J$: $\sup_{\cT_1}|\Phi^J|\le C$ in (A), $\sup_Q|\Phi^J|\le C_Q$ in (A$_\infty$), (T) and (C), uniformly in $J\ge J_0$, whatever subsequence was taken. (The box $\Omega_{J,k}$ depends on $J$, so $\Phi^J$ does; the subsequence may depend on $J$.)

\emph{Step 6 ($J\to\infty$: sandwich on the cone, identity theorem off it).} By the $\cI$ halves of \cref{lem:s07:L,lem:s07:U}, whose hypotheses are \eqref{eq:s07:margin}(i),(ii), for every $k$ and every $\zeta\in U$,
\begin{equation}\label{eq:s08:sandwich}
Z^a_{k,J}(\zeta)\ \le\ \Phi_k(\zeta)\ \le\ Z^a_{k,J+2n}(\zeta).
\end{equation}
Let $k\to\infty$ along the subsequence of Step 5. The outer terms converge along the full sequence (version (C): \cref{thm:s05:C}(b); (T), (A), (A$_\infty$): \cref{hyp:s03:W1}, resp.\ \Wtwo(a), with Montel's theorem), so $Z^a_J\le\Phi^J\le Z^a_{J+2n}$ on $U$ for every subsequential $\Phi^J$. As $J\to\infty$ the outer terms converge on the respective uniqueness set: in (C) on $U=(\mathrm i\R)^{m_F}$ to $Z_\infty(\cdot,0)$ (\cref{thm:s05:C}(b)); in (T) on $U\cap T_\cF$ to $Z^a_{\infty,T}$ (\cref{prop:s08:T}(a); there $J\mapsto Z^a_J(\zeta)$ is real, nondecreasing and bounded by $Z^a_J(0,g_\zeta)\le\exp(\kappa\norm{g_\zeta}_{L^1})$, since $0\le g_\zeta\le\sigma_+/4\le s_1$); in (A) on $U\cap\cT_1$ to $Z^a_\infty$ (\cref{lem:s08:W1}); in (A$_\infty$) on $U$ to $Z^a_\infty$ (\cref{lem:s08:W1} with \Wtwo(b) for every $q$). Hence $\Phi^J\to Z^a_{\infty(,T)}$ pointwise on that set. By Step 5 the family $\{\Phi^J\}_J$ is locally bounded on the corresponding connected domain ($\C^{m}$, $T_\cF$, $\cT_1$, $\C^m$), on which the limit functional is holomorphic (\cref{thm:s05:C}(b), \cref{prop:s08:T}(a), \cref{lem:s08:W1}, \cref{lem:s08:W1} on $\C^m$ respectively); by Montel's theorem every subsequence of $(\Phi^J)$ has a further subsequence converging locally uniformly to a holomorphic function, which agrees with the limit functional on the uniqueness set and hence on the whole domain \cite{Hor90}. So $\Phi^J\to Z^a_{\infty(,T)}$ locally uniformly on the domain, in particular at the finitely many matrix points $\zeta^*_{ij}$ (which lie in $\cT_1$ by construction and in $T_\cF$ by \eqref{eq:s08:ampbound} when every $\amp(\zeta_i)<\sigma_W$), where $\Phi^J(\zeta^*_{ij})\to Z^a_{\infty(,T)}(\theta\zeta_i^*\sqcup\zeta_j)$ because the functionals depend on the family only through the total sources. Positivity passes to the limit:
\begin{equation}\label{eq:s08:midplane}
\sum_{i,j}\bar\gamma_i\gamma_j\,Z^a_{\infty(,T)}(\theta\zeta_i^*\sqcup\zeta_j)\ \ge\ 0\qquad\text{for the plane }\{x_\nu=c\},\ c\in D .
\end{equation}
A side result: the modified-box values converge to the cube state, $\lim_J\lim_{k_m}Z^{a,(3.38)}_{k_m,\Omega_{J,k_m}}(\zeta)=Z^a_{\infty(,T)}(\zeta)$, on the respective domain and along any $k$-subsequence on which the fixed-$J$ limits exist; this is the corrected form of King's passage from the modified boxes to the cube limit \cite[(3.42)]{K2}, with the margin $2n$ in place of one cube step (\cref{lem:s07:R,rem:s07:K8}). The reversal of King's second inequality is proved, weakly at every coupling and strictly at $e=0$ (\cref{lem:s07:R}); no step of the passage to the limit uses that inequality in King's direction.

\emph{Step 7 (every plane).} Let $c\in\R$ be arbitrary and the sources supported in $\{x_\nu>c\}$. If $c\in D$, \eqref{eq:s08:midplane} is the result. Otherwise pick $c'\in D$ and $b\defeq(c'-c)e_\nu$. The translated sources $\tau_b\zeta_i$ are supported in $\{x_\nu>c'\}$, and $\theta_{c'}(\tau_b\zeta)^*=\tau_b(\theta_c\zeta^*)$, because $\rho_{c'}(x+b)=\rho_c(x)+b$ for this $b$ and the linear parts of $\rho_c,\rho_{c'}$ agree. By \eqref{eq:s08:midplane} for the plane $c'$ and the translation invariance of the limit functional,
\[
\sum\bar\gamma_i\gamma_jZ^a_{\infty(,T)}(\theta_c\zeta_i^*\sqcup\zeta_j)=\sum\bar\gamma_i\gamma_jZ^a_{\infty(,T)}\big(\tau_b(\theta_c\zeta_i^*\sqcup\zeta_j)\big)=\sum\bar\gamma_i\gamma_jZ^a_{\infty(,T)}\big(\theta_{c'}(\tau_b\zeta_i)^*\sqcup\tau_b\zeta_j\big)\ \ge\ 0 .
\]
The translation invariance is supplied in (C) by \cref{prop:s06:translations}, and in (T) by \cref{prop:s08:T}(b) on $T_\cF$ (the amplitude $\amp$ is translation invariant, so the matrix points stay in the tube). In the optional extension a plane $c\notin D$ needs a non-lattice translation, which finitely many instances of \Wtwo(b) do not supply; there the step is carried by \cref{prop:s08:T}(b) together with \cref{lem:s08:W1} for $\cF$ and for $\tau_b\cF$ and the identity theorem on $\cT_1$ (\cref{thm:s23:OS2W}).

\subsection{The theorems}

\begin{theorem}[reflection positivity, curvature sector, comparison convention]\label{thm:s08:OS2C}
Assume that the comparison family is convergent at fixed cubes (\cref{def:s03:conv}). Let $\lmu>0$, $\nu\in\{1,2,3\}$, $c\in\R$, and let $\zeta_1,\dots,\zeta_N$ be finite families of real $2$-form sources $f_{il}\in C^\infty_c(\R^3;\Lambda^2)$ supported in the open half-space $\{x_\nu>c\}$, with arbitrary complex coefficients $z_{il}$ and no scalar sources. Then for all $\gamma\in\C^N$
\[
\sum_{i,j=1}^N\bar\gamma_i\gamma_j\,Z_\infty\big(\theta\zeta_i^*\sqcup\zeta_j,\,0\big)\ \ge\ 0 .
\]
Besides the convergence at fixed cubes the proof uses only \cite[Thm.~6.2(1),(2)]{BFS1} (through the first inequality of \cref{lem:s02:cut} and the $\cI$ halves of \cref{lem:s07:L,lem:s07:U}), \cref{thm:s05:C}(b) with \cref{lem:s05:C2}, and \cref{lem:s02:GD}, \cref{lem:s02:alg}(d),(e), \cref{lem:s02:BFS}(e), \cref{lem:s02:baire}; no inequality of class $\cD$ enters.
\end{theorem}

\begin{proof}
At $h\equiv0$ there is no Wick constant and no scheme, and $\cF$ consists of the $f_{il}$ and $\rho^*f_{il}$ with $U=(\mathrm i\R)^{m}$. What the argument of Steps 5--7 needs of the cube functionals is: the existence of $Z_J(\cdot,0)=\lim_kZ_{k,J}(\cdot,0)$ locally uniformly on $\C^m$, entire; the monotonicity of $J\mapsto Z_J(\mathrm if,0)$ on the cone (first inequality of \cref{lem:s02:cut}); the existence of $Z_\infty(\cdot,0)=\lim_JZ_J(\cdot,0)$, entire, locally uniformly on $\C^m$; and translation invariance. The first three are \cref{thm:s05:C}(b), whose local bound is the $\lmu$-free bound of \cref{thm:s05:C}(a) and \cref{lem:s05:C2} rather than a local bound resting on the nonincreasing clause of \cref{lem:s05:TWC} and so on \cite[Thm.~6.2(3),(4)]{BFS1}; the fourth is \cref{prop:s06:translations}, which uses \cref{lem:s02:baire} with $\ln Z_\infty$ finite on all of $E_\rho$ (upper bound from \cref{thm:s05:C}(b), lower bound \cref{lem:s02:alg}(e)) and the identity theorem. The lattice steps \cref{lem:s07:RPbox,lem:s07:L,lem:s07:U} include $g=0$, and only their $\cI$ halves are used; \cref{lem:s08:claim4}(C) is the fixed-$J$ and $J$-uniform bound; Steps 5--7 in version (C) give the claim.
\end{proof}

\begin{theorem}[reflection positivity on the Wick tube, comparison convention]\label{thm:s08:OS2WT}
Let $a$ be an admissible scheme, and assume \Itwo\ and \Ithree\ (the hypotheses of \cref{cor:s11:W2prime}); the $\cD$-class steps use in addition \cref{lem:s02:BFS}(d). Let $\sigma_+=s_1$ be the amplitude of \Itwo\ and $\sigma_W\defeq\sigma_+/8$. Let $\lmu>0$, $\nu$, $c\in\R$, and let $\zeta_1,\dots,\zeta_N$ be finite families of real sources $f_{il}\in C^\infty_c(\R^3;\Lambda^2)$, $h_{il}\in C^\infty_c(\R^3)$ supported in $\{x_\nu>c\}$, with complex coefficients $z_{il},w_{il}$ such that $\amp(\zeta_i)<\sigma_W$ for every $i$; in particular, all purely imaginary $w_{il}$ with arbitrary $z_{il}$ are admitted. Then the values $Z^a_{\infty,T}(\theta\zeta_i^*\sqcup\zeta_j)=\lim_JZ^a_J(\theta\zeta_i^*\sqcup\zeta_j)$ exist (\cref{prop:s08:T}) and
\[
\sum_{i,j=1}^N\bar\gamma_i\gamma_j\,Z^a_{\infty,T}\big(\theta\zeta_i^*\sqcup\zeta_j\big)\ \ge\ 0\qquad\text{for all }\gamma\in\C^N .
\]
Where $Z^a_\infty$ is defined, $Z^a_{\infty,T}=Z^a_\infty$ there. No lower bound on $J$ beyond $J_0$ enters, and \Wtwo\ is not used.
\end{theorem}

\begin{proof}
With the splitting of \cref{def:s08:enlarged} the matrix points lie in $T_\cF$ by \eqref{eq:s08:ampbound}. \Cref{lem:s08:claim4}(T) gives the bounds; Steps 5 and 6 in version (T), with \cref{prop:s08:T}(a) for the limit on $U\cap T_\cF$ and its holomorphy on $T_\cF$, give \eqref{eq:s08:midplane} for $c\in D$; Step 7 with \cref{prop:s08:T}(b) gives every plane. The $\cD$-class steps are the first factor of \eqref{eq:s08:split} (the $\cD$ half of \cref{lem:s07:L} and the second inequality of \cref{lem:s02:cut}) and the nonincreasing factor in \cref{prop:s08:T}(a).
\end{proof}

The arbitrary-amplitude form, reflection positivity of $Z^a_\infty$ at all complex coefficients under \Wtwo, is \cref{thm:s23:OS2W} in the optional extension; its proof is Steps 5--7 in version (A) for $c\in D$, where \Wtwo\ enters only at the finitely many $q_v$, and in version (A) or (A$'$) for $c\notin D$.

\begin{lemma}[reflection positivity of King's printed state]\label{lem:s08:OS2K}
\Cref{thm:s08:OS2C} holds for $\ZKinf$: for $\nu$, $c$, $\zeta_1,\dots,\zeta_N$ as there, $\sum_{i,j}\bar\gamma_i\gamma_j\ZKinf(\theta\zeta_i^*\sqcup\zeta_j,0)\ge0$ for all $\gamma\in\C^N$; and \cref{prop:s08:OS3,prop:s08:poly} below hold for $\ZKinf(\cdot,0)$. The inputs are \cref{hyp:s03:EC}, \cite[Thm.~6.2(1),(2)]{BFS1}, \cref{thm:s05:CK}(a),(b) in place of \cref{thm:s05:C}, \cref{lem:s04:D}(c) and \cref{rem:s04:F} for the two places where the cube convention is seen, and \cref{cor:s06:TK} for the translations; no inequality of class $\cD$ enters.
\end{lemma}

\begin{proof}
\Cref{lem:s07:RPbox,lem:s07:B} concern the modified box state alone and the lattice reflection, and \cref{lem:s08:claim4}(C) bounds $\Phi_k$ on the modified photon form: these are convention-free, as is Step 5 (Montel in $k$; positivity passes). The two places where the Neumann--Dirichlet cube states enter are the ends of the sandwich \eqref{eq:s08:sandwich}, \ie the $\cI$ halves of \cref{lem:s07:L,lem:s07:U}; on King's convention they read $Z^K_{k,J}(\zeta)\le\Phi_k(\zeta)\le Z^K_{k,J+2n}(\zeta)$ for every $k$ and $\zeta\in U=(\mathrm i\R)^m$: the lower end is \cref{lem:s07:L} with Step 1 of \cref{lem:s04:G} prepended (\cref{rem:s04:F}: $\avK{\cdot}\le\ava{\cdot}\le\langle\cdot\rangle^{(3.38)}_{k,\Omega_{J,k}}$ on $\cI$), the upper end is \cref{lem:s07:U} with $\avK{\cdot}$ on $X_{J+2n}$ in place of the comparison state (\cref{lem:s04:D}(b),(c): the box $\Omega_{J,k}$ does not meet the killed layers $T_{J+2n}$, and part (a) of \cref{lem:s07:U} is photon-only), under the hypotheses of \eqref{eq:s07:margin}. Here $J_0$ is the $J_0$ of the enlarged family (the sources and their reflections), which exceeds that of the sources alone when $c\ne0$; \cref{thm:s05:CK}(b) is applied to the enlarged family at $J$ and $J+2n$, $J\ge J_0$. Along Step 5's subsequence the ends converge along the full sequence by \cref{thm:s05:CK}(b), so $Z^K_J\le\Phi^J\le Z^K_{J+2n}$ on $U$ for every subsequential $\Phi^J$; as $J\to\infty$ both ends tend to $\ZKinf$ on $U$, $\{\Phi^J\}_J$ is locally bounded on $\C^m$ by \cref{lem:s08:claim4}(C), and $U$ is a set of uniqueness, so $\Phi^J\to\ZKinf$ locally uniformly on $\C^m$ as in Step 6, in particular at the matrix points; positivity passes to the limit. Step 7 uses \cref{cor:s06:TK} in place of \cref{prop:s06:translations}. \Cref{prop:s08:OS3,prop:s08:poly} use only the existence, holomorphy near $0$ and lattice symmetries of the limit, and the locally uniform convergences just established; their proofs below apply verbatim with $\ZKinf$, the charge-conjugation invariance being exact for $\avK{\cdot}$ as well.
\end{proof}

\subsection{\texorpdfstring{$\OSth$}{OS3} and the polynomial form}

\begin{proposition}[symmetry, realness and charge conjugation]\label{prop:s08:OS3}
Let $Z$ be one of: $Z_\infty(\cdot,0)$ if the comparison family is convergent at fixed cubes (resp.\ $\ZKinf(\cdot,0)$); $Z^a_{\infty,T}$ under the hypotheses of \cref{prop:s08:T}; $Z^a_\infty$ under \Wtwo\ (optional extension). The Schwinger functions
\[
S_{n,p}(f_1,\dots,f_n;h_1,\dots,h_p)\defeq\partial_{z_1}\cdots\partial_{z_n}\partial_{w_1}\cdots\partial_{w_p}\,Z\big(\textstyle\sum_iz_if_i,\sum_lw_lh_l\big)\big|_{0}
\]
(with $p=0$ and $S_n$ in the curvature sector) are (a) symmetric multilinear in the $f$'s and in the $h$'s; (b) real for real test functions; (c) zero for odd $n$. The proof uses only that $Z$ exists and is holomorphic near $0$, with the stated symmetries of the lattice functionals; the existence is the input of the respective sector.
\end{proposition}

\begin{proof}
(a) $Z$ depends on a family only through its total sources (so does every lattice functional, hence both limits), so its Taylor coefficients at $0$ are symmetric multilinear forms in the $f$'s and in the $h$'s; $Z^a_{\infty,T}$ is holomorphic on a tube containing $0$. The definition of $S_{n,p}$ does not depend on the family used to compute it (restrict a larger family to the given sources). (b) For real coefficients every lattice functional is real (real integrand and multiplier, \cref{lem:s02:alg}(a)), so $Z$ is real at the real points near $0$ (for $Z^a_{\infty,T}$ on $\R^m\cap T$, an open piece of $\R^m$ containing $0$); a holomorphic function real on the real points near $0$ has real Taylor coefficients at $0$. (c) Every cube state and every modified box state is invariant under $(A,\vartheta)\mapsto(-A,-\vartheta)$ (the photon terms are even in $A$, the masked divergence included; $\cos(\vartheta_x-e\eps A-\vartheta_y)$ is invariant; $W^a$ depends on $r$ only; Lebesgue measure is preserved), which maps $F_{\Lk}(f)\mapsto-F_{\Lk}(f)$ and fixes $W^a_{\Lk}(h)$. Hence $Z_{k,X}(-z,w)=Z_{k,X}(z,w)$ for all $k$, $X$, and in both limits $Z(-z,w)=Z(z,w)$ (on the $z\mapsto-z$ invariant tube for $Z^a_{\infty,T}$). So $\partial_z^{(n)}\partial_w^{(p)}Z(0)=0$ for odd $n$. This is the argument of \cite[lem:C-OS3lat, thm:C-OS3]{SymPhase} for the cube states.
\end{proof}

\begin{proposition}[the polynomial form of reflection positivity]\label{prop:s08:poly}
Let $P=\sum_\alpha c_\alpha\prod_iF(f_{\alpha i})\prod_lW^a(h_{\alpha l})$ be a polynomial in the smeared fields with real sources supported in $\{x_\nu>c\}$, and $\theta P\defeq\sum_\alpha\bar c_\alpha\prod_iF(\rho^*f_{\alpha i})\prod_lW^a(h_{\alpha l}\circ\rho)$. Then $S(\theta(P)P)\ge0$, where $S$ assigns to a monomial the corresponding Schwinger function of \cref{prop:s08:OS3}: the Taylor coefficient at $0$ of $Z^a_{\infty,T}$ under the hypotheses of \cref{thm:s08:OS2WT}, resp.\ of $Z_\infty(\cdot,0)$ (of $\ZKinf(\cdot,0)$), where at $h\equiv0$ the $W^a$ factors are absent.
\end{proposition}

\begin{proof}
Let all sources of $P$ be supported in $\{x_\nu\ge c+\delta\}$ and let $\cF'$ be the family of all sources of $P$ and their reflections. Its scalar members have no sign, so Steps 4--6 are applied to the split family $\cF'_{\rm split}$ (members $h^\pm_{\alpha l}$, $h^\pm_{\alpha l}\circ\rho$, all $\ge0$); its tube contains $0$ ($k_1(0)=0$), and the unsplit functional is the composition of the split one with the linear map $w_{\alpha l}\mapsto(w^+_{\alpha l},w^-_{\alpha l})=(w_{\alpha l},-w_{\alpha l})$, so by the chain rule the derivatives at $0$ of the unsplit $\Phi_k$, $\Phi^J$, $Z^a_{\infty(,T)}$ are finite linear combinations of those of the split ones, and locally uniform convergence near $0$ passes through the linear map. (i) \emph{Lattice.} For $k\ge k_0$ the sampled $P$ is a gauge-invariant polynomial in the plus-side smeared fields, lies in $L^2$ of the modified box state (Gaussian moments; $\lambda>0$) and in $\mathfrak B_+$ (\cref{def:s07:reflection}), and $\theta(P)=\overline{P\circ\hat\rho}$ by \eqref{eq:s07:thetaTheta}. By \cref{lem:s07:RPbox}(a) (its $L^2$ extension), $\langle\theta(P)P\rangle^{(3.38)}_{k,\Omega_{J,k}}\ge0$. (ii) \emph{Moments as derivatives.} $\langle\theta(P)P\rangle^{(3.38)}=\sum_{\alpha,\beta}\bar c_\alpha c_\beta\langle\prod_iF(\rho^*f_{\alpha i})\prod_lW^a(h_{\alpha l}\circ\rho)\prod_iF(f_{\beta i})\prod_lW^a(h_{\beta l})\rangle^{(3.38)}$, and each mixed moment is a mixed partial derivative at $0$ of $\Phi_k$ for the family $\cF'$ (differentiate under the integral; the scheme multiplier $e^{-wa_k\sS(h)}$ is the $-a_k$ inside $W^a$). (iii) \emph{Limits.} By Steps 5--6 for $\cF'_{\rm split}$, $\Phi_{k_m}\to\Phi^J$ and $\Phi^J\to Z^a_{\infty(,T)}$ locally uniformly on a neighbourhood of $0$, so by Weierstrass' theorem all derivatives at $0$ converge, and $\sum\bar c_\alpha c_\beta\,\partial^{(\alpha,\beta)}Z^a_{\infty(,T)}(0)\ge0$; the left side is $S(\theta(P)P)$ by the family-independence of the Schwinger functions (\cref{prop:s08:OS3}(a)). Every plane: Step 7 (translation of the sources; $S$ is translation invariant by \cref{prop:s08:T}(b), resp.\ \cref{prop:s06:translations} or \cref{cor:s06:TK}). The curvature sector is the case $h\equiv0$ with no splitting. This is the form OS2(b) of \cite[def:OS, thm:C-OS2]{SymPhase} for the cube states.
\end{proof}

\begin{remark}[which form is used, and the amplitude constants]\label{rem:s08:amplitudes}
The main theorem (\cref{thm:s13:main}) states reflection positivity of the Wick sector in the tube form of \cref{thm:s08:OS2WT}, with \cref{prop:s08:OS3,prop:s08:poly} for the Schwinger functions; its inputs for this clause are \cref{hyp:s03:W1}, \Itwo, \Ithree\ and \cite[Thm.~6.2(1)--(4)]{BFS1}, without \Wtwo. This suffices for the consumer: the weak form of $\OSf$ and the reconstruction of \cref{thm:s09:S2K} use reflection positivity for exponentials $e^{\mathrm i\Phi(f)}$ of the smeared fields, whose Wick coefficients are purely imaginary, so amplitude $0$ only. Three positive, $\eps$-free amplitude constants occur in the Wick sector: $\sigma_W=\sigma_+/8$ here; $H/2=\sigma_+/32$, the negative-part radius of the tube $\Dtube$ of \cref{thm:s11:W} ($H=s_1/16$); and $\sigma_0=H/4$ in the bound $\Ezz$ of \cref{thm:s11:W}(ii). On $T_\cF\cap\Dtube$ the objects $Z^a_{\infty,T}$ and $\Ginf$ coincide (\cref{prop:s08:T}(a)). On King's convention the tube statements of this section are read under \TWK\ (\cref{rem:s11:Kconvention}, \cref{prop:s13:WK}); the curvature-sector statement needs no such hypothesis (\cref{lem:s08:OS2K}). The margin $R(\lmu)$ of \cref{lem:s07:U} enters the proof, not the conclusion: reflection positivity is proved at every fixed $\lmu>0$, and in a limit of vanishing photon mass it passes to the limit functionals as a pointwise limit of positive semidefinite matrices, the polynomial form by Weierstrass' theorem.
\end{remark}

% s09 The measure and weak clustering
\section{The measure and weak clustering}

This section builds the measure of King's printed state in the curvature sector and draws from the authors' earlier paper on the torus limits \cite{SymPhase} the consequences of reflection positivity: the ergodic decomposition, clustering, the unique joint vacuum on a component, the growth of the components' moments, and weak \OSf\ on almost every component. The joint field with the Wick square is the subject of the tube sections. The measure is built from the characteristic functional, not from the lattice laws, which converge only locally (\cref{rem:s09:S2K}).

\emph{Conventions.} $\cV\defeq\Lambda^2\R^3$ with its Euclidean fibre norm. For a real $u\in C^\infty_c(\R^3;\Lambda^2)$ put $\cZ^K(u)\defeq\ZKinf(\mathrm iu)$, the one-member functional at the coefficient $\mathrm i$ (\cref{thm:s05:CK}(b), with the total-source convention of \cref{thm:s06:S1K}). Fix an even $r\ge14$. $\cH^{-r}(\R^3;\Lambda^2)$ is the Hermite--Sobolev space, the dual of the completion of $\cS(\R^3;\Lambda^2)$ under $\norm u_r^2=\sum_{\alpha,c}\lambda_\alpha^{r}\ip{h_{\alpha,c}}u^2$, where $(h_{\alpha,c})_{\alpha\in\N^3,\,c\le3}$ is the Hermite basis of $L^2(\R^3;\Lambda^2)$ and $\lambda_\alpha=2|\alpha|+4$ are the eigenvalues of $-\Delta+|x|^2+1$; $\fournorm\cdot$ is the weighted $C^4$ norm of \cref{lem:s02:E}, and $\fournorm u\le C\norm u_{10}$ on $\cS$ \cite[lem:D:hermite(a)]{SymPhase}. $X\defeq H^{-r}_{\mathrm{loc}}(\R^3;\Lambda^2)$ carries the actions $(\tau_a\Phi)(u)=\Phi(\tau_{-a}u)$, $(g\Phi)(u)=\Phi(g^{-1}\cdot u)$ for $g\in\Oh$ (or $O(3)$), $(\theta_j\Phi)(u)=\Phi(\sigma_j\cdot u)$ for the coordinate reflections $\sigma_j$, with the signed transport of $2$-forms of \cref{def:s02:sources}, and the local $\sigma$-algebras $\Sigma(U)$ generated by $\Phi\mapsto\Phi(u)$, $u\in C^\infty_c(U;\Lambda^2)$ real; $\theta_{j,c}\defeq\tau_{ce_j}\theta_j\tau_{-ce_j}$ is the reflection in the plane $x_j=c$ and $\rho_{j,c*}u$ the reflected source (\cref{def:s07:reflection}). These are the space, actions and $\sigma$-algebras of the abstract part of \cite{SymPhase}, restricted from $\Lambda^2\oplus\R$ to the curvature summand. The \emph{lattice curvature field} of King's cube is the random distribution $\Phi^K_{k,J}:u\mapsto F_{\Lk}(u)$ under $\avK\cdot$. $J_0(u)$ is the least $J$ with $\supp u$ at $\ell^\infty$-distance $>1$ from $\R^3\setminus\CJ$.

\begin{theorem}[the measure, the local lattice laws, weak clustering]\label{thm:s09:S2K}
The limits below exist by \cref{hyp:s03:EC}.
\begin{enumerate}[label=\textup{(\alph*)},leftmargin=*]
\item \textup{(The measure; uses \cref{thm:s05:CK,cor:s06:TK} only.)} For every even $r\ge14$ there is exactly one Borel probability measure $\mu^{K,H}$ on $\cH^{-r}(\R^3;\Lambda^2)$ whose mixed moments on $C^\infty_c$ are the Schwinger functions $S^K_n$ of King's printed state. Its characteristic functional on $C^\infty_c$, indeed on $\cS$, is $\cZ^K$ (extended by continuity); its moments on $\cS$ are the continuous extensions of the $S^K_n$ and satisfy the bounds of \cref{thm:s05:CK}(c). Its image $\mu^K_S$ under the restriction map $\iota:\cH^{-r}\to X$ is the unique Borel probability measure on $X$ with characteristic functional $\cZ^K$ on $C^\infty_c$, and also the unique one whose mixed moments on $C^\infty_c$ are the $S^K_n$. \emph{Local convergence of the lattice laws:} for every real $\chi\in C^\infty_c(\R^3)$ and every box $Q$ containing $\supp\chi$ in its interior, the laws $\nu^K_{k,J,\chi}$ of the localised lattice fields $\chi\Phi^K_{k,J}$ on $H^{-r}_{\mathrm{per}}(Q;\Lambda^2)$ converge weakly, along the full sequence $k\to\infty$ at every fixed $J\ge J_0(\chi)$, to a law $\nu^K_{J,\chi}$, and $\nu^K_{J,\chi}\Rightarrow\nu^K_\chi$ as $J\to\infty$, where $\nu^K_\chi$ is the law of $\chi\Phi$ under $\mu^K_S$. Nothing is asserted about the laws of the unlocalised lattice fields at fixed $J$ (\cref{rem:s09:S2K}).
\item \textup{(Invariance and weak clustering; uses also \cref{thm:s06:S1K,lem:s08:OS2K}.)} $\mu^K_S$ satisfies the hypotheses (H1)--(H3) of the abstract part of \cite{SymPhase}, is invariant under $\Cub$ about every point, and the abstract theorems of \cite{SymPhase} apply to it with their proofs: $\mu^K_S=\int\mu_\xi\,d\nu(\xi)$ with components that are translation-invariant and ergodic, $\theta_j$-invariant, reflection positive about every coordinate plane and satisfy (H2) \cite[thm:D:decomp]{SymPhase}; clustering in the mean \cite[thm:D:mean]{SymPhase} and under diagonal separation \cite[thm:D:diag]{SymPhase}; the Osterwalder--Schrader reconstruction in each coordinate direction with a unique joint vacuum on every component \cite[prop:D:recon, prop:D:fixed(iii)]{SymPhase}; the growth bound $|S^\xi_N(u_1,\dots,u_N)|\le\sigma^\xi_N\prod_i\norm{u_i}_{r+2}$, $\sigma^\xi_N\le\alpha_\beta(\xi)(N!)^\beta$ for every $\beta>1$, for the components' Schwartz moments \cite[cor:D:schwinger(iv)]{SymPhase}; and the cubic-invariant decomposition into mixtures of $k(\zeta)\in\{1,2,3,6\}$ ergodic measures \cite[add:Ohdecomp]{SymPhase}. The same holds on $\cH^{-r}$ for $\mu^{K,H}$ \cite[lem:D:hermite(d)]{SymPhase}. $\nu$-almost every component satisfies \OSf\ in the weak form, ergodicity under the translation group $\R^3$ \cite[add:OS4(iii)]{SymPhase}.
\item \textup{(The rotation corollary.)} $\cZ^K$ is $O(3)$-invariant, $\cZ^K(R\cdot u)=\cZ^K(u)$ for all $R\in O(3)$ and real $u$, by \cref{thm:s06:O3} (the rotation theorem \cref{hyp:s03:RC}); the clauses below use only this invariance. Then for $\nu$-a.e.\ $\xi$: $\mu_\xi$ is invariant under $\Euc$; $\E_{\mu_\xi}[\overline{F\circ\theta_{n,c}}\,F]\ge0$ for every unit vector $n$, every $c\in\R$ and every $F\in L^2(\Sigma(H^+_{n,c}),\mu_\xi)$, $H^+_{n,c}=\{x\cdot n>c\}$; and $\E_{\mu_\xi}[G'\,U(a)G]\to\E_{\mu_\xi}[G']\,\E_{\mu_\xi}[G]$ as $|a|\to\infty$ in any manner, for all $G,G'\in L^2(\mu_\xi)$. Consequently $\mu_\xi$ satisfies \OSf\ in each of the forms of \cite[def:OS4, add:OS4]{SymPhase} (the cluster property on the positive-time algebra, ergodicity of the time translations, ergodicity under $\R^3$) for every choice of the time direction, coordinate or not, and in the reconstruction from $\mu_\xi$ in any direction $n$ the joint fixed space is $\ker H_n=\C\Omega_n$ \cite[p:cor:rotcor]{SymPhase}. The moments of $\mu_\xi$ are real and symmetric; the vanishing of its odd curvature moments is \emph{not} claimed for a component.
\end{enumerate}
\end{theorem}

\subsection{Proof of (a)}

\emph{Step 1 (moment bounds, and where they hold).} For real $2$-forms $f_1,\dots,f_n$ at distance $>1$ from $\R^3\setminus C_{J_0}$, \cref{thm:s05:CK}(c) gives, at every $k$ and $J\ge J_0$, $|S^{K,(k,J)}_n|\le(\sqrt2\,e)^n(n!)^{1/2}\prod_j\norm{f_j}_{2,k}$ with $\norm f_{2,k}^2\le3C_0\fournorm f^2$ uniformly in $k$ (\cref{lem:s02:E}), and in the limit the bound of \cref{thm:s05:CK}(c); by \cref{thm:s05:CK}(a), $\avK{F(u)^2}\le2\norm u^2_{2,k}$ and hence $S^K_2(u,u)\le2\norm u_2^2$. The lattice Schwinger functions are the Taylor coefficients at $0$ of $a^K_{k,J}$, that is the moments of $\Phi^K_{k,J}$, and converge to the $S^K_n$. These bounds are uniform in $k$ and $J\ge J_0$ for sources at distance $>1$ inside $C_{J_0}$; at a fixed $J$ they are not available for sources reaching $\partial\CJ$, and neither is the convergence in $k$: \cref{hyp:s03:EC} and \cref{thm:s05:CK}(a) are stated for sources whose plaquettes have their corners in $\XJ$.

\emph{Step 2 (the measure).} (i) \emph{The functional on $\cS$.} For real $u_1,\dots,u_N\in C^\infty_c$ and $c\in\C^N$, at every $(k,J)$ with $J\ge J_0(u)$, $\sum_{l,m}\bar c_lc_m\,Z^K_{k,J}(\mathrm i(u_m-u_l))=\avK{|\sum_mc_me^{\mathrm iF(u_m)}|^2}\ge0$ by the linearity of $F_{\Lk}$ in the source. The limits $k\to\infty$, then $J\to\infty$, exist at these points by \cref{thm:s05:CK}(b) for the family $u$, so $\cZ^K$ is positive definite on $C^\infty_c$ with $\cZ^K(0)=1$. By \cref{cor:s06:TK}'s Lipschitz bound and (E2), $|\cZ^K(u)-\cZ^K(u')|\le\sqrt2\norm{u-u'}_2\le\frac\pi2\fournorm{u-u'}$. Since $\fournorm\cdot$ is a continuous seminorm on $\cS(\R^3;\Lambda^2)$ and $C^\infty_c$ is dense in $\cS$, $\cZ^K$ extends uniquely to a continuous function on $\cS$, still positive definite with $\cZ^K(0)=1$. By the Bochner--Minlos theorem \cite{GJ,Jan97} ($\cS$ is nuclear) there is exactly one Borel probability measure $\mu'$ on $\cS'(\R^3;\Lambda^2)$, with its cylinder $\sigma$-algebra, such that $\int e^{\mathrm i\Phi(u)}\,d\mu'=\cZ^K(u)$ for all real $u\in\cS$.

(ii) \emph{Moments on $C^\infty_c$, and on $\cS$ by uniform integrability.} For real $u_1,\dots,u_N\in C^\infty_c$ the characteristic function of $(\Phi(u_l))_l$ under $\mu'$ is $\varphi(t)=\cZ^K(\sum_lt_lu_l)=Z^{K,u}_\infty(\mathrm it)$, the restriction to $\mathrm i\R^N$ of the entire function $Z^{K,u}_\infty$ (\cref{thm:s05:CK}(b)). A characteristic function real-analytic near $0$ has moments of all orders, with $\int\prod_l\Phi(u_l)\,d\mu'=(-\mathrm i)^N\partial_{t_1}\cdots\partial_{t_N}\varphi(0)$ \cite{Luk70}; the chain rule gives $\partial_{t_1}\cdots\partial_{t_N}\varphi(0)=\mathrm i^N\partial_{z_1}\cdots\partial_{z_N}Z^{K,u}_\infty(0)$, so the factors cancel and
\begin{equation}\label{eq:s09:moments}
\int\prod_{l=1}^N\Phi(u_l)\,d\mu'=S^K_N(u_1,\dots,u_N);\qquad\int\Phi(u)^2\,d\mu'=S^K_2(u,u)\le2\norm u_2^2\le\tfrac{\pi^2}4\fournorm u^2\le C_2''\norm u_{10}^2
\end{equation}
by Step~1, (E2) and $\fournorm u\le C\norm u_{10}$. \emph{Extension to $\cS$.} By \cref{thm:s05:CK}(c) each $S^K_N$ extends to a continuous symmetric multilinear form on $\cS^N$ with $|S^K_N|\le\sigma_N\prod_l\fournorm{u_l}$. For $u_l\in\cS$ take cut-offs $\chi_R=\chi(\cdot/R)$, $\chi\in C^\infty_c$, $\chi=1$ on the unit ball, $0\le\chi\le1$, $R\ge1$: $\chi_Ru_l\to u_l$ in $\cS$, so $\Phi(\chi_Ru_l)\to\Phi(u_l)$ for every $\Phi\in\cS'$; and $\fournorm{\chi_Ru}\le C_\chi\fournorm u$ uniformly in $R\ge1$ with $C_\chi=2^4\norm\chi_{C^4}$, by Leibniz with $|R^{-|\beta|}(D^\beta\chi)(\cdot/R)|\le\norm\chi_{C^4}$ (the norm $\fournorm\cdot$ carries derivatives up to order $4$, so the cut-off does not pass through it with constant $1$). Hence $\int\prod_l\Phi(\chi_Ru_l)^2\,d\mu'=S^K_{2N}(\chi_Ru_1,\chi_Ru_1,\dots)\le\sigma_{2N}C_\chi^{2N}\prod_l\fournorm{u_l}^2$ uniformly in $R$: the products $\prod_l\Phi(\chi_Ru_l)$ are bounded in $L^2(\mu')$, hence uniformly integrable, and converge pointwise, so $\int\prod_l\Phi(u_l)\,d\mu'=\lim_RS^K_N(\chi_Ru_1,\dots)=S^K_N(u_1,\dots,u_N)$ by the continuity of $S^K_N$ on $\cS^N$.

(iii) \emph{The Hermite carrier.} Put $M(\Phi)\defeq\sum_{\alpha,c}\lambda_\alpha^{-r}\Phi(h_{\alpha,c})^2$. By \eqref{eq:s09:moments} with $\norm{h_{\alpha,c}}_{10}^2=\lambda_\alpha^{10}$, $\int M\,d\mu'\le3C_2''\sum_\alpha\lambda_\alpha^{10-r}<\infty$ as soon as $r-10>3$ (there are $(n+1)(n+2)/2$ multi-indices with $|\alpha|=n$), that is for every even $r\ge14$; so $\mu'(M<\infty)=1$. For $\Phi\in\cS'$ and $u\in\cS$ the Hermite expansion $u=\sum_{\alpha,c}\ip{h_{\alpha,c}}u\,h_{\alpha,c}$ converges in $\cS$ \cite[Thm.~V.13]{RS1}, so $\Phi(u)=\sum_{\alpha,c}\ip{h_{\alpha,c}}u\,\Phi(h_{\alpha,c})$ and, on $\{M<\infty\}$, $|\Phi(u)|\le M(\Phi)^{1/2}\norm u_r$ by Cauchy--Schwarz: $\Phi$ extends to $\Psi(\Phi)\in\cH^{-r}$ with $\norm{\Psi(\Phi)}_{-r}^2=M(\Phi)$ and $\Psi(\Phi)=\Phi$ on $\cS$. The map $\Phi\mapsto\Psi(\Phi)$ ($\defeq0$ off $\{M<\infty\}$) is measurable from the cylinder $\sigma$-algebra to the Borel $\sigma$-algebra of the separable Hilbert space $\cH^{-r}$, its coordinates being evaluations. Put $\mu^{K,H}\defeq\Psi_*\mu'$. Then $\int e^{\mathrm i\Psi(u)}\,d\mu^{K,H}=\cZ^K(u)$ and $\int\prod_l\Psi(u_l)\,d\mu^{K,H}=S^K_N(u_1,\dots,u_N)$ for all $u,u_l\in\cS$, by (ii). This is the carrier construction of \cite[lem:D:hermite(b)]{SymPhase} run on $\cS'$.

(iv) \emph{Uniqueness on $\cH^{-r}$.} By \cite[lem:D:hermite(c)]{SymPhase}, $\mu^{K,H}$ is the only Borel probability measure on $\cH^{-r}$ whose mixed moments on $C^\infty_c$ exist and equal the $S^K_n$. The proof there uses only that the moments $m_n$ of $Z=\sum_lt_l\Psi(u_l)$ satisfy $|m_n|\le c_S(n+1)K^nn!$ (then, by Cauchy--Schwarz for the odd absolute moments and the Taylor remainder bound $|y|^n/n!$, the characteristic function of $Z$ is the sum of its Taylor series on every interval of length $<2/K$, hence determined by the moments on all of $\R$) and the Lusin--Souslin theorem on $\cH^{-r}$ \cite{Kec95}. Our moments give the hypothesis with $c_S=1$: $m_n=S^K_n(g,\dots,g)$ for $g=\sum_lt_lu_l$, and $|m_n|\le(\sqrt2\,e)^n(n!)^{1/2}\norm g_2^n\le(n+1)K^nn!$ with $K=\sqrt2\,e\norm g_2$.

(v) \emph{The law on $X$ and its two uniqueness statements.} The map $\iota:\cH^{-r}\to X$, $\iota\Psi\defeq\Psi|_{C^\infty_c}$, is continuous, linear and injective \cite[lem:D:hermite(b)]{SymPhase}: $|\Psi(\chi f)|\le\norm\Psi_{-r}\norm{\chi f}_r\le C_\chi\norm\Psi_{-r}\norm f_{H^r}$ for a cut-off $\chi$, and $C^\infty_c$ is dense in $\cS$, which is dense in $\cH^r$. Put $\mu^K_S\defeq\iota_*\mu^{K,H}$; its characteristic functional on real $C^\infty_c$ is $\cZ^K$. By the uniqueness part of \cite[lem:D:realize]{SymPhase}, which applies to any Borel probability measure on $X$ (for a countable dense $(f_n)\subset C^\infty_c$ the sets $\{\Phi:(\Phi(f_1),\dots,\Phi(f_N))\in B\}$ form a $\pi$-system generating the Borel $\sigma$-algebra, by Lusin--Souslin, and the law of each such vector is determined by $t\mapsto\cZ^K(\sum_kt_kf_k)$), $\mu^K_S$ is the only Borel probability measure on $X$ with this characteristic functional. \emph{Moment uniqueness on $X$:} if $\mu''$ is any Borel probability measure on $X$ whose mixed moments on $C^\infty_c$ are the $S^K_n$, the argument of (iv) applies verbatim on $X$ with this $\pi$-system \cite[prop:discharge]{SymPhase}, so $\mu''=\mu^K_S$. And $\mu^{K,H}$ is the only measure on $\cH^{-r}$ with $\iota_*\mu^{K,H}=\mu^K_S$, since $\iota$ is a Borel isomorphism onto the Borel set $\iota\cH^{-r}$. Nothing in Step~2 sees the lattice: the measure clauses of (a) depend on King's printed state only through $\cZ^K$ and the $S^K_n$.

\emph{Step $2'$ (local convergence of the lattice laws).} Fix a real $\chi\in C^\infty_c(\R^3)$ and a box $Q$ with $\supp\chi$ compact in the interior of $Q$ (as in \cite[prop:C-tight]{SymPhase}), and let $(e_{\alpha,c})_{\alpha\in\Z^3,\,c\le3}$ be the Fourier basis of $L^2(Q;\Lambda^2)$, $|e_{\alpha,c}(x)|=|Q|^{-1/2}$. For a real $Q$-periodic smooth $2$-form $\psi$ put $\ip{\chi\Phi^K_{k,J}}\psi\defeq\Phi^K_{k,J}(\chi\psi)=F_{\Lk}(\chi\psi)$, where $\chi\psi\in C^\infty_c$ is supported in $\supp\chi$ (for the complex basis vectors, used only in (1), take real and imaginary parts); this identifies $\chi\Phi^K_{k,J}$ with an element of the real Hilbert space $H_\chi\defeq H^{-r}_{\mathrm{per}}(Q;\Lambda^2)$, and $\nu^K_{k,J,\chi}$ denotes its law.

(1) \emph{Second moments, uniformly.} For $J\ge J_0(\chi)$ and every $k$, every plaquette carrying $\chi e_{\alpha,c}$ has its centre in $\supp\chi$, hence its corners in $\XJ$, so the second-moment clause of \cref{thm:s05:CK}(a) applies to the real and imaginary parts: $\E^K_{k,J}|F_{\Lk}(\chi e_{\alpha,c})|^2\le2\norm{\chi\Re e_{\alpha,c}}^2_{2,k}+2\norm{\chi\Im e_{\alpha,c}}^2_{2,k}\le4\norm{\chi e_{\alpha,c}}^2_{2,k}$. \emph{Riemann-sum bound:} the plaquettes of one orientation have their centres on a translate of $\eps_k\Z^3$; those carrying $\chi e_{\alpha,c}$ have centres in $\supp\chi$; the cubes of side $\eps_k\le1$ about these centres are disjoint and lie in the $\frac12$-neighbourhood $N_{1/2}(\supp\chi)$; so $\norm{\chi e_{\alpha,c}}^2_{2,k}\le3\,\mathrm{vol}(N_{1/2}(\supp\chi))\,\norm\chi_\infty^2/|Q|=:K_\chi$, uniformly in $\alpha$, $c$, $k$ and $J\ge J_0(\chi)$ (no derivative of $\chi$ enters). Hence, for $s_0>\frac32$,
\begin{equation}\label{eq:s09:secondmoment}
\E^K_{k,J}\norm{\chi\Phi^K_{k,J}}^2_{H^{-s_0}_{\mathrm{per}}(Q)}=\sum_{\alpha,c}(1+|\alpha|^2)^{-s_0}\,\E^K_{k,J}|F_{\Lk}(\chi e_{\alpha,c})|^2\ \le\ 12K_\chi\sum_{\alpha\in\Z^3}(1+|\alpha|^2)^{-s_0}=:C'_\chi<\infty,
\end{equation}
uniformly in $k$ and $J\ge J_0(\chi)$ (the constants $4$ and $12K_\chi$ are valid but not sharp).

(2) \emph{Tightness.} The embedding $H^{-s_0}_{\mathrm{per}}(Q)\hookrightarrow H^{-r}_{\mathrm{per}}(Q)$ is compact for $r>s_0$ (Rellich on the torus); by Chebyshev, $\nu^K_{k,J,\chi}(\norm\cdot_{H^{-s_0}}>\rho)\le C'_\chi/\rho^2$, and the $H^{-s_0}$-ball of radius $\rho$ is compact in $H_\chi$; so $(\nu^K_{k,J,\chi})_{k\ge0,\,J\ge J_0(\chi)}$ is tight on the separable Hilbert space $H_\chi$, and Prokhorov's theorem \cite{Bil99} gives weak limit points along $k$ at fixed $J$.

(3) \emph{Identification at fixed $J$: the full sequence, and where \cref{hyp:s03:EC} enters.} For a real trigonometric polynomial $\psi$, $\E^K_{k,J}e^{\mathrm i\ip{\chi\Phi^K_{k,J}}\psi}=Z^K_{k,J}(\mathrm i\chi\psi)\to\ZKJ(\mathrm i\chi\psi)$ as $k\to\infty$, along the full sequence, by \cref{thm:s05:CK}(b) (that is \cref{cor:s04:G2}(ii)) for the one-member family $(\chi\psi)$ at $J\ge J_0(\chi)\ge J_0(\chi\psi)$: this is \cref{hyp:s03:EC} for that family on King's cube, together with the local bound of \cref{thm:s05:CK}(a) and Montel's theorem. A Borel probability measure on $H_\chi$ is determined by its characteristic functional on the dual $H^r_{\mathrm{per}}(Q;\Lambda^2)$, which is continuous there (dominated convergence), hence by its values on the dense set of real trigonometric polynomials. So every weak limit point of $(\nu^K_{k,J,\chi})_k$ is the one probability measure $\nu^K_{J,\chi}$ on $H_\chi$ with characteristic functional $\psi\mapsto\ZKJ(\mathrm i\chi\psi)$, and $\nu^K_{k,J,\chi}\Rightarrow\nu^K_{J,\chi}$ along the whole sequence $k\to\infty$ by tightness and the uniqueness of the limit point. Without \cref{hyp:s03:EC} the same argument gives the convergence along every continuum subsequence (\cref{thm:s12:LIM,cor:s12:KK}), not along the full sequence; the full-sequence clause is what \cref{hyp:s03:EC} adds for King's printed state.

(4) \emph{$J\to\infty$.} $\norm\cdot^2_{H^{-s_0}_{\mathrm{per}}}=\sum_{\alpha,c}(1+|\alpha|^2)^{-s_0}|\ip\cdot{e_{\alpha,c}}|^2$ is lower semicontinuous on $H_\chi$ (a supremum of finite sums of continuous functions), so by the portmanteau theorem \cite{Bil99} the bound \eqref{eq:s09:secondmoment} passes to $\nu^K_{J,\chi}$: $\int\norm\cdot^2_{H^{-s_0}}\,d\nu^K_{J,\chi}\le C'_\chi$ for every $J\ge J_0(\chi)$; hence $(\nu^K_{J,\chi})_J$ is tight on $H_\chi$ by the argument of (2), and $\ZKJ(\mathrm i\chi\psi)\to\ZKinf(\mathrm i\chi\psi)=\cZ^K(\chi\psi)$ (\cref{thm:s05:CK}(b)) identifies every limit point: $\nu^K_{J,\chi}\Rightarrow\nu^K_\chi$, the measure with characteristic functional $\psi\mapsto\cZ^K(\chi\psi)$.

(5) \emph{$\nu^K_\chi$ is the law of $\chi\Phi$ under $\mu^K_S$.} For $\Phi\in X$, $\chi\Phi\in H^{-r}(\R^3;\Lambda^2)$ is supported in $\supp\chi$, compact in the interior of $Q$, so its periodisation is an element of $H_\chi$ with $\ip{\chi\Phi}\psi=\Phi(\chi\psi)$ for $Q$-periodic smooth $\psi$; the map $\Phi\mapsto\chi\Phi$ is continuous from $X$ to $H_\chi$, since $|\Phi(\chi\psi)|\le\norm{\tilde\chi\Phi}_{H^{-r}}\norm{\chi\psi}_{H^r}\le c_\chi\norm{\tilde\chi\Phi}_{H^{-r}}\norm\psi_{H^r_{\mathrm{per}}(Q)}$ for a cut-off $\tilde\chi=1$ on $Q$. And $\int e^{\mathrm i\ip{\chi\Phi}\psi}\,d\mu^K_S=\int e^{\mathrm i\Phi(\chi\psi)}\,d\mu^K_S=\cZ^K(\chi\psi)$ by Step~2(v).

(6) \emph{Moments.} For every trigonometric polynomial $\psi$ (indeed every $\psi\in H^r_{\mathrm{per}}(Q;\Lambda^2)$, by the uniform-integrability argument of Step~2(ii) with $\chi\psi_n\to\chi\psi$), all moments of $\ip{\chi\Phi^K_{k,J}}\psi=F_{\Lk}(\chi\psi)$ are bounded uniformly in $k$ and $J\ge J_0(\chi)$ (Step~1; $\fournorm{\chi\psi}\le d_\chi\norm\psi_{C^4}$ by Leibniz), so by uniform integrability the moments of $\nu^K_{J,\chi}$ are the Taylor coefficients of $\ZKJ$ at $\chi\psi$, and those of $\nu^K_\chi$ are the $S^K_n(\chi\psi,\dots,\chi\psi)$, consistently with \eqref{eq:s09:moments}. This proves (a). \qed

\begin{remark}[why the convergence clause is local]\label{rem:s09:S2K}
The convergence $Z^K_{k,J}(\mathrm iu)\to\ZKJ(\mathrm iu)$ is available only for sources $u$ whose plaquettes have their corners in $\XJ$, that is for $u$ at distance $>1$ inside $\CJ$, where \cref{hyp:s03:EC} speaks. For a $u$ reaching $\partial\CJ$ the lattice functional is defined only after a convention for the plaquettes with a corner outside $\XJ$ (for instance, drop them), and for it no bound uniform in $k$ and no convergence is asserted: the second-moment bound of \cref{thm:s05:CK}(a) carries the corner hypothesis, and the crude bound $\avK{e^{F(\tilde u)}}\le e^{\frac12\lmu^{-2}\norm{\dstar_k\tilde u}^2_{2,k}}$ (\cref{lem:s02:GD} with $C\le\lmu^{-2}$ as forms) is not uniform in $k$ for a source $\tilde u$ truncated at $\partial\CJ$: $\dstar_k\tilde u$ is of order $\norm u_\infty/\eps_k$ on the $O(\eps_k^{-2}|\partial\CJ|)$ bonds of the truncation layer, so $\norm{\dstar_k\tilde u}^2_{2,k}=O(\eps_k^{-1})$. In particular the fixed-$J$ laws of the unlocalised lattice field on $X$ are not asserted to be tight; nothing below depends on them.
\end{remark}

\subsection{Proof of (b) and (c)}

\emph{Step 3 ((H1)--(H3)).} The hypotheses of the abstract part of \cite{SymPhase} are: (H1) a Polish space with a jointly continuous $\R^3$-action and three involutive homeomorphisms $\theta_j$ with $\theta_j\tau_a=\tau_{\rho_ja}\theta_j$, and $\mu$ invariant under all of them; (H2) local $\sigma$-algebras $\Sigma(U)$ with monotonicity, $\Sigma(U)=\bigvee_{K\Subset U}\Sigma(K)$, $\bigvee_K\Sigma(K)=\cB$, countable generation by bounded continuous functions, and covariance under $U(a)$ and $\Theta_j$; (H3) $\E_\mu[\overline{F\circ\theta_{j,c}}\,F]\ge0$ for $F\in L^2(\Sigma(H^+_{j,c}),\mu)$. On $X$ with the actions of the Conventions, the structural parts of (H1) and all of (H2) are proved in \cite[prop:discharge]{SymPhase} by arguments that use only the space and the push-forward actions (joint continuity of the translations on $H^{-s}_{\mathrm{loc}}$; the $\theta_j$ are linear homeomorphisms; Lusin--Souslin for $\bigvee_K\Sigma(K)=\cB$; countable dense families $\arctan\Phi(f_n)$; $\supp\tau_af=\supp f+a$, $\supp\sigma_j\cdot f=\sigma_j\supp f$); they hold verbatim for our $X$, that paper's space restricted to the curvature summand. \emph{Invariance of $\mu^K_S$:} by the uniqueness in Step~2(v) it suffices that $\cZ^K(\tau_au)=\cZ^K(u)$ and $\cZ^K(\sigma_j\cdot u)=\cZ^K(u)$ for real $u$; the first is \cref{cor:s06:TK} at the coefficient $\mathrm i$, the second \cref{thm:s06:S1K} at $z=\mathrm i$ with $O=\sigma_j$, and the same theorem with $O=\tau_bg\tau_{-b}$ gives the $\Oh$-invariance about every point. \emph{(H3):} for $E_a=e^{\mathrm i\Phi(u_a)}$ with $u_a$ real in $C^\infty_c(\{x_j>c\};\Lambda^2)$, $\overline{E_a\circ\theta_{j,c}}\,E_b=e^{\mathrm i\Phi(u_b-\rho_{j,c*}u_a)}$ and
\[
\E_{\mu^K_S}\big[\overline{E_a\circ\theta_{j,c}}\,E_b\big]=\cZ^K(u_b-\rho_{j,c*}u_a)=\ZKinf(\theta\zeta_a^*\sqcup\zeta_b),\qquad\zeta_a=(\mathrm i,u_a),
\]
the conjugated coefficient $-\mathrm i$ sitting on the reflected source as in \cref{def:s07:reflection}. By \cref{lem:s08:OS2K}, for the plane $x_j=c$ and every $c\in\R$, this Gram matrix is positive semidefinite. The exponential algebra is dense in $L^2(\Sigma(H^+_{j,c}),\mu^K_S)$ \cite[lem:A:mc]{SymPhase} (an algebra of bounded complex functions, closed under conjugation and containing the constants, is dense in $L^2$ of the $\sigma$-algebra it generates), and $\theta_{j,c}$ preserves $\mu^K_S$, so (H3) follows as in \cite[prop:discharge]{SymPhase}. Finally \cite[lem:D:hermite(d)]{SymPhase} carries (H1)--(H3) and the $\Oh$-invariance about every point from $(X,\mu^K_S)$ to $(\cH^{-r},\mu^{K,H})$.

\emph{Step 4 (weak clustering).} The standing hypothesis of the abstract part of \cite{SymPhase} is that $\mu$ is any measure satisfying (H1)--(H3); so \cite[thm:D:decomp]{SymPhase} gives the ergodic decomposition $\mu^K_S=\int\mu_\xi\,d\nu$ with components translation-invariant and ergodic, $\theta_j$-invariant, satisfying (H3) and (H2); \cite[thm:D:mean, thm:D:diag]{SymPhase} give clustering in the mean and under diagonal separation; \cite[prop:D:recon]{SymPhase} is the reconstruction in direction $j$, and on an ergodic component the joint fixed space is $\C\Omega_j$ \cite[prop:D:fixed(iii)]{SymPhase}. The proof of \cite[cor:D:schwinger(iv)]{SymPhase} uses only the $\Ezz$ bound in the form $\E_{\mu_S}Z_{\alpha,c}^{2n}\le c_SC_0^{2n}(2n)!\,\fournorm{h_{\alpha,c}}^{2n}$ on $Z_{\alpha,c}=\Phi(h_{\alpha,c})$, together with \cite[lem:D:hermite(a)]{SymPhase}, Jensen, Markov and Borel--Cantelli; \cref{thm:s05:CK}(c) supplies that bound on $\cS$ through \eqref{eq:s09:moments}, with $(e\pi/2)^{2n}((2n)!)^{1/2}\le(e\pi/2)^{2n}(2n)!$, so $c_S=1$, $C_0=e\pi/2$, and the growth bound of (b) holds for our components with the same proof. \cite[add:Ohdecomp]{SymPhase} assumes in addition that $\Oh$ acts on $X$ by homeomorphisms with $\rho_j=\theta_j$, $g\tau_ag^{-1}=\tau_{ga}$, $\mu$ invariant under every $g$, and $G\mapsto G\circ g^{-1}$ mapping $\Sigma(U)$-measurable functions to $\Sigma(gU)$-measurable ones: structural for the push-forward action ($(g\Phi)(u)=\Phi(g^{-1}\cdot u)$ is a linear homeomorphism of $X$, the relations are identities on test functions, $\supp g\cdot u=g\supp u$) together with the invariance of Step~3; it gives the cubic-invariant decomposition with $k(\zeta)\in\{1,2,3,6\}$. Weak \OSf\ is \cite[add:OS4(iii)]{SymPhase}, ergodicity for the action of $\R^3$, the ergodicity clause of the decomposition for $\nu$-a.e.\ component.

\emph{Step 5 (the rotation corollary).} \cite[p:cor:rotcor]{SymPhase} assumes (H1)--(H3) and that $O(3)$ acts on $X$ by homeomorphisms $x\mapsto Rx$, continuously in $R$ for each $x$, with $\rho_jx=\theta_jx$ for the coordinate reflections, $R\tau_aR^{-1}=\tau_{Ra}$, $G\mapsto G\circ R^{-1}$ mapping $\Sigma(U)$-measurable functions to $\Sigma(RU)$-measurable ones, and $\mu$ invariant under every $R\in O(3)$. The structural parts hold for $(R\Phi)(u)=\Phi(R^{-1}\cdot u)$ on $X$: the action is by linear homeomorphisms, continuous in $R$, and covariant for the local $\sigma$-algebras, as noted in \cite[add:rotcond-rem(1)]{SymPhase} for $H^{-s}_{\mathrm{loc}}(\R^3;\cV)$, and the push-forward respects the restriction to the $\Lambda^2$ summand. The invariance of $\mu^K_S$ is, by Step~2(v), $\cZ^K(R\cdot u)=\cZ^K(u)$ for real $u$, the hypothesis of (c), that is the $O(3)$-invariance of $\ZKinf$ at imaginary coefficients, supplied by \cref{thm:s06:O3}; this section does not derive it. The conclusions are those quoted in (c), proved there from the $R$-invariance and (H3) (reflection positivity about every plane), the argument of \cite[thm:D:decomp]{SymPhase} with a countable dense subgroup (the Euclidean invariance and reflection positivity of a.e.\ component), \cite[thm:D:diag]{SymPhase} along $(1,1,1)/\sqrt3$ after a rotation (clustering in every manner), and \cite[add:OS4rel(a)]{SymPhase} in a frame containing $n$ ($\ker H_n=\C\Omega_n$). They contain nothing on \OSth: the vanishing of the odd curvature moments of a component is not among them, and is not known for a single component.

\emph{Step 6 (curvature alone).} Everything above runs on $X=H^{-r}_{\mathrm{loc}}(\R^3;\Lambda^2)$ for the curvature field alone: the regularity input is \cref{thm:s05:CK}, the reflection positivity input \cref{lem:s08:OS2K}, the invariance input \cref{thm:s06:S1K}, the $O(3)$ input that of (c). The result is an ergodic decomposition of the law of the curvature field; nothing is claimed about the components of a joint law with the Wick square, and the single-axis statements of \cite{SymPhase} that rest on a torus-specific domination are not claimed for King's state. \qed

\begin{remark}[the hypotheses of the cited results, matched]\label{rem:s09:v1match}
For each result of \cite{SymPhase} used above, the step that supplies its printed hypotheses. (H1)--(H3): Step~3 (structure from prop:discharge; invariance from \cref{cor:s06:TK,thm:s06:S1K}; positivity from \cref{lem:s08:OS2K} and lem:A:mc). lem:D:realize, uniqueness part (a measure on $X$ with a given characteristic functional; Lusin--Souslin, lem:A:LS): Step~2(v). prop:discharge (the space and the push-forward actions; ``Determination by $S$''): Steps~3 and 2(v). lem:D:hermite: (a) no hypothesis; (b) the two-point bound $\E\Phi(f)^2\le C\norm f_{10}^2$, Step~2(ii); (c) $|m_n|\le c_S(n+1)K^nn!$, $r\ge14$ even, Step~2(iv) with $K=\sqrt2\,e\norm g_2$, $c_S=1$; (d) (H1)--(H3) on $X$, Step~3. The standing hypothesis ``$\mu$ is any measure satisfying (H1)--(H3)'' of thm:D:decomp, thm:D:mean, thm:D:diag, prop:D:recon, prop:D:fixed: Step~3; the ergodicity needed in prop:D:fixed(iii) holds on a component by thm:D:decomp. add:Ohdecomp (the $\Oh$ action and invariance): Step~4 and \cref{thm:s06:S1K}. cor:D:schwinger(iv) (the $\Ezz$ bound on Hermite functions; lem:D:hermite(a)): \cref{thm:s05:CK}(c) through Step~2(ii), $c_S=1$, $C_0=e\pi/2$. p:cor:rotcor (the $O(3)$ action and invariance): Step~5, the structure from add:rotcond-rem(1), the invariance from the hypothesis of (c). def:OS4, add:OS4(i)--(iii): definitions, Step~5; add:OS4(iii) is the ergodicity clause of thm:D:decomp. prop:C-tight, mechanism only (uniformly bounded second moments of $\Phi(\chi e_\alpha)$): Step~$2'$(1). lem:A:LS, lem:A:mc: Steps~2(iv),(v), 3. The results of \cite{SymPhase} stated for its own limit states (its clustering theorem, its Schwinger-function corollary as stated, its remark on components) are not cited; their abstract content is in the items above, and the odd-moment clause is not transferred to the components (clause (c)).
\end{remark}

\begin{remark}[the comparison family]\label{rem:s09:components}
The same conclusions hold for the limit states of the comparison family $\NDa$: for every continuum subsequence $\cK$ (\cref{def:s12:LIMstate}) the limit state $\ZLIM\cK$ carries a measure with the properties of \cref{thm:s09:S2K}(a) and (b), the local lattice laws converging along $\cK$ instead of the full sequence, and (c) (\cref{thm:s12:LIM}, clauses (v)--(vii)). By \cref{hyp:s03:EC} every such limit state equals $\ZKinf$ (\cref{cor:s04:G2}(iii)), so for King's printed state the two routes to the measure agree. For the iterated limit of the comparison family itself, when it is convergent at fixed cubes (\cref{def:s03:conv}), see \cref{prop:s12:ID}(b); nothing beyond \cref{cor:s04:G2}(iii) is asserted about it here, the two existence clauses not being interchangeable by the inputs of this paper (\cref{cor:s04:G2}(iv)).
\end{remark}

% s10 The Wick tube I: torus inputs and the extensive bound
\section{The Wick tube I: torus inputs and the extensive bound}

The Wick leg controls the joint functional of the curvature and the Wick square on the tube $\Dtube$ of \cref{def:s02:tube}. Its positive leg (sources $q\ge0$) and negative leg (sources $-q$) are proved in the next section from three inputs on \emph{King's tori}: lattice reflection positivity (\cref{prop:s10:RP}), the chessboard estimate it yields (\cref{lem:s10:CB}), and an \emph{extensive bound} \Itwo\ for constant positive scalar sources in the fixed-spacing volume limit (\cref{def:s10:I2}). This section proves the first two and derives the third from the torus inputs \cref{hyp:s03:T1,hyp:s03:T2} and the sign clause \Tthree\ (\cref{prop:s10:U}); a variant with the auxiliary mass on either side of $m$ (\cref{prop:s10:Upm}) serves the discussion of the coupling range; no continuum limit is taken here.

\subsection{King's tori, the two conventions and mid-plane pairs}

\begin{definition}[King's tori]\label{def:s10:torus}
A \emph{King torus} is $T=\R^3/(\ell\Z)^3$ with $\ell\in M\N$ even; $|T|\defeq\ell^3$ is its volume, and $T_J\defeq\R^3/(2MJ\Z)^3$, $J\in\N$, is the torus with fundamental cell $\CJ$. Since $\ell\Z^3\subset\Z^3\subset\Lk$, the lattice $\Lk$ descends to $T$, with $\ell\,5^k$ sites per side, an even number, and $|T|/\eps_k^3$ sites in all. The \emph{periodic state} $\langle\cdot\rangle_{k,T}$ is King's action of Section~2 summed over the finite graph $\Lk/(\ell\Z)^3$, with periodic conditions for both fields \cite[(4.1)]{K2}, \cite[(2.21)]{K1}: Feynman-form gradient terms for every pair of parallel neighbouring bonds, the photon mass $\lmu^2$, hopping terms on every bond, the single-site weights $e^{-\eps^3[\frac12(m^2+\dm)r^2+\lambda r^4]}$. Two normalisations of the counterterm are used:
\begin{enumerate}[label=\textup{(\roman*)},nosep]
\item the \emph{free convention} of Section~2 ($\dm=\dm_\eps$, Wick constant $\sW_\eps$; the convention of every cube state of this paper), with $N^a_{k,T}(0,q)\defeq e^{-(a_k+2\sW_\eps)\sS(q)}\langle e^{\sum_x\eps^3q(x)r_x^2}\rangle_{k,T}$ for a real function $q$ on the sites of $T$, exactly as on boxes (\cref{def:s02:sources}), and
\[
\alpha^a_{k,T}(u)\defeq|T|^{-1}\ln N^a_{k,T}(0,u\mathbf 1_T)\qquad(u\in\R),
\]
the \emph{free energy density} of a constant source;
\item \emph{King's torus convention} \cite[(2.4), (3.33), (3.35)]{K1}: the counterterm $\dm_{k,T}(M)$ and the Wick constant $c_{k,T}(M)$ computed with the periodic propagators of $T$ at scalar mass $M$, together with the field-independent constants $E(M)=E_0(M)+E_1(M)$ of King's torus action. $Z^{(M)}_{k,T}\defeq\int e^{-S}$ is the source-free partition function of this model at scalar mass $M$ (bare coefficient $M^2+\dm_{k,T}(M)$; \cref{hyp:s03:T1} is a statement about these numbers), and $N^{z,\mathrm{tor}}_{k,T}(0,q)$ the normalised source functional of the torus model at mass $m$ in the reference scheme $z$, with $c_{k,T}=c_{k,T}(m)$ in place of $\sW_\eps$.
\end{enumerate}
In either convention the coefficients are the same at every site of $T$. A \emph{mid-plane pair} is $\pi=\{x_\nu=c\}\cup\{x_\nu=c+\ell/2\}$ with $\nu\in\{1,2,3\}$ and $c\in\eps_k(\Z+\frac12)$ (both planes lie between lattice planes, since $\ell5^k$ is even); $\theta$ denotes the reflection of $T$ in $\pi$, which exchanges the two slabs of $T\setminus\pi$ (``the two sides of $\pi$''), acting on the fields as in \cref{def:s07:reflection}.
\end{definition}

\subsection{Reflection positivity of King's tori}

\begin{proposition}[lattice reflection positivity of King's tori]\label{prop:s10:RP}
Let $T$ be a King torus, $k\ge0$, $\pi$ a mid-plane pair and $\theta$ the reflection in $\pi$. For every finite family $G_1,\dots,G_N$ of gauge-invariant Borel functions of the fields on one side of $\pi$ with $\langle|G_i|^2\rangle_{k,T}<\infty$ and every $\gamma\in\C^N$,
\[
\sum_{i,l=1}^N\bar\gamma_i\gamma_l\,\langle\theta(G_i)\,G_l\rangle_{k,T}\ \ge\ 0 .
\]
The periodic state is invariant under $\theta$ and under all lattice translations. This holds in both conventions of \cref{def:s10:torus}: reflection positivity does not see the value of the (site-independent) counterterm.
\end{proposition}

\begin{proof}
The proof is the mechanism of \cref{lem:s07:RPbox} with a pair of planes in place of one and without Dirichlet diagonals, in three steps.

\emph{Step 1: from the Feynman form to the Proca form.} On the torus write $A=Z+d_0\chi$ with $d_0^TZ=0$, $\chi\perp1$ (Hodge decomposition on the finite periodic graph). The Feynman form equals $\frac12\langle Z,(d_1^Td_1+\lmu^2)Z\rangle$ plus a quadratic form in $\chi$, and the Proca form $\frac12\langle A,(d_1^Td_1+\lmu^2)A\rangle$ equals the same $Z$-form plus another $\chi$-form; both $\chi$-forms are positive definite on $1^\perp$ because $\lmu>0$. For gauge-invariant $G$ the matter integral at fixed $A$ is invariant under $A\mapsto A+d_0\chi$ (rotate $\phi$ site by site), so the $\chi$-integrals factor and cancel in the normalised expectation: gauge-invariant expectations in the Feynman form equal those in the Proca form. Fubini's theorem applies by $\lmu>0$, $\lambda>0$ and a growth bound $|G|\le Ce^{a\sum_b|A(b)|}e^{C\sum_xr_x^2}$, satisfied by bounded $G$ and by the observables $\Theta(\zeta)$; so this step uses $\lmu>0$, gauge invariance, $\lambda>0$ and the growth bound, and nothing about the value of the counterterm.

\emph{Step 2: the Proca torus measure is reflection positive about $\pi$.} In the Proca form the action couples the two sides of $\pi$ only through the plaquettes and hopping terms crossing one of the two planes of $\pi$ (no gradient terms between parallel bonds across a plane; the mass term is local). Write each crossing bond variable as $a_b=\alpha_b-\alpha'_b$ with $\alpha_b,\alpha'_b$ independent centred Gaussians of variance $(2\lmu^2\eps^3)^{-1}$, attached to the two sides. Then $\theta$ becomes the swap of the two data sets $\xi=(\omega_+,\alpha)$, $\xi'=(\omega_-^\theta,\alpha')$, whose a priori measures coincide by the $\theta$-invariance of the one-sided terms, and exactly as in Step~2 of \cref{lem:s07:RPbox} every crossing plaquette term is $\frac\eps2X_p(\xi)^2+\frac\eps2X_p(\xi')^2-\eps X_p(\xi)X_p(\xi')$ and every crossing hopping term is $-\eps\sum_{i=1,2}\psi^i_b(\xi)\psi^i_b(\xi')$, with the same real functions on both copies. So the Boltzmann factor is $\cG(\xi)\cG(\xi')\exp(\sum_\iota\gamma_\iota Y_\iota(\xi)Y_\iota(\xi'))$ with all $\gamma_\iota=\eps>0$ and all $Y_\iota$ real, and for $G$ a function of one side $\theta(G)G=\overline{G(\xi')}G(\xi)$. Expanding the exponential, $Z\langle\theta(G)G\rangle$ is a sum with nonnegative coefficients of squared moduli $|\int d\xi\,G(\xi)\cG(\xi)\prod_\iota Y_\iota(\xi)^{n_\iota}|^2$, the interchange of sum and integral being justified as in \cref{lem:s07:RPbox} ($\lambda>0$, $\lmu>0$ and the growth bound). The Hermitian property $\overline{\langle\theta(G')G\rangle}=\langle\theta(G)G'\rangle$ follows from the $\theta$-invariance of the measure, so $(G,G')\mapsto\langle\theta(G')G\rangle$ is a positive semidefinite Hermitian form on bounded gauge-invariant one-sided $G$.

\emph{Step 3: $L^2$ observables.} For gauge-invariant one-sided $G$ with $\langle|G|^2\rangle<\infty$ apply Step~2 to $G_n\defeq G\mathbf 1\{|G|\le n\}$; since $|\theta(G_n)G_n|\le|\theta G||G|$ and $\langle|\theta G||G|\rangle\le\langle|G|^2\rangle$ by Cauchy--Schwarz and $\theta$-invariance, dominated convergence gives the inequality for $G=\sum_i\gamma_iG_i$, hence the form statement.

\emph{Invariance.} Under $\theta$ and under every lattice translation each term of the periodic action, in either convention, goes to a term of the same kind with the same coefficient, and the a priori measure is invariant; the value of the site-independent bare coefficient enters nowhere above.
\end{proof}

\begin{remark}[the torus paper and the literature]
(i) The torus paper \cite{SymPhase} proves \cref{prop:s10:RP} for its lattice model, King's action on the torus with a constant bare coefficient $b_T$, for every even side, every direction, both kinds of plane pair and every $b_T$ \cite[prop:C-RP(a)]{SymPhase}; its printed scope is $\lmu\in(0,\mu_*]$, where $\mu_*=m$ is the photon mass at which that paper freezes its counterterm. The restriction plays no role in the proof, which uses no bound on $\lmu$: the two lemmas behind it there, the passage to the Proca form \cite[lem:C-FP]{SymPhase} and reflection positivity of the Proca torus measure in a pair of planes \cite[lem:C-RPP(ii)]{SymPhase}, are stated for every photon mass $\mu>0$, with the inputs site-independent coefficients, $\lambda>0$ and even sides. So the statement lifts to every $\lmu>0$, as the proof above also shows directly. (ii) \cite[Cor.~5.4]{BFS1} (reflection positivity ``in every gauge'', for covariances $C_F=C_u+dFd^*\ge0$ with $F$ of either sign) is the same mechanism: on the torus the Feynman covariance has this form with $F=-\lmu^{-2}(\Delta_0+\lmu^2)^{-1}\le0$ on functions orthogonal to the constants. It is not a literal instance (the region of \cite[\S5]{BFS1} is a bounded subset of the lattice with a boundary condition, and reflection positivity of periodic conditions is asserted there in a remark), and \cref{prop:s10:RP} does not rest on it. (iii) The chessboard estimate uses \cref{prop:s10:RP} for functions of the $r_x$ only; nothing in this section uses reflection positivity of a cube or of a limit state.
\end{remark}

\subsection{The chessboard estimate on King's tori}

\begin{lemma}[chessboard estimate]\label{lem:s10:CB}
Let $T$ be a King torus at spacing $\eps=\eps_k$, $a$ a scheme, and $q$ a real function on the sites of $T$. In the free convention,
\begin{equation}\label{eq:s10:CB}
\ln N^a_{k,T}(0,q)\ \le\ \sum_{x}\eps^3\,\alpha^a_{k,T}(q(x)),
\end{equation}
and $u\mapsto\alpha^a_{k,T}(u)$ is convex on $\R$ with $\alpha^a_{k,T}(0)=0$. In King's torus convention $u\mapsto|T|^{-1}\ln N^{z,\mathrm{tor}}_{k,T}(0,u\mathbf 1_T)$ is likewise convex with value $0$ at $u=0$.
\end{lemma}

\begin{proof}
We apply the chessboard estimate \cite[Thm.~4.2]{FILS78} (lattice-gauge form: \cite[Thm.~2.10]{BFS1}), whose hypotheses we check: the site set of $T$ is a rectangular lattice box with sides of even length $\ell5^k$; the periodic state is invariant under translations modulo the sides and reflection positive with respect to the reflections in all planes perpendicular to the coordinate axes running midway between neighbouring sites (\cref{prop:s10:RP}; on the torus such a reflection is the reflection in a mid-plane pair, the antipodal plane $c+\ell/2$ being again a mid-plane). Take as the ``spin at $x$'' the gauge-invariant variable $r_x$ and
\[
G_x\defeq\exp\big(\eps^3q(x)(r_x^2-2\sW_\eps-a_k)\big),\qquad\text{so that}\qquad\Big\langle\prod_xG_x\Big\rangle_{k,T}=N^a_{k,T}(0,q).
\]
The $G_x$ are positive with all moments ($\lambda>0$), and the Schwarz inequality for such unbounded gauge-invariant functions, all that \cite[Thm.~4.1]{FILS78} uses beyond cyclic invariance, follows from the bounded case by truncation and monotone convergence. The conclusion is $|\langle\prod_xG_x\rangle|\le\prod_x\langle\prod_yG_x^{(y)}\rangle^{1/|\Lambda|}$, $G_x^{(y)}$ the copy of $G_x$ carried to $y$ by reflections, $|\Lambda|=|T|/\eps^3$ the number of sites. Because $\sW_\eps$, $a_k$ and $\dm_\eps$ are the same at every site, $\prod_yG_x^{(y)}=\exp(q(x)\sum_y\eps^3(r_y^2-2\sW_\eps-a_k))$ has expectation $N^a_{k,T}(0,q(x)\mathbf 1_T)=e^{|T|\alpha^a_{k,T}(q(x))}$ in the same state; taking logarithms gives \eqref{eq:s10:CB}. Convexity: $|T|\alpha^a_{k,T}(u)$ is the logarithm of a Laplace transform in the linear parameter $u$ plus a linear term, and vanishes at $u=0$; the same computation gives the convexity statement in the torus convention.
\end{proof}

\subsection{Constant sources on large tori: the extensive bound}

\begin{definition}[the extensive bound \Itwo]\label{def:s10:I2}
We say that \Itwo\ (one-sided, in the volume limit at fixed spacing; also written \UVS) holds for the scheme $a$ in the free convention if there are $s_1>0$ and $C_U\in[0,\infty)$ such that
\begin{equation}\label{eq:s10:I2}
\sup_{k\ge0}\ \limsup_{J\to\infty}\ \alpha^a_{k,T_J}(s_1)\ \le\ C_U .
\end{equation}
\end{definition}

The positive leg of the next section consumes exactly this, with $\kappa\defeq C_U/s_1$ (by convexity and $\alpha(0)=0$, $\limsup_J\alpha^a_{k,T_J}(u)\le\kappa u$ for $0\le u\le s_1$). Throughout, $M_1\in(0,m)$ is an auxiliary scalar mass, $\alpha^\infty_k(M)=M^2+\dm_k(M)$ is King's free-boundary bare coefficient on $\Lk$ at scalar mass $M$ (\cref{hyp:s03:T2}; it is the fixed-$k$ volume limit of the torus coefficient $M^2+\dm_{k,T_J}(M)$, \cref{lem:s10:V}), and, in the reference scheme $z$,
\[
s^\infty_k(M_1)\defeq\tfrac12\big(\alpha^\infty_k(m)-\alpha^\infty_k(M_1)\big),\qquad
\sigma_1=\sigma_1(M_1)\defeq\liminf_{k\to\infty}s^\infty_k(M_1),
\]
\[
B^z_{k,T}(M)\defeq E(M)-E(m)-\tfrac12\big(2c_{k,T}+z_k\big)\big(\alpha^{\mathrm{tor}}_{k,T}(m)-\alpha^{\mathrm{tor}}_{k,T}(M)\big)|T|,
\]
\[
b^\infty_k(M_1)\defeq\lim_{J\to\infty}|T_J|^{-1}B^z_{k,T_J}(M_1),\qquad
\beta_1=\beta_1(M_1)\defeq\limsup_{k\to\infty}b^\infty_k(M_1),
\]
where $\alpha^{\mathrm{tor}}_{k,T}(M)\defeq M^2+\dm_{k,T}(M)$ is the bare coefficient of King's torus convention and the limit defining $b^\infty_k$ exists by \cref{lem:s10:V}. The constant $\beta_1(M_1)$ is the second constant of \cref{hyp:s03:T2}, which says that $\sigma_1(M_1)$ and $\beta_1(M_1)$ are finite; \Tthree\ says that $\sigma_1(M_1)>0$.

\begin{proposition}[\Itwo\ from the torus inputs]\label{prop:s10:U}
Assume \Tthree\ (\cref{hyp:s03:T3}); the torus inputs are \cref{hyp:s03:T1} at the two masses $m$ and $M_1$ and \cref{hyp:s03:T2}. Then \Itwo\ holds in the reference scheme $z$ for every $s_1\in(0,\sigma_1)$, with a constant $C_U=C_U(s_1)<\infty$ depending on $(e,\lambda,m,\lmu,M)$, on $M_1$ and on $s_1$. For a scheme $a\ne z$ with $\sup_k|z_k-a_k|<\infty$ (an admissible scheme, \cref{def:s03:adm}) it holds with the constant $C_U+s_1\sup_k|z_k-a_k|$.
\end{proposition}

The proof assembles four lemmas on King's torus model at fixed spacing, stated in the form used and with their short proofs where elementary. Notation at fixed $k$ (spacing $\eps$) on $T=T_J$: $f^{\mathrm{tor}}(u)\defeq\ln N^{z,\mathrm{tor}}_{k,T}(0,u\mathbf 1)$ (King's torus model at mass $m$, Wick constant $c_{k,T}$, scheme $z$) and $f^F(u)\defeq\ln N^{z}_{k,T}(0,u\mathbf 1)$ (free convention); both are convex with $f(0)=0$ (\cref{lem:s10:CB}). Write $\alpha^{\mathrm{tor}}\defeq m^2+\dm_{k,T}(m)$, $\alpha^F\defeq m^2+\dm_\eps$, $\sigma_{k,T}\defeq\frac12(\alpha^{\mathrm{tor}}-\alpha^F)$ and $\Delta c_T\defeq\sW_\eps-c_{k,T}$ (the two \emph{convention offsets}), and $s_{k,T}(M_1)\defeq\frac12(\alpha^{\mathrm{tor}}_{k,T}(m)-\alpha^{\mathrm{tor}}_{k,T}(M_1))$.

\begin{lemma}[mass shift]\label{lem:s10:M}
If $M''>0$ satisfies $\alpha^{\mathrm{tor}}_{k,T}(M'')=\alpha^{\mathrm{tor}}_{k,T}(m)-2u$, then
\[
N^{z,\mathrm{tor}}_{k,T}(0,u\mathbf 1)=e^{B}\,Z^{(M'')}_{k,T}\big/Z^{(m)}_{k,T},\qquad B=E(M'')-E(m)-(2c_{k,T}+z_k)\,u\,|T| .
\]
\end{lemma}

\begin{proof}
The source $u\sum_x\eps^3r_x^2$ shifts the bare coefficient of the torus model at mass $m$ by $-2u$, that is to the bare coefficient of the model at mass $M''$; beyond the fixed terms, King's torus action \cite[(2.4)]{K1} contains only the counterterm and the field-independent constants $E_0,E_1$, so the two Boltzmann factors differ by the constant $e^{E(M'')-E(m)}$, and the Wick and scheme multipliers $e^{-(2c_{k,T}+z_k)u|T|}$ of $N^{z,\mathrm{tor}}$ are bookkeeping.
\end{proof}

\begin{lemma}[chord]\label{lem:s10:C}
If $f$ is convex on $\R$ with $f(0)=0$, then $f(u)\le(u/s)f(s)$ for $0\le u\le s$, $s>0$.
\end{lemma}

\begin{proof}
$u=(1-u/s)\cdot0+(u/s)\cdot s$ and $f(0)=0$.
\end{proof}

\begin{lemma}[the volume limit at fixed spacing]\label{lem:s10:V}
Fix $k$. As $J\to\infty$: $\alpha^{\mathrm{tor}}_{k,T_J}(M)\to\alpha^\infty_k(M)$ uniformly for $M\in[M_1,m]$; $s_{k,T_J}(M_1)\to s^\infty_k(M_1)$; $|T_J|^{-1}B^z_{k,T_J}(M_1)\to b^\infty_k(M_1)$, a finite limit; and the offsets vanish, $\sigma_{k,T_J}\to0$ and $\Delta c_{T_J}\to0$.
\end{lemma}

Proved in \moving{existence part: the fixed-spacing torus lemma}; the reason the statement holds: $\dm_{k,T}$ and $E_1$ (\cite[(3.32)--(3.33), (3.35)]{K1}) are finite sums of lattice graphs with translation-invariant massive propagators ($m,\lmu>0$), which on the torus are exact image sums of the propagators of $\eps\Z^3$, each decaying exponentially in $|x|/\eps$ at fixed $\eps$, so dominated convergence over the images gives the limits; $E_0/|T|$ is a Riemann sum of a smooth periodic log-determinant. No rate and no uniformity in the source are claimed or used. The boundedness in $J\ge J_*(k)$ of $\dm_{k,T_J}$ and $c_{k,T_J}$, used in \cref{lem:s10:F}, is a further clause of the same lemma, not one of those listed here. Of the clauses listed, Step~2 of the proof of \cref{prop:s10:U} uses the convergence of $s_{k,T_J}(M_1)$ and the $\limsup_J$ of $|T_J|^{-1}B^z_{k,T_J}(M_1)$, and Step~3 the two offset clauses.

\begin{lemma}[the fixed-spacing extensive bound]\label{lem:s10:F}
Fix $k$ and $v>0$. There is $C_k(v)<\infty$ such that for every King torus $T=T_J$ with $J\ge J_*(k)$, in either convention, $f(\pm v)\le C_k(v)|T|$; hence, by convexity, $|f(w)|\le C_k(v)\,|w|\,|T|/v$ for $|w|\le v$.
\end{lemma}

\begin{proof}
Let $I(b)\defeq\int_\C e^{-\eps^3[\frac12b|\phi|^2+\lambda|\phi|^4]}d^2\phi<\infty$ for real $b$ ($\lambda>0$), $\alpha$ the bare coefficient and $c$ the Wick constant of the convention in use, and $Z_\gamma$ the photon Gaussian integral, finite because $\lmu>0$ (at $\lmu=0$ the constant mode is in the kernel of the Feynman form). The hopping term is $\ge0$, so dropping it bounds the numerator: $\int e^{-S+w\sum\eps^3r^2}\le Z_\gamma I(\alpha-2w)^{|T|/\eps^3}$. The hopping term of a bond is $\le\eps(r_x^2+r_y^2)$, and each site lies on six bonds, so the hopping part of the action is $\le6\eps^{-2}\sum_x\eps^3r_x^2=\sum_x\eps^3\cdot\frac12(12\eps^{-2})r_x^2$ and $\int e^{-S}\ge Z_\gamma I(\alpha+12\eps^{-2})^{|T|/\eps^3}$. Hence $\ln\langle e^{w\sum\eps^3r^2}\rangle\le\eps^{-3}|T|[\ln I(\alpha-2w)-\ln I(\alpha+12\eps^{-2})]$ and
\[
f(w)\le\big\{\eps^{-3}[\ln I(\alpha-2w)-\ln I(\alpha+12\eps^{-2})]+(2|c|+|z_k|)|w|\big\}|T| .
\]
In the free convention $\alpha$ and $c$ do not depend on $T$; in the torus convention they converge as $J\to\infty$ and are bounded for $J\ge J_*(k)$ (the boundedness clause noted after \cref{lem:s10:V}). So $C_k(v)$, the supremum over the sign and over $J\ge J_*(k)$ of the brace at $w=\pm v$, is finite. The last claim is convexity with $f(0)=0$: $f(w)\le(|w|/v)f(\pm v)$ and $f(w)\ge-f(-w)\ge-(|w|/v)f(\mp v)$.
\end{proof}

\begin{proof}[Proof of \cref{prop:s10:U}]
\emph{Step 1: the exact convention bridge.} On the same torus the free-convention action differs from the torus-convention action by $\frac12(\alpha^F-\alpha^{\mathrm{tor}})\sum_x\eps^3r_x^2=-\sigma_{k,T}\sum_x\eps^3r_x^2$ plus field-independent constants. Hence $\langle e^{u\sum\eps^3r^2}\rangle^F=\langle e^{(u+\sigma)\sum\eps^3r^2}\rangle^{\mathrm{tor}}/\langle e^{\sigma\sum\eps^3r^2}\rangle^{\mathrm{tor}}$ with $\sigma=\sigma_{k,T}$, and inserting the Wick and scheme multipliers,
\begin{equation}\label{eq:s10:bridge}
f^F(u)=f^{\mathrm{tor}}(u+\sigma)-f^{\mathrm{tor}}(\sigma)-2\,\Delta c_T\,u\,|T| .
\end{equation}

\emph{Step 2: the bound in the torus convention.} Let $0\le u<\sigma_1$ and $\eta>0$. By the definitions of $\sigma_1$ and $\beta_1$ there is $K_1=K_1(u,\eta)$, which we take at least as large as the tail of \cref{hyp:s03:T1}, such that $s^\infty_k(M_1)>u$ and $b^\infty_k(M_1)<\beta_1+\eta$ for all $k\ge K_1$. Fix such a $k$. By \cref{lem:s10:V} there is $J_*\ge J_*(k)$ with $s\defeq s_{k,T_J}(M_1)>u$ for $J\ge J_*$. The chord (\cref{lem:s10:C}) at this $s$ gives $f^{\mathrm{tor}}(u)\le(u/s)f^{\mathrm{tor}}(s)\le[f^{\mathrm{tor}}(s)]_+$. \Cref{lem:s10:M} with $M''=M_1$ applies exactly, since $\alpha^{\mathrm{tor}}_{k,T_J}(M_1)=\alpha^{\mathrm{tor}}_{k,T_J}(m)-2s$ by the definition of $s$:
\[
f^{\mathrm{tor}}(s)=B^z_{k,T_J}(M_1)+\ln Z^{(M_1)}_{k,T_J}-\ln Z^{(m)}_{k,T_J}\le B^z_{k,T_J}(M_1)+\big(E_+(M_1)+E_-(m)\big)|T_J|
\]
by \cref{hyp:s03:T1}, used as an upper bound at $M_1$ and as a lower bound at $m$. Divide by $|T_J|$ and let $J\to\infty$ (\cref{lem:s10:V} for the first term):
\[
\limsup_{J\to\infty}|T_J|^{-1}f^{\mathrm{tor}}_{k,T_J}(u)\le\big[\beta_1+\eta+E_+(M_1)+E_-(m)\big]_+\qquad(k\ge K_1).
\]
For each of the finitely many $k<K_1$, \cref{lem:s10:F} with any $v\ge u$ gives $\limsup_J|T_J|^{-1}f^{\mathrm{tor}}_{k,T_J}(u)\le C_k(v)u/v<\infty$. Fix $\eta\defeq1$ (any $\eta>0$ serves); then
\[
C^{\mathrm{tor}}(u)\defeq\sup_{k\ge0}\ \limsup_{J\to\infty}|T_J|^{-1}f^{\mathrm{tor}}_{k,T_J}(u)<\infty\qquad\text{for every }u\in[0,\sigma_1).
\]

\emph{Step 3: transfer to the free convention.} Fix $u\in[0,\sigma_1)$ and $p>1$ with $pu<\sigma_1$, and let $p'=p/(p-1)$. Since $u+\sigma=\frac1p(pu)+\frac1{p'}(p'\sigma)$, convexity gives $f^{\mathrm{tor}}(u+\sigma)\le\frac1pf^{\mathrm{tor}}(pu)+\frac1{p'}f^{\mathrm{tor}}(p'\sigma)$, whatever the sign of $\sigma=\sigma_{k,T_J}$. Once $p'|\sigma_{k,T_J}|\le v$, which holds for all large $J$ by \cref{lem:s10:V}, \cref{lem:s10:F} gives $|f^{\mathrm{tor}}(p'\sigma)|,|f^{\mathrm{tor}}(\sigma)|\le C_k(v)p'|\sigma_{k,T_J}||T_J|/v$. Inserting this into \eqref{eq:s10:bridge},
\[
|T_J|^{-1}f^F_{k,T_J}(u)\ \le\ \tfrac1p|T_J|^{-1}f^{\mathrm{tor}}_{k,T_J}(pu)+\tfrac{2C_k(v)p'}{v}|\sigma_{k,T_J}|+2|\Delta c_{T_J}|\,u
\ \xrightarrow[J\to\infty]{\ \limsup\ }\ \tfrac1pC^{\mathrm{tor}}(pu),
\]
because $\sigma_{k,T_J}\to0$ and $\Delta c_{T_J}\to0$ at fixed $k$ (\cref{lem:s10:V}; the factor $\frac1{p'}\le1$ in the middle term has been dropped). Hence $\sup_k\limsup_J\alpha^z_{k,T_J}(u)\le\frac1pC^{\mathrm{tor}}(pu)<\infty$ for every $u\in[0,\sigma_1)$. Take $u=s_1\in(0,\sigma_1)$ and $C_U\defeq[\inf_{p>1:\,ps_1<\sigma_1}\frac1pC^{\mathrm{tor}}(ps_1)]_+$; this is \eqref{eq:s10:I2} in the scheme $z$. For a scheme $a\ne z$, $N^a_{k,T}(0,q)=e^{-(a_k-z_k)\sS(q)}N^z_{k,T}(0,q)$ gives $\alpha^a_{k,T}(u)=\alpha^z_{k,T}(u)+(z_k-a_k)u\le\alpha^z_{k,T}(u)+|z_k-a_k|\,u$, so \eqref{eq:s10:I2} holds with $C_U+s_1\sup_k|z_k-a_k|$.
\end{proof}

\begin{proposition}[the auxiliary mass on either side of $m$]\label{prop:s10:Upm}
Let $M_1>0$, $M_1\ne m$. Assume \cref{hyp:s03:T1} at $m$ and at $M_1$, and \cref{hyp:s03:T2} and \Tthree\ read at this $M_1$: $s^\infty_k(M_1)=\frac12(\alpha^\infty_k(m)-\alpha^\infty_k(M_1))$, of either sign, has $\sigma_1(M_1)=\liminf_ks^\infty_k(M_1)>0$, and $\beta_1(M_1)<\infty$. Then \Itwo\ holds in the scheme $z$ for every $s_1\in(0,\sigma_1(M_1))$, and for $a\ne z$ with the shifted constant, exactly as in \cref{prop:s10:U}.
\end{proposition}

\begin{proof}
The proof of \cref{prop:s10:U} is repeated with \cref{lem:s10:V} stated on the interval $I\defeq[\min(M_1,m),\max(M_1,m)]$, and we check that no step uses $M_1<m$. \Cref{lem:s10:M} is an identity for any $M''$ with $\alpha^{\mathrm{tor}}(M'')=\alpha^{\mathrm{tor}}(m)-2u$, whatever the order of $M''$ and $m$. \Cref{lem:s10:C} needs $0\le u\le s$ only. \Cref{lem:s10:V}: its proof uses only that the propagators are translation-invariant massive lattice covariances at fixed $\eps$ with masses in $I\cup\{\lmu\}$, bounded below by $\min(M_1,m,\lmu)>0$, so it holds on $I$ word for word; Step~2 uses it at the two masses $m$ and $M_1$ only. \Cref{lem:s10:F} carries no mass order, and \eqref{eq:s10:bridge} is an identity. In Step~2 the sign clause is $s^\infty_k(M_1)>u\ge0$, the chord is applied at $s=s_{k,T_J}(M_1)>u\ge0$, \cref{lem:s10:M} at $M''=M_1$ is exact by the definition of $s$, and \cref{hyp:s03:T1} is used as an upper bound at $M_1$ and a lower bound at $m$; none of this uses $M_1<m$. Step~3 carries no mass.
\end{proof}

The two-sided variant matters exactly on the set $E^<_W\setminus E_W$ of couplings at which $m$ minimises $\alpha_\infty$ on $(0,m]$ but not on $(0,\infty)$ (\cref{lem:s14:CW}; for instance when $\alpha'_\infty(m)<0$ and $m$ minimises $\alpha_\infty$ on $(0,m]$): there \Tthree\ fails for every $M_1<m$ while some $M_1>m$ has $\sigma_1(M_1)>0$. The main results use \cref{prop:s10:U} with $M_1\in(0,m)$, as \Tthree\ is posed.

\begin{remark}[the convention of \cref{hyp:s03:T1} and how much of it is used]\label{rem:s10:T1convention}
(i) \Cref{hyp:s03:T1} is posed in King's torus convention (counterterm and Wick constant computed with the periodic propagators of the torus itself), not in the normalisation of Bałaban's own papers. This is what the proof needs: the offsets $\sigma_{k,T_J}$, $\Delta c_{T_J}$ between King's torus convention and the free convention vanish in the volume limit at fixed spacing (\cref{lem:s10:V}), and Step~3 of the proof of \cref{prop:s10:U} transfers the bound through them; between other normalisations they need not vanish. (ii) Only a tail of \cref{hyp:s03:T1} is used: Step~2 of that proof invokes it for $k\ge K_1$ and $J\ge J_*(k)$, the finitely many $k<K_1$ being handled by \cref{lem:s10:F} at fixed spacing, and the constants $E_\pm(M)$ are independent of $k$ and of the torus. (iii) In the scheme $z$, $C_U$ depends on the couplings, on $M$ and on the auxiliary choices $M_1$, $s_1$, and on nothing else; neither it nor any other bound of this section is uniform in $\lmu$.
\end{remark}

% s11 The Wick tube II: the uniform bound and its limits
\section{The Wick tube II: the uniform bound and its limits}

This section proves the $(k,J)$-uniform bound \eqref{eq:s11:W} on the general-source functionals of the comparison cubes on the tube $\Dtube$ from the extensive bound \Itwo\ (\cref{def:s10:I2}): the positive part of the Wick source is controlled through King's tori, the negative part through a tiling into cells of fixed size, each controlled by one instance of \cref{hyp:s03:W1} and the positive bound. Then come the limits on the tube, $\Ezz$, and the continuity and translation invariance of the tube object. States are those of the comparison convention $\NDa$ (\cref{def:s02:acube}) unless King's convention is named; $N^a_{k,X}(0,q)$ is the functional of \cref{def:s02:sources} on a finite $X\subset\Lk$ or a King torus (\cref{def:s10:torus}), the torus with fundamental cell $\CJ$ being written $\mathbb T_J$ here ($\TJ$ is the top-layer set of King's RG cube), and $\alpha^a_{k,\mathbb T}(u)=|\mathbb T|^{-1}\ln N^a_{k,\mathbb T}(0,u\mathbf 1)$. $s_1\in(0,\sigma_1)$ and $C_U$ are the amplitude and constant of \Itwo, $\kappa\defeq C_U/s_1$; the sources $f_1,\dots,f_n$, $h_1,\dots,h_p$ are in $C^\infty_c$ with supports in $Q_R=[-R,R]^3$; $g(\zeta)$, $k(\zeta)=k_+-k_-$ and the norms are as in \cref{def:s02:tube}. The scheme is the reference scheme $z$, or an admissible scheme $a$ (\cref{def:s03:adm}, so that $\sup_k|z_k-a_k|<\infty$); then $C_U$ stands for $C_U+s_1\sup_k|z_k-a_k|$ throughout (\cref{prop:s10:U}).

\subsection{The positive leg on cubes}

\begin{proposition}[cube $\le$ torus at one bare mass]\label{prop:s11:W2}
For every $k$, $J$ and every lattice function $q\ge0$ on $\XJ$: $N^a_{k,\CJ}(0,q)\le N^a_{k,\mathbb T_J}(0,q)$, both in the free convention.
\end{proposition}

\begin{proof}
The periodic state on $\mathbb T_J$ has the sites of $\XJ$ with opposite faces identified; its action is the Neumann--Dirichlet action on $\CJ$ plus, for every wrap-around pair $\{x,y\}$, the photon gradient terms $\eps(A_\alpha(y)-A_\alpha(x))^2$ and the hopping term $\frac\eps2(r_x^2+r_y^2)-\eps r_xr_y\cos(\cdot)$, whose diagonal part is exactly the pair of Dirichlet terms the cube action carries at $x$, $y$ for their outward bonds; so the cube state is the torus state with the wrap-around gradient terms and cosine parts deleted. For the class-$\cI$ observable $e^{\sum\eps^3qr^2}$ ($\Theta_{F,g}$, $F=0$, $g=\eps^3q\ge0$) deleting a gradient term, a quadratic $|\ell(A)|^2$, does not increase the expectation (\cref{lem:s02:BFS}(a)), nor does deleting a cosine part (\cref{lem:s02:BFS}(b)); each intermediate measure is a BFS measure, and $\dm_\eps$, $\sW_\eps$, $a_k$, $\sS(q)$ agree on both sides.
\end{proof}

\begin{proposition}[the positive leg in $L^1$ form, on every Neumann--Dirichlet box]\label{prop:s11:W1prime}
Assume \Itwo\ and \Ithree\ (\cref{prop:s10:RP}, entering through the chessboard estimate \cref{lem:s10:CB}). Let $X\subset\Lk$ be finite, $\langle\cdot\rangle^{\ND}_{k,X}$ the Neumann--Dirichlet state of \cref{def:s02:acube} on $X$, and $0\le q\le s_1$ a lattice function on $X$. Then
\[
\ln N^a_{k,X}(0,q)\ \le\ \kappa\,\norm{q}_{1,k},\qquad\kappa=C_U/s_1 .
\]
\end{proposition}

\begin{proof}
Choose $J'$ with $X$ and all its nearest neighbours in $C_{J'}$. (1) \emph{Box $\le$ cube:} the first inequality of \cref{lem:s02:cut} with $Y=X_{J'}$ and $(f,g)=(0,q)$ (\cref{lem:s02:BFS}(a),(b)); its $X$-state is the state of \cref{def:s02:acube} on $X$, since all six bonds at each site of $X$ are present in $Y$ and no site is killed; the scheme multiplier is common (\cref{lem:s02:alg}(a)). (2) \emph{Cube $\le$ torus:} \cref{prop:s11:W2} on $C_{J'}$, $q$ extended by $0$. (3) \emph{Chessboard:} \cref{lem:s10:CB} on $\mathbb T_{J'}$, $\ln N^a_{k,\mathbb T_{J'}}(0,q)\le\sum_x\eps^3\alpha^a_{k,\mathbb T_{J'}}(q(x))$. (4) \emph{Chord:} $\alpha=\alpha^a_{k,\mathbb T_{J'}}$ is convex with $\alpha(0)=0$, so $\alpha(u)\le(u/s_1)[\alpha(s_1)]_+$ for $0\le u\le s_1$, and (3) is at most $s_1^{-1}[\alpha(s_1)]_+\norm{q}_{1,k}$ (the $\ell^1$ norm of $q$ is the same on $X$, $C_{J'}$, $\mathbb T_{J'}$). (5) The left side does not depend on $J'$; let $J'\to\infty$ at fixed $k$ and use \Itwo: $\limsup_{J'}\alpha^a_{k,\mathbb T_{J'}}(s_1)\le C_U$ for every $k$.
\end{proof}

\begin{corollary}[the positive leg on the cubes]\label{cor:s11:W2prime}
Under the hypotheses of \cref{prop:s11:W1prime}, for $0\le q\le s_1$ on $\XJ$: $\ln N^a_{k,\CJ}(0,q)\le\kappa\norm{q}_{1,k}$, uniformly in $k$ and $J$. In particular \Wtwo(b) holds for $\sup q\le s_1$: for every $0\le q\in C^\infty_c$ with $\sup q\le s_1$ and every $J$ with $\supp q$ at $\ell^\infty$-distance $>1$ from $\R^3\setminus\CJ$, $\lim_kZ^a_{k,J}(0,q)$ exists and $\sup_J\lim_kZ^a_{k,J}(0,q)\le e^{\kappa\norm{q}_{L^1}}$.
\end{corollary}

\begin{proof}
\Cref{prop:s11:W1prime} at $X=\XJ$. The limit exists by \cref{hyp:s03:W1} ($0\le q\le s_1=s_W$); the sampled source satisfies $0\le q\le s_1$ (\cref{lem:s02:E}(E3)) and $\norm{q}_{1,k}\to\norm{q}_{L^1}$.
\end{proof}

\subsection{Cells and the constants of the negative leg}

\begin{definition}[cells, tilings, cutoffs]\label{def:s11:cells}
Let $L_c\defeq\min\{L\in\N:ML\ge5\}$, $\ell_c\defeq2ML_c\ge10$; the \emph{cell} $C_{L_c}=[-ML_c,ML_c)^3$ is a King cube. For $s\in\{0,1\}^3$ let $W_\mu(0)\defeq\ell_c\Z$, $W_\mu(1)\defeq ML_c+\ell_c\Z$, and let the tiling $\cT_s$ consist of the cells $\prod_\mu[w_\mu,w_\mu+\ell_c)$, $w_\mu\in W_\mu(s_\mu)$; every cell is $C_{L_c}+b$ with $b\in ML_c\Z^3\subset\Lk$. The good sets $G_s\defeq\{x:\dist(x_\mu,W_\mu(s_\mu))>2\ \forall\mu\}$ are open and cover $\R^3$ (the $2$-neighbourhoods of $W_\mu(0)$, $W_\mu(1)$ are disjoint as $ML_c\ge5>4$). Fix a smooth partition of unity $0\le\chi_s\le1$, $\sum_s\chi_s=1$, $\supp\chi_s\subset G_s$; so $\chi_s$ vanishes within distance $2$ of every wall of $\cT_s$. Fix $\chi_0\in C^\infty_c$, $0\le\chi_0\le1$, $\chi_0=1$ on $\{\dist_\infty(x,\R^3\setminus C_{L_c})\ge2\}$, $\supp\chi_0\subset\{\dist_\infty(x,\R^3\setminus C_{L_c})>1\}$, and put $\chi_c\defeq\tau_b\chi_0$ for the cell $c=C_{L_c}+b$. For sources supported in $Q_R$ let $J_1(R)\defeq\min\{J\in\N:MJ\ge R+\ell_c\}$; a cell meeting $Q_R$ has $w_\mu\in[-R-\ell_c,R]$, hence lies in $[-R-\ell_c,R+\ell_c)^3\subset\CJ$ for $J\ge J_1(R)$, so every cell meeting $Q_R$ is a full cell inside $\CJ$ and the cells clipped by the faces of $\CJ$ carry no source.
\end{definition}

The constants of the Wick leg are
\begin{equation}\label{eq:s11:constants}
\begin{gathered}
\sigma_+\defeq s_1,\qquad H'\defeq\frac{s_1}2,\qquad H\defeq\frac{H'}8=\frac{s_1}{16},\qquad \kappa\defeq\frac{C_U}{s_1},\qquad B_{\rm cell}\defeq\sup_k\ln N^a_{k,C_{L_c}}(0,-2H'\chi_0),\\
C_{\rm cell}\defeq\tfrac12B_{\rm cell}+\kappa H'\ell_c^3,\qquad C_\Sigma\defeq\frac{(2+6\ell_c^2)^2\pi^2}{\ell_c^3},\qquad K_-\defeq\frac{8\,[C_{\rm cell}]_+\,C_\Sigma}{H'} .
\end{gathered}
\end{equation}
$B_{\rm cell}$ is finite by \cref{hyp:s03:W1} at the single cube $C_{L_c}$ and the single real source $-2H'\chi_0$, the supremum of a convergent sequence of positive numbers ($\supp\chi_0$ is at $\ell^\infty$-distance $>1$ from $\R^3\setminus C_{L_c}$ as \cref{hyp:s03:W1} requires, and the theorem puts no amplitude restriction on a negative Wick source, a position-dependent increase of the scalar mass squared). This instance of \cref{hyp:s03:W1} is \Bcell\ (\cref{hyp:s12:Bcell}) and the only place where it enters the bound \eqref{eq:s11:W}.

\subsection{The negative leg on cubes}

\begin{lemma}[wall insertion]\label{lem:s11:W3}
Let $X_1,\dots,X_N$ be the site sets $c\cap\XJ$ of the cells $c$ of one tiling $\cT_s$ that meet $\CJ$, and $q\ge0$ a lattice function on $\XJ$. Then $N^a_{k,\CJ}(0,-q)\le\prod_{i=1}^NN^a_{k,X_i}(0,-q|_{X_i})$, with $N^a_{k,X_i}$ the functional of the Neumann--Dirichlet state of \cref{def:s02:acube} on $X_i$.
\end{lemma}

\begin{proof}
$e^{-\sum\eps^3qr^2}$ is the class-$\cD$ observable $\Theta^-_{F,g}$, $F=0$, $g=\eps^3q\ge0$. Delete, for every nearest-neighbour pair with ends in different cells, the three photon gradient terms and the cosine part of the hopping term, one at a time: by \cref{lem:s02:BFS}(d) (the inequalities \cite[Thm.~6.2(3),(4)]{BFS1}, proved in Section~2) neither deletion decreases $\langle\Theta^-\rangle$, and each intermediate measure is a BFS measure. The result is a product over the cells (\cref{lem:s02:BFS}(e)) whose factors are the Neumann--Dirichlet states on the $X_i$: the kept diagonal parts of the deleted bonds are their Dirichlet terms, the remaining gradient terms their Neumann terms, and $\dm_\eps$, $\sW_\eps$, $a_k$ are site-independent. The observable and the multiplier $e^{a_k\sS(q)}=\prod_ie^{a_k\sS(q|_{X_i})}$ factorise.
\end{proof}

\begin{lemma}[one cell]\label{lem:s11:W4}
Assume \Itwo\ and \Ithree. Let $c=C_{L_c}+b$ be a cell ($b\in ML_c\Z^3$) and $0\le\tilde q\le H'$ a lattice function supported in $\{x:\dist_\infty(x,\R^3\setminus c)\ge2\}$. Then $\ln N^a_{k,c}(0,-\tilde q)\le C_{\rm cell}$ for every $k$.
\end{lemma}

\begin{proof}
$\chi_c=1$ on $\supp\tilde q$, so $-\tilde q=k_1-k_2$ with $k_1\defeq H'\chi_c-\tilde q\ge0$, $k_2\defeq H'\chi_c$, and \cref{lem:s02:alg}(f) at $f=0$ gives $N^a_{k,c}(0,-\tilde q)\le N^a_{k,c}(0,-2H'\chi_c)^{1/2}N^a_{k,c}(0,2k_1)^{1/2}$. \emph{First factor:} exact lattice translation covariance ($b\in\Lk$; the change of variables $x\mapsto x-b$ transports the boundary conditions, and the action, $\dm_\eps$, $\sW_\eps$, $a_k$ are site-independent) gives $N^a_{k,c}(0,-2H'\chi_c)=N^a_{k,C_{L_c}}(0,-2H'\chi_0)\le e^{B_{\rm cell}}$. \emph{Second factor:} \cref{prop:s11:W1prime} on the box $c\cap\Lk$ with $q'\defeq2k_1$, $0\le q'\le2H'=s_1$, gives at most $\exp(\kappa\norm{q'}_{1,k})\le\exp(2\kappa H'\ell_c^3)$, the cell having volume $\ell_c^3$. Hence $\ln N^a_{k,c}(0,-\tilde q)\le\frac12B_{\rm cell}+\kappa H'\ell_c^3$.
\end{proof}

\begin{lemma}[chord on the cell]\label{lem:s11:W5}
Under the hypotheses of \cref{lem:s11:W4}: $\ln N^a_{k,c}(0,-\tilde q)\le([C_{\rm cell}]_+/H')\sup\tilde q$.
\end{lemma}

\begin{proof}
Let $t\defeq\sup\tilde q/H'\in(0,1]$, $\hat q\defeq\tilde q/t$ ($0\le\hat q\le H'$, same support). $\varphi(u)\defeq\ln N^a_{k,c}(0,-u\hat q)$ is convex (\cref{lem:s02:alg}(d)), $\varphi(0)=0$, $\varphi(1)\le C_{\rm cell}$ (\cref{lem:s11:W4}); so $\varphi(t)\le t\varphi(1)\le t[C_{\rm cell}]_+$.
\end{proof}

\begin{lemma}[eight families and the sum over cells]\label{lem:s11:W67}
\textup{(a)} For a lattice function $q\ge0$ on $\XJ$: $N^a_{k,\CJ}(0,-q)\le\prod_{s\in\{0,1\}^3}N^a_{k,\CJ}(0,-8\chi_sq)^{1/8}$. If $q$ is supported in $Q_R$ and $J\ge J_1(R)$, then for each $s$ every cell of $\cT_s$ meeting $\supp(\chi_sq)$ is a full cell inside $\CJ$, on which $\chi_sq$ is supported at $\ell^\infty$-distance $>2$ from the complement of the cell.
\textup{(b)} For any tiling of $\R^3$ by cubes of side $\ell_c$: $\sum_c\sup_{x\in c}(1+|x|^2)^{-2}\le C_\Sigma$.
\end{lemma}

\begin{proof}
(a) H\"older's inequality with eight exponents $8$ in the one state $\langle\cdot\rangle^{\ND}_{k,\CJ}$ for $e^{-\sum\eps^3qr^2}=\prod_se^{-\sum\eps^3\chi_sqr^2}$; the multipliers match since $\sS(q)=\sum_s\sS(\chi_sq)$ at lattice points. Support: every cell meeting $Q_R$ is a full cell inside $\CJ$ for $J\ge J_1(R)$ (\cref{def:s11:cells}), and a point of $G_s$ in the cell $\prod_\mu[w_\mu,w_\mu+\ell_c)$ has $\dist(x_\mu,\{w_\mu,w_\mu+\ell_c\})>2$ for every $\mu$. (b) For $x,y$ in one cell $|y|\le|x|+\sqrt3\ell_c$, so $1+|y|^2\le(2+6\ell_c^2)(1+|x|^2)$ and $\sup_{x\in c}(1+|x|^2)^{-2}\le(2+6\ell_c^2)^2\ell_c^{-3}\int_c(1+|y|^2)^{-2}dy$; sum over the disjoint cells, with $\int_{\R^3}(1+|y|^2)^{-2}dy=\pi^2$.
\end{proof}

\begin{proposition}[the negative leg]\label{prop:s11:W8}
Assume \Itwo\ and \Ithree. For every $k$, every $J\ge J_1(R)$ and every lattice function $q\ge0$ supported in $Q_R$ with $\sup q\le H$:
\[
\ln N^a_{k,\CJ}(0,-q)\ \le\ K_-\,\wnorm{q}.
\]
\end{proposition}

\begin{proof}
By \cref{lem:s11:W67}(a), then \cref{lem:s11:W3} for each tiling $\cT_s$ (the clipped cells carry no source and contribute $\ln N^a_{k,X_i}(0,0)=0$), then \cref{lem:s11:W5}, applicable since $8\chi_sq\le8\sup q\le H'$ and by the support clause of \cref{lem:s11:W67}(a):
\[
\ln N^a_{k,\CJ}(0,-q)\le\frac18\sum_s\sum_{c\in\cT_s}\ln N^a_{k,c}(0,-8\chi_sq|_c)\le\frac18\sum_s\sum_c\frac{[C_{\rm cell}]_+}{H'}\,8\sup_cq\le\frac{[C_{\rm cell}]_+}{H'}\,\wnorm{q}\sum_s\sum_c\sup_{x\in c}(1+|x|^2)^{-2}\le\frac{8[C_{\rm cell}]_+C_\Sigma}{H'}\,\wnorm{q},
\]
using $\sup_cq\le\wnorm{q}\sup_{x\in c}(1+|x|^2)^{-2}$ (a sampled supremum is at most the continuum one, \cref{lem:s02:E}(E3)) and \cref{lem:s11:W67}(b) for each of the eight tilings.
\end{proof}

The output norm is $\wnorm q$: the chord bound on a cell sees only $\sup_cq$, so no $L^1$ information inside a cell survives, and the tube bound carries $K_-\wnorm{k_-}$ where the torus paper \cite{SymPhase} has an $L^1$ norm.

\subsection{The uniform bound on the tube and its limits}

For a family with supports in $Q_R$ and $\zeta=(z,w)\in\C^{n+p}$ put $G^{(k,\CJ)}(\zeta)\defeq Z^a_{k,J}(\sum_iz_if_i,\sum_jw_jh_j)$, and let $\Dtube$ be the tube of \cref{def:s02:tube} with $\sigma_+=s_1$, $H=s_1/16$. Two facts from several complex variables \cite{Hor90} are used: (T2) a holomorphic function on a connected open $D\subset\C^m$ vanishing on a non-empty relatively open subset of $D\cap\Sigma$, $\Sigma$ an affine totally real subspace of real dimension $m$ (such as $\R^m$ or $(\mathrm i\R)^{m_1}\times\R^{m_2}$), vanishes on $D$ (all its derivatives vanish at such a point); (T3) a locally bounded family of holomorphic functions on $D$ is normal, and if it converges on a set as in (T2) it converges locally uniformly on $D$ to a holomorphic limit (Montel, Vitali).

\begin{theorem}[the tube bound]\label{thm:s11:W}
Assume \Itwo\ (which \cref{prop:s10:U} derives from \Tthree\ and the torus inputs \cref{hyp:s03:T1,hyp:s03:T2}), for the reference scheme $z$ or for an admissible scheme $a$ (\cref{def:s03:adm}), with \Ithree\ $=$ \cref{prop:s10:RP} and \Ifour\ $=$ \cref{lem:s02:BFS}. The constants $\sigma_+,H,\kappa,K_-$ of \eqref{eq:s11:constants} depend on $(e,\lambda,m,\lmu,M,a)$ and on the fixed auxiliary choices $M_1$, $\chi_0$, $s_1\in(0,\sigma_1)$, and on nothing else; no uniformity in $\lmu$ is claimed. For every family with supports in $Q_R$, every $k\ge0$, $J\ge J_1(R)$ and $\zeta\in\Dtube$:
\begin{equation}\label{eq:s11:W}
|G^{(k,\CJ)}(\zeta)|\ \le\ G^{(k,\CJ)}(\Re\zeta)\ \le\ \exp\Big(\norm{g(\zeta)}_{2,k}^2+\kappa\norm{k_+(\zeta)}_{1,k}+K_-\wnorm{k_-(\zeta)}\Big).
\end{equation}
Consequently:
\begin{enumerate}[label=\textup{(\roman*)},leftmargin=*]
\item \textup{(Limits on the whole tube.)} $G_J\defeq\lim_kG^{(k,\CJ)}$ exists locally uniformly on $\Dtube$ for every $J\ge J_1(R)$, $\Ginf\defeq\lim_JG_J$ exists locally uniformly on $\Dtube$ (the limit along $J\ge J_1(R+\frac12)$, hence along any tail), both are holomorphic there, and
\[
|\Ginf(\zeta)|\le\exp\big(\norm{g(\zeta)}_2^2+\kappa\norm{k_+(\zeta)}_{L^1}+K_-\wnorm{k_-(\zeta)}\big)\qquad(\zeta\in\Dtube).
\]
\Wtwo\ is not used; where the full object $Z^a_\infty$ of the optional extension exists, $\Ginf=Z^a_\infty$ on $\Dtube$.
\item \textup{($\Ezz$.)} With $\sigma_0\defeq H/4$, $K'\defeq\kappa\pi^2+K_-$, the mixed Schwinger functions $S_{n,p}(f;h)\defeq\partial^{(n)}_z\partial^{(p)}_w\Ginf(0)$ satisfy, for real test functions,
\[
|S_{n,p}(f_1,\dots,f_n;h_1,\dots,h_p)|\le e^{K'\sigma_0}(\sqrt2e)^n(n!)^{1/2}(e/\sigma_0)^p\,p!\prod_i\norm{f_i}_2\prod_j\wnorm{h_j}\le\sigma_{n+p}\prod_i\fournorm{f_i}\prod_j\fournorm{h_j},
\]
$\sigma_N\defeq e^{K'\sigma_0}C_1^NN!$, $C_1\defeq2\max(e\pi/2,e/\sigma_0)$; for complex test functions the first bound carries the factor $2^{n+p}\le2^N$, $N=n+p$, absorbed in $C_1$. This is $\Ezz$ with $s=4$ and every $\beta>1$; the $S_{n,p}$ are tempered distributions.
\end{enumerate}
\end{theorem}

$\sigma_1$, $C_U$ depend on $(e,\lambda,m,\lmu,M_1)$ through the torus stability constants of \cref{hyp:s03:T1} at $m$, $M_1$, the $c$-number $\beta_1$ of \cref{hyp:s03:T2} and finitely many fixed-spacing constants; $B_{\rm cell}$ depends on everything, including $\lmu$, $M$, $a$. The curvature term of \eqref{eq:s11:W} rests on the constant $2$ of \cref{lem:s05:C1}, numerically nearly sharp but not proved to be (\cref{rem:s05:sharp}).

\begin{proof}
\emph{The bound \eqref{eq:s11:W}.} The first inequality is \cref{lem:s02:alg}(d). Write $g=g(\zeta)$, $k=k(\zeta)$ (real). \emph{Step 1} (curvature off): \cref{thm:s05:C}(a) in the presence of the real scalar source $k$ gives $N^a(g,k)\le e^{\norm g_{2,k}^2}N^a(0,k)$; the supports of the $f_i$ are at distance $>1$ from $\R^3\setminus\CJ$ since $MJ\ge R+\ell_c>R+1$. \emph{Step 2}: $N^a(0,k_+-k_-)\le N^a(0,2k_+)^{1/2}N^a(0,-2k_-)^{1/2}$ (\cref{lem:s02:alg}(f) at $f=0$). \emph{Step 3} (positive leg): $\zeta\in\Dtube$ and (E3) give $0\le2k_+\le s_1$ on the lattice, so $\frac12\ln N^a(0,2k_+)\le\kappa\norm{k_+}_{1,k}$ by \cref{cor:s11:W2prime}. \emph{Step 4} (negative leg): $0\le2k_-\le H$, supported in $Q_R$, $J\ge J_1(R)$, so $\frac12\ln N^a(0,-2k_-)\le\frac12K_-\wnorm{2k_-}=K_-\wnorm{k_-}$ by \cref{prop:s11:W8}.

\emph{(i), $k\to\infty$ at fixed $J\ge J_1(R)$.} Each $G^{(k,\CJ)}$ is entire on $\C^{n+p}$. On a compact $K\subset\Dtube$ the family $\{G^{(k,\CJ)}\}_k$ is uniformly bounded by \eqref{eq:s11:W} and \cref{lem:s02:E}(E1)--(E3) ($\norm g_{2,k}^2\le C_0\fournorm g^2$, $\norm{k_+}_{1,k}\le C_0'\wnorm{k_+}$, $\wnorm{k_\pm}\le\sum_j|\Re w_j|\wnorm{h_j}$). At the real points of $\Dtube$, a non-empty open subset of $\R^{n+p}$, it converges by \cref{hyp:s03:W1} (the positive scalar part is below $\sigma_+/2<s_W$). By (T3) with $\Sigma=\R^{n+p}$, $G^{(k,\CJ)}\to G_J$ locally uniformly on $\Dtube$, $G_J$ is holomorphic, and \eqref{eq:s11:W} passes to the limit with $\norm g_{2,k}\to\norm g_2$, $\norm{k_+}_{1,k}\to\norm{k_+}_{L^1}$ ($k_+$ continuous with compact support) and the sampled $\wnorm{k_-}$ at most the continuum one: $|G_J|$ is bounded on $\Dtube$ by the right side of the display in (i), uniformly in $J\ge J_1(R)$.

\emph{(i), $J\to\infty$ on a sub-tube.} Split $h_j=h_j^+-h_j^-$, $h_j^-\defeq\norm{h_j}_\infty\chi_j$, $h_j^+\defeq h_j+h_j^-\ge0$, with $\chi_j\in C^\infty_c$, $0\le\chi_j\le1$, $\chi_j=1$ on $\supp h_j$, $\supp\chi_j\subset Q_{R+1/2}$. Put $F_{k,J}(z,w^+,w^-)\defeq Z^a_{k,J}(\sum_iz_if_i,\sum_j(w_j^+h_j^++w_j^-h_j^-))$ and $\cD''\defeq\{(z,w^+,w^-)\in\C^{n+2p}:|\Re w^\pm_j|<\delta_j\}$ with $\sum_j\delta_j(\norm{h_j^+}_\infty+\norm{h_j^-}_\infty)<H/2$. On $\cD''$ the split family's scalar source has $\norm{k^{\rm sp}_\pm}_\infty\le\norm{k^{\rm sp}}_\infty<H/2\le\sigma_+/2$, so \eqref{eq:s11:W} holds for it uniformly in $k$ and $J\ge J_1(R+\frac12)$, and the previous paragraph gives $F_{k,J}\to F_J$ locally uniformly on $\cD''$ with the limit bound. On $\Omega''_+\defeq\cD''\cap((\mathrm i\R)^n\times(0,\infty)^{2p})$ the sources lie in the cone $\cK$ (imaginary curvature coefficients, nonnegative total scalar source), so $J\mapsto F_J$ is nondecreasing there (first inequality of \cref{lem:s02:cut} with \cref{lem:s02:alg}(a) at fixed $k$, then $k\to\infty$) and bounded, hence convergent. $\Omega''_+$ is a non-empty open subset of the maximal totally real subspace $(\mathrm i\R)^n\times\R^{2p}$ and $\cD''$ is convex, so (T3) gives $F_J\to F_\infty$ locally uniformly on $\cD''$. Restricting to $w^+_j=w_j$, $w^-_j=-w_j$ gives $G_J\to\Ginf$ locally uniformly on $\cD'\defeq\{(z,w):|\Re w_j|<\delta_j\}\subset\Dtube$ (there $\norm{k_\pm}_\infty\le\norm k_\infty<\sum_j\delta_j\norm{h_j}_\infty<H/2$).

\emph{(i), the whole tube.} $\{G_J\}_{J\ge J_1(R+1/2)}$ is locally uniformly bounded on $\Dtube$ by the $J$-free limit bound; every Montel subsequential limit is holomorphic on $\Dtube$ and agrees with $\Ginf$ on the open set $\cD'$; $\Dtube$ is connected, so all subsequential limits coincide: $G_J\to\Ginf$ locally uniformly on $\Dtube$, and the bound passes to $\Ginf$. Where $Z^a_\infty$ exists it is the same iterated limit of the same $Z^a_J$.

\emph{(ii).} Normalise $\norm{f_i}_2=1$ and let $\rho_j\defeq\sigma_0/(p\wnorm{h_j})$. On the polydisc $|w_j|\le\rho_j$, $z\in\C^n$: $\norm k_\infty\le\sum_j|\Re w_j|\norm{h_j}_\infty\le\sum_j\rho_j\wnorm{h_j}=\sigma_0<H/2\le\sigma_+/2$, so the polydisc lies in $\Dtube$, where $\Ginf$ exists with its bound by (i); the exponent there is $\le\kappa\pi^2\wnorm{k_+}+K_-\wnorm{k_-}\le K'\wnorm k\le K'\sum_j|\Re w_j|\wnorm{h_j}\le K'\sigma_0$ by (E2). Cauchy's estimate for the mixed first $w$-derivatives at $w=0$, $z$ fixed: $|\partial^{(p)}_w\Ginf(z,0)|\le\prod_j\rho_j^{-1}\sup|\Ginf|\le(p/\sigma_0)^p\prod_j\wnorm{h_j}\,e^{K'\sigma_0}e^{\norm{g(z)}_2^2}$, and $p^p\le e^pp!$. The entire function $z\mapsto\partial^{(p)}_w\Ginf(z,0)$ is thus bounded by $Ae^{\norm{g(z)}_2^2}$, and Cauchy's estimate in the $z$'s with radius $(2n)^{-1/2}$, as in the proof of \cref{thm:s05:C}(c), gives $(\sqrt2e)^n(n!)^{1/2}$. Then (E2) and $(n!)^{1/2}p!\le(n+p)!$ give the second bound; the extension to $\cS$, to complex test functions (the factor $2^{n+p}\le2^N$) and to tempered distributions is as in \cref{thm:s05:C}(c), and $C_1^N/(N!)^{\beta-1}\to0$ for every $\beta>1$.
\end{proof}

\begin{lemma}[the tube object: limit, continuity, translations]\label{lem:s11:TWtube}
Assume the bound \eqref{eq:s11:W}, which \cref{thm:s11:W} derives from \Tthree\ and the torus inputs (through \Itwo). Write $\cD=\Dtube$ for the tube of a family with supports in $Q_R$.
\begin{enumerate}[label=\textup{(\roman*)},leftmargin=*,start=3]
\item \textup{(The $J$-limit on the tube.)} For $J\ge J_1(R+\frac12)$, which exceeds the least $J_0$ with the supports at $\ell^\infty$-distance $>1$ from $\R^3\setminus\CJ$ (as $MJ_1(R')-R'\ge\ell_c\ge10>1$), $\Ginf=\lim_JG_J$ exists locally uniformly on $\cD$, is holomorphic on $\cD$, satisfies the bound of \cref{thm:s11:W}\textup{(i)}, is $>0$ at every real point of $\cD$, and $G_J\le\Ginf$ at the cone points of $\cD$ ($z\in(\mathrm i\R)^n$, $w\in\R^p$, $k(\zeta)\ge0$: the real scalar source $k(\zeta)\ge0$ sits in the slot of the cone).
\item \textup{(Continuity on the real tube.)} Let $E_\rho\defeq C^\infty_c(\bar B_\rho;\Lambda^2)\times C^\infty_c(\bar B_\rho)$ with its $C^\infty$ topology, a Fr\'echet space, and $U_\rho\defeq\{(f,h)\in E_\rho:\norm{h_+}_\infty<\sigma_+/2,\ \norm{h_-}_\infty<H/2\}$, open and convex ($h\mapsto\norm{h_\pm}_\infty$ is continuous and convex). Then $\Phi^a(f,h)\defeq\ln\Ginf(f,h)$ is finite, convex and continuous on $U_\rho$.
\item \textup{(Translations on the tube.)} For every family and every $b\in\R^3$: $\Ginf(\tau_b\cdot\zeta)=\Ginf(\zeta)$ for all $\zeta\in\cD$; the tube of the translated family is that of the family, since $k(\tau_b\cdot\zeta)=\tau_bk(\zeta)$ has the same supremum norms.
\end{enumerate}
Besides \cref{hyp:s03:W1} and \eqref{eq:s11:W} the proofs use only the first inequality of \cref{lem:s02:cut}, \cref{lem:s02:alg,lem:s02:baire}, lattice covariance and (T2), (T3): not \Wtwo, and \cref{lem:s02:BFS}(d) only inside \eqref{eq:s11:W}.
\end{lemma}

\begin{proof}
(iii) Existence, holomorphy and the bound are the last three paragraphs of the proof of \cref{thm:s11:W}(i), which use \eqref{eq:s11:W}, \cref{hyp:s03:W1} for the $k$-limits and the first inequality of \cref{lem:s02:cut}. \emph{Positivity at real points.} For real $\zeta\in\cD$ pick $\epsilon>0$ with $-\epsilon\zeta\in\cD$ ($\cD$ is an open neighbourhood of $0$). $\varphi_J(t)\defeq\ln G_J(t\zeta)$ is convex on $[-\epsilon,1]$ (\cref{lem:s02:alg}(d) in the limit) with $\varphi_J(0)=0$, so $0\le\frac{\epsilon}{1+\epsilon}\varphi_J(1)+\frac1{1+\epsilon}\varphi_J(-\epsilon)$, i.e.\ $\varphi_J(1)\ge-\varphi_J(-\epsilon)/\epsilon\ge-C_\epsilon/\epsilon$, $C_\epsilon$ the limit bound of \eqref{eq:s11:W} at $-\epsilon\zeta$; hence $G_J(\zeta)\ge e^{-C_\epsilon/\epsilon}$ uniformly in $J$ and $\Ginf(\zeta)>0$. \emph{Cone points.} There $G_J\le G_{J+1}$ (first inequality of \cref{lem:s02:cut} with \cref{lem:s02:alg}(a) at fixed $k$, then $k\to\infty$; it needs $J\ge J_0$, which the range of (iii) ensures), so $G_J\le\Ginf$.

(iv) \emph{Finite:} (iii) for the two-source family $(f;h)$ at $\zeta=(1,1)$, in $\cD$ by the definition of $U_\rho$; the value is the pointwise iterated limit and does not depend on the family used. \emph{Convex:} for $s,s'\in U_\rho$, $t\in[0,1]$, use the family $(f,f';h,h')$ at the real point $(t,1-t,t,1-t)$, in $\cD$ because $(th+(1-t)h')_\pm\le th_\pm+(1-t)h'_\pm$ pointwise; the convexity of $\ln Z^a_{k,J}$ in the real coefficients (\cref{lem:s02:alg}(d)) passes to the limits. \emph{Borel:} a pointwise limit of functions each depending continuously on finitely many values of the sources. \emph{Continuous:} \cref{lem:s02:baire} on the open convex $U_\rho$.

(v) \emph{Lattice translations at cone points.} Let $b\in5^{-k_0}\Z^3$, $|b|_\infty\le M$, and $\zeta$ a cone point of $\cD$ ($k_-=0$, $q\defeq k(\zeta)\ge0$, $\norm q_\infty<\sigma_+/2$). Steps 1--3 of the proof of \cref{prop:s06:translations} run unchanged with the real scalar source $q$ in the slot of the cone: lattice covariance ($\tau_b$ preserves $\Lk$ and $\sS_{\Lk}$ for $k\ge k_0$) gives $Z^a_{k,\CJ}(\tau_b\cdot\zeta)=Z^a_{k,\CJ-b}(\zeta)$, the interlacing $X_{J-1}\subset(\CJ-b)\cap\Lk\subset X_{J+1}$ and the first inequality of \cref{lem:s02:cut} twice (no killed site) give $Z^a_{k,C_{J-1}}(\zeta)\le Z^a_{k,\CJ}(\tau_b\cdot\zeta)\le Z^a_{k,C_{J+1}}(\zeta)$ for $J\ge J_0+1$. Let $k\to\infty$ by \cref{thm:s11:W}(i) for the three families (the translated family is supported at distance $>1$ from $\R^3\setminus C_{J_0+1}$), then $J\to\infty$ by (iii), the three terms being cone points of their tubes: $\Ginf(\tau_b\cdot\zeta)=\Ginf(\zeta)$ at the cone points of $\cD$. \emph{All $\zeta\in\cD$.} For the split family of the proof of \cref{thm:s11:W}(i) and its translate, $F^{\tau_b}_\infty-F_\infty$ is holomorphic on the connected $\cD''$ and vanishes on $\Omega''_+$, whose points are cone points of the split family's tube and which is an open subset of the maximal totally real subspace $(\mathrm i\R)^n\times\R^{2p}$; by (T2) it vanishes on $\cD''$, and the diagonal gives $G^{\tau_b}_\infty=\Ginf$ on $\cD'$. Then $G^{\tau_b}_\infty-\Ginf$ is holomorphic on the connected $\cD$ ((iii) for both families) and vanishes on the open set $\cD'$, hence on $\cD$. \emph{All $b\in\bigcup_k5^{-k}\Z^3$:} iterate; the tube is the same at every step. \emph{All $b\in\R^3$ at real points.} For real $\zeta\in\cD$ the source pair $s\defeq(\sum_iz_if_i,\sum_jw_jh_j)$ lies in some $U_\rho$, $\tau_bs\in U_{\rho'}$ for $|b|\le\rho'-\rho$, and $b\mapsto\tau_bs$ is continuous into $E_{\rho'}$; by (iv) and the density of $\bigcup_k5^{-k}\Z^3$, $\Phi^a(\tau_bs)=\Phi^a(s)$. \emph{Complex $\zeta\in\cD$.} $G^{\tau_b}_\infty-\Ginf$ is holomorphic on $\cD$ and vanishes on the real points of $\cD$, a non-empty open subset of $\R^{n+p}$; (T2) gives the claim.
\end{proof}

\begin{remark}[King's convention]\label{rem:s11:Kconvention}
On King's RG cube $\NDK_{k,J}$ the negative leg holds as printed: by \cref{lem:s04:E}, the wall insertion of \cref{lem:s11:W3} runs on $\NDK_{k,J}$ with factors that are comparison box states killed on $X_i\cap\TJ$, and every cell on which a cutoff source $8\chi_sq$ does not vanish is a full cell inside $\CJ$ meeting no killed site, because $\chi_s$ vanishes within distance $2$ of every cell wall; so \cref{lem:s11:W4,lem:s11:W5,lem:s11:W67,prop:s11:W8} apply with the printed $K_-$. The direction matters: killing sites of a cell increases a class-$\cD$ functional (\cref{lem:s04:E}(e)), so $B_{\rm cell}$ does not control a source-carrying cell with killed sites. For the positive leg, \cref{rem:s04:F} prepends the interlacing inequality $\avK{\cdot}\le\ava{\cdot}$ on class $\cI$ (Step~1 of \cref{lem:s04:G}) to \cref{prop:s11:W2,prop:s11:W1prime}, so \cref{cor:s11:W2prime} holds for $N^K_{k,\CJ}(0,q)$ at every $(k,J)$ and, for the fixed-cube limits, in the form $\limsup_kZ^{a,K}_{k,J}(0,q)\le e^{\kappa\norm q_{L^1}}$. The assembly of \eqref{eq:s11:W} on $\NDK$ from these pieces, with the check that every step of the proof of \cref{thm:s11:W} runs in the right direction on King's convention, and the existence of the fixed-cube limits of $Z^{a,K}_{k,J}$ on the tube (\cref{thm:s11:W}(i) on $\NDK$) are carried out in \cref{thm:s13:TWK}; the limit statements of this section and of the Wick sector of Sections~8, 9 and~12 then hold on King's convention and the Wick sector of the main theorem is stated for King's printed state (\cref{prop:s13:WK}).
\end{remark}

% s12 Continuum subsequences and LIM states
\section{Continuum subsequences and LIM states}

The axioms of the preceding sections rest on the existence of the fixed-cube limits $\lim_k$: on \cref{hyp:s03:EC} in the curvature sector, on \cref{hyp:s03:W1} on the Wick tube. This section removes that dependence. The $k$-uniform bounds of \cref{lem:s05:C2} and \eqref{eq:s11:W} make the fixed-cube functionals normal families, so one subsequence of spacings serves every cube and every family at once; along it the cube limit $J\to\infty$ exists because the cone functionals are monotone in the cube at fixed $k$ (\cref{lem:s02:cut}). The limit state of such a subsequence (a \emph{LIM state}) satisfies every axiom proved above, with no hypothesis in the curvature sector and, on the Wick tube, under \Tthree\ and two further hypotheses; \cref{hyp:s03:EC} (resp.\ \cref{hyp:s03:W1}) adds exactly the identification with the iterated limit and the independence of the subsequence. The section is written on the comparison convention $\NDa$ (\cref{def:s02:acube}), in which the tube lemmas and the rotation part are written; King's convention is \cref{cor:s12:KK}. Notation is that of \cref{def:s02:sources,def:s02:tube,def:s11:cells} and \eqref{eq:s11:constants}: $a^{(a)}_{k,J}(z)=\ava{e^{\sum_jz_jF(f_j)}}$ for a family $f=(f_1,\dots,f_N)$ of real $2$-forms; $Z^a_{k,J}(\zeta)$ for $\cF=(f_1,\dots,f_n;h_1,\dots,h_p)$ and $\zeta=(z,w)\in\C^{n+p}$ in the scheme $a$; $g(\zeta)$, $k(\zeta)=k_+-k_-$, $\amp(\zeta)$; $U_J\defeq\{x:\dist_\infty(x,\R^3\setminus\CJ)>1\}$, $J_0(f)$ the least $J$ with all supports in $U_J$, and $R(\cF)$ the least $R$ with all supports in $Q_R=[-R,R]^3$. Two standard facts \cite{Hor90} are used by name: the \emph{identity theorem} (a holomorphic function on a connected open $D\subset\C^m$ vanishing on a non-empty relatively open subset of $D\cap\Sigma$, $\Sigma$ a maximal totally real affine subspace such as $\R^m$ or $(\mathrm i\R)^{m_1}\times\R^{m_2}$, vanishes on $D$) and \emph{Montel--Vitali} (a locally bounded family of holomorphic functions on $D$ is normal, and converges locally uniformly on $D$ if it converges on such a set).

\subsection{Definitions and hypotheses}

\begin{definition}[continuum subsequences]\label{def:s12:contsub}
An increasing sequence $\cK=(k_m)_{m\ge1}\subset\N$ of spacings is a \emph{continuum subsequence} (of the comparison family) if for every finite family $f$ of real $2$-forms and every $J\ge J_0(f)$ the sequence $(a^{(a)}_{k_m,J})_m$ converges locally uniformly on $\C^N$; the limit is written $Z^{\cK}_J(\sum_jz_jf_j,0)$ and depends on the family only through the total complex source, as each $a^{(a)}_{k,J}$ does. Fix a scheme $a$. $\cK$ is a \emph{continuum subsequence for the Wick sector} (Form~(B)) if it is a continuum subsequence and, for every finite family $\cF$ of real $2$-forms and scalars and every $J$ with $MJ>R(\cF)+\ell_c$, the sequence $(Z^a_{k_m,J})_m$ converges locally uniformly on the tube $\Dtube$ of $\cF$; the limit is $Z^{a,\cK}_J$. (The strict inequality replaces the threshold $J_1(R)$ of \cref{thm:s11:W} by at most one cube step; only tails in $J$ are used.) Form~(A), convergence on all of $\C^{n+p}$ for the same $J$, is the optional extension of \cref{thm:s23:FormA} and is not used here.
\end{definition}

\begin{definition}[LIM states]\label{def:s12:LIMstate}
The \emph{LIM state} of a continuum subsequence $\cK$ is $\ZLIM{\cK}\defeq\lim_{J\to\infty}Z^{\cK}_J$ on $\C^N$ (curvature sector), resp.\ $Z^{a,\cK}_\infty\defeq\lim_JZ^{a,\cK}_J$ on the tube (Wick sector); \cref{lem:s12:LIM1} gives existence. A subsequence of a continuum subsequence is a continuum subsequence with the same fixed-cube limits, hence the same LIM state.
\end{definition}

\begin{definition}[the limit state of the rotation theorem]\label{def:s12:o3LIM}
For a family $f$ and $J\ge J_0(f)$ let $L^f_J$ be the set of locally uniform limits on $\C^N$ of subsequences of $(a^{(a)}_{k,J})_k$. The rotation theorem (\cref{hyp:s03:RC}) gives an entire $Z^{\mathrm{LIM}}_\infty(\cdot,0)$ to which \emph{every} choice $\zeta_J\in L^f_J$ converges locally uniformly as $J\to\infty$, the subsequence being allowed to depend on $J$. With $J$-dependent subsequences one only has $\liminf_ka^{(a)}_{k,J}(\mathrm iu)\le\zeta_J(\mathrm iu)\le\limsup_ka^{(a)}_{k,J}(\mathrm iu)$, and the $J$-limit is pinned only if the $J$-limits of the two bounds agree, which is the content of that theorem. A LIM state in the sense of \cref{def:s12:LIMstate} uses one $\cK$ common to all $J$ and all families, along which the fixed-cube limits are monotone in $J$ (\cref{lem:s12:LIM1}), and needs no hypothesis; \cref{prop:s12:ID}(a) relates the two.
\end{definition}

The curvature sector needs no hypothesis. On the Wick tube the bound \eqref{eq:s11:W} of \cref{thm:s11:W} is used; it rests on \Tthree\ and the torus inputs (\cref{hyp:s03:T1,hyp:s03:T2,hyp:s03:T3}, which give \Itwo\ with amplitude $s_1\in(0,\sigma_1)$ and constant $C_U$ through \cref{prop:s10:U}; \Ithree\ is \cref{prop:s10:RP} and \Ifour\ is \cref{lem:s02:BFS}) and on one instance of \cref{hyp:s03:W1}, isolated as follows.

\begin{hypothesis}[\Bcell, the cell bound]\label{hyp:s12:Bcell}
In the scheme $a$, $B_{\mathrm{cell}}\defeq\sup_k\ln N^a_{k,C_{L_c}}(0,-2H'\chi_0)<\infty$, for the one cell $C_{L_c}$ (side $\ell_c\ge10$) and the one negative Wick source $-2H'\chi_0$ of \cref{def:s11:cells}.
\end{hypothesis}
This is the boundedness form of \cref{hyp:s03:W1} at one cube and one source, and the only place where it enters the proof of \eqref{eq:s11:W} (through \cref{lem:s11:W4,lem:s11:W5,prop:s11:W8}; $K_-$ is built from $B_{\mathrm{cell}}$). It follows from \cref{hyp:s03:W1} for an admissible $a$.

\begin{hypothesis}[\SCH, the scheme]\label{hyp:s12:SCH}
Either $a=z$, the reference scheme in which \cref{hyp:s03:T2} is computed, or $\sup_k|z_k-a_k|<\infty$; in the second case the constants of \eqref{eq:s11:W} hold with $C_U$ replaced by $C_U+s_1\sup_k|z_k-a_k|$.
\end{hypothesis}
For an admissible scheme (\cref{def:s03:adm}) \SCH\ holds; for LIM states with $a\ne z$ it is a genuine further hypothesis.

\begin{hypothesis}[\UNIQ, independence of the subsequence]\label{hyp:s12:UNIQ}
Any two continuum subsequences $\cK,\cK'$ have the same LIM state: $\ZLIM{\cK}(\cdot,0)=\ZLIM{\cK'}(\cdot,0)$ for every family.
\end{hypothesis}
\UNIQ\ follows from \cref{hyp:s03:RC} and from \cref{hyp:s03:EC} (\cref{prop:s12:ID}(a),(b)). Without them, every LIM state lies on $(\mathrm i\R)^N$ between $\lim_J\liminf_ka^{(a)}_{k,J}$ and $\lim_J\limsup_ka^{(a)}_{k,J}$ (both exist, by monotonicity in $J$), and \UNIQ\ holds if these coincide at every real argument, an entire function being determined on $(\mathrm i\R)^N$; the converse is not claimed. The axioms below hold for each LIM state separately and say nothing about independence of $\cK$, which is among the open problems.

\subsection{Existence of continuum subsequences and the cube limit}

\begin{lemma}[continuum subsequences exist]\label{lem:s12:LIM0}
\begin{enumerate}[label=\textup{(\alph*)},leftmargin=*,nosep]
\item \textup{(Curvature sector; no hypothesis.)} Every increasing sequence of spacings has a continuum subsequence. For every continuum subsequence $\cK$, family $f$ and $J\ge J_0(f)$: $Z^{\cK}_J(\cdot,0)$ is entire, $|Z^{\cK}_J(\sum_jz_jf_j,0)|\le e^{\norm{g(z)}_2^2}$, $Z^{\cK}_J(x,0)$ is real and $\ge1$ at real $x$, and $0\le Z^{\cK}_J(\mathrm it,0)\le1$ at imaginary $\mathrm it$. \emph{Several models at once:} for every countable set $\cM$ of lattice models of King's type (Neumann--Dirichlet cube states with a site-independent real bare scalar coefficient, photon mass $\mu>0$, $\lambda>0$; for instance the models at the photon masses of a countable set, with the counterterm fixed) one $\cK$ is a continuum subsequence for every member of $\cM$.
\item \textup{(Wick sector, Form~(B).)} Assume \Tthree, \Bcell\ and \SCH. Then every increasing sequence of spacings has a continuum subsequence for the Wick sector, and for every such $\cK$, family $\cF$ and $J$ with $MJ>R(\cF)+\ell_c$, $Z^{a,\cK}_J$ is holomorphic on $\Dtube$ and
\begin{equation}\label{eq:s12:Winf}
|Z^{a,\cK}_J(\zeta)|\le\exp\big(\norm{g(\zeta)}_2^2+\kappa\norm{k_+(\zeta)}_{L^1}+K_-\wnorm{k_-(\zeta)}\big)\qquad(\zeta\in\Dtube),
\end{equation}
with the constants of \eqref{eq:s11:constants}. \emph{Several models at once:} for a countable set of models all satisfying \eqref{eq:s11:W} with one set of constants, one $\cK$ serves all.
\end{enumerate}
\end{lemma}

Form~(A), which needs in addition the existence of the limits at every real coefficient vector (\Wtwo(a)), is \cref{thm:s23:FormA}. The subsequence is chosen in Step~2 below and nowhere else (\cref{rem:s12:K9}).

\begin{proof}
\emph{(a).} Inputs: \cref{lem:s05:C2} ($\sup_{k\ge0,\,J\ge J_0(f)}\sup_K|a^{(a)}_{k,J}|<\infty$ for every family and compact $K\subset\C^N$); \cref{lem:s05:C1} at second order ($\ava{F(u)^2}\le2\norm u^2_{2,k}$ for real $u$ with plaquettes in $\XJ$); charge conjugation $(A,\vartheta)\mapsto(-A,-\vartheta)$, a symmetry of every Neumann--Dirichlet state, so $\ava{F(u)}=0$ and $\ava{e^{\mathrm iF(u)}}$ is real.
\emph{Step 1 (a countable core).} Let $\cB$ be the countable set of finite unions of closed boxes with rational corners, for $B\in\cB$ let $S_0(B)$ be a countable $\norm\cdot_\infty$-dense subset of $C^\infty_c(\operatorname{int}B;\Lambda^2)$, and let $\cP$ be the countable set of pairs $(J,u')$, $u'$ a finite family from $\bigcup_BS_0(B)$, $J\ge J_0(u')$.
\emph{Step 2 (Montel and the diagonal).} For $(J,u')\in\cP$, $(a^{u'}_{k,J})_k$ is locally bounded on $\C^N$ uniformly in $k$ (\cref{lem:s05:C2}), hence normal; enumerate $\cP$, extract successively, and let $\cK=(k_m)$ be the diagonal sequence: $(a^{u'}_{k_m,J})_m$ converges locally uniformly for every $(J,u')\in\cP$.
\emph{Step 3 (every family, imaginary coefficients).} Let $f$ be a family, $J\ge J_0(f)$, and $B\in\cB$ with $\bigcup_j\supp f_j\subset\operatorname{int}B\subset B\subset U_{J_0(f)}$ (so every $u'\in S_0(B)$ is admissible at every $J\ge J_0(f)$). For real $2$-forms $u,u',g$ supported in $B$ and every $k$ with $\eps_k\le1$,
\[
|\ava{e^{\mathrm iF(u)}}-\ava{e^{\mathrm iF(u')}}|\le{\ava{F(u-u')^2}}^{1/2}\le\sqrt2\norm{u-u'}_{2,k},\qquad \norm g^2_{2,k}\le V_B\norm g_\infty^2,\quad V_B\defeq3\operatorname{vol}(N_{1/2}(B)),
\]
the first by $|e^{\mathrm ix}-1|\le|x|$ and the second-order bound, the second because the cubes of side $\eps_k$ about the centres of the plaquettes of one orientation on which $g\ne0$ are disjoint and lie in the closed $\ell^\infty$-neighbourhood $N_{1/2}(B)$. Given $t\in\R^N$ and $\eta>0$ pick $u'\in S_0(B)$ with $\norm{\sum_jt_jf_j-u'}_\infty\le\eta$; then $|a^f_{k_m,J}(\mathrm it)-a^f_{k_{m'},J}(\mathrm it)|\le2\sqrt{2V_B}\,\eta+|{\ava{e^{\mathrm iF(u')}}}_{k_m}-{\ava{e^{\mathrm iF(u')}}}_{k_{m'}}|$, and the last term tends to $0$ by Step~2. So $(a^f_{k_m,J}(\mathrm it))_m$ converges for every $t$.
\emph{Step 4 (all complex coefficients).} $(a^f_{k_m,J})_m$ is locally bounded on $\C^N$ and converges on the maximal totally real $(\mathrm i\R)^N$; Montel--Vitali gives locally uniform convergence to an entire $Z^{\cK}_J(\cdot,0)$. The bound $|a^{(a)}_{k,J}(z)|\le e^{\norm{g(z)}^2_{2,k}}$ (\cref{thm:s05:C}(a)) passes to the limit with $\norm g_{2,k}\to\norm g_2$ \eqref{eq:s02:riemann}; realness at real points passes; $a^{(a)}_{k,J}(x)=\ava{e^{F(u_x)}}\ge e^{\ava{F(u_x)}}=1$ by Jensen, and $[1,\infty)$ is closed; $0\le a^{(a)}_{k,J}(\mathrm it)\le1$ by the current representation of \cref{lem:s02:GD}(a), $\ava{e^{\mathrm i\ip vA}}=e^{-\ip v{Cv}/2}\,\E^\sharp\cosh\ip{Cv}{\kappa}\ge0$, and $|e^{\mathrm iF}|=1$. Dependence only on the total source holds at every $(k,J)$, hence in the limit.
\emph{Step 5 (several models).} The constant of \cref{lem:s05:C2}, the bound $e^{\norm g^2_{2,k}}$ and the Lipschitz constant $\sqrt2$ are model-free: \cref{thm:s05:C}(a) is \cref{lem:s02:GD}(a) for a measure of the class of \cref{lem:s02:BFS} plus \cref{lem:s05:C1} for the Neumann--Dirichlet photon covariance, valid for every real site-independent bare coefficient (the matter weight does not enter) and every $\mu>0$. For countable $\cM$ run Step~2 over $\cM\times\cP$ and Steps~3--4 member by member.

\emph{(b).} Inputs: under \Itwo--\Ifour\ and $B_{\mathrm{cell}}<\infty$, for every family with supports in $Q_R$, every $k$, $J\ge J_1(R)$ and $\zeta\in\Dtube$, \eqref{eq:s11:W} gives $|Z^a_{k,J}(\zeta)|\le Z^a_{k,J}(\Re\zeta)\le\exp(\norm g^2_{2,k}+\kappa\norm{k_+}_{1,k}+K_-\wnorm{k_-})$, with \Itwo\ from \Tthree\ by \cref{prop:s10:U} and \SCH\ for the scheme; on compacts of the tube the right side is bounded uniformly in $k$ (\cref{lem:s02:E}).
\emph{Steps 1--2.} As in (a), with $S^{\sW}_0(B)$ a countable $\norm\cdot_\infty$-dense subset of $C^\infty_c(\operatorname{int}B;\Lambda^2)\times C^\infty_c(\operatorname{int}B)$ and $\cP^{\sW}$ the countable set of pairs $(J,\cF')$, $\cF'$ a finite family from $\bigcup_BS^{\sW}_0(B)$, $MJ>R(\cF')+\ell_c$. For such a pair $J\ge J_1(R)$ for every $R\in(R(\cF'),MJ-\ell_c]$, so \eqref{eq:s11:W} holds for $\cF'$ at every $k$ and $(Z^a_{k,J})_k$ is normal on the tube of $\cF'$; the diagonal over $\cP^{\sW}$, together with that of (a), gives $\cK$.
\emph{Step 3 (every family, imaginary coefficients).} Let $\cF=(f_i;h_j)$, $MJ>R(\cF)+\ell_c$, $R_J\defeq MJ-\ell_c>R(\cF)$, so that all supports lie in the open cube $Q^\circ_{R_J}$ and $J\ge J_1(R_J)$; choose $B\in\cB$ with the supports in $\operatorname{int}B\subset B\subset Q^\circ_{R_J}$, so that $(J,\cF')\in\cP^{\sW}$ for every $\cF'$ from $S^{\sW}_0(B)$. For $(s,s')\in\R^{n+p}$ put $\sigma\defeq(\sum_is_if_i,\sum_js'_jh_j)$, so that $Z^a_{k,J}(\mathrm is,\mathrm is')=\ava{e^{\mathrm i\Phi^a_k(\sigma)}}$ with $\Phi^a_k(g,h)\defeq F_{\Lk}(g)+W^a_{\Lk}(h)$. \emph{Lipschitz bound.} For a real pair $q=(g_q,h_q)$ supported in $Q_{R_J}$, $(t_1,t_2)\mapsto Z^a_{k,J}(t_1g_q,t_2h_q)$ is holomorphic, and on the polydisc $|t_1|\le\norm{g_q}_{2,k}^{-1}$, $|t_2|\le\sigma_0/\wnorm{h_q}$ one has $\norm{k_\pm}_\infty\le|\Re t_2|\,\norm{h_q}_\infty\le\sigma_0<H/2<\sigma_+/2$, so the polydisc lies in the tube and \eqref{eq:s11:W} gives $|Z^a_{k,J}|\le e^{1+K'_{\mathrm{lat}}\sigma_0}$ there, $K'_{\mathrm{lat}}\defeq\kappa C_0'+K_-$ ($\norm{k_+}_{1,k}\le C_0'\wnorm{k_+}$ by \cref{lem:s02:E}). The Taylor coefficients at $0$ are the second moments of $F(g_q)$ and $W^a(h_q)$, and Cauchy's estimates give
\[
{\ava{\Phi^a_k(q)^2}}^{1/2}\le C\,(\norm{g_q}_{2,k}+\wnorm{h_q}),\qquad C\defeq\sqrt2\,e^{(1+K'_{\mathrm{lat}}\sigma_0)/2}\max(1,\sigma_0^{-1}),
\]
uniformly in $k$, $J$ and $R_J$; hence $|Z^a_{k,J}(\mathrm i\sigma)-Z^a_{k,J}(\mathrm i\sigma')|\le C(\norm{g_{\sigma-\sigma'}}_{2,k}+\wnorm{h_{\sigma-\sigma'}})$ for real pairs supported in $Q_{R_J}$. \emph{Density.} Given $\eta>0$ pick $\sigma'\in S^{\sW}_0(B)$ with $\norm{g_\sigma-g_{\sigma'}}_\infty,\norm{h_\sigma-h_{\sigma'}}_\infty\le\eta$; with $V_B$ and $\wnorm h\le(1+3R_J^2)^2\norm h_\infty$ for $h$ supported in $Q_{R_J}$, $|Z^a_{k,J}(\mathrm i\sigma)-Z^a_{k,J}(\mathrm i\sigma')|\le C(\sqrt{V_B}+(1+3R_J^2)^2)\eta$ for every $k$, while $Z^a_{k_m,J}(\mathrm i\sigma')$, the core functional at the point $(\mathrm i,\mathrm i)$ of its tube, converges by Step~2. So $(Z^a_{k_m,J})_m$ converges on $(\mathrm i\R)^{n+p}$.
\emph{Step 4 (the tube).} $(Z^a_{k_m,J})_m$ is locally bounded on the connected open $\Dtube$ and converges on the maximal totally real $(\mathrm i\R)^{n+p}\subset\Dtube$; Montel--Vitali gives locally uniform convergence to a holomorphic $Z^{a,\cK}_J$, and \eqref{eq:s11:W} passes to the limit with $\norm g_{2,k}\to\norm g_2$, $\norm{k_+}_{1,k}\to\norm{k_+}_{L^1}$ (Riemann sums; $k_+$ is continuous with compact support) and sampled $\wnorm{k_-}\le$ continuum $\wnorm{k_-}$: this is \eqref{eq:s12:Winf}.
\emph{Step 5 (several models).} With one set of constants in \eqref{eq:s11:W} for all members, $C$ is model-free and Step~2 runs over the product of the model set with $\cP^{\sW}$.
\end{proof}

\begin{lemma}[the cube limit along $\cK$ is a full limit]\label{lem:s12:LIM1}
Let $\cK$ be a continuum subsequence. \textup{(Curvature sector.)} For every family, $Z^{\cK}_J(\cdot,0)\to\ZLIM{\cK}(\cdot,0)$ locally uniformly on $\C^N$ as $J\to\infty$; the limit is entire, $|\ZLIM{\cK}(\sum_jz_jf_j,0)|\le e^{\norm{g(z)}_2^2}$, real and $\ge1$ at real points, in $[0,1]$ at imaginary points; $J\mapsto Z^{\cK}_J(\mathrm it,0)$ is nondecreasing, so $Z^{\cK}_J\le\ZLIM{\cK}$ on $(\mathrm i\R)^N$. \textup{(Wick sector, Form~(B).)} Under the hypotheses of \cref{lem:s12:LIM0}(b), $Z^{a,\cK}_J\to Z^{a,\cK}_\infty$ locally uniformly on $\Dtube$; the limit is holomorphic, obeys \eqref{eq:s12:Winf}, is $>0$ at the real points of the tube, and $Z^{a,\cK}_J\le Z^{a,\cK}_\infty$ at the cone points ($z\in(\mathrm i\R)^n$, $w$ real, $k(\zeta)\ge0$).
\end{lemma}

\begin{proof}
\emph{Curvature sector.} For $t\in\R^N$, $J\ge J_0(f)$ and every $k$: $a^{(a)}_{k,J}(\mathrm it)\le a^{(a)}_{k,J+1}(\mathrm it)\le1$ (first inequality of \cref{lem:s02:cut} with $Y=C_{J+1}\cap\Lk$, $X=\XJ$; the cut-out factor is the Neumann--Dirichlet state on $\CJ$ because $\dm_\eps$ is site-independent; realness by charge conjugation). Both sides converge along $\cK$, so $J\mapsto Z^{\cK}_J(\mathrm it,0)$ is nondecreasing, bounded by $1$, convergent. $\{Z^{\cK}_J\}_J$ is locally bounded on $\C^N$ by the $J$-free bound $e^{\norm{g(z)}_2^2}$ and converges on $(\mathrm i\R)^N$; Montel--Vitali gives the entire limit with the same bound; realness, $\ge1$ at real points and $[0,1]$ at imaginary points pass to the limit.

\emph{Wick sector.} (1) \emph{Split family.} Write $h_j=h_j^+-h_j^-$, $h_j^-\defeq\norm{h_j}_\infty\chi_j$, $h^+_j\defeq h_j+h^-_j\ge0$, with $\chi_j\in C^\infty_c$, $0\le\chi_j\le1$, $\chi_j=1$ on $\supp h_j$, $\supp\chi_j\subset N_{1/2}(\supp h_j)$; put $F_{k,J}(z,w^+,w^-)\defeq Z^a_{k,J}(\sum_iz_if_i,\sum_j(w^+_jh^+_j+w^-_jh^-_j))$ and $\cD''\defeq\{|\Re w^\pm_j|<\delta_j\}$ with $\sum_j\delta_j(\norm{h^+_j}_\infty+\norm{h^-_j}_\infty)<H/2$, a convex subset of the split family's tube. By \cref{lem:s12:LIM0}(b) for the split family, $F_{k_m,J}\to F^{\cK}_J$ locally uniformly on $\cD''$ for $MJ>R(\cF)+\frac12+\ell_c$, with \eqref{eq:s12:Winf}. (2) On $\Omega_+\defeq\cD''\cap((\mathrm i\R)^n\times(0,\infty)^{2p})$ the observable has imaginary curvature coefficients and nonnegative total scalar source, so $J\mapsto F_{k,J}$ is nondecreasing at every $k$ (first inequality of \cref{lem:s02:cut} with \cref{lem:s02:alg}(a), the scheme multiplier being common); along $\cK$, $J\mapsto F^{\cK}_J$ is nondecreasing and bounded on $\Omega_+$, hence convergent. (3) $\Omega_+$ is a non-empty open subset of the maximal totally real $(\mathrm i\R)^n\times\R^{2p}$ and $\cD''$ is connected: $F^{\cK}_J\to F^{\cK}_\infty$ locally uniformly on $\cD''$. (4) At $w^+=w$, $w^-=-w$ the lattice functionals equal $Z^a_{k,J}(\zeta)$ exactly (equal total sources), so $Z^{a,\cK}_J$ converges locally uniformly on $\cD'\defeq\{|\Re w_j|<\delta_j\}\times\C^n\subset\Dtube$; $\{Z^{a,\cK}_J\}_J$ is locally bounded on $\Dtube$ by the $J$-free \eqref{eq:s12:Winf} and converges on the open $\cD'$; Montel--Vitali on the connected tube gives $Z^{a,\cK}_J\to Z^{a,\cK}_\infty$, holomorphic, with \eqref{eq:s12:Winf}. (5) \emph{Positivity at real tube points.} For real $\zeta\in\Dtube$ pick $\epsilon>0$ with $-\epsilon\zeta\in\Dtube$ and let $C_\epsilon\defeq\sup_k(\norm g^2_{2,k}+\kappa\norm{k_+}_{1,k}+K_-\wnorm{k_-})|_{-\epsilon\zeta}<\infty$, the exponent of \eqref{eq:s11:W} at $-\epsilon\zeta$ (bounded in $k$ by \cref{lem:s02:E}, free of $J$). At every $k$, $t\mapsto\ln Z^a_{k,J}(t\zeta)$ is convex on $\R$ (H\"older, \cref{lem:s02:alg}(d)) and vanishes at $t=0$, so $1\le Z^a_{k,J}(\zeta)^{\epsilon/(1+\epsilon)}Z^a_{k,J}(-\epsilon\zeta)^{1/(1+\epsilon)}\le Z^a_{k,J}(\zeta)^{\epsilon/(1+\epsilon)}e^{C_\epsilon/(1+\epsilon)}$, i.e.\ $Z^a_{k,J}(\zeta)\ge e^{-C_\epsilon/\epsilon}$ for every $k$ and $J$; the bound passes along $\cK$ and then in $J$: $Z^{a,\cK}_\infty(\zeta)\ge e^{-C_\epsilon/\epsilon}>0$. (The product inequality of \cref{lem:s02:alg}(e) is not used: the tube is asymmetric, $H=\sigma_+/16$, and for real $\xi\in\Dtube$ with $\norm{k_+(\xi)}_\infty\ge H/2$ the point $-\xi$ lies outside it, where \cref{def:s12:contsub} gives no limit.) (6) At a cone point of the original family $J\mapsto Z^a_{k,J}(\zeta)$ is nondecreasing at every $k$ (first inequality of \cref{lem:s02:cut} with \cref{lem:s02:alg}(a)); pass along $\cK$.
\end{proof}

\subsection{The axioms for LIM states}

\begin{theorem}[LIM states of the comparison family at fixed photon mass]\label{thm:s12:LIM}
Let $e>0$, $\lambda,m,\lmu>0$, $\eps_k=5^{-k}$, $\CJ=[-MJ,MJ)^3$, and the states $\NDa_{k,J}$ of \cref{def:s02:acube}.

\textbf{Curvature sector (no hypothesis).} Every increasing sequence of spacings has a continuum subsequence. For every continuum subsequence $\cK$ and every finite family of real $2$-forms,
\[
\ZLIM{\cK}\Big(\sum_jz_jf_j,0\Big)=\lim_{J\to\infty}\lim_{m\to\infty}\Big\langle e^{\sum_jz_jF_{\Lambda_{k_m}}(f_j)}\Big\rangle^{(a)}_{k_m,J}
\]
exists locally uniformly in $z\in\C^N$, is entire, and depends on the family only through the total source. Its Schwinger functions $S^{\cK}_n(f_1,\dots,f_n)\defeq\partial_{z_1}\cdots\partial_{z_n}\ZLIM{\cK}(\sum_jz_jf_j,0)|_{z=0}$ satisfy:
\begin{enumerate}[label=\textup{(\roman*)},leftmargin=*,nosep]
\item \textup{(\OSz.)} $|\ZLIM{\cK}(\sum_jz_jf_j,0)|\le e^{\norm{g(z)}_2^2}$ (no dependence on $\lmu,e,\lambda,m$) and $\le\exp(\half\lmu^{-2}\norm{\sum_j(\Re z_j)\dstar f_j}_2^2)$; $|S^{\cK}_n|\le(\sqrt2e)^n(n!)^{1/2}\prod_j\norm{f_j}_2\le(e\pi/2)^n(n!)^{1/2}\prod_j\fournorm{f_j}$; the $S^{\cK}_n$ are tempered distributions, and \Ezz\ holds with $s=4$ and every $\beta>\half$.
\item \textup{(\OSo.)} $S^{\cK}_n$ is invariant under all translations and under $\Oh$ about every point, $F$ transforming as a $2$-form.
\item \textup{(\OSt.)} Reflection positivity about every coordinate plane $\{x_\nu=c\}$: for exponentials $E_a=e^{\sum_lz_{al}F(f_{al})}$ with arbitrary complex $z_{al}$ and $f_{al}$ supported in $\{x_\nu>c\}$, $(S(\overline{\theta_cE_a}E_b))_{ab}$ is positive semidefinite, and $S(\theta_c(P)P)\ge0$ for polynomials $P$ in such fields.
\item \textup{(\OSth.)} $S^{\cK}_n$ is real, symmetric, and zero for odd $n$.
\item \textup{(The measure.)} For every even $r\ge14$ there is exactly one Borel probability measure $\mu^{\cK}_F$ on $\cH^{-r}(\R^3;\Lambda^2)$ whose moments on $C^\infty_c$ are the $S^{\cK}_n$; its characteristic functional is $\ZLIM{\cK}(\mathrm i\,\cdot,0)$ extended by continuity to $\cS$, and its image on $H^{-r}_{\mathrm{loc}}$ is the unique measure there with this characteristic functional and the unique one with these moments. For every $\chi\in C^\infty_c$ and box $Q$ containing $\supp\chi$ in its interior, the laws of $\chi F_{\Lambda_{k_m}}$ under $\NDa_{k_m,J}$ converge weakly on $H^{-r}_{\mathrm{per}}(Q;\Lambda^2)$, \emph{first as $m\to\infty$ along $\cK$} at every fixed $J\ge J_0(\chi)$ and then as $J\to\infty$, to the law of $\chi\Phi$ under $\mu^{\cK}_F$.
\item \textup{(Weak \OSf.)} $\mu^{\cK}_F$ is $\Cub$-invariant and satisfies the clauses of \cref{thm:s09:S2K}(b) with $\mu^{\cK}_F$ in place of $\mu^K_S$: it is an integral of translation-ergodic, reflection-positive, $\theta_j$-invariant components clustering in the mean and under diagonal separation, with a unique joint vacuum per coordinate direction and the component growth bound; almost every component satisfies \OSf\ in the weak form; the $\Oh$-invariant decomposition into mixtures of $k\in\{1,2,3,6\}$ ergodic measures exists.
\item \textup{(Euclidean clause.)} $\ZLIM{\cK}$ does not depend on $\cK$, equals the limit state $Z^{\mathrm{LIM}}_\infty$ of \cref{def:s12:o3LIM}, is $\Euc$-invariant for all complex coefficients (\cref{hyp:s03:RC}), and the Euclidean forms of (ii) (full $\Euc$), (iii) (reflection positivity about every plane) and (vi) (almost every component $\Euc$-invariant, reflection positive about every plane, clustering in any manner, with a unique vacuum in every direction; the vanishing of a component's odd moments is not claimed) hold.
\end{enumerate}

\textbf{Wick sector, Form~(B).} Assume \Tthree, \Bcell\ and \SCH. Every increasing sequence of spacings has a continuum subsequence for the Wick sector, and for every such $\cK$ and family the LIM state $Z^{a,\cK}_\infty$ exists locally uniformly on $\Dtube$, is holomorphic there, obeys \eqref{eq:s12:Winf}, and (i)--(vi) hold for the joint Schwinger functions $S^{\cK}_{n,p}(f;h)=\partial^{(n)}_z\partial^{(p)}_wZ^{a,\cK}_\infty(0)$ with: in (i) the mixed bounds $|S^{\cK}_{n,p}|\le e^{K'\sigma_0}(\sqrt2e)^n(n!)^{1/2}(e/\sigma_0)^p\,p!\prod_i\norm{f_i}_2\prod_j\wnorm{h_j}$, $\sigma_0=H/4$, $K'=\kappa\pi^2+K_-$, and \Ezz\ with $s=4$ and every $\beta>1$; in (ii) $W^a$ a scalar, the identity holding on the tube; in (iii) exponentials with arbitrary complex curvature coefficients and Wick coefficient vectors $\zeta_1,\dots,\zeta_N$ of \emph{amplitude $\amp(\zeta_i)<H/2$} (in particular all purely imaginary Wick coefficients), and the polynomial form; in (iv) $S^{\cK}_{n,p}=0$ for odd $n$; in (v) the measures on $\cH^{-r}(\R^3;\Lambda^2\oplus\R)$ and $H^{-r}_{\mathrm{loc}}$, the localised lattice fields $\chi(F_{\Lk},W^a_{\Lk})$ converging along $\cK$ and then in $J$; (vi) unchanged; and \emph{no Euclidean clause}. \emph{Schemes:} if $a,a'$ both satisfy the hypotheses, every continuum subsequence has a further subsequence $\cK'$ along which $a'_k-a_k\to\kappa^{\cK'}$, and $Z^{a',\cK'}_\infty(\zeta)=e^{-\kappa^{\cK'}\sum_jw_j\int h_j}Z^{a,\cK'}_\infty(\zeta)$: the two LIM states differ by the shift $W\mapsto W-\kappa^{\cK'}\int(\cdot)$ of the Wick component, and joint cumulants of order $\ge2$ are scheme-free.
\end{theorem}

\begin{proof}
Each clause is the corresponding proof for the iterated limit (\cref{thm:s05:C,prop:s06:translations,thm:s06:S1K,thm:s08:OS2C,thm:s08:OS2WT,prop:s08:OS3,thm:s09:S2K}) with exactly one of three replacements: ($\alpha$) a fixed-$(k,J)$ identity or inequality between Neumann--Dirichlet functionals of explicit finite families is passed along $\cK$ (both sides converge by \cref{lem:s12:LIM0}, which covers \emph{every} family); ($\beta$) Montel on a $k$-uniform bound, the identification of the limit by convergence at real points being replaced by the choice of $\cK$; ($\gamma$) a step about the limit object only, unchanged and not repeated here. The only uniqueness used is that of the monotone cube limit along $\cK$ (\cref{lem:s12:LIM1}).

\emph{(i).} \cref{lem:s12:LIM0,lem:s12:LIM1} give existence, entireness and the first bound; the second, $\exp(\half\lmu^{-2}\norm{\dstar_kg}^2_{2,k})$ at every $(k,J)$ (\cref{thm:s05:C}(a)), passes along $\cK$ and $J$ with $\norm{\dstar_kg}_{2,k}\to\norm{\dstar g}_2$. The Schwinger bounds, the extension to $\cS$, the kernel theorem and \Ezz\ are the arguments of \cref{thm:s05:C} (Wick sector: \cref{thm:s11:W}(ii), on the polydisc inside the tube) run on the limit function with its bound ($\gamma$).

\emph{(ii), translations.} For $b\in5^{-k_0}\Z^3$, $|b|_\infty\le M$, $k\ge k_0$: $Z^a_{k,\CJ}(\tau_b\cdot\zeta)=Z^a_{k,\CJ-b}(\zeta)$ (exact lattice covariance) and $C_{J-1}\cap\Lk\subset(\CJ-b)\cap\Lk\subset C_{J+1}\cap\Lk$, so for a cone point $s$ supported in $U_{J_0}$ and $J\ge J_0+1$, $Z^a_{k,C_{J-1}}(s)\le Z^a_{k,\CJ}(\tau_b\cdot s)\le Z^a_{k,C_{J+1}}(s)$ (first inequality of \cref{lem:s02:cut}, twice). The three terms are functionals of the families $s$ and $\tau_b\cdot s$ (supports in $Q_{R+M}$, the same tube), so they converge along $\cK$ ($\alpha$), and \cref{lem:s12:LIM1} sends both ends to $\ZLIM{\cK}(s)$: $\ZLIM{\cK}(\tau_b\cdot s)=\ZLIM{\cK}(s)$ at cone points, with no uniqueness of the $k$-limit. All coefficients ($\gamma$): the difference of the two limit functions is entire and vanishes on $(\mathrm i\R)^N$ (curvature), resp.\ for the split family is holomorphic on the connected $\cD''$ and vanishes on $\Omega_+$, so the identity theorem, the diagonal (4) of \cref{lem:s12:LIM1} and the identity theorem on the connected tube give it on $\Dtube$. All $b\in\bigcup_k5^{-k}\Z^3$ by iteration; all $b\in\R^3$ at imaginary coefficients because the Lipschitz bounds of \cref{lem:s12:LIM0}, Step~3, pass to the limits ($|\ZLIM{\cK}(\mathrm iu,0)-\ZLIM{\cK}(\mathrm iu',0)|\le\sqrt2\norm{u-u'}_2$; $|Z^{a,\cK}_\infty(\mathrm i\sigma)-Z^{a,\cK}_\infty(\mathrm i\sigma')|\le C(\norm{g_{\sigma-\sigma'}}_2+\wnorm{h_{\sigma-\sigma'}})$) and $\tau_{b_n}s\to\tau_{b_0}s$ in $L^2$ and in $\wnorm\cdot$ for $s\in C^\infty_c$; complex coefficients by the identity theorem.

\emph{(ii), the cubic group.} Coordinate permutations are exact at every $(k,J)$ (\cref{lem:s02:pass}); pass along $\cK$ and $J$ ($\alpha$). For the reflection $\sigma_3$ and a cone point $(f,g)$ supported in $U_{J-1}$, at every $k$, $Z^a_{k,C_{J-1}}(\mathrm if,g)\le Z^a_{k,\CJ}(\mathrm i\sigma_3\cdot f,\sigma_3\cdot g)\le Z^a_{k,C_{J+1}}(\mathrm if,g)$ (the frame flip of \cref{lem:s02:pass} with the one-site shift $X_+=(\CJ+\eps_ke_3)\cap\Lk$ and two deletions in the cone direction, \cref{lem:s02:BFS}(a),(b)); the middle term is the functional of the reflected family (same tube), so ($\alpha$) and \cref{lem:s12:LIM1} give $\ZLIM{\cK}(\mathrm i\sigma_3\cdot f,\sigma_3\cdot g)=\ZLIM{\cK}(\mathrm if,g)$ on the cone (all of it at $g=0$; in Form~(B) at cone points with $\norm g_\infty<\sigma_+/2$): the reflections are not exact lattice symmetries; the sandwich replaces exactness and survives along $\cK$. Conjugation $\sigma_\nu=P\sigma_3P^{-1}$, group closure on the cone, the identity theorem off the cone (on $\C^N$ with $(\mathrm i\R)^N$; on $\cD''$ with $\Omega_+$), invariance about every point by the translations, and differentiation are ($\gamma$), as in \cref{thm:s06:S1K}.

\emph{(iii), reflection positivity.} Fix $\nu$, a mid-plane $c\in D$ (\cref{def:s07:reflection}) and the enlarged family $\cF$ of \cref{def:s08:enlarged} with the modified-box functionals $\Phi_k(\zeta)=Z^{a,(3.38)}_{k,\Omega_{J,k}}(\zeta)$, $\zeta\in\C^m$, and support radius $R_\cF$. The fixed-$k$ steps (reflection positivity of the modified box, \cref{lem:s07:RPbox}; the sandwich lemmas \cref{lem:s07:L,lem:s07:U}) concern the photon form and the inequalities of \cref{lem:s02:BFS} only and are unchanged. \emph{Curvature sector.} At $h=0$, $|\Phi_k(z)|\le e^{\norm{f_{\Re z}}^2_{2,k}}$ by \cref{lem:s02:GD}(a) and the modified photon form, uniformly in $k$ and $J$. ($\beta$) Extract from $\cK$ a further $\cK'$ along which $\Phi_k\to\Phi^J$ locally uniformly on $\C^m$; positive semidefiniteness of the Gram matrices at every $k\ge k_0(c,\delta)$ passes to $\Phi^J$. ($\alpha$) On $U$ (imaginary curvature coefficients) $a^{(a)}_{k,\CJ}(\zeta)\le\Phi_k(\zeta)\le a^{(a)}_{k,C_{J+2n}}(\zeta)$ at every $k$ (\cref{lem:s07:L,lem:s07:U}, cone halves); along $\cK'$ the ends converge to $Z^{\cK}_J,Z^{\cK}_{J+2n}$, so $Z^{\cK}_J\le\Phi^J\le Z^{\cK}_{J+2n}$ on $U$, and \cref{lem:s12:LIM1} sends both ends to $\ZLIM{\cK}$. $\{\Phi^J\}_J$ is locally bounded on $\C^m$ ($J$-free bound) and $U$ is a set of uniqueness, so $\Phi^J\to\ZLIM{\cK}$ locally uniformly, in particular at the matrix points, where $\Phi^J(\zeta^*_{ij})\to\ZLIM{\cK}(\theta\zeta^*_i\sqcup\zeta_j)$ (dependence through the total source only); positivity passes. $\cK'$ may depend on $J$; this is harmless, since the sandwich ends converge along all of $\cK$ and the bound on $\Phi^J$ is $J$-free, so neither $\lim_J\Phi^J$ nor the positivity depends on the choice. Every plane: translate by $(c'-c)e_\nu$ and use (ii). The polynomial form: Weierstrass along $\cK'$ and $J$ (\cref{prop:s08:poly}). \emph{Wick sector, amplitude $<H/2$, without \cref{hyp:s03:W1}.} Let $\zeta_1,\dots,\zeta_N$ be families of real sources supported in $\{x_\nu>c\}$ with $\amp(\zeta_i)<H/2$; the $(i,j)$ Gram entry is $\Phi_k(\zeta^*_{ij})$, with total scalar source $\sum_l\bar w_{il}\,h_{il}\circ\rho+\sum_lw_{jl}h_{jl}$. The bound of \cref{lem:s08:claim4}, which depends on \cref{hyp:s03:W1}, is replaced by \cref{lem:s12:claim4prime} below, after which the argument is completed.

\emph{(iv).} Symmetry: Taylor coefficients of a function of the total source. Realness: real at real points at every $(k,J)$, closed under limits (Form~(B): at the real points of the tube). Charge conjugation gives $Z^a_{k,J}(-z,w)=Z^a_{k,J}(z,w)$ at every $(k,J)$, hence for $Z^{\cK}_J$ and $\ZLIM{\cK}$ (the tube is invariant under $z\mapsto-z$); odd $z$-derivatives at $0$ vanish.

\emph{(v), (vi).} The moment bounds at $(k,J)$ (\cref{thm:s05:C}(a), resp.\ \eqref{eq:s11:W}, with Cauchy's estimates) are uniform in $k$ and $J$ above the thresholds, and the lattice Schwinger functions converge along $\cK$ then $J$ (Weierstrass). The measure is built from the limit functional ($\gamma$): $\cZ^{\cK}(f)\defeq\ZLIM{\cK}(\mathrm if)$ on real $f\in C^\infty_c$ is normalised, positive definite (at every $(k,J)$ the matrix $(Z^a_{k,J}(\mathrm i(f_m-f_l)))_{lm}$ is a Gram matrix; the entries converge along $\cK$ then $J$) and Lipschitz for $\sqrt2\norm\cdot_2$, resp.\ $C(\norm\cdot_2+\wnorm\cdot)$, hence continuous on $\cS$; Bochner--Minlos, moments by Lukacs \cite{Luk70} and uniform integrability, the Hermite carrier on $\cH^{-r}$, $r\ge14$ even, and the two uniqueness statements are the steps of \cref{thm:s09:S2K}(a) with \cite[lem:D:hermite(c), lem:D:realize, prop:discharge]{SymPhase}. \emph{Local laws along $\cK$} (($\beta$); the one weakened clause): tightness of the laws of $\chi\Phi^a_{k,J}$ on $H^{-r}_{\mathrm{per}}(Q;\cV)$ is uniform in $k$ and $J\ge J(\chi)$ from the second-moment bounds (Prokhorov \cite{Bil99}); the characteristic functionals $Z^a_{k,J}(\mathrm i\chi\psi)$ converge \emph{along $\cK$} to $Z^{a,\cK}_J(\mathrm i\chi\psi)$; a probability measure on $H^{-r}_{\mathrm{per}}$ is determined by its characteristic functional on trigonometric polynomials; so the laws converge along $\cK$, then in $J$ (\cref{lem:s12:LIM1}), to the law of $\chi\Phi$ under the measure just built (the periodisation identity of \cref{thm:s09:S2K}). (vi): $\Cub$-invariance of the measure is (ii) at imaginary coefficients plus uniqueness of the characteristic functional; (H3) of \cite{SymPhase} is (iii) at purely imaginary Wick coefficients ($\amp=0<H/2$) with the density of the exponential algebra \cite[lem:A:mc]{SymPhase}; then \cite[thm:D:decomp, thm:D:mean, thm:D:diag, prop:D:recon, cor:D:schwinger(iv), add:Ohdecomp, add:OS4(iii)]{SymPhase}, stated for a measure satisfying (H1)--(H3) and matched as in \cref{rem:s09:v1match}.

\emph{Schemes.} $\sup_k|a'_k-a_k|\le\sup_k|a'_k-z_k|+\sup_k|z_k-a_k|<\infty$ by \SCH\ for both. For a continuum subsequence $\cK$ for $a$, Bolzano--Weierstrass gives $\cK'\subset\cK$ with $a'_k-a_k\to\kappa^{\cK'}\in\R$. At every $(k,J)$, $Z^{a'}_{k,J}(\zeta)=e^{-\sum_jw_j(a'_k-a_k)\sS_{\Lk}(h_j)}Z^a_{k,J}(\zeta)$ (\cref{rem:s02:scheme}) with $\sS_{\Lk}(h_j)\to\int h_j$; so along $\cK'$ the $a'$-functionals converge wherever the $a$-functionals do, $\cK'$ is a continuum subsequence for $a'$ with $Z^{a',\cK'}_J=e^{-\kappa^{\cK'}\sum_jw_j\int h_j}Z^{a,\cK'}_J$, and then $J\to\infty$. (For admissible $a,a'$, $a'_k-a_k$ converges along the full sequence and $\kappa^{\cK'}=\kappa$.) The cumulant statement is that of \cref{thm:s13:main}. \Bcell\ transfers between schemes with bounded difference: $B^{a'}_{\mathrm{cell}}\le B^a_{\mathrm{cell}}+\sup_k|a'_k-a_k|\,\norm{2H'\chi_0}_{1,k}$.

\emph{(vii).} By \cref{prop:s12:ID}(a), $\ZLIM{\cK}=Z^{\mathrm{LIM}}_\infty$, which is $\Euc$-invariant for all $z\in\C^N$ by \cref{hyp:s03:RC}. (ii) in Euclidean form follows by differentiation; for (iii), a plane $P$ is $E\{x_\nu=c\}$ for a Euclidean motion $E$, $\theta_P=E\theta_{\nu,c}E^{-1}$, and $\ZLIM{\cK}(\theta_P\zeta^*_i\sqcup\zeta_j)=\ZLIM{\cK}(\theta_{\nu,c}(E^{-1}\zeta_i)^*\sqcup E^{-1}\zeta_j)$ by invariance at complex coefficients, the $E^{-1}\zeta_i$ being supported in $\{x_\nu>c\}$; (vi) is \cite[p:cor:rotcor]{SymPhase} with $\cZ^{\cK}$ $O(3)$-invariant, as in \cref{thm:s09:S2K}(c), with its proviso on the odd moments of a component.
\end{proof}

\begin{lemma}[a $k$- and $J$-uniform bound for the modified-box functionals on the tube of the enlarged family]\label{lem:s12:claim4prime}
Under \Tthree, \Bcell, \SCH, for every $k$, every $J\ge\max(J_0,J_1(R_\cF))$ and every $\zeta\in\cD^{\cF}_{\sigma_+,H}$ (the tube of $\cF$ in the sense of \cref{def:s02:tube}),
\[
|\Phi_k(\zeta)|\le\exp\big(2\norm{g(\zeta)}^2_{2,k}+\kappa\norm{k_+(\zeta)}_{1,k}+K_-\wnorm{k_-(\zeta)}\big).
\]
\end{lemma}
The curvature constant is $2$, not the $1$ of \eqref{eq:s11:W}, because the Cauchy--Schwarz split below doubles the curvature source and \cref{thm:s05:C}(a) is quadratic in it; nothing depends on its value.

\begin{proof}
$|\Phi_k(\zeta)|\le\Phi_k(\xi)$, $\xi\defeq\Re\zeta$ (\cref{lem:s02:alg}(d)). Let $k(\xi)$ be the total real scalar source at $\xi$ and $k_1\defeq k(\xi)_+$, $k_2\defeq k(\xi)_-$ its \emph{pointwise} parts, sampled on the lattice. Cauchy--Schwarz at the lattice level (\cref{lem:s02:alg}(f); the multiplier $e^{-a_k\sS(k_1-k_2)}$ splits by linearity of $\sS$): $\Phi_k(\xi)\le Z^{a,(3.38)}(2f_\xi,-2k_2)^{1/2}Z^{a,(3.38)}(0,2k_1)^{1/2}$. \emph{First factor} (class $\cD$): the $\cD$-half of \cref{lem:s07:L}, which needs only that the sources vanish off $\XJ$ ($J\ge J_0$), gives $\le Z^a_{k,\CJ}(2f_\xi,-2k_2)$; \cref{thm:s05:C}(a) with the real scalar source $-2k_2$ gives $\le e^{\norm{2f_\xi}^2_{2,k}}N^a_{k,\CJ}(0,-2k_2)$ (the same scheme multiplier on both sides); \cref{prop:s11:W8} with $q\defeq2k_2\ge0$, supported in $Q_{R_\cF}$, $J\ge J_1(R_\cF)$ and $\sup q=2\norm{k(\xi)_-}_\infty<H$ (this is $\zeta\in\cD^\cF_{\sigma_+,H}$; the proposition allows $\sup q\le H$) gives $\ln N^a_{k,\CJ}(0,-2k_2)\le K_-\wnorm{2k_2}$. Hence $\half\ln(\text{first factor})\le2\norm{g(\zeta)}^2_{2,k}+K_-\wnorm{k_-(\zeta)}$, as $f_\xi=g(\zeta)$ and $k_2=k_-(\zeta)$. \emph{Second factor} (class $\cI$): the $\cI$-half of \cref{lem:s07:U} with $B=C_{J+2n}\cap\Lk$ gives $\le Z^a_{k,C_{J+2n}}(0,2k_1)$, and \cref{cor:s11:W2prime}, uniform in $k$ and in the cube, with $0\le2k_1\le s_1$ on the lattice ($\norm{k_+(\zeta)}_\infty<\sigma_+/2$, sampled $\le$ continuum) gives $\ln Z^a_{k,C_{J+2n}}(0,2k_1)\le\kappa\norm{2k_1}_{1,k}$, so $\half\ln(\text{second factor})\le\kappa\norm{k_+(\zeta)}_{1,k}$. Add. On compacts of the tube the right side is bounded uniformly in $k$ (\cref{lem:s02:E}).
\end{proof}

The lemma is \eqref{eq:s11:W} (with curvature constant $2$) for the modified-box functional, transported through \cref{lem:s07:L,lem:s07:U}; the two factors live on $\CJ$ and $C_{J+2n}$, which is harmless as \cref{prop:s11:W8} and \cref{cor:s11:W2prime} are uniform in the cube, and the amplitude $H/2$ is exactly the tube condition $\norm{k(\xi)_-}_\infty\le H/2$ on the total source. \emph{The matrix points lie in the tube:} at $\xi^*_{ij}=\Re\zeta^*_{ij}$ the two sums of the total scalar source have disjoint supports, so $\norm{k(\xi^*)_\pm}_\infty\le\max(\amp(\zeta_i),\amp(\zeta_j))<H/2\le\sigma_+/2$, and by continuity a neighbourhood of the finitely many matrix points lies in $\cD^\cF_{\sigma_+,H}$. \emph{End of the proof of (iii), Wick sector.} ($\beta$) By \cref{lem:s12:claim4prime}, $(\Phi_k)_k$ is locally bounded on the open connected $\cD^\cF_{\sigma_+,H}$; extract $\cK'\subset\cK$ with $\Phi_k\to\Phi^J$ locally uniformly there; Gram positivity passes ($\cK'$ may depend on $J$; harmless as in the curvature sector). ($\alpha$) Let $U$ be the set of coefficient vectors with purely imaginary curvature coefficients and positive real scalar coefficients on the nonnegative $\tilde h_j$; $U\cap\cD^\cF\ne\emptyset$ ($k_-=0$, $k_+$ small), and it is a relatively open subset of the maximal totally real $(\mathrm i\R)^{m_F}\times\R^{m_h}$, a set of uniqueness in $\cD^\cF$. On $U$, $Z^a_{k,\CJ}(\zeta)\le\Phi_k(\zeta)\le Z^a_{k,C_{J+2n}}(\zeta)$ at every $k$ (\cref{lem:s07:L,lem:s07:U}, $\cI$-halves); the ends are functionals of $\cF$ at points of its tube, converging along $\cK'$ (\cref{lem:s12:LIM0}(b), $J$ above the threshold of $\cF$) to $Z^{a,\cK}_J,Z^{a,\cK}_{J+2n}$, which tend to $Z^{a,\cK}_\infty$ (\cref{lem:s12:LIM1}, locally uniformly on the tube of $\cF$, in particular on $U\cap\cD^\cF$). $\{\Phi^J\}_J$ is locally bounded on $\cD^\cF$ ($J$-free), so $\Phi^J\to Z^{a,\cK}_\infty$ locally uniformly there, in particular at the matrix points, where the limit is $Z^{a,\cK}_\infty(\theta\zeta^*_i\sqcup\zeta_j)$; positivity passes. Every plane by (ii) (the amplitude is translation invariant); the polynomial form by Weierstrass.

\begin{remark}[amplitudes in $[H/2,\sigma_+/8)$: the route through the fixed-cube limits, for a further subsequence]\label{rem:s12:amplitudes}
The proof of \cref{lem:s08:claim4} bounds the first factor on the Wick tube $T_{\cF'}$ of \cref{def:s08:enlarged} by the vertex argument at the cube $C_{J_0}$, which needs the fixed-cube limits in their boundedness form only. For amplitudes in $[H/2,\sigma_+/8)$ the matrix points lie in $T_{\cF'}$ but in general outside $\cD^{\cF}_{\sigma_+,H}$, the only domain on which \cref{def:s12:contsub} gives convergence along $\cK$; what follows is therefore stated here only for a further subsequence. Assume the boundedness in $k$ of the functionals of the class-$\cD$ families $(2f_\xi,-2k_2(\xi))$ at the one cube $C_{J_0}$ (with the coefficient-sign split of \cref{def:s08:enlarged}). The vertex argument then gives local boundedness of $(Z^a_{k,J})_k$ and $(\Phi_k)_k$ on $T_{\cF'}$ (first factors at $C_{J_0}$ by that boundedness, second factors by \cref{cor:s11:W2prime}); Montel, with a diagonal over the countably many $J$, yields a further subsequence $\cK'\subset\cK$, depending on the family, along which $Z^a_{k,J}\to Z^{a,\cK'}_J$ and $\Phi_k\to\Phi^J$ locally uniformly on $T_{\cF'}$ (on $T_{\cF'}\cap\cD^\cF$ the first limit is $Z^{a,\cK}_J$). The sandwich step runs on $T_{\cF'}$ with the uniqueness set $U\cap T_{\cF'}$, the $J$-free local bound of \cref{prop:s08:T}(a) on $T_{\cF'}$ serving the Vitali step in $J$, and \cref{prop:s08:T}(b) along $\cK'$ for the remaining planes; so \OSt\ holds up to $\amp(\zeta_i)<\sigma_W=\sigma_+/8$ for the LIM state of such a $\cK'$ on $T_{\cF'}$. Whether the convergence holds along $\cK$ itself on $T_{\cF'}$ is not claimed here; under the hypotheses of Form~(A), \cref{thm:s23:FormA} gives \OSt\ along $\cK$ with no amplitude restriction. The remark is used nowhere else.
\end{remark}

\subsection{Identification, and King's printed state}

\begin{proposition}[identification]\label{prop:s12:ID}
\begin{enumerate}[label=\textup{(\alph*)},leftmargin=*,nosep]
\item For every continuum subsequence $\cK$ and every family, $\ZLIM{\cK}(\cdot,0)=Z^{\mathrm{LIM}}_\infty(\cdot,0)$, the limit state of the rotation theorem (\cref{hyp:s03:RC}); in particular all LIM states coincide, \UNIQ\ holds, and they are $\Euc$-invariant. \Cref{hyp:s03:RC} does \emph{not} give \cref{hyp:s03:EC}.
\item The following are equivalent: (i) the comparison family is convergent at fixed cubes (\cref{def:s03:conv}), that is the existence in $(0,\infty)$ of the fixed-cube limits of the comparison family at real coefficients; (ii) for every $J$ and family, $a^{(a)}_{k,J}\to Z_J$ locally uniformly on $\C^N$ along the full sequence; (iii) every increasing sequence of spacings is a continuum subsequence. Under them $Z^{\cK}_J=Z_J$ and $\ZLIM{\cK}=Z_\infty\defeq\lim_JZ_J$ for every $\cK$: every LIM state is the iterated limit of the comparison family.
\item Assume \Tthree\ and let $a$ be admissible. Then \SCH\ holds (\cref{def:s03:adm}) and \Bcell\ holds (\cref{hyp:s03:W1} for $a$ at the one cell and source), every increasing sequence of spacings is a continuum subsequence for the Wick sector, and for every $\cK$, $Z^{a,\cK}_\infty=\Ginf$ of \cref{thm:s11:W} on $\Dtube$, which equals $Z^a_{\infty,T}$ of \cref{prop:s08:T} on $T_{\cF'}\cap\Dtube$ (\cref{prop:s08:T} lives on the Wick tube $T_{\cF'}$, a different domain). Once \SCH\ and \Bcell\ are given, the argument uses \cref{hyp:s03:W1} for $a$ alone.
\end{enumerate}
\end{proposition}

\begin{proof}
(a) Let $\cK$ be a continuum subsequence, $f$ a family, $J\ge J_0(f)$. By \cref{def:s12:contsub}, $Z^{\cK}_J(\cdot,0)$ is the locally uniform limit of $(a^{(a)}_{k_m,J})_m$, so $Z^{\cK}_J(\cdot,0)\in L^f_J$. By \cref{hyp:s03:RC} every choice $\zeta_J\in L^f_J$, $J\ge J_0(f)$, converges locally uniformly to $Z^{\mathrm{LIM}}_\infty(\cdot,0)$; take $\zeta_J\defeq Z^{\cK}_J(\cdot,0)$, which by \cref{lem:s12:LIM1} converges to $\ZLIM{\cK}(\cdot,0)$, and limits are unique. Hence any two continuum subsequences have the same LIM state, and invariance is the second clause of \cref{hyp:s03:RC}. That theorem pins the $J$-limit of the subsequential limits, not the fixed-cube limits themselves, so it does not give \cref{hyp:s03:EC}.

(b) (i)$\Rightarrow$(ii): $(a^{(a)}_{k,J})_k$ is locally bounded on $\C^N$ (\cref{lem:s05:C2}) and converges at every real point; $\R^N$ is a set of uniqueness for entire functions, and Montel--Vitali gives locally uniform convergence of the whole sequence to an entire $Z_J$. (ii)$\Rightarrow$(iii): every subsequence converges to the same limit. (iii)$\Rightarrow$(i): the full sequence is itself an increasing sequence, hence a continuum subsequence, so $a^{(a)}_{k,J}$ converges locally uniformly, and at real $x$ the limit lies in $[1,\infty)\subset(0,\infty)$ (\cref{lem:s12:LIM0}, Step~4). The last clause is \cref{lem:s12:LIM1}.

(c) Admissibility gives $\sup_k|z_k-a_k|<\infty$, hence \SCH; \cref{hyp:s03:W1} for $a$ at the cell gives $B_{\mathrm{cell}}<\infty$; so \eqref{eq:s11:W} holds. $(Z^a_{k,J})_k$ is locally bounded on the tube by \eqref{eq:s11:W} and converges at the real points of the tube by \cref{hyp:s03:W1}, a non-empty open subset of $\R^{n+p}$ and a set of uniqueness in the connected tube; Montel--Vitali: the whole sequence converges locally uniformly on the tube (\cref{thm:s11:W}(i)), so every increasing sequence is a continuum subsequence for the Wick sector with $Z^{a,\cK}_J=G_J$, and $Z^{a,\cK}_\infty=\Ginf$ by \cref{lem:s12:LIM1}; on $T_{\cF'}\cap\Dtube$ this is the same $J$-limit of the same $G_J$ as $Z^a_{\infty,T}$ (\cref{prop:s08:T}(a)).
\end{proof}

\begin{corollary}[the iterated limits of the comparison family]\label{cor:s12:K}
Assume that the comparison family is convergent at fixed cubes (\cref{def:s03:conv}). Then every increasing sequence of spacings is a continuum subsequence, all LIM states coincide with the iterated limit $Z_\infty(f,0)=\lim_J\lim_k\ava{e^{F(f)}}$ of the comparison family, and the curvature sector (i)--(vi) of \cref{thm:s12:LIM} holds for it with the lattice laws of (v) converging along the \emph{full} sequence of spacings; (vii) holds as well. For general sources assume \Tthree\ and let $a$ be admissible; then every increasing sequence is a continuum subsequence for the Wick sector, every LIM state is $\Ginf$ ($=Z^a_{\infty,T}$ on $T_{\cF'}\cap\Dtube$), $\kappa^{\cK}=\kappa$ is independent of $\cK$, and Form~(B) of \cref{thm:s12:LIM} holds for $\Ginf$ with the full-sequence lattice clause.
\end{corollary}

\begin{proof}
Curvature sector: \cref{prop:s12:ID}(b) makes every increasing sequence a continuum subsequence with $\ZLIM{\cK}=Z_\infty$; (i)--(vi) are those of \cref{thm:s12:LIM}, and in (v) the local laws converge along the full sequence because the characteristic functionals now do (step ($\beta$) with the full sequence); (vii) is \cref{prop:s12:ID}(a). Wick sector: \cref{prop:s12:ID}(c); the clauses are Form~(B) of \cref{thm:s12:LIM} for $\Ginf$, with \OSt\ at amplitude $<H/2$ by the route of \cref{lem:s12:claim4prime}, which does not use \cref{hyp:s03:W1}, and at amplitude $<\sigma_W=\sigma_+/8$ by the route of \cref{thm:s08:OS2WT}; both hold ($H/2=\sigma_+/32<\sigma_+/8$: the route without the fixed-cube limits gives the smaller amplitude, and what it improves is the list of inputs); $\kappa^{\cK}=\kappa$ by the scheme clause.
\end{proof}

\Cref{hyp:s03:W1} is the scheme-$a$ analogue, with scalar sources, of the existence clause of King's theorem; it is not that theorem, which is false as printed with scalar sources in King's Wick ordering. With \Wtwo\ in addition one obtains Form~(A) for the full object $Z^a_\infty$ and, under \RW, its Euclidean clause; this is the optional extension of \cref{thm:s23:FormA}.

\begin{corollary}[King's convention]\label{cor:s12:KK}
Define continuum subsequences and LIM states of the K-family $a^K_{k,J}(z)=\avK{e^{\sum_jz_jF(f_j)}}$ (\cref{def:s02:Kcube}) by \cref{def:s12:contsub,def:s12:LIMstate} with $\avK\cdot$ in place of $\ava\cdot$.
\begin{enumerate}[label=\textup{(\alph*)},leftmargin=*,nosep]
\item \textup{(No hypothesis.)} Every increasing sequence of spacings has a continuum subsequence of the K-family; for every such $\cK$, family and $J\ge J_0(f)$ the fixed-cube limits $Z^{K,\cK}_J$ and the LIM state $Z^{K,\cK}_\infty=\lim_JZ^{K,\cK}_J$ exist, are entire and obey $|Z^{K,\cK}_J|,|Z^{K,\cK}_\infty|\le e^{\norm{g(z)}_2^2}$, with the other clauses of \cref{lem:s12:LIM0}(a) and \cref{lem:s12:LIM1}. Nothing beyond this is asserted for a K-family LIM state from the bounds alone.
\item \textup{(With \cref{hyp:s03:EC}.)} Every increasing sequence of spacings is a continuum subsequence of the K-family, every K-family LIM state is King's printed state $\ZKinf$, every LIM state of the comparison family is $\ZKinf$ as well (\cref{cor:s04:G2}(iii)), and the curvature-sector axioms of \cref{thm:s13:main} hold for it, with the lattice laws of \cref{thm:s09:S2K}(a) converging along the full sequence.
\end{enumerate}
With scalar sources, the corresponding statements on King's convention are supplied by \cref{thm:s13:TWK,prop:s13:WK}.
\end{corollary}

\begin{proof}
(a) The proof of \cref{lem:s12:LIM0}(a) uses only the $k$- and $J$-uniform local bound of \cref{lem:s05:C2} (Step~2), the second-order bound $\langle F(u)^2\rangle\le2\norm u^2_{2,k}$ of \cref{lem:s05:C1} (Step~3), and the current representation and Jensen's inequality (Step~4); that of \cref{lem:s12:LIM1} uses in addition the first inequality of \cref{lem:s02:cut}. All of these hold on King's convention: the bounds by \cref{thm:s05:CK} (the second-order bound is photon-only and transfers by the covariance identity of \cref{lem:s04:G}), the cut inequality by \cref{cor:s04:G1}. So the two lemmas hold for the K-family by the same proofs; the argument makes no use of \cref{hyp:s03:EC}. (b) By \cref{hyp:s03:EC}, \cref{cor:s04:G2}(ii)($\alpha$) gives, for every $J\ge J_0$ and family, locally uniform convergence of $a^K_{k,J}$ along the full sequence, so every increasing sequence is a continuum subsequence of the K-family, and by ($\beta$) every K-family LIM state is $\ZKinf$; \cref{cor:s04:G2}(iii) is the statement for the comparison family; the axioms are the curvature sector of \cref{thm:s13:main}, and the full-sequence lattice laws are \cref{thm:s09:S2K}(a).
\end{proof}

\begin{remark}[on the construction of the subsequences]\label{rem:s12:K9}
Every subsequence of spacings used in this paper is constructed in Step~2 of the proof of \cref{lem:s12:LIM0} from the $k$-uniform bounds of \cref{lem:s05:C2} and \eqref{eq:s11:W} and Montel's theorem, or is a further extraction from one so constructed, and the identification of its limits uses only the monotone cube limit of \cref{lem:s12:LIM1}; no subsequence statement is cited from \cite{K2}.
\end{remark}


\part*{Rotation invariance}
\section{The rotation part}
This part proves the rotation theorem (\cref{hyp:s03:RC}): King's cubes are compared with periodic tori carrying a smooth mass wall, the walled tori on two lattices whose orientations differ by $\arctan(3/4)$ are compared, and King's state is identified as the lower envelope of the walled tori through a heavy-shell decoupling estimate on King's block-spin flow; the limit state of the comparison family is then $O(3)$-invariant for all complex coefficients, hence invariant under every Euclidean motion.

\moving{rotation part: seven sections in preparation}

\part*{The main theorem}
% s13 The main theorem for King's printed state
\section{The main theorem for King's printed state}

We state the main theorem in full and prove it by assembling Sections~4--12. The object is King's printed state $\ZKinf(f,0)=\lim_J\lim_k\avK{e^{F(f)}}$ on King's RG cube (\cref{def:s02:Kcube}), block factor $5$, at a fixed photon mass $\lmu>0$ and all positive couplings $e,\lambda,m$ (\cref{lem:s14:CW}). The curvature sector (Leg~C) is proved on King's convention $\NDK$. The Wick sector (Leg~W) is proved in its tube form, Form~(B), on the comparison convention $\NDa$ (\cref{def:s02:acube}), on which \cref{thm:s11:W,lem:s11:TWtube,thm:s08:OS2WT} are written and whose functionals $Z^a_{k,J}$ are the ones of \cref{hyp:s03:W1}; it is carried to King's convention by \cref{thm:s13:TWK,prop:s13:WK}: the two tube statements that the interlacing of \cref{lem:s04:G} does not transfer, the bound \eqref{eq:s11:W} and the fixed-cube limits on the tube for the K-family functionals $Z^{a,K}_{k,J}$, are proved there directly. The torus inputs are \cref{hyp:s03:T1,hyp:s03:T2}, and the coupling condition is the sign clause \Tthree\ (\cref{hyp:s03:T3}), proved on the explicit regions of \cref{prop:s14:T3}; $\cK=\{(\mathrm if,g):g\ge0\}$ is the cone. \Cref{thm:s01:main} is the theorem below.

\begin{theorem}[the axioms for King's printed state at fixed photon mass]\label{thm:s13:main}
Let $e>0$, $\lambda,m,\lmu>0$, $\eps_k=5^{-k}$, $\CJ=[-MJ,MJ)^3$ with King's block integer $M\ge1$, and let $\NDK_{k,J}$ be King's cube state.

\smallskip\noindent\textbf{A. The object.} For every finite family $f=(f_1,\dots,f_N)$ of real $2$-forms in $C^\infty_c(\R^3;\Lambda^2)$ the limit
\[
\ZKinf\Big(\sum_jz_jf_j,0\Big)=\lim_{J\to\infty}\lim_{k\to\infty}\avK{e^{\sum_jz_jF_{\Lk}(f_j)}}
\]
exists locally uniformly in $z\in\C^N$, is entire, depends on the family only through the total source, is real and $\ge1$ at real $z$, lies in $[0,1]$ at purely imaginary $z$, and satisfies $|\ZKinf(\sum_jz_jf_j,0)|\le e^{\norm{g(z)}_2^2}$, $g(z)=\sum_j(\Re z_j)f_j$, with no dependence on $\lmu$, $e$, $\lambda$, $m$ or $M$.

\smallskip\noindent\textbf{B. Leg~C (the curvature field).} The Schwinger functions $S^K_n(f_1,\dots,f_n)=\partial_{z_1}\cdots\partial_{z_n}\ZKinf(\sum_jz_jf_j,0)|_{z=0}$ satisfy:
\begin{enumerate}[label=\textup{(\roman*)},leftmargin=*]
\item \emph{\OSz.} $|\ZKinf(\sum_jz_jf_j,0)|\le e^{\norm{g(z)}_2^2}$ for all $z\in\C^N$, and also $\le\exp(\frac12\lmu^{-2}\norm{\sum_j(\Re z_j)\dstar f_j}_2^2)$; the same two bounds, with the lattice norms $\norm{\cdot}_{2,k}$ and $\dstar_k$, hold at every $(k,J)$ for the lattice functionals $\avK{e^{\sum_jz_jF_{\Lk}(f_j)}}$. For real test functions
\[
|S^K_n(f_1,\dots,f_n)|\le(\sqrt2\,e)^n(n!)^{1/2}\prod_{i=1}^n\norm{f_i}_2\le(e\pi/2)^n(n!)^{1/2}\prod_{i=1}^n\fournorm{f_i};
\]
each $S^K_n$ extends to a continuous symmetric multilinear form on $\cS(\R^3;\Lambda^2_\C)^n$, a tempered distribution, the bound carrying the factor $2^n$ for complex test functions; so \Ezz\ holds with $s=4$, $\sigma_n=(e\pi)^n(n!)^{1/2}$ and every $\beta>\frac12$.
\item \emph{\OSo.} The $S^K_n$ are invariant under all translations of $\R^3$ and under the cubic group $\Oh$ about every point, $F$ transforming as a $2$-form. Moreover $\ZKinf$ is invariant under the Euclidean group $\Euc$ and the $S^K_n$ are Euclidean covariant.
\item \emph{\OSt.} Reflection positivity about every coordinate plane $\{x_\nu=c\}$, $\nu\in\{1,2,3\}$, $c\in\R$: for finite families $\zeta_1,\dots,\zeta_N$ of real $2$-form sources supported in $\{x_\nu>c\}$ with arbitrary complex coefficients, $\sum_{i,j}\bar\gamma_i\gamma_j\ZKinf(\theta\zeta_i^*\sqcup\zeta_j,0)\ge0$ for all $\gamma\in\C^N$ (the matrix form for exponentials), and $S(\theta(P)P)\ge0$ for every polynomial $P$ in the smeared curvature with sources supported in $\{x_\nu>c\}$ (the polynomial form). The same holds about every plane $\{x\cdot n=c\}$, $n$ a unit vector.
\item \emph{\OSth.} $S^K_n$ is real for real test functions, symmetric, and zero for odd $n$.
\item \emph{The measure.} For every even $r\ge14$ there is exactly one Borel probability measure $\mu^{K,H}$ on $\cH^{-r}(\R^3;\Lambda^2)$ whose moments on $C^\infty_c$ are the $S^K_n$; its characteristic functional is $\ZKinf(\mathrm i\,\cdot,0)$, extended by continuity to $\cS$; its image $\mu^K_S$ on $H^{-r}_{\rm loc}(\R^3;\Lambda^2)$ is the unique Borel probability measure there with this characteristic functional, and the unique one with these moments. For every real $\chi\in C^\infty_c(\R^3)$ and every box $Q$ containing $\supp\chi$ in its interior, the laws of the localised lattice curvature fields $\chi\Phi^K_{k,J}$ under $\NDK_{k,J}$ on $H^{-r}_{\rm per}(Q;\Lambda^2)$ converge weakly, along the full sequence $k\to\infty$ at every fixed $J\ge J_0(\chi)$ and then as $J\to\infty$, to the law of $\chi\Phi$ under $\mu^K_S$. Nothing is asserted about the laws of the unlocalised lattice fields at fixed $J$.
\item \emph{Weak \OSf.} $\mu^K_S$ is $\Cub$-invariant and is an integral $\int\mu_\xi\,d\nu(\xi)$ of translation-ergodic, reflection-positive, $\theta_j$-invariant components which cluster in the mean and under diagonal separation, have a unique joint vacuum in each coordinate direction, and satisfy the factorial growth bound $|S^\xi_N(u_1,\dots,u_N)|\le\sigma^\xi_N\prod_i\norm{u_i}_{r+2}$, $\sigma^\xi_N\le\alpha_\beta(\xi)(N!)^\beta$ for every $\beta>1$, with component constants (no tube bound); there is an $\Oh$-invariant decomposition into mixtures of $k\in\{1,2,3,6\}$ ergodic measures; and $\nu$-almost every component satisfies \OSf\ in the weak form (ergodicity under $\R^3$). Moreover $\nu$-almost every component is $\Euc$-invariant, reflection positive about every plane, and satisfies \OSf\ in each of its forms for every time direction, with $\ker H_n=\C\Omega_n$ in every direction $n$; its moments are real and symmetric. The vanishing of a component's odd moments is not claimed.
\end{enumerate}

\smallskip\noindent\textbf{C. Leg~W, Form~(B), on the comparison convention.} Let the couplings satisfy the sign clause \Tthree\ (\cref{hyp:s03:T3}; in particular, let them lie in one of the explicit regions of \cref{prop:s14:T3}), and fix an admissible scheme $a$ (\cref{def:s03:adm}; $z=a^\sharp$ the reference scheme). Let $\sigma_+=s_1$, $H=s_1/16$, $\kappa=C_U/s_1$, $K_-$ be the constants of \eqref{eq:s11:constants} ($K_-$ built from $B_{\rm cell}$; for $a\ne z$ with $C_U$ replaced by $C_U+s_1\sup_k|z_k-a_k|$), $\sigma_0=H/4$ and $K'=\kappa\pi^2+K_-$. Then for every finite family $\cF=(f_1,\dots,f_n;h_1,\dots,h_p)$ of real $2$-forms and real scalars the tube object
\[
\Ginf(\zeta)=\lim_{J\to\infty}\lim_{k\to\infty}Z^a_{k,J}\Big(\sum_iz_if_i,\sum_jw_jh_j\Big),\qquad\zeta=(z,w)\in\Dtube=\{\norm{k_+(\zeta)}_\infty<\sigma_+/2,\ \norm{k_-(\zeta)}_\infty<H/2\},
\]
exists locally uniformly on the tube and is holomorphic there; on the Wick tube of \cref{def:s08:enlarged} the limit $Z^a_{\infty,T}$ of \cref{prop:s08:T} exists and equals $\Ginf$ where both are defined. The joint Schwinger functions $S_{n,p}(f_1,\dots,f_n;h_1,\dots,h_p)=\partial^{(n)}_z\partial^{(p)}_w\Ginf(0)$ of the pair $(F,W^a)$ satisfy:
\begin{enumerate}[label=\textup{(\roman*)},leftmargin=*]
\item \emph{\OSz.} $|\Ginf(\zeta)|\le\exp(\norm{g(\zeta)}_2^2+\kappa\norm{k_+(\zeta)}_{L^1}+K_-\wnorm{k_-(\zeta)})$ on $\Dtube$, and the same bound with lattice norms holds for $Z^a_{k,J}$ uniformly in $k\ge0$ and $J\ge J_1(R)$ when the supports lie in $Q_R$; for real test functions
\[
|S_{n,p}(f_1,\dots,f_n;h_1,\dots,h_p)|\le e^{K'\sigma_0}(\sqrt2\,e)^n(n!)^{1/2}(e/\sigma_0)^p\,p!\prod_{i=1}^n\norm{f_i}_2\prod_{j=1}^p\wnorm{h_j},
\]
with the factor $2^{n+p}$ for complex test functions; the $S_{n,p}$ are tempered distributions and \Ezz\ holds with $s=4$ and every $\beta>1$. The constants $\sigma_+,H,\kappa,K_-$ depend on $(e,\lambda,m,\lmu,M,a)$, on the fixed auxiliary choices $M_1$, $\chi_0$ and the amplitude $s_1\in(0,\sigma_1)$ of \Itwo, and on nothing else; they are not uniform in $\lmu$.
\item \emph{\OSo.} $\Ginf(\tau_bR\cdot\zeta)=\Ginf(\zeta)$ on $\Dtube$ for every $b\in\R^3$ and $R\in\Oh$, $F$ transforming as a $2$-form and $W^a$ as a scalar (the transformed family has the same tube); hence the $S_{n,p}$ are invariant under translations and under $\Oh$ about every point. No Euclidean clause is asserted.
\item \emph{\OSt.} For finite families $\zeta_1,\dots,\zeta_N$ of real sources supported in $\{x_\nu>c\}$ with complex coefficients of amplitude $\amp(\zeta_i)<\sigma_W\defeq\sigma_+/8$ (in particular all purely imaginary Wick coefficients with arbitrary complex curvature coefficients), $\sum_{i,j}\bar\gamma_i\gamma_jZ^a_{\infty,T}(\theta\zeta_i^*\sqcup\zeta_j)\ge0$ about every coordinate plane; and the polynomial form $S(\theta(P)P)\ge0$ for polynomials in the smeared $F$ and $W^a$ with sources in $\{x_\nu>c\}$.
\item \emph{\OSth.} The $S_{n,p}$ are symmetric multilinear in the $f$'s and in the $h$'s, real for real test functions, and zero for odd $n$.
\item[\textup{(v),(vi)}] \emph{The measure and weak \OSf.} Clauses (v) and (vi) of Leg~C hold for the joint field on $\cH^{-r}(\R^3;\Lambda^2\oplus\R)$ and $H^{-r}_{\rm loc}(\R^3;\Lambda^2\oplus\R)$, every even $r\ge14$, with moments $S_{n,p}$, characteristic functional $(u,v)\mapsto\Ginf(\mathrm iu,\mathrm iv)$ on real test functions (purely imaginary coefficients have $k(\zeta)=0$ and lie in the tube), the component growth bound with the constants of (i), and the local convergence of the laws of the scheme-shifted lattice fields $\chi(F_{\Lk},W^a_{\Lk})$ along the full sequence $k\to\infty$ at every fixed $J$ with $MJ>R(\chi)+\ell_c$, then $J\to\infty$. No Euclidean clause is asserted. Two admissible schemes $a,a'$ give objects differing by the shift $W\mapsto W-\kappa_{a'a}\int(\cdot)$ of the Wick component, $\kappa_{a'a}=\lim_k(a'_k-a_k)$; joint cumulants of order $\ge2$ are scheme-free.
\end{enumerate}
No bound at positive scalar sources beyond the cone $\sup q\le s_1$ of \cref{cor:s11:W2prime} is assumed. The arbitrary-amplitude form (Form~(A)), which carries the only Euclidean clause of the Wick sector, is the optional extension of \cref{thm:s23:FormA}.

\smallskip\noindent\textbf{D. Leg~W on King's convention.} Under the conditions of~C, clauses (i)--(vi) of~C hold, with the same constants, for the tube object $G^K_\infty=\lim_J\lim_kZ^{a,K}_{k,J}$ of King's printed state on $\NDK$ (\cref{thm:s13:TWK,prop:s13:WK}).
\end{theorem}

\begin{proof}
\emph{A.} \Cref{thm:s05:CK}(b): by \cref{hyp:s03:EC} the two limits exist locally uniformly on $\C^N$ and are entire, $|\ZKJ|,|\ZKinf|\le e^{\norm{g(z)}_2^2}$, $J\mapsto\ZKJ(\mathrm it)$ is nondecreasing with values in $[0,1]$, and $\ZKinf(x)\ge1$ at real $x$; realness at real $z$ is the symmetry $(A,\vartheta)\mapsto(-A,-\vartheta)$ of every $\NDK_{k,J}$, closed under limits. Each lattice functional depends on the family only through its total source, hence so do the limits. The bound is absolute (\cref{rem:s05:absolute}).

\emph{B.} (i) is \cref{thm:s05:CK}: the finite-volume bound (a) with $\min\{\norm{g}^2_{2,k},\frac12\lmu^{-2}\norm{\dstar_kg}^2_{2,k}\}$ at every $(k,J)$; the two limit bounds (b), the second being (b)'s second bound, carried from the lattice by $\norm{\dstar_kg}_{2,k}\to\norm{\dstar g}_2$, \eqref{eq:s02:riemann}; the Schwinger bounds, the extension to $\cS$, the factor $2^n$ and \Ezz\ (c). (ii): translations are \cref{cor:s06:TK} (the lattice sandwich \cref{lem:s06:TK}, the two limits, the identity theorem, density); the cubic group about every point is \cref{thm:s06:S1K}; the Euclidean clause is \cref{thm:s06:O3}, that is the rotation theorem \cref{hyp:s03:RC} with \cref{cor:s04:G2} and \cref{cor:s06:TK}. (iii): the matrix form about every coordinate plane and the polynomial form are \cref{lem:s08:OS2K} with \cref{prop:s08:poly}; a plane $\{x\cdot n=c\}$ is the image of a coordinate plane under a Euclidean motion $\gamma$, and by (ii) $\ZKinf(\theta_\gamma\zeta_i^*\sqcup\zeta_j,0)=\ZKinf(\gamma^{-1}\cdot(\theta_\gamma\zeta_i^*\sqcup\zeta_j),0)$ is a Gram entry about the coordinate plane for the sources $\gamma^{-1}\cdot\zeta_i$, supported in the corresponding coordinate half-space. (iv) is \cref{prop:s08:OS3} for $\ZKinf(\cdot,0)$ (through \cref{lem:s08:OS2K}; charge conjugation gives the vanishing for odd $n$). (v) is \cref{thm:s09:S2K}(a), including the local convergence of the lattice laws along the full sequence, with \cref{rem:s09:S2K} for the unlocalised fields; (vi) is \cref{thm:s09:S2K}(b) and (c), whose $O(3)$ input is (ii)'s Euclidean clause; the disclaimer on odd moments is that of \cref{thm:s09:S2K}(c).

\emph{C.} Under \Tthree, \cref{prop:s10:U} gives \Itwo\ with amplitude $s_1\in(0,\sigma_1)$ and constant $C_U$ from the torus inputs \cref{hyp:s03:T1,hyp:s03:T2}; for $a\ne z$ the difference $z_k-a_k$ converges by admissibility, so $\sup_k|z_k-a_k|<\infty$ and $C_U$ is replaced as stated. \Ithree\ is \cref{prop:s10:RP}, and the correlation inequalities are \cref{lem:s02:BFS}. So \cref{thm:s11:W} applies: the bound \eqref{eq:s11:W} holds uniformly in $k\ge0$, $J\ge J_1(R)$, with the stated dependence of the constants and no uniformity in $\lmu$. \emph{(i)}: \cref{thm:s11:W}(i) gives the fixed-cube limits $G_J=\lim_kZ^a_{k,J}$ locally uniformly on $\Dtube$ (local boundedness by \eqref{eq:s11:W}, convergence at the real points of the tube by \cref{hyp:s03:W1}, Vitali on the connected tube), and \cref{lem:s11:TWtube}(iii) the limit $\Ginf=\lim_JG_J$, locally uniformly on $\Dtube$, holomorphic, with the bound in the limit; \cref{thm:s11:W}(ii) gives the Schwinger bounds, the factor $2^{n+p}$ and \Ezz. On the Wick tube $T_{\cF'}$ of \cref{def:s08:enlarged}, \cref{prop:s08:T}(a) gives $Z^a_{\infty,T}$, equal to $\Ginf$ on $T_{\cF'}\cap\Dtube$ as the same $J$-limit of the same $G_J$. \emph{(ii)}: translations are \cref{lem:s11:TWtube}(v). For the cubic group, \cref{prop:s12:ID}(c) shows that for an admissible $a$ under \Tthree\ the hypotheses \Bcell\ (\cref{hyp:s03:W1} at the one cell and source of \cref{hyp:s12:Bcell}) and \SCH\ (\cref{hyp:s12:SCH}) hold, every increasing sequence of spacings is a continuum subsequence for the Wick sector (\cref{def:s12:contsub}), and the LIM state $Z^{a,\cK}_\infty$ of every continuum subsequence $\cK$ is $\Ginf$ on $\Dtube$; the Wick sector of \cref{thm:s12:LIM} under \Tthree, \Bcell, \SCH\ gives $Z^{a,\cK}_\infty(\tau_bR\cdot\zeta)=Z^{a,\cK}_\infty(\zeta)$ on the tube with $W^a$ a scalar. Hence $\Ginf$ is $\Cub$-invariant on $\Dtube$, and the $S_{n,p}$ are by differentiation at $0$. \emph{(iii)}: \cref{thm:s08:OS2WT} under \cref{hyp:s03:W1}, \Itwo, \Ithree\ and \cref{lem:s02:BFS}, whose Gram entries exist by \cref{prop:s08:T}(a) and whose passage to the planes $c\notin D$ is \cref{prop:s08:T}(b); the polynomial form is \cref{prop:s08:poly}. \emph{(iv)} is \cref{prop:s08:OS3} for $Z^a_{\infty,T}$. \emph{(v),(vi)}: the Wick sector of \cref{thm:s12:LIM}, clauses (v) and (vi), for the LIM state of the full sequence, which is $\Ginf$ by \cref{prop:s12:ID}(c); the lattice-law clause then holds along the full sequence (\cref{cor:s12:K}). The Gram points of (v)'s characteristic functional have purely imaginary coefficients and lie in the tube, and the reflection positivity consumed there is (iii) at amplitude $0$ (\cref{rem:s13:amplitudes}). For the scheme clause, the shift constant of \cref{thm:s12:LIM} is $\lim_k(a'_k-a_k)$, which exists by admissibility, so no further subsequence is needed and it is independent of $\cK$ (\cref{cor:s12:K}).

\emph{D.} \Cref{thm:s13:TWK,prop:s13:WK}.
\end{proof}

\begin{remark}[the Euclidean clause]\label{rem:s13:rotation}
The Euclidean clauses of B(ii), (iii), (vi) rest on the rotation theorem \cref{hyp:s03:RC}, proved in the rotation part \moving{rotation part}, and on the identification of its limit state with $\ZKinf$ through \cref{cor:s04:G2}; the Wick sector carries no Euclidean clause (\cref{thm:s23:FormA}).
\end{remark}

\begin{theorem}[the Wick tube on King's convention]\label{thm:s13:TWK}
Let the couplings satisfy \Tthree\ and let $a$ be an admissible scheme. Let $\sigma_+=s_1$, $H$, $\kappa$ be the constants of \eqref{eq:s11:constants} and let $K_-$ be the constant there, computed from the cell constant $B^{(a)}_{\rm cell}$ of the comparison convention (\cref{hyp:s12:Bcell}), which is finite. Then for every finite family of real sources with supports in $Q_R$, every $k\ge0$, every $J\ge J_1(R)$ and every $\zeta\in\Dtube$,
\begin{equation}\label{eq:s13:WK}
|Z^{a,K}_{k,J}(\zeta)|\ \le\ Z^{a,K}_{k,J}(\Re\zeta)\ \le\ \exp\big(\norm{g(\zeta)}_{2,k}^2+\kappa\norm{k_+(\zeta)}_{1,k}+K_-\wnorm{k_-(\zeta)}\big),
\end{equation}
the bound \eqref{eq:s11:W} for the K-family functionals with the constants of the comparison convention; and the fixed-cube limits $G^K_J(\zeta)=\lim_{k\to\infty}Z^{a,K}_{k,J}(\zeta)$ exist for every $\zeta\in\Dtube$, locally uniformly, so that \cref{thm:s11:W}(i),(ii) hold on $\NDK$: $G^K_J$ and $G^K_\infty=\lim_JG^K_J$ exist locally uniformly on $\Dtube$ and are holomorphic, the bound \eqref{eq:s13:WK} passes to the limits, and the Schwinger bounds and \Ezz\ of \cref{thm:s13:main}C(i) hold for $G^K_\infty$.
\end{theorem}

\begin{proof}[Sketch]
The proof of \cref{thm:s11:W} is run on $\NDK$; every step of it that sees the states only as members of the class of \cref{lem:s02:BFS} is unchanged, and three steps are convention-sensitive.

(1) \emph{The negative leg and its direction.} By \cref{lem:s04:E} every source-carrying cell of the wall insertion of \cref{lem:s11:W3} has all its sites alive, so \cref{lem:s11:W4,lem:s11:W5,lem:s11:W67,prop:s11:W8} hold for $\avK{\cdot}$ as printed, with $K_-$ computed from the cell constant of the comparison convention (\cref{rem:s11:Kconvention}). That constant is finite: killing the sites of a cell increases a class-$\cD$ functional (\cref{lem:s04:E}(e)), so $B^{(a)}_{\rm cell}\le B^K_{\rm cell}$, and $B^K_{\rm cell}<\infty$ by \cref{hyp:s03:W1} at the one cell and the one negative source of \cref{def:s11:cells}.

(2) \emph{The positive leg.} The interlacing inequality $\avK{\cdot}\le\ava{\cdot}$ on class $\cI$ (Step~1 of \cref{lem:s04:G}, \cref{rem:s04:F}), prepended to \cref{prop:s11:W1prime}, gives $\ln N^K_{k,J}(0,q)\le\kappa\norm q_{1,k}$ for every lattice function $0\le q\le s_1$ vanishing on $\TJ$ (\cref{cor:s11:W2prime} on $\NDK$); \Itwo\ comes from \Tthree\ and the torus inputs through \cref{prop:s10:U}, \Ithree\ is \cref{prop:s10:RP} and \Ifour\ is \cref{lem:s02:BFS}, none of which sees the cube convention.

(3) \emph{Assembly and limits.} The splitting of the Wick source into its positive and negative parts and the H\"older step of the proof of \cref{thm:s11:W} use only the correlation inequalities and run unchanged, giving \eqref{eq:s13:WK} at every $(k,J)$ with the sources' supports at distance $>1$ from $\R^3\setminus\CJ$, so that the positive part vanishes on $\TJ$. The family $(Z^{a,K}_{k,J})_k$ is then locally bounded on the connected tube and converges at the real points of the tube by \cref{hyp:s03:W1}, hence locally uniformly (Vitali), as in \cref{thm:s11:W}(i); the $J$-limit, the holomorphy and the Schwinger bounds follow as in \cref{prop:s13:WK}(a) below. Kill steps outside the scope of \cref{lem:s04:G}, for an observable that does not depend on the killed sites, are the monotonicity of \cref{lem:s02:BFS} under the replacement of a single-site measure by a point mass.
\end{proof}

\begin{proposition}[the Wick sector on King's convention]\label{prop:s13:WK}
Let the couplings satisfy \Tthree\ and let $a$ be an admissible scheme. Then, with the constants of \cref{thm:s13:main}C:
\begin{enumerate}[label=\textup{(\alph*)},leftmargin=*,nosep]
\item $G^K_J\defeq\lim_kZ^{a,K}_{k,J}$ exists locally uniformly on $\Dtube$ for every $J\ge J_1(R)$, $G^K_\infty\defeq\lim_JG^K_J$ exists locally uniformly on $\Dtube$, both are holomorphic, $|G^K_\infty(\zeta)|\le\exp(\norm{g(\zeta)}_2^2+\kappa\norm{k_+(\zeta)}_{L^1}+K_-\wnorm{k_-(\zeta)})$, and the Schwinger bounds and \Ezz\ of \cref{thm:s13:main}C(i) follow.
\item $G^K_\infty(\tau_bR\cdot\zeta)=G^K_\infty(\zeta)$ on $\Dtube$ for all $b\in\R^3$, $R\in\Oh$, $W^a$ a scalar.
\item Reflection positivity about every coordinate plane in the matrix form at amplitude $<\sigma_W$ and in the polynomial form, and \OSth, for $G^K_\infty$ and its Schwinger functions.
\item The measure and weak \OSf\ of \cref{thm:s13:main}C(v),(vi) for the joint field under $\NDK$, with the local convergence of the laws of $\chi(F_{\Lk},W^a_{\Lk})$ under $\NDK_{k,J}$ along the full sequence $k\to\infty$ at every fixed $J\ge\max(J_0(\chi),J_1(R_\chi))$, then $J\to\infty$.
\end{enumerate}
\end{proposition}

\begin{proof}
The proofs of \cref{thm:s11:W}(i), \cref{lem:s11:TWtube}, \cref{thm:s06:S1K} with scalar sources, \cref{thm:s08:OS2WT} with \cref{prop:s08:T}, and \cref{thm:s09:S2K} in the Wick sector are run on $\NDK$ with the replacements below; every other step of theirs is convention-free, since it sees the states only as members of the class of \cref{lem:s02:BFS} or concerns the limit object alone.

(a) By \cref{thm:s13:TWK} the bound \eqref{eq:s11:W} holds for $Z^{a,K}_{k,J}$ uniformly in $k\ge0$, $J\ge J_1(R)$, and $\lim_kZ^{a,K}_{k,J}(\zeta)$ exists for $\zeta\in\Dtube$; so $(Z^{a,K}_{k,J})_k$ is locally bounded and convergent on the connected tube, and $G^K_J$ exists locally uniformly and is holomorphic (Vitali), with the bound. The $J$-limit is the argument of \cref{lem:s11:TWtube}(iii), whose one convention-sensitive input is the monotonicity of the cone functionals in the cube; on $\NDK$ that is \cref{cor:s04:G1} (for $J\ge J_1(R+\frac12)$, which exceeds the $J_0$ of that corollary). The split family, the uniqueness set, Vitali on the split tube and the extension to $\Dtube$ are as there; the Schwinger bounds are Cauchy's estimates on a polydisc in the tube, as in \cref{thm:s11:W}(ii).

(b) The steps of \cref{thm:s06:S1K} with scalar sources. Coordinate permutations are exact on $\NDK$ for $(f,h)$ by \cref{lem:s06:PK}. The reflection sandwich is \cref{lem:s04:C}(b) at the cone points of the tube: for $(f,g)\in\cK$ with $\norm g_\infty<\sigma_+/2$ the point $(\mathrm i,1)$ lies in $\Dtube$, so (a) supplies the three $k$-limits and the two $J$-limits of the two-step sandwich $C^K_{J\mp2}$. Conjugation and group closure on the cone are unchanged. Off the cone, the identity theorem on the split family's tube of \cref{lem:s11:TWtube}, where (a) gives the limits, extends the identity from the cone points to the connected tube; so \cref{lem:s04:C} is needed on $\cK$ only, and no class-$\cD$ form of it enters (the proof splits $h=h^+-h^-$ and never applies the sandwich to a negative scalar source). Translations are \cref{lem:s06:TK}, stated on $\cK$ and applied at the cone points of the tube with (a) for the limits, then the continuity route of \cref{cor:s06:TK} with the Wick-sector second-moment bound, which Cauchy's estimates give from \eqref{eq:s11:W} on $\NDK$; the translated family needs $J\ge\max(J_0+1,J_1(R+M+\frac12))$. Invariance about every point and differentiation concern the limit object.

(c) \Cref{thm:s08:OS2WT} and \cref{prop:s08:T} on $\NDK$: in the sandwich of the $J$-limit the upper box is \cref{lem:s04:D} (the printed boxes $\CJ$, $C_{J+2n}$) and the lower end carries the kill step of \cref{rem:s04:F}; the positive-leg bound \cref{cor:s11:W2prime} holds on $\NDK$ by \cref{rem:s04:F} prepended to \cref{prop:s11:W1prime} (\cref{rem:s11:Kconvention}); the cube monotonicity in both classes is \cref{cor:s04:G1}; the limits are (a) and the translations to the planes $c\notin D$ are (b). \OSth\ and the polynomial form use only the existence and holomorphy of the limit near $0$ and charge conjugation (\cref{prop:s08:OS3,prop:s08:poly}).

(d) The steps of \cref{thm:s09:S2K} on $\cV=\Lambda^2\oplus\R$ with the scheme-shifted field: \eqref{eq:s11:W} on $\NDK$ in the moment bounds, (a) for the limit functionals and the lattice laws, (b) and (c) for the invariance and the reflection positivity consumed by the abstract chain of \cite{SymPhase}. The second moments of the lattice field under $\NDK_{k,J}$ come from \eqref{eq:s11:W} on $\NDK$ by Cauchy's estimates, so tightness and the identification of the limit law run as in the curvature sector, with the fixed-cube limits of (a) along the full sequence, for $J\ge\max(J_0(\chi),J_1(R_\chi))$; the remaining steps concern the limit measure.
\end{proof}

\begin{remark}[the three amplitude constants]\label{rem:s13:amplitudes}
Three amplitudes occur in the Wick sector, all built from $s_1$: the matrix form of \OSt\ is proved at amplitude $\amp(\zeta_i)<\sigma_W=\sigma_+/8$ on the Wick tube (\cref{thm:s08:OS2WT}); the tube $\Dtube$ has half-widths $\sigma_+/2$ for the positive and $H/2=\sigma_+/32$ for the negative part of the real Wick source, and the form of \OSt\ proved for LIM states (\cref{thm:s12:LIM}) takes $\amp(\zeta_i)<H/2<\sigma_W$, so that every Gram point lies in $\Dtube$, where those limits exist; and $\sigma_0=H/4$ is the radius of the polydisc in Cauchy's estimates for \Ezz. The consumer of \OSt\ in clauses (v),(vi), the reflection positivity of exponentials $e^{\mathrm i\Phi(u)}$ with real $u$ in the clustering theorems of \cite{SymPhase}, needs purely imaginary Wick coefficients only, that is amplitude $0$ (\cref{rem:s08:amplitudes}).
\end{remark}

\begin{remark}[the photon mass and the couplings]\label{rem:s13:mu0}
The photon mass $\lmu>0$ is fixed throughout \cref{thm:s13:main}. No relation $\lmu\le m$ is used or needed: the one input whose printed proof in the torus paper \cite{SymPhase} carries the restriction $\lmu\in(0,\mus]$, the lattice reflection positivity of King's tori, holds for all $\lmu>0$, since that proof uses no bound on $\lmu$ and the paper's own lemma on the finite-volume photon covariance lifts the scope (\cref{prop:s10:RP}, \cref{lem:s14:CW}(d)). Reflection positivity of the cube limit needs the margin $R(\lmu)=\pi\sqrt3/(2\lmu)$ of \eqref{eq:s07:margin} inside the upper half of the sandwich (\cref{lem:s07:U}) and nowhere in the conclusion. The constants of Leg~C are absolute; those of Leg~W depend on $\lmu$ and are not uniform in it (a tube bound uniform in the photon mass is the first step of the separate work on the limit of vanishing photon mass). The couplings range over all of $e>0$, $\lambda,m,\lmu>0$: their quantifier is the sources' (\cref{lem:s14:CW}(e), \cref{rem:s14:quantifier}), the stopping scale of the lattice constructions behind \cref{hyp:s03:EC,hyp:s03:W1,hyp:s03:T1} being chosen after the point and the datum, with the mass cap of \cite[Prop.~3.1]{Ba2} one more one-sided cap on that scale (\cref{rem:s14:BALm}); so the theorem contains no sentence ``$e,\lambda$ sufficiently small'', and the only coupling condition of the Wick sector is the one-sided sign clause \Tthree\ (\cref{lem:s14:CW}(b)), proved on the explicit regions of \cref{prop:s14:T3}.
\end{remark}

% s14 The coupling range
\section{The coupling range}

\providecommand{\sxivNE}{\ensuremath{(\mathrm{NE})}}
\providecommand{\sxivNEl}{\ensuremath{(\mathrm{NE}^{<})}}
\providecommand{\sxivCT}{\ensuremath{(\mathrm{CT}_{\infty})}}

This section collects what the sources behind the inputs of \cref{thm:s13:main} say about the couplings and shows that the only coupling condition anywhere in this paper is the one-sided sign clause \Tthree\ of the Wick leg, which \cref{prop:s14:T3} proves on explicit regions.

\medskip\noindent\textbf{King's range, units, two letters $J$.} Let $e>0$ and $\lambda,m,\lmu>0$. Every constant of the cited constructions, including the number $J_{\rm K}$ of renormalisation steps after which the running couplings are small and the stopping scale $5^{-J_{\rm K}}$ (a member of the family $\{5^{-k}\}$), may depend on $(e,\lambda,m,\lmu)$, on the large-block integer $M$ and on the cube or torus; no numerical restriction is placed on the couplings themselves. That a stopping scale exists at every positive point is part of the cited constructions (\cite[Thm.~3.1]{K1}, \cite{Ba1,Ba2,Ba3,Ba4}, and the existence and rotation parts of this paper), under the mass cap of \cite[Prop.~3.1]{Ba2}, a one-sided cap on the stopping scale (\cref{rem:s14:BALm}). This is what ``King's range'' means here: all positive $(e,\lambda,m,\lmu)$, with a stopping scale chosen after them. The cube index $J$ of $\CJ$ is ours; King's stopping index, the ``fixed integer $J$'' of \cite[Thm.~3.1]{K1}, is written $J_{\rm K}$. Couplings are measured in the units in which the coarsest spacing is $1$; in $d=3$ all four of $e^2,\lambda,m,\lmu$ have the dimension of a mass, rescaling to the unit lattice at spacing $\eps$ replacing $m^2,\lambda,e$ by $m^2\eps^2,\lambda\eps,e\eps^{1/2}$ \cite[Sect.~0]{K1} and $\lmu^2$ like $m^2$.

\medskip\noindent\textbf{The relative bare coefficient.} Let $\alpha^\infty_k(M)=M^2+\dm_k(M)$, $s^\infty_k(M_1)=\frac12(\alpha^\infty_k(m)-\alpha^\infty_k(M_1))$ and $\sigma_1(M_1)=\liminf_ks^\infty_k(M_1)$ be as in \cref{hyp:s03:T2}; the theorem is stated through $\sigma_1$ alone. Describing \emph{where} \Tthree\ fails uses one further statement about the continuum limit of this coefficient, which is of the same kind as \cref{hyp:s03:T2} (\cref{rem:s03:T2}) and which we name rather than use tacitly:

\smallskip\noindent\sxivCT\ \emph{For every $M>0$ the limit
\[
\alpha_\infty(M)-\alpha_\infty(m)\defeq\lim_{k\to\infty}\big[\alpha^\infty_k(M)-\alpha^\infty_k(m)\big]
\]
exists and is a polynomial of degree $\le2$ in $(e^2,\lambda)$ with finite coefficients, differentiable in $M^2$ at $M=m$.}

\smallskip\noindent The divergent parts of the counterterm ($1/\eps$ at one loop, $\ln(1/\eps)$ at two loops) do not depend on the scalar mass and cancel in the difference. \sxivCT\ is not a hypothesis of \cref{thm:s13:main}; it enters only the description of the failure set. Under \sxivCT\ put
\[
s_\infty(M_1)\defeq\tfrac12\big(\alpha_\infty(m)-\alpha_\infty(M_1)\big),\qquad s_*\defeq\sup_{M>0}s_\infty(M)\in[0,\infty],
\]
\[
E_W\defeq\{s_*=0\}=\{m\text{ is a global minimiser of }M\mapsto\alpha_\infty(M)\text{ on }(0,\infty)\},\qquad
E^<_W\defeq\{m\text{ minimises }\alpha_\infty\text{ on }(0,m]\}\supseteq E_W,
\]
subsets of the coupling space ($s_*\ge0$ because $M=m$ gives $0$; the value $+\infty$ is harmless below). The two sign clauses are
\begin{itemize}[nosep,leftmargin=*]
\item[\sxivNE] $m$ is not a global minimiser of $\alpha_\infty$ on $(0,\infty)$, i.e.\ $s_*>0$, i.e.\ $(e,\lambda,m,\lmu)\notin E_W$;
\item[\sxivNEl] some $M_1\in(0,m)$ has $\alpha_\infty(M_1)<\alpha_\infty(m)$, i.e.\ $(e,\lambda,m,\lmu)\notin E^<_W$ (the one-sided form; it implies \sxivNE).
\end{itemize}
Throughout, $'$ denotes $d/dM^2$, $N$ is the number of real components of the scalar field, and ``one-loop order'' means the orders $e^2$ and $\lambda$.

\begin{lemma}[the coupling range]\label{lem:s14:CW}
Let $e>0$ and $\lambda,m,\lmu>0$.
\begin{enumerate}[label=\textup{(\alph*)},leftmargin=*]
\item \emph{The curvature sector.} The existence theorem of the curvature sector, \cref{hyp:s03:EC} (the existence clause of \cite[Thm.~2.1]{K2} at zero scalar source), carries no condition on $e,\lambda,m,\lmu$. The correlation inequalities it uses, \cref{lem:s02:BFS} (Theorem~6.2(1),(2) of \cite{BFS1}), have the standing hypotheses of a finite single-site measure, positive hopping coefficients and a positive photon covariance, which for King's action are $\lambda>0$ and $\lmu>0$; the margin $R(\lmu)=\pi\sqrt3/(2\lmu)$ of \cref{lem:s07:U} needs $\lmu>0$ only; and no lemma proved in this paper for the curvature sector carries a condition on the couplings beyond those of its inputs. Hence the curvature sector of \cref{thm:s13:main} holds on all of King's range, with no exceptional set. The one proof-internal proviso in the Ba{\l}aban antecedents of \cref{hyp:s03:EC}, the mass clause of \cite[Prop.~3.1]{Ba2}, is a cap on the dimensionless scalar mass at the stopping scales and no condition on the bare couplings (\cref{rem:s14:BALm}). Of the bounds of the paper only the absolute bounds of \cref{thm:s05:CK} are uniform in $\lmu$.
\item \emph{The Wick sector, Form~(B).} The inputs of Form~(B) beyond \cref{hyp:s03:W1,hyp:s03:T1,hyp:s03:T2} and \Tthree\ --- \cref{prop:s10:RP,lem:s10:CB,lem:s02:BFS,lem:s10:F,lem:s10:V} and the tube statements \cref{thm:s11:W,lem:s11:TWtube,thm:s08:OS2WT}, whose conditions on the Wick amplitudes are not conditions on the couplings --- print no condition on the bare couplings beyond positivity ($\lambda>0$, $\lmu>0$); \cref{hyp:s03:T1} is used at the two masses $m$ and $M_1$. The only coupling condition of Form~(B) is therefore \Tthree, which \cref{prop:s14:T3} proves on explicit regions. Under \sxivCT:
\begin{enumerate}[label=\textup{(\roman*)},nosep,leftmargin=*]
\item $\sigma_1(M_1)=\lim_ks^\infty_k(M_1)=s_\infty(M_1)$ for every $M_1>0$. Hence \Tthree\ holds if and only if \sxivNEl: it fails exactly on $E^<_W$. If \cref{hyp:s03:T1,hyp:s03:T2} are available also at auxiliary masses $M_1>m$, then by \cref{prop:s10:Upm} the two-sided clause \sxivNE\ replaces \Tthree, with failure set exactly $E_W$; the branch $M_1>m$ is needed exactly on $E^<_W\setminus E_W$, for instance when $\alpha'_\infty(m)<0$ and $m$ minimises $\alpha_\infty$ on $(0,m]$.
\item To one-loop order, for every $M_1>0$,
\[
s_\infty(M_1)=\tfrac12(m^2-M_1^2)+\frac{\lambda(N+2)(m-M_1)}{2\pi}-e^2\,\frac{(m-M_1)\,(m\lmu+mM_1+M_1\lmu)}{8\pi(\lmu+M_1)(m+\lmu)},
\]
so that
\begin{equation}\label{eq:s14:alphaprime}
\alpha'_\infty(m)=1+\frac{\lambda(N+2)}{2\pi m}-e^2\,\frac{m+2\lmu}{8\pi(m+\lmu)^2},\qquad \frac{m+2\lmu}{8\pi(m+\lmu)^2}\le\frac1{4\pi(m+\lmu)},
\end{equation}
and to one-loop order \Tthree\ holds if and only if
\begin{equation}\label{eq:s14:ec2}
e^2<e_c^2(\lambda,m,\lmu)\defeq\max\Big\{\frac{8\pi(m+\lmu)^2}{m+2\lmu}\Big(1+\frac{\lambda(N+2)}{2\pi m}\Big),\ 4\pi(m+\lmu)\Big(1+\frac{\lambda(N+2)}{\pi m}\Big)\Big\},
\end{equation}
the first entry being the larger if and only if $\lambda(N+2)\lmu\le\pi m^2$. So the one-loop failure set of \Tthree\ is the half-line $e^2\ge e_c^2$. The first entry is the zero $e_*^2\ge4\pi(m+\lmu)$ of \eqref{eq:s14:alphaprime}; at $m=\lmu=1$, $\lambda=0$ it equals $32\pi/3\approx33.5$ (for orientation only). Nothing finer than one-loop order is claimed about $E^<_W$.
\item For the two-sided set, $E_W\subset\{\alpha'_\infty(m)=0\}$; for each $(\lambda,m,\lmu)$ at which the polynomial $e^2\mapsto\alpha'_\infty(m)$, of degree $\le2$, is not identically zero, this is at most two values of $e^2$. At $\lmu=m$ the window $0<\lambda,e^2<0.1m$ of the torus paper \cite{SymPhase} lies outside $E^<_W$, by (d).
\end{enumerate}
\item \emph{The existence and rotation parts.} The existence theorems \cref{hyp:s03:EC,hyp:s03:W1} and the rotation theorem \cref{hyp:s03:RC} are proved for $e>0$, $\lambda,m,\lmu>0$ in King's range, and none carries a condition on the bare couplings beyond positivity: every smallness condition of \cite{Ba1,Ba2,Ba3,Ba4} that they use (listed in the proof) is a one-sided cap on a stopping scale $\beta\le1$, imposed last in the schedule of constants, and an admissible $\beta$ exists at every positive point. \RW\ (optional extension) carries no coupling threshold: the amplitude threshold $s_*$ restricts the source, not the couplings, and vanishes exactly on $E_W$.
\item \emph{Comparison with the torus paper.} The window $0<\lambda<0.1m$, $0<e^2<0.1m$, $\lmu\in(0,m]$ of \cite{SymPhase} is a hypothesis of its torus theorem only and is not inherited by King's state. Nor is $\lmu\le m$: the one King-side input printed under it, the torus paper's reflection-positivity proposition behind \cref{prop:s10:RP} (stated for $\lmu\in(0,\mus]$, $\mus=m$), is lifted to all $\lmu>0$ by a proof that uses only site-independent coefficients, $\lambda>0$ and even sides. At $\lmu=m$ the torus paper's certificate gives $s_\infty(0.7m)\ge0.1595\,m^2$ on the whole window (identifying the two torus counterterm schemes at $\lmu=m$, which compute the same graphs), so \Tthree\ holds there with $M_1=0.7m$ and an explicit margin; \cref{prop:s14:T3} sharpens this and gives explicit regions at every $\lmu\in(0,m]$ (\cref{rem:s14:torus-window}).
\item \emph{The quantifier.} For all $e>0$, $\lambda,m,\lmu>0$ and each datum there are an integer $J_{\rm K}$ (\cite[Thm.~3.1]{K1}) and a stopping scale $\beta\le1$ such that the running couplings $(e\beta^{1/2},\lambda\beta,m\beta,\lmu\beta)$ satisfy finitely many one-sided caps on $\beta$ and on three of the four of them ($\lmu\beta$ is uncapped in the sources); no condition is placed on the bare couplings, and the theorems consumed from \cite{K1,K2} carry no coupling hypothesis. A sentence ``$e,\lambda$ sufficiently small, $m,\lmu>0$ fixed'', read as a hypothesis on the bare couplings, is strictly stronger and is printed in none of the sources (\cref{rem:s14:quantifier}).
\end{enumerate}
\end{lemma}

\begin{proof}
\emph{(a).} Theorem~2.1 of \cite{K2} asserts, for $g,h$ supported in $C_J$ and all $z,w\in\C$, the existence of $\lim_K\langle\Theta(zg,wh)\rangle^{\eps_K}_{N,D,C_J}$; no hypothesis on the couplings precedes it. Its proof rests on \cite[Thm.~3.1]{K1} (the convergence of the effective action, referred there to \cite{Ba1,Ba2,Ba3,Ba4}), stated ``for some fixed integer $J$'', our $J_{\rm K}$; neither it nor \cite[Thms.~2.1, 3.4]{K1} carries a coupling hypothesis, the smallness in \cite{K1} being of the fields and of the stopping spacing. So the one proviso of \cite[Prop.~3.1]{Ba2} is upstream of \cref{hyp:s03:EC}. For King's action the single-site measure $e^{-\eps^3[\frac12\alpha r^2+\lambda r^4]}r\,dr$ is finite iff $\lambda>0$ (or $\lambda=0$, $\alpha>0$), the hopping weights are positive, and the Feynman-form photon covariance is positive iff $\lmu>0$; the ``$g>0$'' of Theorem~6.2 of \cite{BFS1} is the sign of the scalar source, and \cref{lem:s02:BFS} covers $g=0$. The margin of \cref{lem:s07:U} concerns the photon Gaussians of the two boxes alone, so $\lmu$ is its only parameter (no fixed margin works at $\lmu=0$). The lemmas of this paper see only the form of the action.

\emph{The torus inputs.} \Cref{hyp:s03:T1} is Ba{\l}aban's stability theorem \cite[(1.14)]{Ba1}, \cite[(1.3)]{Ba2}, with constants independent of $\eps$ and of the torus, under $\lambda>0$, $\lmu^2>0$, $m^2>0$. Every printed smallness in \cite{Ba1,Ba2,Ba3,Ba4} is of one of three kinds. \emph{Running charge:} the bounds hold ``for $e(L^k\eps)$ sufficiently small'', $e(L^k\eps)=e\,(L^k\eps)^{(4-d)/2}$ \cite[(2.23)]{Ba1}; \cite{Ba4} proves them for a unit-lattice model whose charge is this running charge (its theorem is stated there as Proposition~2.1 of \cite{Ba1}), so ``$e$ sufficiently small'' there is a cap on $e\beta^{1/2}$ at the stopping scale $\beta$. \emph{Scale:} the stopping scale $\eps_0$ with $L^K\eps\le\eps_0<L^{K+1}\eps$ is ``some fixed sufficiently small number'' \cite{Ba3}, and ``$L^{k+1}\eps\le1$'', ``$L^k\eps$ sufficiently small'' \cite{Ba1,Ba2} are caps on it; Ba{\l}aban's units have $\lambda=1$ (in $d=3$: spacing $\lambda\eps$, mass squared $m^2/\lambda^2$, charge $e\lambda^{-1/2}$, the running charge invariant), so his ``$\eps$ small'' is ``$\lambda\eps$ small'' in ours, and his constants may depend on $e\lambda^{-1/2}$, $m/\lambda$, $\lmu/\lambda$: a dependence, not a condition. \emph{Mass:} Proposition~3.1 of \cite{Ba2} carries the sentence ``we assume $m^2\le O(1)$ also''; this is the mass cap of that proposition (\cref{rem:s14:BALm}). \Cref{hyp:s03:T2} carries no coupling condition: the bound on the relative bare coefficient holds for every mass in a compact interval containing $m$, on either side of $m$, with a constant depending on the couplings (\cref{rem:s03:T2}; the power-counting theorem of \cite{Rei88} that enters it has none). \Cref{lem:s10:F} needs a finite single-site integral ($\lambda>0$) and a finite photon Gaussian on the torus ($\lmu>0$; at $\lmu=0$ the constant mode lies in the kernel); \cref{lem:s10:V} needs massive propagators; \cref{prop:s10:RP} needs finite single-site moments and a positive Proca covariance and does not see the counterterm. The tube statements restrict Wick amplitudes only. So the only coupling condition of Form~(B) is the sign clause of \cref{prop:s10:U}, i.e.\ \Tthree.

\emph{(b)(i).} By \cref{lem:s10:V}, $s^\infty_k(M_1)$ is the fixed-$k$ volume limit of the torus coefficient, the relative bare coefficient on $\Lk$ in the free convention; by \sxivCT\ its limit $s_\infty(M_1)$ exists, hence $\sigma_1(M_1)=\liminf_k=\lim_k=s_\infty(M_1)$. Thus some $M_1\in(0,m)$ has $\sigma_1(M_1)>0$ iff \sxivNEl, and some $M_1>0$ has $\sigma_1(M_1)>0$ iff \sxivNE. \Cref{prop:s10:Upm} gives \Itwo\ from \cref{hyp:s03:T1,hyp:s03:T2} at an auxiliary mass on either side of $m$, so with those inputs \sxivNE\ suffices; on $E^<_W\setminus E_W$ every $M_1<m$ fails and some $M_1>m$ works.

\emph{(b)(ii).} The closed form is the one-loop evaluation of King's counterterm \cite[(3.32)--(3.33)]{K1} (tadpole and zero-momentum photon-exchange graphs; the $1/\eps$ divergence cancels in the difference). It equals $\frac12\int_{M_1^2}^{m^2}\alpha'_\infty(\mu)\,d\mu^2$ with $\alpha'_\infty(\mu)=1+\frac{\lambda(N+2)}{2\pi\mu}-e^2\frac{\mu+2\lmu}{8\pi(\mu+\lmu)^2}$, an orientation-free definite integral, valid on either side of $m$; at $\mu=m$ this is \eqref{eq:s14:alphaprime}, whose inequality is $m\ge0$. For $M_1\in(0,m)$ write $s_\infty(M_1)=\frac{m-M_1}2\beta(M_1)$ with
\[
\beta(M_1)=(m+M_1)+c-e^2\,\frac{m\lmu+(m+\lmu)M_1}{4\pi(m+\lmu)(\lmu+M_1)},\qquad c\defeq\frac{\lambda(N+2)}{\pi};
\]
$s_\infty(M_1)>0$ iff $\beta(M_1)>0$ iff $e^2<e_c^2(M_1)\defeq\big[(m+M_1)+c\big]\frac{4\pi(m+\lmu)(\lmu+M_1)}{m\lmu+(m+\lmu)M_1}$. With $\mathcal N(M_1)\defeq(m+M_1+c)(\lmu+M_1)$ and $\mathcal D(M_1)\defeq m\lmu+(m+\lmu)M_1$, the derivative of $e_c^2(M_1)$ has the sign of $\mathcal N'\mathcal D-\mathcal N\mathcal D'=2m\lmu M_1+(m+\lmu)M_1^2-c\lmu^2$, positive on $(0,m)$ when $c=0$ and, when $c>0$, negative then positive with one zero. So $e_c^2(\cdot)$ is increasing, or decreasing then increasing, on $(0,m)$, and its supremum over the open interval is the larger of the two end values, not attained: $M_1\to m^-$ gives $\frac{8\pi(m+\lmu)^2}{m+2\lmu}(1+\frac c{2m})$, the zero $e_*^2$ of \eqref{eq:s14:alphaprime}, and $M_1\to0^+$ gives $4\pi(m+\lmu)(1+\frac cm)$. Hence to one-loop order \Tthree\ holds iff $e^2<e_c^2(\lambda,m,\lmu)$ as in \eqref{eq:s14:ec2}, and fails for every $e^2\ge e_c^2$ (then $\beta\le0$ for every $M_1$); the first entry is the larger iff $(m+\lmu)(2m+c)\ge(m+c)(m+2\lmu)$, i.e.\ $c\lmu\le m^2$.

\emph{(b)(iii).} Under \sxivCT, $M^2\mapsto\alpha_\infty(M)$ is differentiable at $m$, and a global minimum at the interior point $m$ is a critical point: $E_W\subset\{\alpha'_\infty(m)=0\}$. A polynomial of degree $\le2$ in $e^2$ that is not identically zero has at most two zeros. The clause at $\lmu=m$ is \cref{rem:s14:torus-window}.

\emph{(c).} The existence and rotation parts fix $e>0$, $\lambda,m,\lmu>0$ in King's range and the block factor $5$, and fix their constants after the couplings. In their schedule of constants, $\beta$ is fixed last as the minimum of finitely many caps, each one-sided and monotone in $\beta$ with a right side that is a positive finite function of the couplings and of the entries fixed before it; so an admissible $\beta$ exists at every positive point. Ba{\l}aban's own thresholds and King's $J_{\rm K}$ are imposed on $\beta$ after $M$ is fixed.

\emph{(d).} The torus paper states that its numerical window enters through its counterterm-gap input alone and that none of its other conditions is a condition on $e$ or $\lambda$. Its reflection-positivity proposition for King's tori is stated for $\lmu\in(0,\mus]$, $\mus=m$, while the two lemmas on the periodic photon covariance that carry its proof are stated for every $\mu>0$; the proof of \cref{prop:s10:RP} uses only site-independent coefficients, $\lambda>0$ and even sides, and so holds for every $\lmu>0$. The certificate is \cref{rem:s14:torus-window}.

\emph{(e).} The quantifier is read off \cite[Thm.~3.1]{K1} (``for some fixed integer $J$, for all $K\ge J$, and all $k\le K-J$'') and the stopping-scale caps listed above; the last sentence is \cref{rem:s14:quantifier}.
\end{proof}

\begin{remark}[the mass clause of Ba{\l}aban's upper bound]\label{rem:s14:BALm}
Proposition~3.1 of \cite{Ba2} asserts an inequality (3.26), summed over the scales $k=0,\dots,K$, with a constant $\gamma_0=\gamma_0(d,a)$ and mass terms $(L^k\eps)^dm^2|\phi_k|^2$; rescaled to the unit lattice, its $k$-th term is the per-scale inequality (3.29) with the mass term $(L^k\eps)^2m^2|\phi'_k|^2$, so $(L^k\eps)^2m^2$ is the dimensionless mass at scale $k$. The proposition carries the sentence ``we assume $m^2\le O(1)$ also'', $m^2$ being the bare scalar mass squared in Ba{\l}aban's units ($\lambda=1$), i.e.\ $m^2/\lambda^2$ in ours. The clause under which this paper cites the proposition concerns the scales at which it is applied:
\begin{quote}
\emph{Inequality (3.29) of \cite{Ba2} holds, with $\gamma_0=\gamma_0(d,a)$, at every scale $k\le K$ at which $(L^k\eps)^2m^2\le c_m(d,a)$; hence (3.26) holds whenever this is so for every $k\le K$.}
\end{quote}
Here $c_m(d,a)$ is a constant depending on $d$ and $a$ only. The clause belongs to the statement of the proposition: the mass that enters its printed proof, referred to \cite[Sect.~4]{Ba4}, is the unit-lattice mass at scale $k$, the text of \cite{Ba2} makes no further use of the sentence, and the proposition is cited here under the clause. It is then one more one-sided cap on the stopping scale $\beta=\eps_0$ ($(L^k\eps)^2m^2\le\beta^2m^2$ at every scale used), of the kind every other smallness condition of \cite{Ba1,Ba2,Ba3,Ba4} has; a further monotone cap on $\beta$ leaves the existence of an admissible stopping scale unchanged, and nothing is imposed on the bare mass.

Read instead as an absolute bound on $m^2/\lambda^2$, the sentence would exclude small $\lambda$ at fixed $m$, including all of the window $\lambda<0.1m$ of \cite{SymPhase} ($m^2/\lambda^2>100$), and would be an unprinted hypothesis of the theorem (1.3) of \cite{Ba2}, whose only printed condition on $m$ is $m^2>0$; that is not the reading under which the proposition is proved, and it is used nowhere. The clause sits upstream of \cref{hyp:s03:T1} at both masses, of \cref{hyp:s03:EC,hyp:s03:W1} (through \cite[Thm.~3.1]{K1}) and of the steps of the rotation part that rest on \cite{Ba1,Ba2}, as every cap on the stopping scale does; it is not a condition of this paper.
\end{remark}

\begin{remark}[the torus paper's window and its certificate at $\lmu=m$]\label{rem:s14:torus-window}
(i) \emph{The certificate.} The torus paper's coefficient $a_{K,T}(M)=M^2+\dm_{K,T}(\eps;M)$ is the counterterm of \cite[(3.32)--(3.33)]{K1} (orders $2\le\alpha+2\beta\le4$, zero-momentum normalisation) computed on the periodic torus with the periodic scalar propagator of mass $M$ and the periodic photon propagator of the reference mass $\mus$, without block averaging; its theorem on this coefficient gives, for $m>0$, $\mus=m$, $M_1=0.7m$, uniformly in $\eps\le1/m$ and in every periodic torus of side $\ge96/m$, the explicit two-loop bound $a_{K,T}(m)-a_{K,T}(0.7m)\ge0.319\,m^2$. King's torus-convention coefficient $\alpha^{\mathrm{tor}}_{k,T}(M)$ of \cref{def:s10:torus} is the same polynomial over the same orders, as translation-invariant lattice graphs on the torus with the photon of mass $\lmu$ and the same normalisation, the counterterm graphs of \cite{K1} being defined with the full lattice propagators, without block averaging, as we read that paper. So at $\lmu=m=\mus$ the two schemes compute the same graphs and $a_{K,T}(M)=\alpha^{\mathrm{tor}}_{k,T}(M)$ (an identification of ours). Then $s_{k,T}(0.7m)=\frac12(a_{K,T}(m)-a_{K,T}(0.7m))\ge0.1595\,m^2$ for all $\eps_k\le1/m$ and every $T_J$ with $2MJ\ge96/m$; \cref{lem:s10:V} ($J\to\infty$ at fixed $k$) gives $s^\infty_k(0.7m)\ge0.1595\,m^2$ for every such $k$, hence $\sigma_1(0.7m)\ge0.1595\,m^2>0$: at $\lmu=m$, \Tthree\ holds on the whole window with $M_1=0.7m$ and an explicit margin, without \sxivCT\ and without $C_4$; under \sxivCT\ the window lies outside $E^<_W$.

(ii) \emph{Away from $\lmu=m$.} The torus paper's counterterm is computed at the photon mass $\mus=m$, King's at $\lmu$; the two coefficients differ by a term of order $e^2|m-\lmu|$ (the zero-momentum photon-exchange graph at the two photon masses), so the certificate does not transfer to $\lmu\ne m$. There the comparison is through one-loop order: by \eqref{eq:s14:ec2}, $e_c^2\ge4\pi(m+\lmu)>0.1m$, so the one-loop form of \Tthree\ holds on the torus paper's window at every $\lmu>0$; beyond one loop, \cref{prop:s14:T3} places the window inside the region where \Tthree\ holds for every $\lmu\in[0.03m,m]$, and a box shrinking like $\sqrt{\lmu}$ in $e^2$ for every $\lmu\in(0,m]$; the exact position of $E^<_W$ is not determined.
\end{remark}

\begin{proposition}[explicit coupling regions on which \Tthree\ holds]\label{prop:s14:T3}
Let $N=2$ and $M_1=0.7m$, and write $\hat\lambda=\lambda/m$, $\hat e^2=e^2/m$, $\hat\mu=\lmu/m$. Let $\dm_k(M)$ be King's counterterm \cite[(3.32)--(3.33)]{K1} on $\Lk$ at the photon mass $\lmu$ (the orders $2\le\alpha+2\beta\le4$ in $e^\alpha\lambda^\beta$, zero-momentum normalisation), and let $\sigma_1(M_1)=\liminf_ks^\infty_k(M_1)$ be as in \cref{hyp:s03:T2}.
\begin{enumerate}[label=\textup{(\alph*)},leftmargin=*]
\item \emph{($\lmu=m$.)} $\sigma_1(0.7m)\ge m^2\,\sigma'(\hat\lambda,\hat e^2)$, where
\begin{equation}\label{eq:s14:sigmaprime}
\begin{split}
\sigma'(\hat\lambda,\hat e^2)\defeq{}&0.255+0.190986\,\hat\lambda-0.008426\,\hat e^2\\
&-0.144555\,\hat\lambda^2-4.665214\,\hat e^2\hat\lambda-0.851892\,\hat e^4 .
\end{split}
\end{equation}
Hence \Tthree\ holds on the region $\{\sigma'>0\}$, which contains the box $0\le\lambda,e^2\le0.1m$ (on which $\sigma'\ge0.2166$), the segment $\{\lambda=0,\ e^2<0.5422\,m\}$ and the segment $\{e=0,\ \lambda<2.1440\,m\}$.
\item \emph{(Every $\lmu\in(0,m]$; closed form.)} \Tthree\ holds whenever
\begin{equation}\label{eq:s14:sigmastar}
\begin{split}
\sigma^*(\hat\lambda,\hat e^2;\hat\mu)\defeq{}&0.255\,\big(1-0.0468103\,\hat e^2\big)+0.1909859\,\hat\lambda-0.1445549\,\hat\lambda^2\\
&-6.5937745\,\hat e^2\hat\lambda-\tfrac12\Big(4.7629+\frac{0.77120851}{\hat\mu}\Big)\hat e^4\ >\ 0 ;
\end{split}
\end{equation}
in particular, for every $\lmu\in(0,m]$, on each of the boxes $\{\lambda\le\lambda_1m,\ e^2\le c_1\sqrt{\lmu m}\}$ with $(\lambda_1,c_1)\in\{(0,0.3014),(0.05,0.2533),(0.1,0.2152),(0.2,0.1616),(0.5,0.0884),(1,0.0447)\}$.
\item \emph{(Every $\lmu\in(0,m]$; interval bounds.)} With the two-loop coefficients bounded by computer-assisted interval bounds (rounded outward) in place of the closed form of (b), \Tthree\ holds on a region $R(\lmu)$, non-decreasing in $\lmu$, which contains the box $[0,0.2m]^2$ for $\lmu\in[0.5m,m]$, $[0,0.15m]^2$ for $\lmu\in[0.1m,m]$, $[0,0.1m]^2$ for $\lmu\in[0.03m,m]$ and $[0,0.075m]^2$ for $\lmu\in[0.01m,m]$, the segment $\{e=0,\ \lambda<2.1440\,m\}$ for every $\lmu$, and at $\lmu=m$ the box $[0,0.239m]^2$ and the segment $\{\lambda=0,\ e^2<0.5618\,m\}$.
\item \emph{($e=0$; exact.)} At $e=0$, for every $M_1\in(0,m)$, $\lim_ks^\infty_k(M_1)=\frac12(m^2-M_1^2)+\frac{(N+2)\lambda(m-M_1)}{2\pi}-\frac{(N+2)\lambda^2}{\pi^2}\ln\frac m{M_1}$ exactly, and \Tthree\ holds if and only if $\lambda<\frac\pi4\big(1+\sqrt{1+\frac{16}{N+2}}\big)m$, which is $\frac\pi4(1+\sqrt5)\,m\approx2.5416\,m$ for $N=2$; at $M_1=0.7m$ the exact window is $\lambda<2.1440\,m$.
\end{enumerate}
Outside these regions the sign clause \Tthree\ is not decided here. On each of them, under \sxivCT, the two-sided clause \sxivNE\ holds as well, with no use of \cref{prop:s10:Upm}.
\end{proposition}

\begin{proof}[Sketch]
By \cref{lem:s10:V}, $s^\infty_k(M_1)$ is the volume limit at fixed $k$ of the torus coefficient, and King's counterterm has the five orders $(\alpha,\beta)\in\{(0,1),(2,0),(0,2),(2,1),(4,0)\}$ (the odd orders in $e$ vanish), so
\[
s^\infty_k(M_1)=A+\lambda S_{01,k}+e^2S_{20,k}+\lambda^2S_{02,k}+e^2\lambda S_{21,k}+e^4S_{40,k},\qquad A=\tfrac12(m^2-M_1^2),
\]
with coupling-free coefficients $S_{\alpha\beta,k}=\frac12\big[\dm_{(\alpha\beta)}(\eps_k;m)-\dm_{(\alpha\beta)}(\eps_k;M_1)\big]$. The three low orders are lattice integrals evaluated exactly: $S_{01,k}=2(N+2)(m^2-M_1^2)\int_BG_mG_{M_1}>0$, with limit $\frac{(N+2)(m-M_1)}{2\pi}$; $S_{20,k}=-A\,I_{\eps_k}$ with $0<I_\eps\le\frac{\pi^2}{32(m+M_1)}$ and $I_\eps\to I_\infty=\frac{m\lmu+mM_1+M_1\lmu}{4\pi(m+M_1)(\lmu+M_1)(m+\lmu)}\le\frac1{4\pi(m+M_1)}$, which decreases in $\lmu$; $S_{02,k}=-16(N+2)\sum_x\eps_k^3\big[C_{M_1}(x)^3-C_m(x)^3\big]<0$, with limit $-\frac{N+2}{\pi^2}\ln\frac m{M_1}$. These give the one-loop form of \cref{lem:s14:CW}(b)(ii) and, at $e=0$, the exact limit of (d); its sign analysis ($M_1\mapsto g_\lambda(M_1)$ is strictly concave with $g_\lambda(m)=0$, so some $M_1<m$ has $g_\lambda(M_1)>0$ if and only if $g_\lambda'(m)<0$) gives the threshold in (d). The two-loop coefficients are bounded below for every $k$ with $\eps_k\le1/m$: at $\lmu=m$ by the four upper sums of the two-loop certificate of \cite{SymPhase} (the lemmas of its Appendix~G1), stated there on tori and carried to $\Lk$ at fixed $\eps$ by the convergence of the periodic propagators to the infinite-lattice ones, which give $e^2\lambda S_{21,k}+e^4S_{40,k}\ge-\frac{9.330427}2\,e^2\lambda-\frac{1.703784}2\,e^4$ in units $m=1$, hence (a); for $\lmu<m$ by the same majorants re-evaluated with the photon propagator at mass $\lmu$ and the torus floor removed, in interval arithmetic over a partition of $(0,m]$, which gives (c); and by $\lmu$-free closed forms of the same majorants, whose only $\lmu$-dependence is the zero-momentum photon propagator and produces the $1/\hat\mu$ term, which gives (b). In (a)--(c) the lower bound is concave in $\lambda$ at fixed $e^2$ and in $e^2$ at fixed $\lambda$, so its minimum over a box is attained at a corner; this locates the boxes and the segments. The two-sided clause follows as in \cref{lem:s14:CW}(b)(i), the relative limit existing under \sxivCT.
\end{proof}

\begin{remark}[the quantifier over the couplings]\label{rem:s14:quantifier}
The theorems of this paper print the sources' quantifier of \cref{lem:s14:CW}(e), in the units of the coarsest spacing, with the cube family $\CJ$ not rescaled and $\beta=5^{-J_{\rm K}}$. A sentence ``$e,\lambda$ sufficiently small, $m,\lmu>0$ fixed'' describes \emph{where} the smallness lives (the running couplings, through Ba{\l}aban's $\eps_0$), but as a hypothesis on the bare couplings it is strictly stronger than what the sources need. (i) All four couplings have mass dimension one, so its dimensionless content is a bound on $e^2/m$ and $\lambda/m$, while the sources' hypothesis on the bare couplings is empty beyond positivity: the point $(e^2,\lambda,m,\lmu)=(100,100,1,1)$ is admitted by every source (with a large $J_{\rm K}$) and excluded by the phrase in every system of units. (ii) In the units in which the stopping scale is $1$, $m$ and $\lmu$ are not fixed but rescaled to $m\beta$ and $\lmu\beta$, of which $m\beta$ is capped (by the mass cap of \cite[Prop.~3.1]{Ba2} and by a further cap of the schedule) and $\lmu\beta$ is not: in those units the theorem reads ``three couplings small, the fourth rescaled''. (iii) Even at a fixed stopping scale $\beta_0$ the admissible set is not a product neighbourhood of $(0,0)$ in $(e,\lambda)$: one cap of the schedule of constants of the rotation part has the form $e\le\lambda^{1/4}c(\beta_0)$, a cusp that fails as $\lambda\to0$ at fixed $e$. So the phrase names no set exactly, and the sources' form survives because $\beta$ can be taken to $0$.
\end{remark}

% s23 Open problems
\section{Open problems}

Throughout, $K$ denotes King's printed state of \cref{thm:s13:main}: the Neumann--Dirichlet cube limit on King's RG cube at block factor $5$, fixed photon mass $\lmu>0$, with an admissible scheme $a$ in the Wick-square sector; $S$ denotes a torus limit state of the authors' earlier paper \cite{SymPhase} (block factor $2$, torus counterterm scheme with reference photon mass $m$, $0<\lambda,e^2<0.1m$, $\lmu\in(0,m]$). No identification between the two is proved. ``Symmetric phase'' is, for both, the expectation that the states lie above the surface predicted by one-loop perturbation theory; it is neither a hypothesis nor a conclusion of any theorem here, and \cref{thm:s13:main} holds on all of King's coupling range (\cref{lem:s14:CW}), the surface included. The problems below are extensions of the main theorem, not conditions on it.

\subsection{Relation to the torus states}

Three hypotheses bear on the relation of $K$ to $S$: $(\mathrm{FAM}_L)$, King's cube construction exists on the block-factor-$2$ family in the torus counterterm scheme; $(\mathrm{KI}_L)$, the fixed-torus continuum limit exists in that family; $(\mathrm{CE})$, a convergent unit-scale cluster expansion with sources on the cubic tori, their tilted copies and their cuts along King's cubes, uniform in the lattice spacing. Under $(\mathrm{FAM}_L)$ and a sign or decay hypothesis on the difference of the two counterterms, every torus state dominates the King state at a matched bare mass on the observables $e^{\mathrm iF(f)+\Wick{k}}$, $k\ge0$, a one-sided inequality carrying no information about rotations. Under $(\mathrm{CE})$ at a point of a massive window ($\lambda m^2\ge K_0e^6$, $K_0\ge1/(18\pi^2)$, small $e^2,\lambda$) every torus limit state equals the cube limit of the same family and scheme in every order of limits; with $(\mathrm{FAM}_L)$ or $(\mathrm{KI}_L)$ it is unique, exponentially clustering and ergodic; with King's cube limits on both families it equals $K$ at an explicitly shifted bare mass. Whether $(\mathrm{CE})$ holds anywhere at $e>0$ is open. Two obstacles are structural: the torus paper computes its mass counterterm with the photon propagator frozen at $m$ while King's free-boundary counterterm uses the model's own $\lmu$, a difference of order $e^2|m-\lmu|$, not small (\cref{rem:s14:torus-window}); and King's construction needs block factor $5$ for the rotation by $\arctan(3/4)$, while no lattice family admits both a dyadic structure and a rotation of infinite order. On the one-loop surface $\lambda m^2=e^6/(72\pi^2)(1+O(e^2/m))$ the two constructions are expected to select different phases (periodic tori the broken one, Dirichlet faces the symmetric one); this is a heuristic, not a theorem, and no statement of this paper uses the phase.

\subsection{Uniqueness and ergodicity of \texorpdfstring{$K$}{K}}

Uniqueness of King's cube limit is part of the construction (monotonicity in the cube), but uniqueness of the infinite-volume state, in the sense of King's expected identity between cube limits and periodic-box limits, is not proved, and neither is ergodicity of $K$ under translations: \cref{thm:s13:main} gives clustering only on the ergodic components. The usual route, a cluster expansion with a mass gap, is $(\mathrm{CE})$. This is distinct from \UNIQ\ (\cref{hyp:s12:UNIQ}), the independence of a LIM state from its continuum subsequence, which follows from \cref{hyp:s03:EC} and from \cref{hyp:s03:RC} (\cref{prop:s12:ID}, \cref{cor:s04:G2}). Cubic symmetry of $K$ is a theorem (\cref{thm:s06:S1K}): it needs only the uniqueness of the monotone cube limit and the reflection sandwich (\cref{rem:s06:K4}).

\subsection{Non-triviality}

For the torus states the truncated correlation of $F(g)^2$ with the Wick square is non-positive, and strictly negative at small couplings under the two hypotheses of \cite[thm:NT-cond]{SymPhase}: the one-loop statement (C) on the continuum coefficient and a remainder bound on it, which is unproved and, as printed, is expected to fail below the first-order surface. The analogous statements for King's cube, and a non-vanishing two-point function, are left to later work. For $K$ nothing is proved here beyond the sign inequalities that hold on every lattice. Whether $K$ is non-Gaussian, or even not decoupled, is open.

\subsection{The limit of vanishing photon mass}

Every constant of the Wick leg of \cref{thm:s13:main} depends on $\lmu$ and is not uniform in it (\cref{rem:s13:mu0}); the absolute bounds of the curvature leg are. The limit $\lmu\downarrow0$, taken along a frozen family (counterterm and scheme frozen at a reference mass, so that the members are monotone in $\mu$ on the cone by the correlation inequalities), is left to later work: in the curvature sector it needs the existence theorem for every member and nothing else; in the Wick sector it needs a cell bound uniform in the photon mass, together with the torus inputs at the reference mass and the transfer of \cref{thm:s13:TWK} for the members on King's convention; the identification of the members with King's models at shifted scalar masses, masslessness proper (the vanishing of the photon mass in the limit theory) and the two-point function are further questions there. None of these enters the present paper.

\subsection{The optional extension: \texorpdfstring{\Wtwo}{(W2)} at arbitrary amplitudes, Form (A), and the Euclidean clause of the Wick sector}

Form~(B) of \cref{thm:s13:main} lives on the tube $\Dtube$: the only positive-source bound it uses is the one proved on the cone $\sup q\le s_1<\sigma_1(M_1)$ (\cref{cor:s11:W2prime}), and the only existence statement is \cref{hyp:s03:W1}, whose positive scalar part is below the same threshold. Above the threshold, the existence of the fixed-cube limits and the uniform positive-source bound are the hypothesis \Wtwo\ (\cref{hyp:s03:W2}), which is open.

\begin{theorem}[Form (A): the Wick leg at arbitrary amplitudes]\label{thm:s23:FormA}
Let the couplings satisfy \Tthree, and let $a$ be an admissible scheme satisfying \Wtwo\ for the comparison convention $\NDa$. Then for every finite family $(f_1,\dots,f_n;h_1,\dots,h_p)$ of real $2$-forms and real scalars the full object $Z^a_\infty(\zeta)=\lim_J\lim_kZ^a_{k,J}(\sum_iz_if_i,\sum_jw_jh_j)$ exists for every $\zeta=(z,w)\in\C^{n+p}$, locally uniformly, $Z^a_\infty$ is entire, and $Z^a_\infty=\Ginf$ on $\Dtube$. Its Schwinger functions are those of Form~(B), and:
\begin{enumerate}[label=\textup{(\roman*)},leftmargin=*,nosep]
\item \OSz\ as in Form~(B)(i): the tube bound on $\Dtube$ and the bounds on the $S_{n,p}$; no quantitative bound on $Z^a_\infty$ off the tube is asserted.
\item \OSo: $Z^a_\infty(\tau_bR\cdot\zeta)=Z^a_\infty(\zeta)$ for every $\zeta\in\C^{n+p}$, $b\in\R^3$, $R\in\Oh$ ($F$ a $2$-form, $W^a$ a scalar).
\item \OSt\ in matrix form about every coordinate plane for exponentials with arbitrary complex coefficients (\cref{thm:s23:OS2W}), and in polynomial form (\cref{prop:s08:poly}).
\item \OSth\ as in Form~(B)(iv).
\item[\textup{(v),(vi)}] As in Form~(B), the characteristic functional $(u,v)\mapsto Z^a_\infty(\mathrm iu,\mathrm iv)$ being the restriction of the entire $Z^a_\infty$ to the imaginary axes.
\end{enumerate}
\emph{Euclidean clause.} If in addition \RW\ holds for the comparison family, $Z^a_\infty$ is invariant under every Euclidean motion, its mixed Schwinger functions are $\Euc$-covariant, and $\nu$-almost every ergodic component of the measure of (v) is $\Euc$-invariant and reflection positive about every plane (no vanishing of odd moments is asserted for a component). This is the only Euclidean clause of the Wick sector.

\emph{On King's convention.} The same statements for the K-family object $Z^{a,K}_\infty$ under \Wtwo\ read for $\NDK$ are the content of \FAK\ (\cref{hyp:s03:FAK}); nothing is derived for it from the comparison-family statements, since \Wtwo\ on $\NDK$ does not give \Wtwo\ for the comparison family (\cref{cor:s04:G2}(iv)).
\end{theorem}

\begin{proof}[Indications]
Existence at all complex coefficients, locally uniform convergence and entireness are \cref{lem:s08:W1} under \Wtwo\ (version (A$_\infty$) of \cref{lem:s08:claim4}); $Z^a_\infty=\Ginf$ on $\Dtube$ as the same $J$-limit of the same $Z^a_J$ (\cref{thm:s11:W}). (i) is Form~(B)(i). (ii): Form~(B)(ii) is the identity on $\Dtube$, an open neighbourhood of the origin of $\C^{n+p}$, and both sides are entire. (iii) is \cref{thm:s23:OS2W} with \cref{prop:s08:poly}; (iv) is \cref{prop:s08:OS3} for $Z^a_\infty$; (v),(vi): the measure of Form~(B) is built from $\Ginf$ on the imaginary axes, which lie in $\Dtube$. The Euclidean clause is \RW\ itself; the component clause follows through the ergodic decomposition of Form~(B)(vi) as in the curvature sector (\cref{thm:s09:S2K}(c)).
\end{proof}

For a LIM state of a continuum subsequence (\cref{thm:s12:LIM}) the corresponding Form~(A) requires both clauses of \Wtwo: the existence of the limits at every real coefficient vector and the uniform bound at every amplitude.

\begin{theorem}[reflection positivity at arbitrary complex coefficients]\label{thm:s23:OS2W}
Let $a$ be an admissible scheme satisfying \Wtwo, let $\nu\in\{1,2,3\}$, $c\in\R$, and let $\zeta_1,\dots,\zeta_N$ be finite families of real sources $f_{il}\in C^\infty_c(\R^3;\Lambda^2)$, $h_{il}\in C^\infty_c(\R^3)$ supported in $\{x_\nu>c\}$, with arbitrary complex coefficients $z_{il},w_{il}$. Then
\[
\sum_{i,j=1}^N\bar\gamma_i\gamma_j\,Z^a_\infty\big(\theta\zeta_i^*\sqcup\zeta_j\big)\ \ge\ 0\qquad\text{for all }\gamma\in\C^N .
\]
\emph{Refinement.} Let $B_1$, $B''$, the vertex sources $q_v$ ($v\in\vertex(B'')$) and $q_\zeta$ be those of \cref{def:s08:enlarged} for the plane $\{x_\nu=c\}$. (a) If $c\in D$, the proof uses \Wtwo(b) only at the finitely many sources $q_v$, and it suffices to assume it at the single source $(1+\eta)q_\zeta$ for any one $\eta>0$ (\cref{rem:s08:oneinstance}, one direction only); all other steps use \Wtwo(a) and the inequalities of \cref{lem:s02:BFS}, and the tube machinery is not needed. (b) If $c\notin D$, with $c'\in D$ and $b\defeq(c'-c)e_\nu$: the finitely many instances $q_v$ and $\tau_bq_v$, $v\in\vertex(B'')$, together with the tube inputs \Itwo, \Ithree\ of \cref{thm:s08:OS2WT}. Such a plane needs a non-lattice translation, which finitely many instances of \Wtwo(b) do not supply; it is carried by the translation invariance of the tube object (\cref{prop:s08:T}(b)).
\end{theorem}

\begin{proof}
Let $c\in D$ and let $\cF$, $\Phi_k$, $\zeta^*_{ij}$, $\cT_1$, $T_\cF$, $U$ be as in \cref{def:s08:enlarged}. By \cref{lem:s08:W1}, under \Wtwo(a) and \Wtwo(b) at the $q_v$, $\{Z^a_J\}_{J\ge J_0}$ is uniformly bounded on $\cT_1$, $Z^a_J\to Z^a_\infty$ locally uniformly on $\cT_1$, $Z^a_\infty$ is holomorphic there and $Z^a_J\le Z^a_\infty$ on $U\cap\cT_1$. By \cref{lem:s08:claim4}(A), $\limsup_k\sup_{\cT_1}|\Phi_k|\le C$ with $C$ independent of $J\ge J_0$, and at fixed $J$ the $\Phi_k$ are locally bounded on $\C^m$ (\Wtwo(a) only). Montel's theorem gives a subsequence along which $\Phi_k\to\Phi^J$ locally uniformly, $\Phi^J$ entire; the matrices $(\Phi_k(\zeta^*_{ij}))_{ij}$ are positive semidefinite for large $k$ (\cref{lem:s07:RPbox}(b)), hence so is their limit, and $\sup_{\cT_1}|\Phi^J|\le C$ uniformly in $J$. On $U$ every lattice value is real and the scheme multipliers of the two states agree, so the $\cI$ halves of \cref{lem:s07:L,lem:s07:U} give $Z^a_{k,J}\le\Phi_k\le Z^a_{k,J+2n}$, hence $Z^a_J\le\Phi^J\le Z^a_{J+2n}$ on $U$. On $U\cap\cT_1$, $Z^a_J\uparrow Z^a_\infty$ (\cref{lem:s08:W1}), so $\Phi^J\to Z^a_\infty$ there; by the uniform bound on $\cT_1$, Montel's theorem and the uniqueness set $U\cap\cT_1$ (\cref{def:s08:enlarged}), $\Phi^J\to Z^a_\infty$ locally uniformly on the connected $\cT_1$, in particular at the $\zeta^*_{ij}$, and positivity passes to the limit. The only instances of \Wtwo(b) used are the $q_v$ of \cref{lem:s08:W1,lem:s08:claim4}; \cref{rem:s08:oneinstance} replaces them by $(1+\eta)q_\zeta$.

Let $c\notin D$, $c'\in D$, $b\defeq(c'-c)e_\nu$. The $\tau_b\zeta_i$ are supported in $\{x_\nu>c'\}$ and $\theta_{c'}(\tau_b\zeta)^*=\tau_b(\theta_c\zeta^*)$, so the display at $c$ follows from the display at $c'$ for the translated families once $Z^a_\infty$ is translation invariant at the points in question. The enlarged family of the translated problem at $c'$ is $\tau_b\cF$, with the same coefficient space, matrix points, $B_1$, $B''$, and vertex sources $\tau_bq_v$; the argument above for $\tau_b\cF$ gives the positivity of $(Z^a_\infty(\tau_b\cdot\zeta^*_{ij}))_{ij}$, and $\zeta\mapsto Z^a_\infty(\tau_b\cdot\zeta)$, $\zeta\mapsto Z^a_\infty(\zeta)$ are holomorphic on $\cT_1$ by \cref{lem:s08:W1}. On the non-empty open set $\cT_1\cap T_\cF$ both equal the tube limits $Z^a_{\infty,T}(\tau_b\cdot\zeta)$, $Z^a_{\infty,T}(\zeta)$, which agree by \cref{prop:s08:T}(b); by the identity theorem they agree on $\cT_1$, in particular at the $\zeta^*_{ij}$.
\end{proof}

\begin{hypothesis}[\Bminus{u}: stability of the shifted torus model]\label{hyp:s23:Bminus}
Let $u>0$. In the free convention and the scheme $a$,
\[
\sup_{k\ge0}\ \limsup_{J\to\infty}\ \alpha^a_{k,T_J}(u)<\infty,\qquad\alpha^a_{k,T}(u)=|T|^{-1}\ln N^a_{k,T}(0,u\mathbf 1),
\]
the free-energy density of the torus functional at the constant scalar source $u$; that is, \Itwo\ of \cref{def:s10:I2} at the amplitude $s_1=u$.
\end{hypothesis}

For $u$ below the threshold $s_*$ of \cref{lem:s14:CW}, \Bminus{u} is \cref{prop:s10:U} (or \cref{prop:s10:Upm}) at an auxiliary mass $M_1$ with $\sigma_1(M_1)>u$, so Form~(B), whose constants are built from \Itwo\ at one such amplitude, already contains every amplitude at which the positive-source bound is known. For $u\ge s_*$ the hypothesis is the exact obstruction. By \cref{lem:s10:M}, exact at fixed $(k,T)$, $\alpha^z_{k,T}(u)$ is, up to explicit $c$-numbers, $|T|^{-1}$ times the difference of the free energies of King's torus model with bare scalar coefficient $\alpha^{\rm tor}_{k,T}(m)-2u$ and with $\alpha^{\rm tor}_{k,T}(m)$, both with the counterterms of the mass $m$: the second is \cref{hyp:s03:T1} at $m$, the first the ultraviolet stability of the shifted model. For $u<s_*$ the shifted coefficient is that of King's model at a positive mass $M_1$; for $u\ge s_*$ it lies, asymptotically in $k$, below $\inf_{M>0}\alpha^\infty_k(M)$, and no positive scalar mass realises it. The printed proofs of \cite{Ba1,Ba2} expand about the Gaussian at the positive scalar mass $m$ and use its mass term in the small-field bounds; no printed theorem covers a bare coefficient below every positive-mass value. Conversely \Bminus{u} is sufficient for \Wtwo(b): \cref{prop:s11:W1prime} at $s_1=u$, whose steps do not depend on the amplitude, gives $\ln N^a_{k,X}(0,q)\le(C_U/u)\norm q_{1,k}$ for $0\le q\le u$ on every Neumann--Dirichlet box, hence \Wtwo(b) in $L^1$ form on $\sup q\le u$. It is not asserted that \Wtwo\ implies \Bminus{u}: the inequalities available bound a constant-source density from below by local sources, never from above, so \Bminus{u} is sufficient and not shown necessary. Also open: an $L^1$ bound for the negative part of the Wick source on cubes, as on the torus paper's tori (the tube bound \eqref{eq:s11:W} carries the weighted norm $\wnorm{k_-}$).

\subsection{Rotation invariance of the Wick sector}

\Cref{thm:s13:main} gives invariance under $\Euc$ for the curvature sector; for the Wick sector it gives invariance under $\Cub$ on the tube and no Euclidean clause. Euclidean invariance of the joint field is the hypothesis \RW\ of the optional extension (\cref{thm:s23:FormA}), and on King's convention \FAK; whether the heavy-shell comparison of the rotation part extends to the Wick square, on the tube or at arbitrary amplitudes, is open. Cubic symmetry is a theorem in every case.

\subsection{The coupling condition}

The sign clause \Tthree\ is proved on the explicit regions of \cref{prop:s14:T3} (at $\lmu=m$ on a region containing $0\le\lambda,e^2\le0.1m$; at every $\lmu\in(0,m]$ on explicit boxes; exactly for $\lambda<\frac\pi4(1+\sqrt5)\,m$ at $e=0$), and the Wick leg is stated there. Outside those regions the exact position of the failure set $E^<_W$ is not determined; the two-sided clause that $m$ is not a global minimiser of $\alpha_\infty$, with the auxiliary mass on either side of $m$ (\cref{prop:s10:Upm}, \cref{lem:s14:CW}(b)), would enlarge the region. The three hypotheses \Bcell, \SCH, \UNIQ\ of the statements on continuum subsequences (\cref{hyp:s12:Bcell,hyp:s12:SCH,hyp:s12:UNIQ}) are proved for the objects of the main theorem (\cref{prop:s12:ID}) and are open only for LIM states of schemes that are not admissible. The mass clause of \cite[Prop.~3.1]{Ba2}, a cap on the dimensionless scalar mass at the stopping scales where that proposition is used, belongs to the statement of the proposition, which is cited under it, and imposes nothing on the bare mass (\cref{rem:s14:BALm}).


% biblio.tex -- thebibliography block, input at the end of main.tex (before \end{document}).
% Keys are fixed; writers cite only these keys. No names of the present authors anywhere.
\begin{thebibliography}{FILS78}

\bibitem{Ba1} T.~Ba{\l}aban, $(\mathrm{Higgs})_{2,3}$ quantum fields in a finite volume. I. A lower bound, Commun.\ Math.\ Phys.\ \textbf{85} (1982), 603--626. DOI 10.1007/BF01403506.

\bibitem{Ba2} T.~Ba{\l}aban, $(\mathrm{Higgs})_{2,3}$ quantum fields in a finite volume. II. An upper bound, Commun.\ Math.\ Phys.\ \textbf{86} (1982), 555--594. DOI 10.1007/BF01214890.

\bibitem{Ba3} T.~Ba{\l}aban, $(\mathrm{Higgs})_{2,3}$ quantum fields in a finite volume. III. Renormalization, Commun.\ Math.\ Phys.\ \textbf{88} (1983), 411--445. DOI 10.1007/BF01213217.

\bibitem{Ba4} T.~Ba{\l}aban, Regularity and decay of lattice Green's functions, Commun.\ Math.\ Phys.\ \textbf{89} (1983), 571--597. DOI 10.1007/BF01214744.

\bibitem{BIJ88} T.~Ba{\l}aban, J.~Imbrie, A.~Jaffe, Effective action and cluster properties of the abelian Higgs model, Commun.\ Math.\ Phys.\ \textbf{114} (1988), 257--315. DOI 10.1007/BF01225037.

\bibitem{BFS1} D.~Brydges, J.~Fr\"ohlich, E.~Seiler, On the construction of quantized gauge fields. I. General results, Ann.\ Physics \textbf{121} (1979), 227--284.

\bibitem{Bil99} P.~Billingsley, \emph{Convergence of Probability Measures}, 2nd ed., Wiley, New York, 1999.

\bibitem{FILS78} J.~Fr\"ohlich, R.~Israel, E.~H.~Lieb, B.~Simon, Phase transitions and reflection positivity. I. General theory and long range lattice models, Commun.\ Math.\ Phys.\ \textbf{62} (1978), 1--34.

\bibitem{GJ} J.~Glimm, A.~Jaffe, \emph{Quantum Physics. A Functional Integral Point of View}, 2nd ed., Springer, New York, 1987.

\bibitem{Hor90} L.~H\"ormander, \emph{An Introduction to Complex Analysis in Several Variables}, 3rd ed., North-Holland, Amsterdam, 1990.

\bibitem{Jan97} S.~Janson, \emph{Gaussian Hilbert Spaces}, Cambridge Tracts in Mathematics \textbf{129}, Cambridge University Press, 1997.

\bibitem{Kec95} A.~S.~Kechris, \emph{Classical Descriptive Set Theory}, Graduate Texts in Mathematics \textbf{156}, Springer, New York, 1995.

\bibitem{K1} C.~King, The $U(1)$ Higgs model. I. The continuum limit, Commun.\ Math.\ Phys.\ \textbf{102} (1986), 649--677. DOI 10.1007/BF01221651.

\bibitem{K2} C.~King, The $U(1)$ Higgs model. II. The infinite volume limit, Commun.\ Math.\ Phys.\ \textbf{103} (1986), 323--349. DOI 10.1007/BF01206942.

\bibitem{LL10} G.~F.~Lawler, V.~Limic, \emph{Random Walk: A Modern Introduction}, Cambridge Studies in Advanced Mathematics \textbf{123}, Cambridge University Press, Cambridge, 2010. DOI 10.1017/CBO9780511750854.

\bibitem{Luk70} E.~Lukacs, \emph{Characteristic Functions}, 2nd ed., Griffin, London, 1970.

\bibitem{OS73} K.~Osterwalder, R.~Schrader, Axioms for Euclidean Green's functions, Commun.\ Math.\ Phys.\ \textbf{31} (1973), 83--112.

\bibitem{OS75} K.~Osterwalder, R.~Schrader, Axioms for Euclidean Green's functions. II, Commun.\ Math.\ Phys.\ \textbf{42} (1975), 281--305.

\bibitem{Ran95} T.~Ransford, \emph{Potential Theory in the Complex Plane}, London Mathematical Society Student Texts \textbf{28}, Cambridge University Press, Cambridge, 1995. DOI 10.1017/CBO9780511623776.

\bibitem{Rei88} T.~Reisz, A power counting theorem for Feynman integrals on the lattice, Commun.\ Math.\ Phys.\ \textbf{116} (1988), 81--126. DOI 10.1007/BF01239027.

\bibitem{Rei88b} T.~Reisz, Renormalization of Feynman integrals on the lattice, Commun.\ Math.\ Phys.\ \textbf{117} (1988), 79--108. DOI 10.1007/BF01228412.

\bibitem{RS1} M.~Reed, B.~Simon, \emph{Methods of Modern Mathematical Physics. I. Functional Analysis}, revised and enlarged ed., Academic Press, New York, 1980.

\bibitem{RotO3} Rotation invariance of the three-dimensional abelian Higgs model: the curvature sector, companion paper (2026).

\bibitem{Sei82} E.~Seiler, \emph{Gauge Theories as a Problem of Constructive Quantum Field Theory and Statistical Mechanics}, Lecture Notes in Physics \textbf{159}, Springer, Berlin, 1982.

\bibitem{SymPhase} The symmetric phase of the three-dimensional abelian Higgs model: torus limits and the Osterwalder--Schrader axioms with cubic symmetry, companion paper (2026).

\end{thebibliography}


\end{document}
